\documentclass[11pt]{article}
\usepackage[margin=1in]{geometry}
\usepackage[T1]{fontenc}
\usepackage[utf8]{inputenc}
\usepackage{textcomp}
\usepackage{lmodern}
\usepackage{microtype}
\usepackage{xcolor}
\usepackage{amsmath,amssymb,amsthm,mathtools}
\usepackage{enumitem}
\usepackage{setspace}
\usepackage{xurl}
\usepackage[hidelinks]{hyperref}
\providecommand{\tightlist}{\setlength{\itemsep}{0pt}\setlength{\parskip}{0pt}}
\setstretch{1.18}
\setlength{\parindent}{0pt}
\setlength{\parskip}{6pt plus 2pt minus 1pt}
\setlist{itemsep=2pt,topsep=4pt}
\setcounter{secnumdepth}{-\maxdimen}
\emergencystretch=3em
\sloppy
\begin{document}
\section*{On the invariant theory of forms in $n$ variables}

In the projective invariant theory of forms in $n$ variables, the principal problem has been solved: the finiteness of the system of forms has been proved. A number of further questions, however, have not been treated, namely those that concern methods for investigating the relations among forms. I should like briefly to communicate the results that I have obtained with respect to such questions, but for the sake of clarity I shall first sketch the questions in the ternary case, where they are already known.

I think first of the two fundamental theorems of the symbolic method: the theorem that all invariant formations can be represented symbolically, and the second theorem that all relations among invariants are obtained by successive application of a small number of symbolic identities; in other words, that the symbolic method gives all relations without leaving the domain of invariants.

Built upon these are the processes for generating forms: the symbolic contraction process, the non-symbolic differentiation processes, the theory of expansions into series, and so on. In this connection I should also mention a theorem proved in my dissertation: all contractions can be generated by successive application of the two possible \emph{cogredient} contractions; by this I mean, for a symbolic product $s,t,u,v$, the two contractions $(stu)$ and $(vx)$. Upon this theorem one can build the theory of reducentia and the reduction of systems of forms, analogously to the binary case, as I showed in my dissertation.

When we pass to the $n$-ary case, already the first theorem of the symbolic method holds only in a qualified sense. As is well known, for quaternary forms there already appear, besides point and plane coordinates $x$ and $u$, line coordinates $p_\rho$, namely the minors obtained from two rows of point coordinates. Analogously, in the $n$-ary case we have rows of variables $p_\rho$ whose individual elements are the minors taken from $\rho$ rows $x$, and symbolic rows $S_\rho$ whose elements behave like minors taken from $\rho$ rows $s$, the latter being cogredient to the $u$.

The symbolic representation by determinants is known to hold only when the symbolic rows $S_\rho$ are resolved into the rows $s$ and these are distributed among different determinants; then only those aggregates of determinants have a real meaning that depend on the $S_\rho$. Which determinants depend only on the $S_\rho$ can, to be sure, be recognized by differential equations; but in that way one obtains no survey of the totality of invariant formations, nor of the processes by which they are generated.

At this point the theorem on corresponding matrices becomes decisive. To every matrix product
\[
(v^{(1)}\cdots v^{(\rho)} \mid y^{(1)}\cdots y^{(\rho)}),
\]
where the rows of the second matrix are contragredient to those of the first, there corresponds a determinant; and conversely every determinant can be transformed into a matrix product with a prescribed number $\rho$ of rows. Hence determinant and matrix product appear as completely equivalent, and the first theorem of the symbolic method takes the following form:

\begin{quote}
All invariant formations can be represented symbolically by matrix products.
\end{quote}

From two symbolic rows $S_\rho,T_\tau$ one can now, with the help of rows of variables, form very easily a matrix product that contains these symbolic rows explicitly and therefore represents an invariant formation of the type $(S_\rho\mid p_\rho)(T_\tau\mid p_\tau)$, namely
\[
(G_{n-\rho-\lambda}S_\rho T_\tau \mid p_{\rho+\lambda}p_{\tau-\lambda}),
\qquad
\lambda = 1,2,\dots,\tau \text{ or } n-\rho.
\tag{1}
\]
More important is the converse statement: all invariant formations can be represented explicitly by matrix products, so that the above process is the genuine extension of the binary and ternary contraction process. To prove this converse it is necessary to transfer the second theorem of the symbolic method to the $n$-ary case, namely to determine the totality of the independent indecomposable identities in which all rows $S_\rho,p_\rho$ are treated as indecomposable. These identities arise from the known identities containing only rows $x$ and $u$ by replacing the product theorem for determinants by a system of product theorems for matrices and then combining several matrix products into a single one by means of these same product theorems. In this way one arrives at two dual formulas, of which I cite only one:
\[
\sum_{i=1}^{k} (S_{\rho_1}\cdots S_{\rho_k}\mid p_\rho x^{(i)})(S_\rho\mid x^{(i)})
=
(S_{\rho_1}\cdots S_{\rho_k}S_\rho\mid p_\rho).
\tag{2}
\]

The only specialization contained in these formulas is that a row $x$ appears. For completely general symbolic and variable rows one reaches a \emph{decomposition identity}, that is, an identity in which a row $p_\rho$ appears as though resolved into the individual rows $x$. Clebsch already used the simplest special case of the decomposition identity in his ``fundamental problem of invariant theory'' in order to derive the elementary covariants occurring in series expansions.

With the aid of these identities, which mediate the passage from the non-explicit determinant expressions to matrix products, one can in fact prove that matrix representation yields the desired explicit symbolic representation.

The same formalism then gives the differentiation processes --- for a row of variables gives rise to a row of symbols by replacing its entries with differential symbols --- and by means of the decomposition identity it becomes possible to transfer Mertens's proof almost literally to the $n$-ary case.

\clearpage
\section*{The finiteness theorem for the invariants of finite groups}

Let the finite group $G$ consist of the $h$ linear transformations (with nonvanishing determinant)
\[
A_1,\dots,A_h,
\]
where $A_k$ is represented by the linear transformation
\[
x^{(k)}_\nu=\sum_{\mu=1}^n a^{(k)}_{\nu\mu}x_\mu
\qquad (\nu=1,\dots,n),
\]
or, more briefly, by
\[
(x^{(k)})=A_k(x).
\]
Thus the group $G$ carries the row $(x)$ with elements $x_1,\dots,x_n$ into the rows $(x^{(k)})$ with elements $x^{(k)}_1,\dots,x^{(k)}_n$. Since the identity must occur among $A_1,\dots,A_h$, the row $(x)$ itself also occurs among the rows $(x^{(k)})$.

By an \emph{integral rational (absolute) invariant} of the group I mean an integral rational function of $x_1,\dots,x_n$ that remains identically unchanged under application of $A_1,\dots,A_h$; thus for such an invariant $f(x)$ one has
\[
f(x)=f(x^{(1)})=f(x^{(2)})=\cdots=f(x^{(h)}).
\tag{1}
\]

\subsection*{1. A full invariant system from the Galois resolvent}

Formula \textup{(1)} expresses that $f(x)$ is an integral rational symmetric function of the rows $(x^{(1)}),\dots,(x^{(h)})$ --- in fact, since each summand $f(x)$ contains only the single row $(x)$, this is the simplest case, usually called the \emph{uniform} case. By the known theorem on symmetric functions of rows, $f(x)$ can therefore be expressed integrally and rationally through the elementary symmetric functions of these rows, that is, through the coefficients $G_{\alpha_1\dots\alpha_n}(x)$ of the ``Galois resolvent''
\[
\Phi(z,u)=\prod_{k=1}^h\bigl(z-(u_1x^{(k)}_1+u_2x^{(k)}_2+\cdots+u_nx^{(k)}_n)\bigr)
      = z^h+\sum G_{\alpha_1\dots\alpha_n}(x)\,u_1^{\alpha_1}\cdots u_n^{\alpha_n} z^{\alpha_0},
\tag{2}
\]
where the $G_{\alpha_1\dots\alpha_n}(x)$ are invariants of degree $\alpha_1+\cdots+\alpha_n$ in the variables $x$.

Thus we have proved:

\medskip
\noindent
\textbf{Theorem.}
\emph{The coefficients $G_{\alpha_1\dots\alpha_n}(x)$ of the Galois resolvent form a complete system of invariants of the group, in such a way that every invariant can be expressed integrally and rationally in terms of these finitely many invariants.}

\subsection*{2. A second complete system}

Starting from \textup{(1)} one can also make the following still more elementary observation, which at the same time leads to a second complete system.

Suppose
\[
f(x)=a+b\,x_\nu+\cdots+c\,x_\nu^m+\cdots,
\]
where $a,b,\dots,c$ are constants. Then from \textup{(1)} one has
\[
h\cdot f(x)
 = h\cdot a
   + b\sum_{k=1}^h x^{(k)}_\nu
   + \cdots
   + c\sum_{k=1}^h \bigl(x^{(k)}_\nu\bigr)^m
   + \cdots.
\]
Hence every invariant can be assembled integrally and linearly from the special quantities
\[
s_m=\sum_{k=1}^h \bigl(x^{(k)}_\nu\bigr)^m
\qquad (m=0,1,2,\dots),
\]
so it is enough to prove the theorem for these special functions.

Up to a numerical factor, $s_m$ is the coefficient of $u_1^{\alpha_1}\cdots u_n^{\alpha_n}$ in the expression
\[
h\cdot s_m=\sum_{k=1}^h\bigl(u_1x^{(k)}_1+\cdots+u_nx^{(k)}_n\bigr)^m,
\]
where $u=u_1+\cdots+u_n$ is written only as a convenient abbreviation, and this expression represents the $u$-power sum of the $h$ linear forms.

Now it is well known that the infinitely many power sums $S_m$ are integral rational functions of $S_1,S_2,\dots,S_h$; their coefficients are determined by the invariants
\[
S_m=\sum_{k=1}^h\bigl(x^{(k)}_1+\cdots+x^{(k)}_n\bigr)^m
\qquad (m\le h).
\]
Thus we obtain a second complete system:

\medskip
\noindent
\textbf{Theorem.}
\emph{A complete invariant system of the group is given by all invariants $J_{\alpha_1\dots\alpha_n}$ for which}
\[
\alpha_1+\cdots+\alpha_n\le h,
\]
\emph{where $h$ is the order of the group.}

The connection with part~1 is given by the remark that the power sums $S_1,\dots,S_h$ may be replaced by the elementary symmetric functions of $x^{(k)}_1,\dots,x^{(k)}_n$, which leads back to the complete system there given by the $G_{\alpha_1\dots\alpha_n}(x)$.

Both results show that all invariants can be expressed integrally and rationally through those whose degree in the $x$ does not exceed the order $h$ of the group.

\subsection*{3. Rational representation}

From what has just been proved one obtains consequences for rational representation. Every rational absolute invariant can --- as is seen immediately by multiplying numerator and denominator by the conjugates of the denominator with respect to the group --- be represented as the quotient of two integral rational absolute invariants, not necessarily relatively prime. It follows that every rational absolute invariant can be expressed rationally through the coefficients $G_{\alpha_1\dots\alpha_n}(x)$ of the Galois resolvent $\Phi(z,u)$.

This theorem can also be proved in various easy ways without passing through the finiteness theorem; it is already formulated in Weber II, \S 58.

\subsection*{4. Final remark}

Finally, I should mention that the results derived here are already implicitly contained in my paper \emph{Körper und Systeme rationaler Funktionen}: namely, the complete system given in part~1 appears there in Theorems VI and VII, while an argument for rational representation independent of the present finiteness theorem appears there in Theorem III. Both the proof and the result become, however, considerably simpler in the special situation treated here than in the general theory.

I was led to apply those general investigations to the invariants of finite groups by conversations with E.~Fischer, to whom the proof given in part~2 is due in substance as well. In the general case the complete system arising there is not rationally definable, and the remark in part~3 likewise goes back to him.

\begin{flushright}
Erlangen, May 1915.
\end{flushright}

\clearpage
\section*{On integral rational representation of the invariants of a system of arbitrarily many basic forms}

At the end of his contribution to the Schwarz Festschrift, \emph{On the invariants of a system of arbitrarily many basic forms}, Hilbert expressed the conjecture that these invariants can be represented integrally and rationally by the finitely many invariants of the system \((J,PJ)\), where \(J\) denotes a complete system of invariants for the linearly independent basic forms and \(PJ\) denotes the invariants obtained from \(J\) by polar processes.

In what follows I show, in Part I, that this conjecture is indeed correct, and in fact as a simple consequence of the reduction theorem: every form in arbitrarily many rows of \(n\) variables can be represented as a sum of polars of forms that involve only \(n\) fixed rows.

In Part II I give a more conceptual proof of the reduction theorem by interpreting the question as a problem of equivalence of linear families of forms. I owe this point of view, and likewise an essential part of the reasoning employed, to Fischer's paper \emph{On algebraic systems of modules} and to unpublished propositions of E. Fischer, whose use he kindly permitted me. Finally, in Part III I indicate how analogous considerations also lead to the most general series expansions.

\section*{I. The invariant-theoretic consequence of the reduction theorem}

Let \(\Theta\) be a form homogeneous in each of the \(N>n\) rows
\[
A_1,A_2,\dots,A_N,
\]
where in general the row \(A_i\) consists of the \(n\) elements
\[
a_1^{(i)},\dots,a_n^{(i)}.
\]
The coefficients of \(\Theta\) may be indeterminates or quantities belonging to a given field of rationality. Let \(P\) denote an integral rational function---with rational numerical coefficients---of the polar operators
\[
P_{ih}
=
a_1^{(i)}\frac{\partial}{\partial a_1^{(h)}}
+
a_2^{(i)}\frac{\partial}{\partial a_2^{(h)}}
+\cdots+
a_n^{(i)}\frac{\partial}{\partial a_n^{(h)}},
\]
where by a product such as \(P_{ih}P_{jk}\Theta\) one understands the action of \(P_{ih}\) on \(P_{jk}\Theta\).

Then the reduction theorem referred to above is expressed by the identity
\[
\Theta = \sum P\Theta_0,
\tag{1}
\]
where the forms \(\Theta_0\) involve only the rows
\[
A_1,A_2,\dots,A_n
\]
and are obtained from \(\Theta\) by polar processes \(P\).

Now consider the basic system \((F)\) consisting of the \(N\) forms of the same order and the same number of variables,
\[
F_1,F_2,\dots,F_N,
\]
whose coefficients are taken respectively to be the \(N\) rows
\[
A_1,A_2,\dots,A_N.
\]
Let \((J)\) denote a complete system of invariants of
\[
F_1,\dots,F_n,
\]
that is, a finite set of invariants such that every invariant of
\[
F_1,\dots,F_n
\]
can be expressed integrally and rationally through the invariants \((J)\). Hilbert's conjecture to be proved is then the following: for the simultaneous invariants \(S\) of
\[
F_1,\dots,F_N
\]
a complete system is given by the invariants \((J,PJ)\).

Applying the reduction theorem \((1)\) to a simultaneous invariant \(S\)---which can itself be regarded as a form in the \(N\) rows \(A_i\)---one sees at once that every simultaneous invariant is representable as a sum of polars of invariants that involve only the \(n\) rows
\[
A_1,\dots,A_n.
\]
But such forms are invariants of
\[
F_1,\dots,F_n,
\]
and hence can be expressed integrally and rationally through the invariants \((J)\). On the other hand, by the very definition of the polar processes, every form obtained from an invariant of
\[
F_1,\dots,F_n
\]
by a polar process is again a simultaneous invariant of
\[
F_1,\dots,F_N.
\]
Since the polar process preserves integral rational expressibility, it follows that every simultaneous invariant \(S\) of
\[
F_1,\dots,F_N
\]
can be represented integrally and rationally through the invariants \((J,PJ)\).

Thus we have proved:

\medskip
\noindent\textbf{Theorem.}
Every simultaneous invariant of a system of arbitrarily many basic forms is integrally and rationally expressible through the finite family \((J,PJ)\), where \(J\) is a complete invariant system for \(n\) linearly independent basic forms and \(PJ\) denotes all invariants obtained from \(J\) by polar operations.

\medskip
This establishes Hilbert's conjecture.

\bigskip
\noindent In the next batch I will continue with Noether's more conceptual proof of the reduction theorem via equivalence of linear families of forms.

\clearpage
\section*{On integral rational representation of the invariants of a system of arbitrarily many basic forms}

\section*{II. Linear families and the proof of the reduction theorem}

For the proof of the reduction theorem, we first insert some preliminaries about linear families of forms. By a \emph{linear family of forms} over \(K\) we mean a system of forms of the same dimension so complete that, together with \(\theta_1\) and \(\theta_2\), every combination
\[
c_1\theta_1+c_2\theta_2
\]
again belongs to the system, where \(c_1,c_2\) lie in a prescribed field \(K\), and \(K\) is also the coefficient field of the \(\theta\). Among these forms there is always a finite number --- say \(\varrho\) --- that are linearly independent over \(K\), whereas every \(\varrho+1\) forms satisfy a linear relation with coefficients in \(K\); the family has rank \(\varrho\) with respect to \(K\).

We shall use the following evident fact: if \(L\) and \(T\) are two linear families over \(K\) of the same rank \(\varrho\), and if \(T\) is a subfamily of \(L\), then \(L\) and \(T\) are identical (equivalent).

With this understood, we return to the reduction theorem. The identity proved in Section I,
\[
\theta=\sum P\cdot Z,
\tag{1}
\]
remains valid in analogous form if the same polar process \(P\) is applied to both sides. Thus the reduction theorem simultaneously gives a representation of \emph{all} polars of \(\theta\) by sums of the special polars \(PZ\). On the other hand, these special polars certainly belong to the family generated by \(\theta\). The problem is therefore to prove the equivalence of suitable linear families.

For simplicity let \(\theta(x,y,z)\) be of dimensions \(\alpha,\beta,\gamma\) in the three binary rows \(x_1,x_2\), \(y_1,y_2\), \(z_1,z_2\); the coefficients of \(\theta\) are at first regarded as indeterminates \(a\) (or, if one prefers, rational numbers). We then consider the following three linear families.

\medskip
\noindent\textbf{1. The family \(L\).}
This consists of all forms \(f(x,y,z)\) obtained from \(\theta\) by polar processes and having, in each of the rows \(x,y,z\), the same dimensions as \(\theta\). In particular, \(L\) contains \(\theta\) itself. If one views the \(f(x,y,z)\) as forms in \(x,y,z\) and the indeterminates \(a\), then the coefficient field is the field \(R\) of rational numbers; \(K\) must likewise be taken to be \(R\), since only for rational \(c_1,c_2\) does \(c_1P_1+c_2P_2\) remain a polar process. Let the rank of \(L\) over \(R\) be \(\varrho\).

\medskip
\noindent\textbf{2. The family \(G\).}
This consists of all forms \(g(\Lambda;x,y)\) defined by
\[
g(\Lambda;x,y)=f(x,y,\Lambda_1x+\Lambda_2y),
\]
or, written through polar processes,
\[
g(\Lambda;x,y)=\sum_{i=0}^{\gamma}\binom{\gamma}{i}\Lambda_1^{\gamma-i}\Lambda_2^i\, f(x,y,z)^{(i)},
\]
where
\[
f(x,y,z)^{(i)}=
\frac{1}{i!(\gamma-i)!}
\left(y\frac{\partial}{\partial z}\right)^{\gamma-i}
\left(x\frac{\partial}{\partial z'}\right)^i
f(x,y,z),
\]
and \(f(x,y,z)^{(i)}\) arises successively from \(f(x,y,z)\) by the process \((yz')\frac{\partial^2}{\partial z\,\partial z'}\).
The coefficients \(g_i(x,y)\) occurring in \(g\) are therefore forms \(Z\) of the identity \((1)\). Let \(G\) have rank \(\delta\) over \(R\).

\medskip
\noindent\textbf{3. The family \(T\).}
This is obtained from \(G\) by the inverse polar construction by which \(G\) arose from \(L\). Its forms \(h(x,y,z)\) are defined by
\[
h(x,y,z)=\sum_{i=0}^{\gamma}\binom{\gamma}{i}z_1^{\gamma-i}z_2^i\, g_i(x,y),
\]
with
\[
g_i(x,y)=\binom{\gamma}{i}f(x,y,z)^{(i)}.
\]
The individual \(h_i\), and therefore also \(h\), are thus forms of the type \(PZ\) and sums of such forms. Let the rank of \(T\) be \(r\).

As the construction shows, every \(h(x,y,z)\) of \(T\) is obtained from \(\theta\) by polar processes and has, in the rows \(x,y,z\), the same dimensions as \(\theta\). Hence \(T\) is a subfamily of \(L\). To conclude equivalence it remains only to prove that the ranks are equal.

\medskip
\noindent\textbf{Lemma a).}
From
\[
f(x,y,\Lambda_1x+\Lambda_2y)=0
\]
(identically in \(\Lambda\)) it follows necessarily that
\[
f(x,y,z)=0.
\]

This follows immediately from the identically valid relation among three binary rows
\[
(xy)^2 z+(yz)^2 x+(zx)^2 y=0,
\]
which yields
\[
f\bigl(x,y,(xy)^2z+(yz)^2x+(zx)^2y\bigr)=0.
\]
Since \((xy)^2=(x_1y_2-x_2y_1)^2\not\equiv 0\), the substitution \(\Lambda_1=(yz)^2\), \(\Lambda_2=(zx)^2\) in the identity \(f(x,y,\Lambda_1x+\Lambda_2y)=0\) gives
\[
f(x,y,z)\,(xy)^2=0,
\]
hence \(f(x,y,z)=0\).

Thus there is a one-to-one correspondence between the forms of \(G\) and those of \(L\); moreover every linear relation in \(G\) induces, and is induced by, a linear relation in \(L\). Therefore
\[
\varrho=\delta.
\]

\medskip
\noindent\textbf{Lemma b).}
From
\[
h(x,y,z)=0
\]
it follows necessarily that
\[
g(\Lambda;x,y)=0.
\]

Indeed, \(h(x,y,z)=0\) means that each of the terms
\[
h_i(x,y)\,z_1^{\gamma-i}z_2^i
\]
vanishes, and therefore so does every linear combination of these terms. In particular,
\[
g(\Lambda;\xi,\eta)=\sum_{i=0}^{\gamma}\binom{\gamma}{i}\Lambda_1^{\gamma-i}\Lambda_2^i h_i(\xi,\eta)=0
\]
for all \(\xi,\eta\), and consequently \(g(\Lambda;x,y)=0\). This proves
\[
\delta=r.
\]

Since \(T\subseteq L\) and \(\varrho=\delta=r\), the two families are equivalent. Therefore every form in \(L\), and in particular \(\theta\) itself, can be represented as a linear combination of linearly independent forms \(h^{(\nu)}(x,y,z)\):
\[
\theta=\sum c_\nu h^{(\nu)}(x,y,z)=\sum P\cdot Z.
\]
Because this identity was derived for indeterminate coefficients \(a\), it remains valid after replacing those indeterminates by elements of an arbitrary rationality field. Hence the reduction theorem is proved in the case of three binary rows.

\medskip
The general proof, for \(N\) rows each of \(n\) variables, proceeds in completely parallel fashion. One again considers three linear families \(L\), \(G\), and \(T\): the family \(L\) of all forms obtained from \(\theta(A)\) by polar processes; the family \(G\) obtained from the forms of \(L\) by substituting
\[
A_{n1}=\Lambda_1A_{11}+\cdots+\Lambda_nA_{1n},\quad \dots,\quad
A_{nn}=\Lambda_1A_{n1}+\cdots+\Lambda_nA_{n,n-1};
\]
and the family \(T\) consisting of sums of forms \(PZ\). The same two lemmas remain valid --- in particular because among any \(n+1\) rows there is always an identically valid linear relation --- and therefore \(L\) and \(T\) have the same rank. This yields once again the identity
\[
\theta=\sum P\cdot Z,
\]
first for indeterminate coefficients, hence also for coefficients lying in an arbitrary rationality field. Thus the reduction theorem is established in full generality.

\section*{III. General series expansions}

We now indicate briefly how analogous considerations also produce the most general series expansions (for \(N<n\)). Let \(Z\) be a separately homogeneous form in the \(n\) rows \(A_1,\dots,A_n\), each consisting of \(n\) variables, and let \(\Delta\) denote the determinant of these rows. Then we first prove the identity corresponding to \((1)\):
\[
Z=\sum P\cdot H \pmod{\Delta},
\tag{2}
\]
where the forms \(H\) contain only the rows \(A_1,\dots,A_{n-1}\) and are derived from \(Z\) by processes \(P\).

The proof consists in converting the relations derived in II into congruences modulo \(\Delta\). For simplicity let us again consider the case of three ternary rows \(x,y,z\). The three linear families \(L,G,T\) are defined as before; write \(\varrho,\delta,r\) for their ordinary ranks, and \(\varrho^*,r^*\) for the ranks of \(L\) and \(T\) modulo \(\Delta\).

One obtains, in place of Lemma a), the following.

\medskip
\noindent\textbf{Lemma a').}
From
\[
f(x,y,\Lambda_1x+\Lambda_2y)=0
\]
it follows necessarily that
\[
f(x,y,z)=0\pmod{\Delta}.
\]

Indeed, for any four ternary rows there is an identically valid relation
\[
\Delta_x x-\Delta_y y+\Delta_z z=\Delta_t t,
\]
and therefore, modulo \(\Delta\), always a linear relation among three rows,
\[
\Lambda_x x-\Lambda_y y+\Lambda_z z\equiv 0 \pmod{\Delta}.
\]
From this one deduces, exactly as before, that
\[
f(x,y,\Lambda_1x+\Lambda_2y)\equiv 0 \pmod{\Delta}
\quad\Rightarrow\quad
f(x,y,z)\equiv 0 \pmod{\Delta}.
\]
Hence \(\varrho^*=\delta\).

Lemma b) from Section II remains valid without change, so \(\delta=r^*\). To show in addition that \(r=r^*\), Noether introduces another lemma.

\medskip
\noindent\textbf{Lemma c).}
From
\[
h(x,y,z)\equiv 0 \pmod{\Delta}
\]
it follows necessarily that
\[
h(x,y,z)=0.
\]

The reason is that \(h(x,y,z)\), being a polar form, vanishes under the familiar \(\Omega\)-process. Thus one has a differential equation of the form
\[
\Omega h=\sum_i \frac{\partial^2 h}{\partial x_i\,\partial y_i}=0.
\]
If now
\[
h(x,y,z)=\Delta(x,y,z)\,k(x,y,z),
\]
then after differentiating one obtains
\[
\Omega h=(xy)\,k+\cdots,
\]
and from this it follows that \(h\) itself must vanish. Consequently
\[
\varrho^*=\delta=r=r^*,
\]
and because \(T\) is a subfamily of \(L\), identity \((2)\) is proved.

Applying identity \((2)\) to the multiplier of \(\Delta\) then leads to an expansion
\[
Z=\varphi+\Delta\varphi_1+\Delta^2\varphi_2+\cdots+\Delta^r\varphi_r,
\]
where each \(\varphi_i\) vanishes under the \(\Omega\)-process. Repeated application of \(\Omega\) and the elementary identity
\[
\Omega(\Delta\psi)=\psi+\Delta\,\Omega(\psi)
\]
then yields
\[
\varphi_i=\sum P\cdot H_i,
\]
and therefore the general expansion
\[
Z=
\sum P\cdot H
+\Delta\sum P\cdot H_1
+\Delta^2\sum P\cdot H_2
+\cdots
+\Delta^r\sum P\cdot H_r.
\tag{3}
\]

This is the most general series expansion of a form in \(n\) rows of \(n\) variables in terms of polars of forms involving only \(n-1\) rows. The corresponding expansion for forms in \(N\) rows of \(n\) variables (\(N<n\)) is then obtained, as Noether remarks, by a simple substitution argument.

She concludes that, together with the reductions and their combinations, these expansions exhaust the known expansion formulas for forms in arbitrarily many rows of \(n\) cogredient variables.

\medskip
\noindent
Erlangen, 5 January 1915.

\clearpage
\section*{An algebraic criterion for absolute irreducibility}

A polynomial with coefficients in an arbitrarily given field is called \emph{absolutely irreducible} if it remains irreducible in an algebraically closed field extending the coefficient field. I shall show below that, in contrast with irreducibility relative to a prescribed field, absolute irreducibility can be formulated algebraically. More precisely, the following theorem holds.

\medskip
\noindent\textbf{Theorem.}
To every polynomial in $n>2$ variables with indeterminate coefficients one can attach a whole rational function of those coefficients and of further indeterminates --- the \emph{reducibility form} --- in such a way that for every specialized polynomial of the same degree the necessary and sufficient condition for absolute irreducibility is the nonvanishing of the reducibility form. In other words: for every specialized system of coefficients for which the reducibility form is not identically zero in the auxiliary indeterminates, the polynomial is absolutely irreducible. Equivalently, the polynomial remains irreducible modulo every prime ideal of every finite algebraic field containing the coefficient field, with at most finitely many exceptional prime ideals. In particular, the containing field may be the rational numbers.

Further applications to the theory of relatively integral functions will be given elsewhere.

\section*{\S 1. The universal product relation}

We first show that every polynomial in more than two variables with indeterminate coefficients is absolutely irreducible. Write
\begin{align}
F(x) &= \sum A_{i_1\dots i_n}x_1^{i_1}\cdots x_n^{i_n}, \qquad i_1+\cdots+i_n=l, \tag{0.1}\\
G(x) &= \sum B_{j_1\dots j_n}x_1^{j_1}\cdots x_n^{j_n}, \qquad j_1+\cdots+j_n=m, \tag{0.2}\\
H(x) &= F(x)G(x)=\sum C_{k_1\dots k_n}(A,B)x_1^{k_1}\cdots x_n^{k_n},
\qquad k_1+\cdots+k_n\le l+m. \tag{0.3}
\end{align}
Here the $A$ and $B$ are indeterminates, and the coefficients $C(A,B)$ are integral bilinear expressions in the $A$'s and $B$'s.

Consequently the quotients
\[
D = \frac{C(A,B)}{C_{0\dots 0}(A,B)}
\]
are rational functions of the quotients $A/A_{0\dots 0}$ and $B/B_{0\dots 0}$. But the number of such functions is larger than the number of their arguments: namely, the three respective counts are
\[
\binom{l+n}{n}-1,\qquad \binom{m+n}{n}-1,
\qquad \binom{l+m+n}{n}-1,
\]
and one has the inequality
\[
\binom{l+n}{n}+\binom{m+n}{n}
\ge \frac{(l+m+n)!}{n!(l+m)!}
= \binom{l+m+n}{n}.
\]
Therefore there is at least one algebraic relation among the $D(A,B)$, and hence also among the coefficients $C(A,B)$:
\[
\Phi(Z)\not\equiv 0 \quad \text{(identically in $Z$)},
\qquad \Phi(C(A,B))=0 \quad \text{(identically in $A,B$)}.
\tag{0.4}
\]
Here $\Phi$ denotes an integral polynomial.

Hence a polynomial with indeterminate coefficients,
\[
E(x)=\sum Z_{k_1\dots k_n}x_1^{k_1}\cdots x_n^{k_n},
\qquad k_1+\cdots+k_n\le l+m,
\tag{0.5}
\]
is necessarily absolutely irreducible for $n>2$.

\section*{\S 2. The prime ideal of reducibility relations}

The totality of all such polynomials $\Phi(Z)$ that vanish identically after the substitution $Z=C(A,B)$ forms an ideal in the polynomial ring; more precisely it forms a \emph{prime ideal} $\mathfrak B$. Indeed, if a product $\Phi_1\Phi_2$ vanishes identically under the substitution, then at least one factor vanishes identically.

Since relation \((0.4)\) holds identically in $A,B$, it remains valid for every specialized choice of $A,B$, where $A,B$, and therefore also $\Gamma=C(A,B)$, may be elements of an arbitrary field. Thus the coefficients $\Gamma$ of every polynomial that becomes reducible in an extension field satisfy the relations
\[
\Phi(\Gamma)=0;
\]
in other words, they are zeros of the prime ideal $\mathfrak B$, equivalently of finitely many basis polynomials for $\mathfrak B$. What remains is to prove the converse.

For this it is important to note that all relations among $A,B,C$ remain valid if in the polynomials \((0.1)\)--\((0.3)\) we homogenize by adding a new variable $y$ and pass to homogeneous forms $F(x,y)$, $G(x,y)$, $H(x,y)$ of dimensions $l$, $m$, $l+m$. Thus coefficients $\Gamma$ may occur among the zeros of $\mathfrak B$ for which, say, $F(x,y)$ degenerates to a power of $y$; in passing back to inhomogeneous polynomials this corresponds to a lowering of degree for $H(x)$ rather than to an actual factorization. It will be shown that, apart from this degree-lowering phenomenon, the converse always holds. In particular, for homogeneous forms the converse holds without exception, provided not all coefficients of $F$ vanish; that is, every zero of $\mathfrak B$ is produced by a parameter representation
\[
\Gamma=C(A,B)
\]
with coefficients $A,B$ in a suitable extension field.

Noether proves this by constructing finitely many special functions $\Phi(Z)$ directly. She associates to the polynomials \((0.1)\)--\((0.3)\) corresponding one-variable polynomials by means of Kronecker substitution, shows that gaps occur among the resulting exponents, and from the norms of the corresponding coefficients obtains the desired functions $\Phi(Z)$ and therefore the criterion.

\section*{\S 3. Kronecker substitution}

Set
\[
d=l+m+1.
\tag{0.6}
\]
Then the Kronecker substitution associates to the polynomials $F,G,H$ the one-variable polynomials
\begin{align}
f(\xi) &= \sum A_{i_1\dots i_n}\,\xi^{i_1+i_2d+i_3d^2+\cdots+i_nd^{n-1}}, \tag{0.7}\\
g(\xi) &= \sum B_{j_1\dots j_n}\,\xi^{j_1+j_2d+j_3d^2+\cdots+j_nd^{n-1}}, \tag{0.8}\\
h(\xi) &= \sum C_{k_1\dots k_n}(A,B)\,\xi^{k_1+k_2d+\cdots+k_nd^{n-1}},
\qquad k_1+\cdots+k_n\le l+m. \tag{0.9}
\end{align}
This correspondence is one-to-one, because from an induced exponent one can recover the exponents in the original monomial: the base-$d$ expansion is unique, and since $d=l+m+1$ no carrying occurs among the exponents that arise from multiplication.

Thus one has simply
\[
h(\xi)=f(\xi)g(\xi),
\tag{0.10}
\]
and conversely, whenever a relation of the form \((0.10)\) holds in such a way that no exponents other than the induced ones occur in $f$ and $g$, one may recover the original multivariable factorization
\[
H(x)=F(x)G(x).
\tag{0.11}
\]

Moreover, genuine gaps do appear among the induced exponents in $f(\xi)$ and $g(\xi)$. The highest exponent in $f(\xi)$ is $l d^{n-1}$, while the second highest is $(l-1)d^{n-1}+d^{n-2}$. Their difference is
\[
d^{n-2}(d-1)>1,
\]
because $d=l+m+1>3$ and $n>2$. The same holds for $g(\xi)$.

The relations \((0.10)\) and \((0.11)\), first established for the indeterminates $A,B$, remain valid after every specialization of $A,B$. To preserve the degree one homogenizes \((0.10)\) and \((0.11)\) by means of a new variable $\eta$ (and likewise $y$ in the multivariable side). Thus a homogeneous form $H(x,y)$ factors in an algebraic extension field if and only if the corresponding binary form $h(\xi,\eta)$ splits into two factors in such a way that only the induced exponents appear in the factors.

This is the point at which Noether passes from the mere existence of algebraic relations to the construction of the reducibility form via norms. The later sections produce the explicit criterion from these one-variable relations.

\clearpage
\section{An algebraic criterion for absolute irreducibility: completion}

\subsection*{4. Norms of the auxiliary linear form}

Let
\[
  p = l d^{n-1}, \qquad q = m d^{n-1},
\]
so that these are the degrees of the one-variable polynomials $f(\xi)$ and $g(\xi)$ obtained above by Kronecker substitution. Write
\[
  \prod (\xi-a_i\eta)=\xi^p+r_1\xi^{p-1}\eta+\cdots+r_p\eta^p,
\]
\[
  \prod (\xi-b_j\eta)=\xi^q+s_1\xi^{q-1}\eta+\cdots+s_q\eta^q,
\]
and
\[
  \prod (\xi-c_k\eta)=\xi^{p+q}+t_1\xi^{p+q-1}\eta+\cdots+t_{p+q}\eta^{p+q}.
\]
Thus $r,s,t$ denote the homogenized elementary symmetric functions in the rows $a$ and $b$; equivalently, they are the symmetric functions, linear in each separate row, from which all integral symmetric functions that are homogeneous and of equal degree in each row are composed uniquely and integrally.

Introduce further indeterminates $u_{\mu\nu}$ and form the linear expression
\begin{equation}
  \ell(u,a,b)=\sum u_{\mu\nu} a^\mu b^\nu,                         \tag{9}
\end{equation}
where the summation is extended over prescribed index pairs $(\mu,\nu)$. Multiply $\ell(u,a,b)$ by all its conjugates obtained from permutations of the rows $a,b$, omitting repetitions. This product is the norm $N(\ell(u))$ of $\ell(u)$ with respect to the field of the $t(a,b)$, when $\ell(u)$ is regarded as a function in the field of the $r(a,b)$ and $s(a,b)$. It is an integral symmetric function of all rows $a,b$, homogeneous of the same degree in each row. Hence, by the preceding discussion, it is a uniquely determined integral polynomial in the $t(a,b)$ and the $u_{\mu\nu}$:
\begin{equation}
  N(\ell(u,a,b)) = T(t(a,b),u),                                      \tag{10}
\end{equation}
identically in $a,b,u$.

The same argument applied to $z-\ell(u)$ gives
\[
 N(z-\ell(u))=z^\delta+T_{\delta-1}(t,u)z^{\delta-1}+\cdots+T_0(t,u),
\]
where all $T_i$ are integral polynomials in $t$ and $u$. Both sides vanish for $z=\ell(u)$. Because the elementary symmetric functions $r(a,b)$ and $s(a,b)$ are algebraically independent, the identity remains an identity in $r,s,u$ once $t(a,b)$ is replaced by the corresponding integral bilinear expressions $w(r,s)$ and $\ell(u)$ is interpreted as a function of $r$ and $s$. Thus one obtains
\begin{equation}
  \ell(u)^\delta+T_{\delta-1}(w(r,s),u)\ell(u)^{\delta-1}+\cdots+T_0(w(r,s),u)=0,        \tag{11}
\end{equation}
identically in $r,s,u$.

The identities (10) and (11) persist under every specialization of the systems $a,b$, and hence of $r,s$. Suppose in particular that the specialized systems are chosen so that all products $a^\mu b^\nu$ that occur in (9) vanish. Then $\ell(u)$ vanishes under the substitution. If $w(\rho,\sigma)=\tau$, equation (11) gives
\begin{equation}
  T_0(\tau,u)=0                                                    \tag{12}
\end{equation}
identically in $u$.

Conversely, let $\tau_1,\ldots,\tau_{p+q}$ be a system, not all zero, for which (12) holds. In an algebraic extension field there is a factorization
\[
  \xi^{p+q}+\tau_1\xi^{p+q-1}\eta+\cdots+\tau_{p+q}\eta^{p+q}
  =\prod_i(\xi-\alpha_i\eta)\prod_j(\xi-\beta_j\eta),
\]
where in no factor do both $\alpha$ and $\beta$ vanish simultaneously, although some individual $\alpha$ or $\beta$ may vanish. Substitution in (10), using (12), shows that $\ell(u,\alpha,\beta)$ or one of its conjugate factors vanishes identically in $u$. After a suitable numbering of the $\alpha$'s and $\beta$'s, putting $r(\alpha,\beta)=\rho$ and $s(\alpha,\beta)=\sigma$ means exactly that
\begin{equation}
  \tau(\xi,\eta)=\rho(\xi,\eta)\sigma(\xi,\eta),                    \tag{13}
\end{equation}
with all products that occur in $\ell(u)$ vanishing, but not all coefficients of $\rho$ or of $\sigma$ vanishing.

\subsection*{5. Construction of the reducibility equations}

Now take $p$ and $q$ to be the degrees of $f(\xi)$ and $g(\xi)$ in the Kronecker substitution. Split the indices $\mu,\nu$ into two classes. In the first class, $\mu$ runs through the exponents missing from $f(\xi)$ and $\nu$ through the exponents occurring in $g(\xi,\eta)$; in the second class, $\nu$ runs through the exponents missing from $g(\xi)$ and $\mu$ through the exponents occurring in $f(\xi,\eta)$.

Let
\[
  e(\xi)=\sum Z_k\xi^k
\]
be the one-variable polynomial corresponding to the original polynomial $E(x)$ under Kronecker substitution. Define $S(Z,u)$ to be the integral polynomial in $Z$ and $u$ obtained from the right hand side $T(t,u)$ of (10), with summation restricted to the index classes just described, by replacing the $t$ by the coefficients $Z$ of $e(\xi)$ and by zeros in the positions of exponents absent from $e$.

If now one substitutes the coefficients $A,B$ of $F,G$ and the corresponding zeros of $f$ and $g$ for $r,s$, then the linear form $\ell(u)$ vanishes by the chosen summation. Moreover $t=w(r,s)$ becomes the coefficient system $C(A,B)$, together with the corresponding zeros of $h(\xi)$. Hence $T(t,u)$ becomes $S(C(A,B),u)$. From (11) one obtains
\begin{equation}
  S(C(A,B),u)=0                                                   \tag{14}
\end{equation}
identically in $A,B,u$.

Since (14) remains true under specialization, the coefficient systems
\[
  \Gamma=C(A,B)
\]
of all specialized forms $H(x)$, and more generally $H(x,y)$, that split into two factors in an extension field, satisfy
\begin{equation}
  S(\Gamma,u)=0                                                   \tag{15}
\end{equation}
identically in $u$.

Conversely, suppose that the coefficients $\Gamma$, not all zero, of a polynomial $H(x)$ or a form $H(x,y)$ satisfy (15). This is equivalent, by the preceding construction, to $T(\Gamma,u)=0$ for the coefficients of the one-variable form $h(\xi,\eta)$. By (13) there is, in an algebraic extension field of the coefficient field, a factorization
\[
  h(\xi,\eta)=f(\xi,\eta)g(\xi,\eta),
\]
where, by the restricted choice of indices, no exponents occur in $f$ and $g$ except those induced by Kronecker substitution. Therefore this one-variable factorization comes from a factorization of the original homogeneous form,
\[
  H(x,y)=F(x,y)G(x,y).
\]
If no factor becomes of degree zero when $y=1$, in particular if the coefficients $\Gamma$ do not lower the degree of $H(x)$, one obtains also
\[
  H(x)=F(x)G(x).
\]
Thus (15) is necessary and sufficient for $H(x,y)$, and, after exclusion of degree lowering, also $H(x)$, to split in an extension field.

It follows further that $S(Z,u)$ cannot vanish identically in $Z$ and $u$. Section 1 showed that for $n>2$ the polynomial with indeterminate coefficients, and therefore the corresponding homogeneous form in one further variable, is absolutely irreducible. Consequently
\[
  S(Z,u)=\sum_i \Phi_i(Z)U_i,
\]
where the $U_i$ are monomials in the $u$'s and, by (14), each $\Phi_i(Z)$ belongs to the prime ideal $\mathfrak B$ of all reducibility relations. We have therefore shown that $\mathfrak B$ has no zeros other than those represented by the parameterization $\Gamma=C(A,B)$ with $A,B$ in an algebraic extension field. The converse asserted at the end of section 2 is proved.

\subsection*{6. The reducibility form}

The condition for absolute irreducibility can now be formulated directly. Decompose the degree of $E(x)$ in every possible way as a sum of two positive integers $l+m$. For each such decomposition form the corresponding $S(Z,u)$. The product over all decompositions,
\begin{equation}
  R(Z,u)=\prod S(Z,u),                                            \tag{16}
\end{equation}
is the reducibility form of $E(x)$. Indeed, by the preceding result every splitting of $E(x)$, respectively of $E(x,y)$, makes one factor of (16) vanish, and conversely. This proves the main theorem stated in the introduction, and at the same time shows how the reducibility form may actually be constructed in finitely many steps.

\subsection*{7. Ostrowski's theorem}

The existence of the reducibility form (16) immediately gives the theorem of Ostrowski cited in the introduction. Let
\[
  E(x)=\sum \Gamma_i x^i
\]
be an absolutely irreducible polynomial whose coefficients $\Gamma_i$ are algebraic integers. Let $K$ be a finite algebraic number field containing all the $\Gamma_i$, and let $\mathfrak p$ be a prime ideal of $K$.

For the exact degree of $E(x)$ the reducibility form $R(Z,u)$ does not vanish after the substitution $Z=\Gamma$; say
\[
  R(\Gamma,u)=\sum_j P_j U_j\ne 0
\]
identically in $u$. The ideal $(P_1,P_2,\ldots)$ is divisible only by finitely many prime ideals of $K$. Hence, except for these finitely many $\mathfrak p$, the reduction of $R(\Gamma,u)$ modulo $\mathfrak p$ is not identically zero. Since the residue classes modulo $\mathfrak p$ again form a field, it follows that $E(x)$ modulo $\mathfrak p$ is irreducible, indeed absolutely irreducible. The same conclusion holds for the homogenized form $E(x,y)$, so that no lowering of degree occurs modulo $\mathfrak p$.

If the coefficients of $E(x)$ are algebraic fractions rather than algebraic integers, one must also exclude the finitely many prime ideals that divide a common denominator. For every remaining prime ideal one may replace $E(x)$ by a scalar multiple with algebraic integral coefficients, and the preceding argument applies. This proves the theorem.

\newpage

\section{Formal calculus of variations and differential invariants}

The generation of all differential invariants and the derivation of the principal theorems are most uniformly carried out by the formal methods of the calculus of variations. This principle goes back to Riemann and received its principal development, in form-theoretic terms, through Lipschitz.

Begin with the variational problem belonging to a homogeneous form $f(dx)$,
\begin{equation}
  \delta \int f(dx)=0.                                            \tag{140}
\end{equation}
By partial integration one is led to the invariant system of equations
\[
  \psi_i=0,
\]
where the Lagrange expressions $\psi_i$ are contragredient to the differentials $dx_i$. Formally this is seen most simply by defining the $\psi_i$ through the integration-free identity corresponding to partial integration, the Lagrange central equation,
\[
  df-dh=-\sum_i D\psi_i\,dx_i.
\]
Here both $f$ and $df$ are invariants, and consequently the right hand side is invariant as well.

For a quadratic form one obtains in particular
\[
  \sum g_{ij}dx_i dx_j -2\sum g_{ij}dx_iDdx_j=-2\sum_i \psi_i(d,d)Dx_i.
\]
If $d$ is replaced by the cogredient expression $d+\lambda\delta$ and powers of $\lambda$ are compared, one obtains the corresponding system of identities
\begin{align*}
  \sum Dg_{ij}dx_i dx_j -2\sum g_{ij}dx_iDdx_j &=-2\sum_i\psi_i(d,d)Dx_i,\tag{141a}\\
  \sum Dg_{ij}dx_i\delta x_j-\sum g_{ij}\delta x_iDdx_j-\sum g_{ij}dx_iD\delta x_j &=-2\sum_i\psi_i(d,\delta)Dx_i,\tag{141b}\\
  \sum Dg_{ij}\delta x_i\delta x_j-2\sum g_{ij}\delta x_iD\delta x_j &=-2\sum_i\psi_i(\delta,\delta)Dx_i.\tag{141c}
\end{align*}
It follows that $\psi_i(d,\delta)$, and in particular $\psi_i(d,d)$ and $\psi_i(\delta,\delta)$, are contragredient vectors. From (141) one computes
\[
  \psi_i(d,\delta)=\sum_j g_{ij}D\delta x_j+\sum_{j,k}\Gamma_{ijk}\,dx_j\delta x_k,
\]
and from this the cogredient vector
\begin{equation}
  p_i(d,\delta)=\sum_j g^{ij}\psi_j(d,\delta).                   \tag{142}
\end{equation}
Setting $p(d,d)=0$ gives, as is evident, the equations of the geodesics.

Knowledge of the cogredient vector $p(d,\delta)$ leads directly to covariant differentiation. If, for example, $h(d,\delta)$ is a form in the two rows $dx$ and $\delta x$, then
\begin{equation}
  \nabla h(dx,\delta x)=dh-\sum_i \frac{\partial h}{\partial dx_i}p_i(d,\delta)-\sum_i\frac{\partial h}{\partial \delta x_i}p_i(\delta,\delta)            \tag{143}
\end{equation}
(as written in Noether's notation) is a difference of invariants and hence itself an invariant; it is formed in such a way that second differentials cancel. Thus $\nabla h(dx,dx)$ is a form involving only first differentials: the covariant derivative of $h$. The same applies when $h$ contains more rows of differentials.

The construction of the curvature form in Riemann and Lipschitz rests on the same principle. One passes from $\delta^2 f$ to the ``normal form of the second variation''
\begin{equation}
  Q=\sum D^2g_{ij}dx_i dx_j-2\sum d\delta g_{ij}dx_i\delta x_j+\sum \delta^2 g_{ij}\delta x_i\delta x_j,             \tag{144}
\end{equation}
which, being an aggregate of invariants, is itself an invariant, and in which the third differentials have now been eliminated, generalizing the Lagrange central equation. The second differentials are eliminated in analogy with covariant differentiation by forming
\begin{equation}
  K=Q-2\sum_i\{p_i(d,\delta)\psi_i(d,\delta)-p_i(\delta,d)\psi_i(d,d)\}.       \tag{145}
\end{equation}
This may also be written in the form
\begin{equation}
  K=dD\psi(d,d)-2D\psi(d,\delta)dx-2\{p(d,\delta)\psi(d,\delta)-p(d,d)\psi(d,\delta)\}.       \tag{146}
\end{equation}
The law of construction of $K$ may be phrased as follows: in $\delta^2 f$ the second differentials are eliminated by setting $p(d,d)$ and $p(d,\delta)$, equivalently the right hand sides of (141), equal to zero. This is Riemann's definition of the curvature form. It must be explicitly noted, however -- and (146) makes this especially clear -- that this zero substitution may be made only after the differentiation has been carried out. Lipschitz's formal construction avoids this difficulty. Correspondingly, the covariant derivative (143) may be defined by eliminating the second differentials in $dh$ through setting the expressions $p(d,\delta)$ and the analogous ones equal to zero.

From the curvature form, by successive formation of covariant derivatives, one obtains a series of forms, Christoffel's series $[\Phi_0,\Phi_1,\ldots]$. After adjoining $f$, their projective invariants exhaust, according to Christoffel and Ricci, the totality of the differential invariants of the quadratic form; and their equivalence with respect to linear transformations, without further integrability conditions, expresses the equivalence of two quadratic differential forms. The inner reason for this reduction theorem lies in the existence of the Riemann normal coordinates arising from the variational problem (140): they transform the geodesics through a point into straight lines and therefore transform linearly under an arbitrary transformation of the variables $x$. The formation of all invariants and the equivalence question are thereby reduced to the corresponding question for the homogeneous components in the expansion of $f$ in normal coordinates; one shows that these components are linear combinations of $f,\Phi_0,\Phi_1,\ldots$. For quadratic forms this is carried out in detail by Vermeil, following a note of Noether. According to that note, the theorems and methods remain valid when an arbitrary differential expression is taken as the starting point. Thus the reach of the methods of the formal calculus of variations, as compared with those of pure elimination theory, becomes apparent.

Levi-Civita gave a geometric interpretation of setting $p(d,\delta)=0$ as the parallel displacement of a vector. This concept, later axiomatized by Weyl, is essentially restricted to quadratic forms, but it allows a generalization of another kind. It is preserved under enlargement of the group, namely multiplication of the $g_{ij}$ by an arbitrary function, provided that, besides the quadratic form, a linear form is also taken as given. Then $p(d,\delta)$ is uniquely determined, constructions of invariants analogous to the preceding ones can be carried out, and geodesic coordinates can be defined. The equation $p(d,d)=0$, however, no longer arises from a variational problem. Questions of reduction theorems and equivalence have not yet been attacked in full generality in this setting; only for invariants of lowest order has Weitzenböck given a reduction theorem.

Finally, one should mention the general ``invariant variational problems'', that is, those in which the simple or multiple integral $J$ is invariant under a group in Lie's sense. Special cases of physical interest were studied by Hamel, Herglotz, Lorentz, Fokker, Weyl, and Klein. The principled formulation by Noether shows that invariance of $J$ under a finite group $G_\rho$ with $\rho$ essential parameters corresponds to $\rho$ linearly independent divergences; invariance under an infinite group containing $\rho$ arbitrary functions and their derivatives up to order $\sigma$ corresponds to $\rho$ identities among the Lagrange expressions and their derivatives up to order $\sigma$. In both cases the converse holds. Since the Lagrange expressions become relative invariants of the group, this gives at the same time a process for generating invariants.

\newpage

\section{K. Hentzelt, revised by E. Noether: On polynomial ideals and resultants}

The following concerns the general problem of elimination theory: to determine the common zeros of all polynomials in an ideal of a polynomial ring, or equivalently the zeros of any finite number of polynomials forming a basis of the ideal. At the same time the multiplicities that occur will receive a conceptual interpretation. The coefficients of the polynomials may belong to any field; the zeros then belong to the algebraically closed field derived from it. Moreover the ideal may, without loss of generality, be assumed to be transformed, in the sense introduced below. The essential result is then the following.

To every polynomial ideal $\mathfrak m$ there may be attached a resultant form
\begin{equation}
  R_{\mathfrak m}=R^{(1)}(x_1,\ldots,x_n)\cdots R^{(i)}(x_i,\ldots,x_n)\cdots R^{(n)}(x_n)\equiv 0(\mathfrak m),       \tag{A}
\end{equation}
so that $R_{\mathfrak m}$ vanishes for all zeros of $\mathfrak m$ and only for them. If $\mathfrak n$ is a divisor of $\mathfrak m$ and the resultant forms $R_{\mathfrak n}$ and $R_{\mathfrak m}$ agree, then the ideals $\mathfrak n$ and $\mathfrak m$ themselves agree.

The second part of the theorem says that the resultant form gives the zeros with their characteristic multiplicity. It is analogous to the theorem that two ideals in an algebraic number field, one of which divides the other, are equal if and only if their norms agree. The inner reason for both theorems is the same.

Indeed $\mathfrak m$ may be considered as a module $M_{n-1}$ of linear forms: each polynomial is regarded as a linear form in the monomials in $x_1,\ldots,x_{n-1}$, and $x_n$ is adjoined to the coefficient domain. The resultant factor $R^{(n)}$ then becomes the norm of the corresponding fundamental module $G_{n-1}$ after $M_{n-1}$. This at once gives an interpretation of the multiplicity -- the product of degree and exponent of a factor of $R^{(n)}$ -- as the number of linearly independent residue classes. This fact was previously known only in the special case in which the resultant can be defined for indeterminate coefficients.

The derivation proceeds as follows. Sections 1 and 2 give theorems on modules of linear forms whose coefficients are either rational integers or polynomials in one indeterminate. Section 3 introduces ideals of polynomials; the isomorphism theorem of section 4 connects these ideals with the preceding module theory. Section 5 proves the second part of the main theorem. Section 7 gives the theorem about zeros, after section 6 has developed a general elimination theory. In section 6, by adjoining new indeterminates, the decomposition of the resolvent into linear forms in these indeterminates is also proved. This supplies, for the elimination method used here, an unambiguous proof of a claim made by Kronecker.

\subsection{1. Modules of linear forms}

In the first two sections the phrase ``integral quantities'' means rational integers or polynomials in one indeterminate with coefficients in an arbitrary abstract field $P$. By a linear form $a(\xi)$ one understands a linear expression in finitely many indeterminates $\xi_1,\ldots,\xi_k$ with integral quantities as coefficients.

A module $M$ of linear forms is defined by the following closure conditions: together with $a(\xi)$ and $b(\xi)$ it contains $a(\xi)-b(\xi)$; together with $a(\xi)$ it contains $c\,a(\xi)$, where $c$ is any integral quantity.

The rank of $M$ is, as usual, the maximum number of linearly independent linear forms in $M$. The $h$-th determinant divisor $d_h$ of $M$ is the greatest common divisor of all determinants of order $h$ formed from coefficient matrices of $h$ forms of $M$. If $\rho$ is the rank, then the non-zero determinant divisors are $d_1,\ldots,d_\rho$, and the elementary divisors are
\[
  e_1=d_1,
  \quad e_2=d_2/d_1,
  \ldots,
  \quad e_\rho=d_\rho/d_{\rho-1}.
\]
The module is finite: it has a module basis consisting of exactly $\rho$ linear forms $a_1(\xi),\ldots,a_\rho(\xi)$ from which every form of $M$ is represented linearly with integral coefficients. After an unimodular substitution $\xi=\Theta(\eta)$ the elementary divisor theorem gives the module equality
\begin{equation}
  M=(e_1\eta_1,e_2\eta_2,\ldots,e_\rho\eta_\rho).                  \tag{1}
\end{equation}

\textbf{Definition I.} A module $G$ is called a \emph{fundamental module} if it has no proper divisor of the same rank.

Equivalently, $G$ is a fundamental module precisely when from
\begin{equation}
  c\,g(\xi)\equiv 0(G),\qquad c\ne 0,                              \tag{2}
\end{equation}
there follows $g(\xi)\equiv 0(G)$.

The relation between an arbitrary module and a fundamental module is given by the following theorem.

\textbf{Theorem I.} Every module $M$ is divisible by one and only one fundamental module $G$ of the same rank. This $G$ is the smallest divisor of $M$ of the same rank and, in suitable coordinates, is represented by
\begin{equation}
  G=(\eta_1,\eta_2,\ldots,\eta_\rho).                              \tag{3}
\end{equation}
It will be called the fundamental module of $M$.

The assertion that $G$ is the smallest divisor of the same rank means that $G$ consists exactly of those linear forms $g(\xi)$ for which there is a relation
\begin{equation}
  b\,g(\xi)
  \equiv 0(M),\qquad b\ne 0.                                      \tag{4}
\end{equation}
From this description $G$ is plainly a module. It is fundamental: if $c g(\xi)\equiv0(G)$ with $c\ne0$, then multiplying (4) by $c$ gives $bcg(\xi)\equiv0(M)$, whence again $g(\xi)\equiv0(G)$. No second fundamental divisor of the same rank can exist, for it would have to divide $G$ and, being fundamental, could not be a proper divisor. The representation (3) is fundamental, since any proper divisor of it would have to contain an additional indeterminate and would therefore have higher rank.

A second concept needed later is Dedekind's quotient of two modules. If $M$ and $B$ are modules of linear forms, define the quotient $M/B$ to be the set of integral quantities $c$ satisfying
\[
  cB\equiv0(M).
\]
This is an ideal of integral quantities, hence a principal ideal in the present coefficient domains. Similarly, if $\mathfrak b$ is an ideal of integral quantities, define $M/\mathfrak b$ to be the set of linear forms $c(\xi)$ satisfying
\[
  c(\xi)\mathfrak b\equiv0(M).
\]
This is again a module of linear forms, and, provided $\mathfrak b\ne0$, it has the same rank as $M$.

The quotient concept gives the reciprocal relation between a module's fundamental module and its highest elementary divisor.

\textbf{Theorem II.} If $e_\rho$ is the highest elementary divisor of $M$ and $G$ its fundamental module, then
\begin{equation}
  (e_\rho)=M/G,
  \qquad
  G=M/(e_\rho).                                                    \tag{5}
\end{equation}
Indeed, since every elementary divisor divides $e_\rho$, (1) and (3) give
\begin{equation}
  e_\rho G\equiv0(M),
  \qquad (e_\rho)G\equiv0(M).                                     \tag{6}
\end{equation}
Thus $(e_\rho)$ is divisible by the quotient $M/G$. Conversely, from $bG\equiv0(M)$ one obtains $b\eta_i\equiv0(M)$ for each $i$, hence $b$ is divisible by $e_\rho$.

For later use it is essential that the latter quotient representation of $G$ can also be proved without recourse to the elementary divisors. This gives it greater generality.

\textbf{Theorem III.} If the integral quantities are polynomials in several indeterminates, then still
\begin{equation}
  G=M/(d),                                                         \tag{7}
\end{equation}
where $d$ is any non-identically vanishing determinant of order $\rho$ from $M$.

The notions of rank, fundamental module, and module quotient remain valid in this enlarged setting, although a basis representation need no longer be available. Let $a_1(\xi),\ldots,a_\rho(\xi)$ be linearly independent forms of $M$ and choose the indeterminates so that $d=|a_{ij}|\ne0$. Then $M$ contains special forms of the shape
\begin{equation}
  d\xi_i+\sum_{j>\rho} b_{ij}\xi_j,
  \qquad i=1,\ldots,\rho.                                        \tag{8}
\end{equation}
It cannot contain a non-zero form depending only on $\xi_1,\ldots,\xi_\rho$, for then a determinant of order $\rho+1$ would be non-zero. Since $M/(d)$ is a divisor of the same rank as $M$, it is divisible by $G$. Conversely, if $g\in G$, then $c g\in M$ for some $c\ne0$; multiplying by $d$ and using (8) shows that $dg\in M$, whence $g\in M/(d)$. This proves (7).

\subsection{2. Decomposition theorems for norms and modules}

We return to modules of linear forms over rational integers or over polynomials in one indeterminate over $P$.

\textbf{Definition III.} Group all linear forms congruent modulo $M$ into residue classes. The symbol $G/M$ denotes the system of residue classes modulo $M$ that are represented by elements of $G$. The norm of $G$ after $M$, written $N(G/M)$, is the determinant of the transition substitution from $G$ to $M$, that is, the determinant of a substitution expressing any linearly independent module basis of $M$ through one of the fundamental module $G$ of $M$.

From (1) and (3), or from Theorem II, one obtains
\begin{equation}
  N(G/M)=e_1e_2\cdots e_\rho,
  \qquad
  G=M/(N(G/M)).                                                     \tag{9}
\end{equation}
If the coefficients are rational integers, $N(G/M)$ is the number of residue classes of $G$ modulo $M$. If the coefficients are polynomials in one indeterminate, the number of residue classes is generally infinite; the degree of $N(G/M)$ is then the number of residue classes linearly independent over $P$.

If $B$ is a divisor of $M$ of the same rank, then $B$ contains either all representatives of a residue class or none. It is obtained by adjoining finitely many residue classes and their linear combinations to $M$. Consequently:

\textbf{Theorem IV.} If $B$ is a divisor of $M$ of the same rank, so that $B$ and $M$ have the same fundamental module $G$, and if
\[
  N(G/B)=N(G/M),
\]
then $B=M$.

\textbf{Theorem V.} Suppose
\begin{equation}
  N(G/M)=rs,\qquad (r,s)=1.                                      \tag{10}
\end{equation}
Then there is one and only one module $R$, and one and only one module $S$, both divisors of $M$ of the same rank, such that
\[
  r=N(G/R),\qquad s=N(G/S).
\]
They are
\begin{equation}
  R=M/(s),
  \qquad S=M/(r),                                                  \tag{11}
\end{equation}
and $M$ is the least common multiple of $R$ and $S$ in the module sense. Moreover, if $r_i=(r,e_i)$ and $s_i=(s,e_i)$, then the elementary divisors of $R$ and $S$ are respectively the $r_i$ and $s_i$.

The proof is the usual elementary-divisor argument. In coordinates where $M=(e_1\eta_1,\ldots,e_\rho\eta_\rho)$ and $G=(\eta_1,\ldots,\eta_\rho)$, the coprimeness of $r$ and $s$ separates the elementary divisors into relatively prime parts. Conversely, any modules satisfying the norm conditions have the same elementary divisors and the same fundamental module as the constructed modules; Theorem IV then forces equality.

\subsection{3. Ideals of polynomials}

Let $\mathfrak m$ be an ideal of polynomials in $n$ indeterminates $x_1,\ldots,x_n$ with coefficients in the field $P$. Such an ideal is, in the usual sense, a set closed under subtraction and under multiplication by arbitrary polynomials.

An ideal $\mathfrak m$ is called \emph{transformed} when it remains unchanged, in the relevant sense, under a sufficiently general triangular linear transformation of the variables. By adjoining new indeterminates to the field $P$ one may always pass to a transformed ideal. The purpose of doing so is to avoid exceptional positions in elimination: the successive resultants then have the expected degrees and the associated factors separate correctly.

The basic constructions are these. Regard the polynomials of $\mathfrak m$ as linear forms in the monomials in $x_1,\ldots,x_i$ with coefficients in the polynomial ring $P[x_{i+1},\ldots,x_n]$. This gives a module $M_i$. Its fundamental module is denoted $G_i$. The highest elementary divisor of $M_i$ is denoted $E^{(i+1)}$; in the elimination setting these elementary divisors are the individual resultant factors. Their product over the successive steps gives the resultant form $R_{\mathfrak m}$.

The highest elementary divisor is also characterized as a greatest common divisor of the maximal determinants of the coefficient matrix of the module. Thus the construction is independent of the chosen basis of the ideal. This independence is crucial: the resultant form belongs to the ideal and not to an arbitrarily selected system of generators.

For a transformed ideal, the fundamental ideals $\mathfrak g_i$ of the successive stages are again transformed. This is a special case of the following assertion: if $\mathfrak m$ is transformed, then the fundamental ideal obtained from the associated module is transformed. The proof is a direct application of the determinant criterion and of the invariance of the determinant divisors under the admissible transformations.

\subsection{4. The isomorphism theorem}

The relation between polynomial ideals and modules of linear forms is governed by an isomorphism theorem. Let $M_i$ be the module obtained from $\mathfrak m$ by treating the polynomials as linear forms in the monomials in $x_1,
\ldots,x_i$, with coefficients in the remaining variables. Then the residue class system of the module agrees with the corresponding residue class system of the ideal after adjoining the remaining variables. In symbols, the construction gives an isomorphism of the form
\[
  G_i/M_i \simeq \mathfrak g_i/\mathfrak m
\]
within the appropriate coefficient domain.

This theorem is proved by induction on $i$. For $i=1$ the assertion is immediate, because one is dealing with ordinary linear forms. In the induction step one first writes the module at stage $i$ in terms of the previous stage, then separates coefficients with respect to the new variable. The determinant divisors and the fundamental module behave functorially under this passage. Hence the quotient system of the polynomial ideal is faithfully reflected by the quotient system of the associated module.

The practical consequence is that statements about norms, elementary divisors, and residue classes of modules can be transported to polynomial ideals. Thus, when one says that a resultant factor is the norm of a fundamental module, this is simultaneously a statement about the residue classes of the corresponding polynomial ideal.

\subsection{5. Norms, elementary factors, and equality of ideals}

We now prove the second part of the main theorem. Let $\mathfrak n$ be a divisor of $\mathfrak m$ in the ideal sense, and suppose that the resultant forms of $\mathfrak n$ and $\mathfrak m$ agree. Passing successively to the associated modules, the equality of the resultant factors says precisely that the corresponding norms agree at each stage. By Theorem IV the modules at each stage are equal. By the isomorphism theorem the associated fundamental ideals and residue class systems are therefore equal. Running the construction back through the variables gives
\[
  \mathfrak n=\mathfrak m.
\]

Equivalently, the resultant form does not merely locate the zero set; it records the characteristic multiplicity with which the zero set occurs. If a factor occurs with exponent $a$ and degree $d$, then the product $ad$ is the number of linearly independent residue classes represented by the corresponding component. This is the precise sense in which the resultant form supplies the multiplicities.

\subsection{6. Successive elimination}

The elimination theory is now developed by eliminating the variables one after another. Start from a transformed ideal $\mathfrak a$ and form, at each stage, the greatest common divisor of those polynomials that remain after viewing the preceding variables as coefficients. The successive greatest common divisors are denoted, in Noether's notation, by expressions such as $D^{(i)}$ and $H^{(i)}$; they are primitive in the new indeterminates and regular in the variable to be eliminated.

If a polynomial of the ideal has a zero $\alpha=(\alpha_1,\ldots,\alpha_n)$, then the first elimination polynomial vanishes at the corresponding value of the eliminated variable. Conversely, because the ideal is transformed and because the determinant divisors are chosen in the general position secured by the transformation, every linear factor of the elimination polynomial extends in one and only one way to a complete zero system of the ideal after adjoining the remaining parameters.

In concrete terms, after introducing new indeterminates $u_1,\ldots,u_n$ and setting
\[
  z=u_1x_1+\cdots+u_nx_n,
\]
one obtains a resolvent polynomial whose factors have the form
\[
  z-(u_1\alpha_1+
\cdots+u_n\alpha_n).
\]
Noether's argument shows that this factorization exhausts the resolvent: no residual factor remains which is irreducible in $z$ and the $u$'s. If such a factor existed, it would acquire a linear factor in an algebraic extension; the corresponding zero would have to coincide with one already obtained from the ideal, contradicting the assumption that a further factor remained.

Thus the explicit decomposition of the resolvent is
\begin{equation}
  H(z,u)=\prod\bigl(z-(u_1\alpha_1+
\cdots+u_n\alpha_n)\bigr)^{e_\alpha},        \tag{36}
\end{equation}
where the exponents are determined by the corresponding residue-class multiplicities. The construction differs from Kronecker's method in that it depends only on the given ideal and not on a choice of basis of that ideal.

\subsection{7. Resultant form and elimination theory}

To connect the resultant form with elimination theory, apply the successive elimination just described to the quotient of the ideal by the fundamental ideal of stage $n-1$:
\[
  \mathfrak a=\mathfrak m/\mathfrak g_{n-1}.
\]
First one must verify that this quotient is again transformed. Since $\mathfrak m$ is transformed, and since the fundamental ideals of transformed ideals are transformed, the quotient inherits the transformation property.

The determinant construction shows that the highest elementary divisor $E^{(n)}(x)$ divides the corresponding elimination polynomial $D^{(n)}(x)$. Conversely, the divisibility properties of the quotient and the regularity conditions show that $D^{(n)}(x)$ divides the resultant factor, up to the multiplicities already accounted for by the elementary divisors. The same reasoning remains valid after the auxiliary variables $u$ are introduced, so that the factorizations of the resolvent and of the resultant agree term by term.

The conclusion is:

\textbf{Theorem XIII.} If the resultant form $R^{(n)}(x)$ is decomposed in an algebraic extension field of $P(u,x_{i+1},\ldots,x_n)$ into linear factors
\[
  x_i-\alpha_i,
\]
then each $\alpha_i$ can be completed in one and only one way to a complete zero system
\[
  \alpha_1,
\ldots,\alpha_n
\]
of $\mathfrak m/\mathfrak g_{n-1}$, and hence of $\mathfrak m$. The finitely many zeros of $\mathfrak m$ arising in this way are summarized by the explicit decomposition
\begin{equation}
  R(z,u)=\prod\bigl(z-(u_1\alpha_1+
\cdots+u_n\alpha_n)\bigr)^{e_\alpha}.           \tag{37}
\end{equation}
This decomposition remains valid, because of the regularity of the resultant in $z$, after the adjoined parameters are specialized to arbitrary values in $P(u)$ or in its algebraic closure.

This also separates the zeros of $\mathfrak m$ according to the individual factors of the resultant form. The number of indeterminates that must be adjoined to $P(u)$ is called the dimension of the corresponding algebraic set.

Finally group, in the decomposition of $R(z,u)$, those linear factors that are conjugate over $P(u,x_{i+1},\ldots,x_n)$. One obtains a decomposition of the resultant into powers of irreducible polynomials
\[
  R^{(n)}(x)=S_1(x)^{\lambda_1}\cdots S_r(x)^{\lambda_r}.
\]
To this decomposition corresponds a representation of the relevant fundamental ideal as an intersection, or least common multiple in the ideal sense, of components. Each $S_j$ is the norm of the corresponding component, and the exponent $\lambda_j$ has the residue-class meaning already described: if $\gamma_j$ is the degree of $S_j$, then the multiplicity $\lambda_j\gamma_j$ is the number of residue classes linearly independent over the appropriate coefficient field.

When the coefficient field $P$ has characteristic zero, the auxiliary indeterminates $u_i$ may be specialized to elements of $P$ so that the required regularity of the determinants is preserved. One may even choose the specialization so that the resultant remains the greatest common divisor of finitely many maximal determinants. Then the whole elimination theory is retained after specialization; the factorization is obtained simply by substituting the chosen numerical values for the $u_i$ in the general formula.

\newpage

\section{Algebraic and differential invariants}

If one is to report on the development of algebraic and differential invariants, it is natural in both areas to restrict attention to the critical period which, according to Hilbert's remark, follows the naive and formal period. This critical period is characterized, for algebraic invariants, by Hilbert himself; for differential invariants, by Riemann. In substance, it is characterized in the first case by the arithmetic methods of algebra, which developed essentially in connection with these questions and here acquired their sharpest form; in the second, by the methods of the formal calculus of variations.

I shall therefore discuss the arithmetic methods and their significance beyond the special subject, while in the second part I shall restrict myself to the subordination of differential invariants to algebraic invariants, precisely as achieved in connection with the methods of the calculus of variations.

\subsection{1. The algebraic invariant problem}

The algebraic invariants are based, as is well known, on the general linear group in $x_1,\ldots,x_n$:
\begin{equation}
  x_i=\sum_j s_{ij}x'_j, \qquad \text{or shortly } x=S(x'),          \tag{1}
\end{equation}
where the $s_{ij}$ are indeterminates and $\Delta=|s_{ij}|$ is non-zero. If $f(a,x)$, $g(b,x),\ldots$ are forms in $x$ of dimensions $\alpha,\beta,\ldots$ with indeterminate coefficient rows $a,b,\ldots$, then (1) induces transformations of $a,b,\ldots$ by substituting into the forms and comparing coefficients. The transformed coefficients $a',b',\ldots$ are linear homogeneous functions of $a,b,\ldots$, their coefficients depending on the $s_{ij}$.

An invariant is a whole rational function of the indeterminates $a,b,\ldots,x,y,\ldots$ which, under the transformations of the variables and of the coefficient rows, reproduces itself up to a factor, namely a power of the substitution determinant:
\begin{equation}
  J(a',b',\ldots,x',y',\ldots)=\Delta^\omega J(a,b,\ldots,x,y,\ldots).              \tag{3}
\end{equation}
The coefficients of $J$ may be assumed to be rational numbers, or even integers, because every invariant with coefficients in an arbitrary number field $P$ can be expressed linearly, with coefficients in $P$, through finitely many such special invariants. Likewise it is no restriction to assume $J$ separately homogeneous in each row; any invariant is a finite sum of such components.

Conversely, the requirement (3) determines the transformations (1) and the induced transformations of the coefficients. Ostrowski and Schur have recently shown that the induced transformations can also be characterized as the totality of separate linear transformations in each row which, with certain explicit exceptions, leave a given invariant not involving the variables $x,y,\ldots$ invariant.

The product of two invariants is again an invariant. A sum is an invariant if and only if the weights, the exponents $\omega$ in (3), agree. If one restricts to substitutions of determinant one, however, arbitrary sums are again invariants, and these exhaust the invariants of the restricted group. Thus one obtains an integral domain, indeed a homogeneous integral domain: if a polynomial belongs to it, then its homogeneous components belong to it individually.

Here arises the fundamental problem at which invariant theory and arithmetic algebra as a whole developed: is this integral domain finite? That is, does it possess a finite integral basis $I_1,\ldots,I_k$ such that every invariant can be represented rationally and integrally by them,
\[
  I=G(I_1,\ldots,I_k)?
\]
Hilbert's solution of the problem rests on reducing the question of an integral basis to that of an ideal basis. Gordan's proof in the binary case did not extend because of the intractable symbolic calculations, and the Mertens-Hilbert proof relied on decomposing a binary form into linear factors. I sketch Hilbert's train of thought, emphasizing the arithmetic ideas, and then mention three related questions: actual construction of the basis in finitely many steps, finiteness questions for subgroups, and finiteness questions with integrality of coefficients taken into account.

Recall the ideal definition in a finite algebraic number field. A system $\mathfrak a$ of numbers of the field is called an ideal if it belongs to the domain $\mathfrak o$ of all algebraic integers of the field; if with $\alpha$ and $\beta$ it contains $\alpha-\beta$; and if with $\alpha$ it contains $A\alpha$ for every $A\in\mathfrak o$. Every ideal has an ideal basis,
\[
  \mathfrak a=(\alpha_1,\ldots,\alpha_r),
\]
meaning that every element of $\mathfrak a$ is of the form $A_1\alpha_1+\cdots+A_r\alpha_r$.

This definition uses only the fact that $\mathfrak o$ is an integral domain. It therefore carries over word for word when $\mathfrak o$ is replaced by any integral domain, as do the notions of ideal and ideal basis.

Hilbert's finiteness proof decomposes into three statements.
\begin{enumerate}
  \item In the ring of all polynomials in finitely many indeterminates over a rationality domain, every ideal has a finite ideal basis.
  \item A homogeneous integral domain of polynomials is finite if and only if the ideal-basis theorem holds in it.
  \item For the invariant domain the ideal-basis theorem holds in the sharpened form that every ideal basis formed relative to the full polynomial ring is at the same time an ideal basis relative to the invariant domain.
\end{enumerate}
In Hilbert's proof the last assertion follows from a known theorem concerning the $\Omega$-process. E. Fischer later showed that the assertion follows already from the characteristic property of the general linear group that together with each transformation it contains the contragredient, respectively the conjugate-complex one.

\subsection{2. Function fields and actual construction}

The actual construction of the basis depends on a further arithmetic transformation of the finiteness concept, using function-field theory.

An integral domain $\mathcal J$ of polynomials is finite if and only if it contains a finite subdomain $\mathcal H$ such that $\mathcal J$ depends algebraically and integrally on $\mathcal H$. Any integral basis of $\mathcal H$ is then supplemented by a fundamental system of the algebraic integral elements to give a basis of $\mathcal J$.

In the case of homogeneous domains $\mathcal J$, such as invariant domains, for which the ideal-basis theorem holds in the sharpened form above, such a finite subdomain $\mathcal H$ can be obtained from finitely many polynomials $J_1,\ldots,J_s$ with the property that the vanishing of these polynomials forces the vanishing of all polynomials of $\mathcal J$. This is again based on a general theorem of ideal theory: to every ideal of polynomials one can attach a finite set of functions that controls its zeros and hence the integral dependence of other functions over the generated subdomain.

From this upper bound K. Hentzelt obtained an upper bound for the degrees of the invariants needed for a basis. This makes the construction finite in principle, although not necessarily practical in the older symbolic sense.

\subsection{3. Subgroups and integrality}

The second adjoining question concerns invariants of subgroups. If one replaces the full linear group by a subgroup, then the invariant domain need not automatically satisfy Hilbert's original argument. Nevertheless, for many groups the arithmetic method persists. In particular, if the invariant domain is integral over a finite polynomial subdomain, then the same finiteness conclusion follows. This applies to finite groups and to many of the classical groups appearing in invariant theory.

The third question concerns integrality. One asks not only for rational representation of invariants by a finite basis, but for integral rational representation, with integer coefficients when the original invariant has integer coefficients. This requires replacing mere field arguments by ideal-theoretic divisibility and congruence arguments. Noether's own work on integral rational representation shows that, after suitable choice of the basic invariants and of an appropriate finite subring, integral dependence and ideal bases give the required representation. Thus the arithmetic method not only proves finiteness over fields; it also explains the integrality phenomena in invariant theory.

\subsection{4. Differential invariants}

I turn now to differential invariants and to the question raised at the beginning: their reduction to algebraic invariant theory.

The underlying group is, as usual, given by transformations
\begin{equation}
  x_i=x_i(y_1,\ldots,y_n).                                       \tag{4}
\end{equation}
The differentials transform by
\[
  dx_i=\sum_j \frac{\partial x_i}{\partial y_j}dy_j,
\]
and the higher differentials transform by formulas obtained by differentiating repeatedly. If one starts with a differential expression
\[
  f(x,dx,d^2x,\ldots,d^rx),
\]
then a differential invariant is an expression that reproduces itself, possibly up to a relative factor, under the induced transformations of all the differential quantities.

The key observation is that these transformation formulas are algebraic in the derivatives of the coordinate transformation. If the derivatives of $x$ with respect to $y$ are treated as independent algebraic variables subject only to the non-vanishing of the Jacobian, then the problem of differential invariants becomes an algebraic invariant problem for a suitable algebraic group acting on the finite jet variables up to a prescribed order.

For each order $r$ one obtains a finite algebraic problem. The higher differentials that occur in differentiating the invariant are eliminated by the formal variational methods described earlier. Thus the theory of differential invariants is subordinated to algebraic invariant theory: the differential invariants up to a fixed order are algebraic invariants of a finite system of forms or tensors built from the jets.

This reduction theorem shows its full force for a homogeneous differential expression $f(x,dx)$. The sequence of forms obtained from the curvature form and its covariant derivatives supplies the algebraic data whose projective invariants exhaust the differential invariants. Conversely, equivalence of the corresponding algebraic systems expresses the equivalence of the differential forms, subject to the integrability conditions already built into their construction.

Therefore the critical development of differential invariant theory parallels that of algebraic invariant theory. The old formal computations are replaced by structural, arithmetic-algebraic methods: finite generation, integral dependence, ideals, and elimination. In this way the theory of differential invariants enters the same conceptual circle as the theory of algebraic invariants.

\newpage

\section{Elimination theory and general ideal theory}

\subsection*{Introduction}

In what follows the task is to place elimination theory, in the arithmetical form that I gave to Hentzelt's exposition and that may be called an ideal theory in a polynomial domain, inside general ideal theory. At the same time the parts of the theory concerning zeros will be founded anew. The basic notions of the two theories, together with those of field theory, are collected first, so that the later argument can refer back to them.

The ordering and the new proof are grouped around four questions. First, there is an arithmetical formulation of the notion of dimension. Secondly, there is the theory of zeros. Thirdly, there is the connection between the decomposition theorems for norms and elementary divisors and the decomposition theorems of general ideal theory. Fourthly, there is the characterization of prime ideals and primary ideals by the norm and by the elementary-divisor form.

The arithmetical form of the dimension concept is given by chains of prime ideals. It is the arithmetical equivalent of the geometric fact that surfaces contain curves and curves contain points. This arithmetical formulation is shown to be identical with the parameter definition used in elimination theory. With a slight modification it can be transferred to arbitrary rings; there, together with the corresponding transfer of parts of the other questions, it gives a precise view of the structure of ideals.

For zeros Hentzelt and Noether attack the given ideal directly, as is usual, by going back to successive elimination. This does not completely avoid the character of a verification. The method followed here is the one always used for polynomials in one indeterminate. One writes the given polynomial as a product of powers of prime functions, or primary functions, and for each single prime function constructs the field of zeros as a field isomorphic to the residue-class field. The degree of the prime function gives the degree of this field; the degree of the primary function, whose zeros are the same as those of the prime function, gives the degree of the residue-class ring modulo that primary function.

Passing for each prime function to the Galois field gives the decomposition into linear factors. Thus the problem of zeros, linear factorization, and multiplicity is settled for the original polynomial. Quite analogously, for a prime ideal the system of residue classes is enlarged by forming quotients to the residue-class field, and a zero-field is constructed that is isomorphic to it. This zero-field contains a subfield generated by adjoining certain indeterminates. The degree of the norm gives the degree relative to this subfield; and by passing to the Galois field one obtains the decomposition of the elementary-divisor form, and therefore of the norm, into linear factors. For the corresponding primary ideal the zeros are the same; the decomposition is already supplied, since the norm becomes a power of the elementary-divisor form of the prime ideal. The degree of the norm gives the degree of the residue-class ring of the primary ideal relative to the subfield derived from the indeterminates.

The question for an arbitrary ideal is thereby reduced to the connection between the decomposition theorems for norms and elementary divisors and those of general ideal theory. The fact that the zeros of an ideal are assembled from the zeros of its associated prime ideals corresponds to the decomposition of the norm into greatest primary factors. These greatest primary factors become powers of the elementary-divisor forms of the individual associated prime ideals. This gives, in general, the factorization of the norm into linear factors. The multiplicity receives its interpretation, in Hentzelt--Noether, by means of fundamental ideals and their components. The degree of a greatest primary factor of the norm is equal to the number of linearly independent residue classes, after the indicated adjunction of variables, of the fundamental ideal of level $i-1$ modulo a component of the fundamental ideal of level $i$.

The fundamental ideal of level $i-1$ will be identified with the isolated component of the decomposition that is defined by all prime ideals of dimension higher than $n-i$. A component of the fundamental ideal of level $i$ will be identified with the isolated component defined by the prime ideal corresponding to the factor of the norm and by all prime ideals of higher dimension. The degree of the corresponding factor of the norm is then the degree of that subring of residue classes of this component which consists only of classes of the fundamental ideal of level $i-1$.

The characterization of prime ideals and proper primary ideals is obtained as a direct consequence of the consideration of residue-class fields. If the original coefficient domain is a perfect field, then to prime ideals there correspond, as norm and elementary-divisor form, prime functions; and to proper primary ideals there correspond proper primary functions, and conversely. If the coefficient field is imperfect, one can state only necessary conditions or sufficient conditions. Examples show that the norm of a prime ideal may become properly primary, and that the elementary-divisor form of a proper primary ideal may be a prime function. More precisely, in characteristic $p$ the norm of a prime ideal becomes the $p^f$-th power, $f\geq 0$, of its elementary-divisor form. For proper primary ideals whose elementary-divisor form is a prime function, the norm may become an arbitrary power, as examples again show.

Finally absolute prime ideals are considered, that is, ideals which remain prime ideals after extension of the coefficient domain to an algebraically closed field. For absolute prime ideals with coefficients in a finite algebraic number field, the characterization gives a transfer of a theorem first formulated by A. Ostrowski for absolute prime functions: the property of being a prime ideal is preserved modulo every prime ideal of the number field, with at most finitely many exceptions.

The last question is analogous to ideal theory in algebraic number fields. As the elementary-divisor form of an ideal of such a field one may take the least rational integer divisible by the ideal. For a prime ideal this is always a prime number; conversely an ideal is prime as soon as its norm is prime. The proofs run in parallel. The examples show that over an imperfect field there occur analogues of prime ideals of higher degree and of ramification ideals, whereas over a perfect field there are only prime ideals of first degree and no ramification ideals.

\subsection{1. Basic notions from elimination theory, general ideal theory, and field theory}

Let $\mathfrak S$ denote the integral domain of all polynomials in variables $y_1,\ldots,y_n$ with coefficients in an abstractly given field $P$. The variables $y$ are subjected to a linear transformation with indeterminate coefficients $u_{rs}$,
\[
  y_r=\sum_s u_{rs}x_s,
\]
with inverse transformation
\[
  x_s=\sum_r t_{sr}y_r.
\]
The $u_{rs}$, or equivalently the $t_{sr}$, are adjoined to $P$, giving the coefficient field $P(u)=P(t)$. Let $\mathfrak S'$ denote the polynomial domain in $x_1,\ldots,x_n$ with coefficients in this enlarged coefficient field. The domains $\mathfrak S$ and $\mathfrak S'$ are the domains underlying elimination theory. Ideals in these domains are denoted by $\mathfrak a,\mathfrak b,\ldots$ and $\mathfrak a',\mathfrak b',\ldots$. In any ring an ideal is defined in the usual way: together with $\alpha$ and $\beta$ it contains $\alpha-\beta$, and together with $\alpha$ it contains $\lambda\alpha$ for arbitrary ring element $\lambda$. The unit ideal is denoted by $\mathfrak o$.

An ideal $\mathfrak m'$ in $\mathfrak S'$ is called transformed if it arises from an ideal $\mathfrak m$ in $\mathfrak S$ by the transformation and by adjoining the $u$'s, or equivalently the $t$'s. Thus $\mathfrak m'$ consists of the linear combinations, with coefficients in the enlarged polynomial domain, of the transformed polynomials $f(x)=f(y)$ belonging to $\mathfrak m$. If $f(y)$ runs through all polynomials of $\mathfrak m$, then $\mathfrak m'$ contains all corresponding $f(x)$, and because the coefficients are symmetric under the relevant changes of rows and columns, it also contains the forms obtained from them by the corresponding permutations of the variables. Unless the contrary is expressly stated, every ideal occurring below is to be assumed transformed.

If an ideal is not transformed, one may make it transformed by composing the preceding transformation with an additional linear substitution with new indeterminate coefficients. For an already transformed ideal this amounts only to replacing the old transformation coefficients by bilinear combinations of the old and new coefficients.

Throughout, expressions such as $f^{(i)},g^{(i)},a^{(i)},b^{(i)}$ denote polynomials in the transformed domain that are free of $x_1,\ldots,x_{i-1}$.

To each ideal $\mathfrak m$ there are associated $n$ fundamental ideals
\[
  \mathfrak g_0,\mathfrak g_1,\ldots,\mathfrak g_{n-1}.
\]
The fundamental ideal of level $i-1$, $\mathfrak g_{i-1}$, consists precisely of those polynomials $G(x)$ for which there is a non-zero polynomial $B^{(i)}$ such that
\[
  B^{(i)}G(x)\equiv 0 \pmod{\mathfrak m}.
\]
The polynomial $B^{(i)}$ may depend on $G$, but the existence of a finite ideal basis shows that one may choose a single such $B^{(i)}$ that works for all $G(x)$ under consideration. The fundamental ideals of transformed ideals are again transformed ideals. Equivalently, $\mathfrak g_i$ is the ideal obtained from $\mathfrak m$ by adjoining $x_{i+1},\ldots,x_n$ to the coefficient field and then restricting attention to polynomials integral in $x_1,\ldots,x_i$. In particular the last fundamental ideal $\mathfrak g_n$ is always the unit ideal.

Every polynomial $f(x)$ can be regarded as a linear form in the monomials of $x_1,\ldots,x_{i-1}$, with coefficients that are polynomials in $x_i,\ldots,x_n$. In this way the ideal $\mathfrak m$ becomes a module $M_{i-1}$ of linear forms over the coefficient domain of these coefficients: with two linear forms it contains their difference, and with one linear form it contains every multiple by a coefficient. The fundamental module $G_{i-1}$ consists of the linear forms $g(\xi)$ for which some non-zero coefficient $B^{(i)}$ satisfies $B^{(i)}g(\xi)\equiv0\pmod{M_{i-1}}$.

The elementary divisors of these modules are understood in the usual determinantal sense. If $M$ has rank $\rho$, the elementary divisors $E_1,\ldots,E_\rho$ are obtained from the determinant divisors; all but finitely many of them are equal to one in the situations at hand. The highest elementary divisor $E^{(i)}$ is the greatest common divisor, in the polynomial sense, of the appropriate highest minors. It may be chosen integral and primitive in the remaining variables, and regular in the leading variable.

The product $R^{(i)}$ of the non-unit elementary divisors occurring in this module representation is called the norm of $G_{i-1}$ modulo $M_{i-1}$, or the individual norm of the ideal at level $i$. Symbolically,
\[
  R^{(i)}=N(G_{i-1}\mid M_{i-1})=N(\mathfrak g_{i-1}\mid \mathfrak g_i).
\]
After adjoining the remaining variables, $R^{(i)}$ is the determinant of the transition substitution from $G_{i-1}$ to $M_{i-1}$. Its degree gives the number of residue classes of $\mathfrak g_{i-1}$ modulo $\mathfrak g_i$ that are linearly independent over the relevant coefficient field.

The product of the highest elementary divisors
\[
  E_{\mathfrak m}=E^{(1)}E^{(2)}\cdots E^{(n)}
\]
is called the elementary-divisor form of $\mathfrak m$. The product
\[
  R_{\mathfrak m}=R^{(1)}R^{(2)}\cdots R^{(n)}
\]
is called the norm, or resultant form, of $\mathfrak m$. Both the elementary-divisor form and the norm are divisible by the ideal. The norm is divisible by the elementary-divisor form, and a power of the elementary-divisor form is divisible by the norm. More generally, successive products of the individual factors give congruences modulo the successive fundamental ideals, and hence finally modulo $\mathfrak m$.

If $\mathfrak n$ divides $\mathfrak m$ and their norms $R_{\mathfrak n}$ and $R_{\mathfrak m}$ agree, then $\mathfrak n=\mathfrak m$. Indeed, for a divisor of $\mathfrak m$, equality of the successive individual norms implies equality of the corresponding fundamental ideals; and since the last fundamental ideal is the unit ideal, the induction proceeds backwards. Equivalently, if $\mathfrak n$ is a proper divisor of $\mathfrak m$, then the first individual norm of $\mathfrak n$ that differs from the corresponding one of $\mathfrak m$ is a proper divisor.

A prime function means a polynomial irreducible over $P(u)$; a primary function is a power of a prime function; a proper primary function is a power higher than the first. A decomposition of a polynomial into greatest primary factors is a decomposition in which every factor is primary and the product of any two factors is no longer primary. To the decomposition of an individual norm into greatest primary factors corresponds a representation of a fundamental ideal as a least common multiple of ideals
\[
  \mathfrak g_i=[\mathfrak t_1,\ldots,\mathfrak t_s].
\]
The factors of the individual norm determine the components $\mathfrak t_j$ of the fundamental ideals. The highest elementary divisor of such a component is the greatest common divisor of the corresponding factor and the corresponding elementary-divisor factor; it is therefore a greatest primary factor of the elementary-divisor form. These components are uniquely determined by their behavior with respect to the fundamental ideals and by the requirement that their individual norm, or their highest elementary divisor, be the prescribed greatest primary factor.

The highest dimension of an ideal is defined as follows. If $R^{(i)}$ is the first individual norm different from one, then the ideal is said to have highest dimension $n-i$. The unit ideal is assigned dimension $-1$. If only one individual norm different from one occurs, the highest dimension is simply called the dimension of the ideal. If $\mathfrak m$ contains a non-zero polynomial free of $x_1,\ldots,x_{i-1}$, then the ideal has highest dimension at most $n-i$. If $\mathfrak m$ has highest dimension $n-i$ and $\mathfrak n$ divides $\mathfrak m$, then $\mathfrak n$ has highest dimension at most $n-i$.

The zeros of an ideal are all systems of values in a suitable algebraic extension field of $P(u;x_{i+1},\ldots,x_n)$ for which every polynomial of the ideal vanishes. To each factor $R^{(i)}$ of the norm there correspond, after the indicated adjunction, as many zeros of $\mathfrak m$ as the degree of $R^{(i)}$ gives. Writing the factor in bilinear combinations of the transformation coefficients and further indeterminates gives an explicit decomposition, whose factors correspond to associated systems of zeros. Thus among the zeros of an ideal of highest dimension $n-i$ there are zeros depending on $n-i$ parameters, but no zeros depending on more parameters. This proves the agreement between the arithmetical dimension just defined and the usual parameter definition of dimension in elimination theory.

The ring used in general ideal theory is a commutative ring in which the finite chain condition holds: every chain of ideals
\[
  \mathfrak a_1\supset \mathfrak a_2\supset\cdots
\]
in which each is a proper divisor of the preceding one stops after finitely many steps. This condition is equivalent to the condition that every ideal possess a finite basis; in particular it holds in the polynomial domains just considered.

In such a ring an ideal $\mathfrak p$ is prime if divisibility of a product by $\mathfrak p$ implies divisibility of at least one factor by $\mathfrak p$. An ideal $\mathfrak q$ is primary if divisibility of a product by $\mathfrak q$ implies divisibility of one factor by $\mathfrak q$ or divisibility of a power of the other factor by $\mathfrak q$:
\[
  ab\equiv0\pmod{\mathfrak p}\Rightarrow a\equiv0\pmod{\mathfrak p}\text{ or }b\equiv0\pmod{\mathfrak p},
\]
\[
  ab\equiv0\pmod{\mathfrak q}\Rightarrow a\equiv0\pmod{\mathfrak q}\text{ or }b^r\equiv0\pmod{\mathfrak q}\text{ for some }r.
\]
To every primary ideal $\mathfrak q$ there is a unique associated prime ideal $\mathfrak p$: it divides $\mathfrak q$, and some power of it is divisible by $\mathfrak q$; moreover $\mathfrak p$ consists of all elements of the ring some power of which is divisible by $\mathfrak q$. Every prime ideal is primary; proper primary ideals are those not themselves prime.

The decomposition theorem says that every ideal admits a representation
\[
  \mathfrak a=[\mathfrak q_1,\ldots,\mathfrak q_s]
\]
as a least common multiple of primary ideals. It may be chosen shortest, meaning that no $\mathfrak q_i$ occurs in the least common multiple of the others. Then the primary components $\mathfrak q_i$ need not be uniquely determined, but their associated prime ideals are determined by $\mathfrak a$. Components belonging to isolated prime ideals, that is, to associated prime ideals which contain no other associated prime ideal, are themselves uniquely determined.

Finally, from Steinitz's field theory, recall that a field is perfect if every prime function with coefficients in it splits in a suitable extension field into distinct linear factors. Otherwise it is imperfect. Every field of characteristic zero is perfect; a field of characteristic $p$ is perfect exactly when together with every element it contains its $p$-th root. Over an imperfect field a prime function can be a function of $x^{p^f}$ and split into $p^f$-th powers of distinct linear factors. The integer $f$ is the exponent. A prime function is of the first kind if the exponent is zero, i.e. if it splits into distinct linear factors in a suitable extension. An extension is of the first kind if every element is of the first kind. In a finite extension of an imperfect field there is a subfield of first kind; its degree is the reduced degree, and the exponent measures the inseparable $p$-power part.

\subsection{2. Elementary-divisor form and norm of prime and primary ideals}

We now work throughout with the polynomial domain defined above. The first lemmas state that passage to the transformed domain does not change primeness or primaryness.

\textbf{Lemma I.} The transform of a prime ideal is prime; the transform of a primary ideal is primary. In other words, the property of being prime or primary is preserved when the indeterminates $u_{rs}$ are adjoined to the coefficient field.

For prime ideals the proof is based on the fact that the residue classes modulo a prime ideal form a domain without zero divisors, a property preserved by adjoining indeterminates. If $a(x)b(x)$ is divisible by the transformed ideal and if neither factor is so divisible, then, after writing the factors in terms of the inverse transformation, the highest non-vanishing coefficients with respect to a lexicographic order give a product not divisible by the original prime ideal, a contradiction.

For primary ideals the argument is parallel, but uses the Dedekind--Mertens theorem. If $a$ is not divisible by $\mathfrak q$ and no power of $b$ is divisible by $\mathfrak q$, one chooses transformed representatives $a(x),b(x)$ with the same property. The ideals generated by their coefficients then satisfy the Dedekind--Mertens relation for a sufficiently high exponent; this prevents the coefficient ideal of the product from being divisible by $\mathfrak q$, and therefore the product itself is not divisible by the transformed ideal.

\textbf{Lemma II.} If
\[
  \mathfrak m=[\mathfrak c_1,\ldots,\mathfrak c_s]
\]
is a representation in the original polynomial domain, then the transformed ideal satisfies
\[
  \mathfrak m'=[\mathfrak c_1',\ldots,\mathfrak c_s'].
\]
Consequently, in a shortest representation by greatest primary components, all components may be assumed transformed.

The proof is direct. Every transformed polynomial of $\mathfrak m$ is, after multiplication by a suitable denominator from the transformation coefficients, a linear combination of transformed polynomials from the original ideal. Conversely, a polynomial divisible by every transformed component is divisible by the transformed least common multiple. Applying Lemma I to a primary decomposition in the original domain gives a primary decomposition in the transformed domain.

The first characterization of prime and primary ideals by elementary-divisor form and norm is then the following.

\textbf{Theorem I.} The elementary-divisor form of a prime ideal is a prime function; its norm is a power of this prime function. For a primary ideal, elementary-divisor form and norm are primary functions, and indeed powers of the same prime function, namely the elementary-divisor form of the associated prime ideal. For prime and primary ideals the highest dimension agrees with the dimension in the ordinary sense, and this dimension is the same for a primary ideal and its associated prime ideal.

For the unit ideal the assertion is clear. Suppose $\mathfrak p$ is a prime ideal of highest dimension $n-i$. The congruences defining the fundamental ideals show that $\mathfrak p$ coincides with the appropriate fundamental ideal. All individual norms before level $i$ are therefore one, and the elementary-divisor form is the single factor $E^{(i)}$. If $E^{(i)}$ split into a product of two non-units, the defining greatest-common-divisor property would imply that one factor already gives divisibility by $\mathfrak p$; hence $E^{(i)}$ is a prime function. The norm is a power of it because a power of the elementary-divisor form is divisible by the norm and the norm is divisible by the elementary-divisor form.

For a primary ideal $\mathfrak q$ the same argument gives a single elementary-divisor form at the relevant level. If $\mathfrak p$ is the associated prime ideal, then some power of $\mathfrak p$ is divisible by $\mathfrak q$, and $\mathfrak q$ divides $\mathfrak p$. The comparison of dimensions follows from these mutual divisibilities; the elementary-divisor form of $\mathfrak q$ and its norm are powers of the elementary-divisor form of $\mathfrak p$.

The dimension criterion for prime ideals is stated next.

\textbf{Theorem II.} A prime ideal has no proper divisor of the same highest dimension; conversely, every ideal with this property is prime.

The proof uses the following elementary fact.

\textbf{Lemma III.} If a prime ideal in a polynomial domain over a field has only finitely many residue classes linearly independent over that field, then it has no proper divisor different from the unit ideal. Equivalently, its residue-class system is already a field.

The hypothesis gives a finite number of residue classes such that every residue class is linearly dependent on them. Multiplication by a non-zero residue class is then an injective linear transformation of a finite-dimensional vector space over the ground field, hence also surjective. Thus every non-zero class has an inverse; the residue classes form a field.

To prove Theorem II, localize the situation at the stage determined by the highest dimension. A proper divisor of a prime ideal of the same highest dimension would, after adjoining the remaining parameters, give a proper divisor of a zero-dimensional prime ideal, contradicting Lemma III. Conversely, if an ideal $\mathfrak q$ has no proper divisor of the same highest dimension and nevertheless is not prime, choose $a,b$ with $ab\equiv0\pmod{\mathfrak q}$ but neither factor divisible by $\mathfrak q$. The ideal $[\mathfrak q,a]$ is then a proper divisor; the dimension comparison, together with the finite-chain condition, contradicts the defining property. Hence $\mathfrak q$ is prime.

\subsection{3. Residue-class fields and fields of zeros}

Let $\mathfrak p$ be a prime ideal. The residue classes modulo $\mathfrak p$ form a domain. By forming quotients one obtains the residue-class field $(\mathfrak R)$ of $\mathfrak p$. We shall construct a field of zeros isomorphic to $(\mathfrak R)$ and use this to connect the norm with the usual splitting into linear factors.

Let $\mathfrak p$ be of dimension $n-i$. The residue classes of the last $n-i$ variables may be taken as parameters. After adjoining them, the remaining residue classes are algebraic over the field $P(u,x_{i+1},\ldots,x_n)$. The elementary-divisor form $P^{(i)}$ is the prime function of a suitable class, say the class of a fundamental linear form. The residue class of this form is algebraic over the parameter field and generates, together with the parameters, a field isomorphic to the quotient field of the residue-class domain.

\textbf{Theorem III.} To the residue-class field of a prime ideal $\mathfrak p$ there corresponds an isomorphic field of zeros. If $\mathfrak p$ has dimension $n-i$, this zero-field is a finite algebraic extension of the parameter field $P(u,x_{i+1},\ldots,x_n)$. Conversely every zero-field of $\mathfrak p$ of this transcendence degree is isomorphic to the residue-class field.

A zero $\alpha=(\alpha_1,\ldots,\alpha_n)$ is obtained by sending each residue class $(x_j)$ to the corresponding element $\alpha_j$ in an extension field. Since exactly the congruences of $\mathfrak p$ become zero, this gives an isomorphism of the residue-class field with the field $P(\alpha_1,\ldots,\alpha_n)$. Conversely any zero-field gives such a residue-class homomorphism and hence the same abstract field.

Passing from the zero-field to the corresponding Galois field over the parameter field gives the factorization of the elementary-divisor form. If $\alpha_1$ is a chosen primitive element, then the associated prime function in $\alpha_1$ splits in the Galois field into linear factors. Returning from the transformed variables $x$ to the original variables $y$, the element $\alpha_1$ becomes the fundamental form
\[
  u_1\beta_1+\cdots+u_n\beta_n,
\]
where $\beta=(\beta_1,\ldots,\beta_n)$ runs through systems of zeros in the original variables depending algebraically on the parameters. Thus the linear factors of the norm are precisely the factors corresponding to the conjugate systems of zeros.

In modern terms, the residue-class field is the function field of the irreducible component determined by $\mathfrak p$, and the zero-field is a concrete algebraic realization of that function field in which the coordinates of the generic point are displayed. Noether's proof keeps the older elimination terminology, but the construction is this one.

\textbf{Lemma IV.} The elementary-divisor form of a prime ideal is the prime function attached to a primitive element of its residue-class field; its norm is the product of the elementary divisors obtained in the module representation, and hence a power of that prime function. If the coefficient field is perfect, the power is the first power.

In an imperfect field the prime function may be inseparable. Then the norm records the full degree, while the elementary-divisor form records the reduced degree. This is the source of the possible power in Theorem I.

\textbf{Theorem IV.} In the Galois field determined by a prime ideal $\mathfrak p$, the norm $N(\mathfrak p)$ decomposes into linear factors corresponding to the conjugate zeros of $\mathfrak p$. The degree of the norm equals the number of residue classes linearly independent over the parameter field, counted with the inseparable exponent when the coefficient field is imperfect.

For a primary ideal $\mathfrak q$ associated with $\mathfrak p$ the zero-field is the same as that of $\mathfrak p$, because $\mathfrak p$ and $\mathfrak q$ have the same zeros. The norm of $\mathfrak q$ is a power of the elementary-divisor form of $\mathfrak p$; the exponent measures the length of the residue-class ring of $\mathfrak q$ over the residue-class field of $\mathfrak p$ at the corresponding level.

\subsection{4. Dimension}

We now compare the arithmetical dimension defined by the first non-unit norm with the usual parameter dimension.

\textbf{Theorem V.} The dimension of a prime ideal is given by the transcendence degree of its residue-class field over $P(u)$. If $\mathfrak p$ is of dimension $n-i$, then its residue-class field contains $n-i$ algebraically independent parameters and is algebraic over the field generated by these parameters.

The proof follows from the construction of the zero-field. The last $n-i$ variables can be chosen so that their residue classes are algebraically independent, while every further residue class is algebraic over them. No set of $n-i+1$ residue classes can be algebraically independent, because the relevant elementary-divisor form gives an algebraic equation. Conversely the selected parameter classes are independent because otherwise a non-zero polynomial in them would belong to $\mathfrak p$, forcing a non-unit individual norm at an earlier level.

This proves the equivalence of the dimension in elimination theory with the arithmetical chain definition. Noether's formulation is deliberately independent of geometric language: an ideal has dimension $n-i$ when the first non-trivial elimination norm occurs at level $i$; it has the same dimension when its residue-class field has transcendence degree $n-i$; and the same number is characterized by the maximum length of chains of prime ideals lying above it.

\textbf{Theorem VI.} Let
\[
  \mathfrak p_0\supset \mathfrak p_1\supset\cdots\supset \mathfrak p_r
\]
be a chain of prime ideals. Then the dimensions strictly decrease along a proper chain. Conversely, a prime ideal of dimension $n-i$ can be placed in a chain of prime ideals of length $n-i$ ending at a maximal ideal. Thus the dimension equals the maximum length of a chain from the given prime ideal down to a maximal ideal, counted in the appropriate direction.

The geometric translation is that passing from a variety to a proper subvariety lowers the number of independent parameters. In the polynomial domain, the arithmetical proof uses divisibility of norms: a proper divisor of a prime ideal cannot have the same highest dimension by Theorem II; hence the dimension must decrease. Conversely, the existence of parameters and of algebraic equations in the remaining variables allows one successively to specialize a parameter and obtain a proper prime divisor at each step until a zero-dimensional prime ideal is reached.

\textbf{Theorem VII.} A prime ideal $\mathfrak p$ is of dimension $n-i$ if and only if it contains a prime chain of exactly $i$ algebraic elimination steps, or equivalently if the first non-unit individual norm attached to $\mathfrak p$ occurs at level $i$.

This theorem supplies the promised arithmetical formulation of the dimension concept. It is invariant under the transformations used in elimination theory and does not depend on a distinguished basis of the ideal.

\subsection{5. Connection with the decomposition theorems of general ideal theory}

The components introduced in Hentzelt--Noether through fundamental ideals are now identified with the isolated components of the primary decomposition of general ideal theory. This is the main comparison theorem.

Let the associated prime ideals of $\mathfrak m$ be arranged by dimension:
\[
  \mathfrak p_{0,1},\ldots;\quad \mathfrak p_{1,1},\ldots;\quad \mathfrak p_{2,1},\ldots;\quad\ldots,
\]
where $\mathfrak p_{\lambda,\mu}$ denotes a prime ideal of dimension $n-\lambda$. Let $\mathfrak q_{\lambda,\mu}$ be the corresponding primary ideals in a fixed primary decomposition of $\mathfrak m$.

For a given level $i$, the fundamental ideal of level $i$ is, by definition, the unit ideal for those primary components whose associated prime ideals have dimension less than the corresponding threshold; for components of higher dimension it remains the given primary component. Thus the fundamental ideals are least common multiples of precisely those primary components whose associated primes are sufficiently large in dimension.

\textbf{Lemma V.} The fundamental ideal of level $i$ is the isolated component of the primary decomposition defined by all associated prime ideals of dimension greater than $n-i$. More precisely, it is the least common multiple of those primary components whose associated prime ideals have dimension greater than $n-i$.

This follows from the defining congruence for the fundamental ideals. Multiplication by a suitable polynomial in the eliminated variables removes the components in smaller dimensions, while components of higher dimension remain. Since isolated components are uniquely determined in general ideal theory, the fundamental ideals are independent of the chosen primary decomposition.

\textbf{Theorem VIII.} If all associated prime ideals of $\mathfrak m$ have the same dimension $n-i$, then the norm of $\mathfrak m$ decomposes into greatest primary factors corresponding one-to-one to the associated prime ideals. Each factor is a power of the elementary-divisor form of the corresponding prime ideal.

In this unmixed case the relevant fundamental ideals reduce directly to the primary components. The individual norm at level $i$ is the only non-unit norm. Its factorization into greatest primary factors therefore mirrors the primary decomposition of $\mathfrak m$.

\textbf{Theorem IX.} In the general case, where associated prime ideals of different dimensions occur, the decomposition of the norm of $\mathfrak m$ is obtained by grouping the contributions according to dimension. The greatest primary factors of the norm at the first non-trivial level correspond to the isolated components of greatest dimension; after their removal the next individual norm gives the components of the next dimension, and so on.

This is precisely the arithmetical form of the statement that the zeros of an ideal are the union of the zeros of its associated prime ideals. The norm does not merely vanish at those zeros; its factorization displays the associated components in the order of their dimension.

\textbf{Theorem X.} The degree of a greatest primary factor of the norm corresponding to an associated prime ideal $\mathfrak p$ is equal to the number of residue classes, linearly independent over the parameter field, of the component of the fundamental ideal at level $i-1$ modulo the component of the fundamental ideal at level $i$ determined by $\mathfrak p$.

This theorem gives the multiplicity. Instead of assigning multiplicity by a geometric convention at the zeros, Noether assigns it by a length, or number of independent residue classes, in a quotient of isolated components. For polynomials in one indeterminate this is exactly the ordinary multiplicity of a root: the primary component is the power of the corresponding prime ideal, and the number of residue classes is the exponent. In higher-dimensional elimination theory, the same algebraic length replaces the root multiplicity.

The theorem also explains the role of the ``basic ideals'' in Hentzelt--Noether. They are not additional arbitrary constructions; they are the isolated components selected by dimension in the primary decomposition. Their quotient lengths are the multiplicities recorded by the norm.

\textbf{Theorem XI.} The decomposition of an ideal into primary components and the decomposition of its norm into greatest primary factors determine each other at the level of associated prime ideals and multiplicities. In particular, two ideals one of which divides the other and which have the same norm must be equal.

This recovers, in the polynomial domain, the arithmetical theorem analogous to the theorem on norms of ideals in algebraic number fields.

\subsection{6. Characterization of prime ideals and proper primary ideals by elementary-divisor form and norm}

The preceding theorems characterized primary ideals in a way that did not yet separate prime ideals from proper primary ideals. We now obtain the separation.

\textbf{Theorem XII.} A necessary condition for an ideal to be prime is that its elementary-divisor form be a prime function; a sufficient condition is that its norm be a prime function. Equivalently, the elementary-divisor form of a prime ideal is always a prime function, and the norm of a proper primary ideal is always a proper primary function.

The first assertion has already been proved in Theorem I. For the sufficient condition, suppose the norm of an ideal is a prime function. If the ideal had a proper divisor of the same highest dimension, then the first individual norm in which the two ideals differ would be a proper divisor of the given prime function, impossible unless it is a unit. Thus every proper divisor has lower highest dimension, and Theorem II shows that the ideal is prime.

A second proof uses the following lemma.

\textbf{Lemma VII.} The residue-class system represented by polynomials modulo the elementary-divisor form $E^{(i)}$ of an ideal whose associated prime ideals all have dimension $n-i$ is isomorphic to a subsystem of the residue-class system modulo the ideal. The degree of the elementary-divisor form is the degree of this subsystem over the quotient field $P(u,x_{i+1},\ldots,x_n)$.

Indeed, classes represented by the same polynomial correspond to each other. The correspondence is injective because exactly those polynomials divisible by the elementary-divisor form represent the zero class in the subsystem. If the norm is a prime function and hence agrees in degree with the elementary-divisor form, this subsystem exhausts all residue classes. Since a residue-class system modulo a prime function has no zero divisors, the ideal is prime.

In general Theorem XII cannot be sharpened to necessary and sufficient conditions. Over imperfect fields examples show two failures: the norm of a prime ideal can be a proper power, and the elementary-divisor form of a proper primary ideal can be a prime function. However the perfect-field case is clean.

\textbf{Theorem XIII.} If the coefficient field $P$ is perfect, then the norm and the elementary-divisor form of a prime ideal are prime functions and are identical. The norm and the elementary-divisor form of a proper primary ideal are proper primary functions. Thus, over a perfect field, the condition that the elementary-divisor form be a prime function is necessary and sufficient for primeness; the same condition may be stated using the norm.

It remains only to show that over a perfect field the ideal is prime whenever the elementary-divisor form is a prime function, and that the norm then equals this prime function. Since $P$ is perfect, a prime elementary-divisor form is of first kind. The residue-class field obtained by adjoining a root is therefore separable; by the usual Lagrange formula all the necessary residue classes lie in the subsystem described in Lemma VII. That subsystem exhausts the residue classes of the ideal. It has no zero divisors, so the ideal is prime. The equality of degrees gives equality of norm and elementary-divisor form.

For imperfect coefficient fields the following sharper statements hold.

\textbf{Theorem XIV.} For a prime ideal over a field of characteristic $p$, the norm is a $p^f$-th power of its elementary-divisor form, where $f$ is the exponent of the finite extension of residue-class fields. If the elementary-divisor form of a proper primary ideal is a prime function, then its norm is a power of that prime function; the exponent need not be a power of $p$ in general.

The proof follows the field-theoretic decomposition into reduced degree and inseparable exponent. The elementary-divisor form measures the reduced degree; the norm measures the full degree. Their quotient is therefore the inseparable $p^f$ part for prime ideals. For proper primary ideals the nilpotent, or primary, multiplicity contributes an additional exponent.

Noether gives explicit examples in characteristic $2$ and $3$. If $P$ is obtained by adjoining independent indeterminates to a finite field, then an ideal such as
\[
  (y_1^2+\lambda,\; y_2^2+\mu)
\]
can be prime, because its residue-class field is purely inseparable of degree four; its elementary-divisor form has reduced degree while its norm is the square or fourth power accordingly. Specializing parameters can produce a proper primary ideal whose elementary-divisor form is still prime. In characteristic $3$, a similar construction shows that the norm exponent for a proper primary ideal need not equal the characteristic exponent predicted for prime ideals.

These examples demonstrate exactly where imperfectness enters. In a perfect coefficient field, no inseparable power is hidden in the field extension, and the elementary-divisor form and norm detect primeness without ambiguity.

\subsection{7. Absolute prime ideals}

A prime ideal with coefficients in a field $P$ is called an absolute prime ideal if it remains prime after extending the coefficient field to the algebraically closed field $A$ belonging to $P$. Correspondingly, a prime function is absolute if it remains irreducible, or prime, over $A$. The characterization is the following.

\textbf{Theorem XV.} A necessary and sufficient condition for a prime ideal to be absolute is that its elementary-divisor form, chosen integral and primitive in the transformation coefficients $u$, be an absolute prime function as a polynomial in the $x$ and $u$. The elementary-divisor form then coincides with the norm, so the condition may also be stated for the norm.

The proof starts from the observation that norm and elementary-divisor form are preserved under algebraic extension of the coefficient field. Each individual factor is defined as a greatest common divisor in a polynomial ring, and this property is unchanged by algebraic extension. Passing from $P$ to the algebraically closed field $A$ passes in particular to a perfect field. By Theorem XIII, over $A$ primeness is equivalent to the elementary-divisor form being a prime function. Hence the original prime ideal is absolute exactly when that form is absolutely prime. Since over $A$ the norm equals the elementary-divisor form, and equality descends back to $P$ by the preceding invariance, the assertion follows.

For absolute prime functions with coefficients in a finite algebraic number field, Ostrowski proved that the function remains absolute prime modulo every prime ideal of the number field with only finitely many exceptions. The transfer to ideals uses the following finite-denominator statement.

\textbf{Theorem XVI.} Let the coefficient domain be the ring $\mathfrak o^*$ of algebraic integers in a finite algebraic number field, rather than a field. Then the norm and elementary-divisor form of a polynomial ideal remain the norm and elementary-divisor form after reduction modulo every prime ideal of $\mathfrak o^*$, with at most finitely many exceptions.

The proof modifies the polynomial domain. Its coefficients now lie in $\mathfrak o^*$, while the transformation coefficients $u$ are adjoined polynomially. In forming the module representation for the individual norms, one reduces powers of the eliminated variable modulo a polynomial regular in that variable. The possible denominators that occur are products of finitely many polynomials in the $u$ with coefficients in $\mathfrak o^*$. Their coefficients generate an ideal of $\mathfrak o^*$ divisible by only finitely many prime ideals. Avoiding these prime ideals preserves the module representation after reduction.

The same is done for the greatest-common-divisor property defining the norm and elementary-divisor form. One chooses the norm and elementary-divisor form so that they are divisible by the ideal over $\mathfrak o^*[u]$. For all but finitely many prime ideals of $\mathfrak o^*$ the relevant ranks of the modules, the relevant determinants, and the greatest-common-divisor relations survive reduction. Thus, modulo such a prime ideal, the reduced forms are still exactly the norm and elementary-divisor form of the reduced ideal.

Combining Theorems XV and XVI with Ostrowski's theorem gives the final result.

\textbf{Theorem XVII.} If $\mathfrak p$ is an absolute prime ideal with coefficients that are algebraic integers in a finite algebraic number field, then $\mathfrak p$ remains a prime ideal, more precisely an absolute prime ideal, modulo every prime ideal of the number field, with at most finitely many exceptions.

Choose the norm $R_{\mathfrak p}$ as in Theorem XVI so that it is divisible by $\mathfrak p$ over $\mathfrak o^*[u]$. Then
\[
  R_{\mathfrak p}=T(u)P(x),
\]
where $P(x)$ may be chosen integral and primitive in the $u$. By Theorem XV, regarded as a polynomial in $u$ and $x$ with coefficients in the number field, $P$ is an absolute prime function. Ostrowski's theorem says that $P$ remains absolute prime modulo all but finitely many prime ideals. Avoiding also the finite exceptional set in Theorem XVI, this reduced $P$ is the norm of the reduced ideal. Theorem XV then implies that the reduced ideal is again an absolute prime ideal.

This completes the transfer of the arithmetic theorem on absolute prime functions to absolute prime ideals.

\newpage

\section{Elimination theory and ideal theory}

I want to show how elimination theory can be built up in complete parallel with the theory of zeros for polynomials in one indeterminate. This construction rests on a connection between the elimination theory given by Hentzelt and the methods of field theory as developed by Steinitz, and of ideal theory as I have developed it.

First consider the case of polynomials in one indeterminate. The coefficient domain is a fixed field $P$, and notions such as prime function are to be understood relative to $P$. The field may be abstractly defined in Steinitz's sense; in special cases it may be the field of rational numbers or the field of all complex numbers, as was formerly usual in elimination theory. For polynomials in one indeterminate the theory of zeros falls into four steps.

First comes the decomposition theorem, which reduces the question to prime functions, that is, to polynomials irreducible over $P$. From a factorization
\[
  a(x)=q_1(x)q_2(x)\cdots q_s(x),
\]
where the underlying $p_i(x)$ are distinct prime functions and the $q_i(x)$ are their powers, hence greatest primary functions, it follows that $a(x)$ vanishes exactly at the zeros of the prime functions $p_i(x)$.

Second, one constructs the zero-field for a prime function $p(x)$. Such a field
\[
  K=P(\alpha), \qquad p(\alpha)=0,
\]
can always be constructed, by Steinitz, as an extension field of $P$ isomorphic to the residue-class field of $p(x)$. Conversely, every zero-field lying in an extension field of $P$ is isomorphic to that residue-class field. Thus zero-field and residue-class field are abstractly identical.

Third, one decomposes $p(x)$ into linear factors in the Galois field determined by $p(x)$. Repetition of the construction just described leads to an essentially uniquely determined Galois field $P(\alpha_1,\ldots,\alpha_r)$ large enough for $p(x)$ to split into linear factors.

Fourth, one interprets multiplicity for the original polynomial $a(x)$. By the preceding steps the zeros and the corresponding linear factorization are known for an arbitrary polynomial. The multiplicity may be taken to be the degree of the corresponding primary function $q(x)$, equivalently the number of residue classes modulo $q(x)$ that are linearly independent over $P$. If $P$ is algebraically closed and the prime functions are linear, this number is exactly the usual multiplicity of the zero.

Now pass to general elimination theory. Let $\mathfrak S$ be the domain of all polynomials in $x_1,\ldots,x_n$ with coefficients in $P$. Elimination theory may then be defined as the theory of the zeros of an ideal in this polynomial domain. An ideal $\mathfrak a$ is a system of polynomials such that with $f$ and $g$ it contains $f-g$, and with $f$ it contains $Af$ for every polynomial $A$ in $\mathfrak S$. A zero of $\mathfrak a$ is a system of values $\alpha_1,\ldots,\alpha_n$ in an extension field of $P$ for which every polynomial $f$ in $\mathfrak a$ vanishes.

If $f_1,\ldots,f_s$ is an ideal basis, so that every $f\in\mathfrak a$ has the form
\[
  f=A_1f_1+\cdots+A_sf_s,
\]
then the zeros of $\mathfrak a$ are exactly the common zeros of $f_1,\ldots,f_s$. Conversely, any finite system of polynomials may be regarded as a basis of the ideal it generates. Thus the ideal-theoretic definition of zeros is the usual one, but avoids the arbitrariness of singling out a particular basis. In one variable this difficulty did not appear, because only principal ideals occur and the basis polynomial is determined up to a unit. Even there, however, the ideal viewpoint gives a clearer interpretation.

The theory again falls into four steps.

I. The decomposition theorem represents an ideal as a least common multiple of greatest primary components and thereby reduces the question to prime ideals. An ideal $\mathfrak p$ is prime if divisibility of a product by $\mathfrak p$ implies divisibility of one factor; an ideal $\mathfrak q$ is primary if divisibility of a product by $\mathfrak q$ implies divisibility of one factor or of some power of the other factor. To every primary ideal $\mathfrak q$ there belongs a unique associated prime ideal $\mathfrak p$ which divides $\mathfrak q$ and some power of which is divisible by $\mathfrak q$. One has a representation
\[
  \mathfrak a=[\mathfrak q_1,\ldots,\mathfrak q_s],
\]
where the $\mathfrak q_i$ are greatest primary components. Although the primary components themselves need not be uniquely determined, their associated prime ideals are uniquely determined by $\mathfrak a$. Since every associated prime ideal divides $\mathfrak a$, and conversely a product of powers of the associated primes is divisible by $\mathfrak a$, the zeros of the ideal agree with the zeros of its associated prime ideals.

II. The residue classes modulo a prime ideal form a ring without zero divisors. If the prime ideal is not the unit ideal, this ring contains a non-zero element, and by adjoining quotients it becomes the residue-class field of the prime ideal. There is always an extension field of $P$ isomorphic to this residue-class field; it becomes the zero-field of the prime ideal. Thus
\[
  K=P(\alpha_1,\ldots,\alpha_n),
\]
where $(\alpha_1,\ldots,\alpha_n)$ is a zero of $\mathfrak p$. The field has transcendence degree $k$, $0\leq k<n$, over $P$: it contains $k$ algebraically independent elements over $P$, while every $k+1$ elements are algebraically dependent. Conversely, every zero-field of transcendence degree $k$ is isomorphic to the residue-class field. Therefore zero-field and residue-class field are abstractly the same object.

III. The norm of $\mathfrak p$ is decomposed into linear factors in the Galois field determined by $\mathfrak p$. If the variables $x_1,\ldots,x_n$ are obtained from original variables $y_1,\ldots,y_n$ by a linear transformation with indeterminate coefficients, then the zero-field of transcendence degree $k$ can be viewed as a finite algebraic extension of $P(x_{k+1},\ldots,x_n)$. Passing to its Galois field, the prime function for a primitive element splits into linear factors. This prime function plays, after returning to the $y$ variables, the role of the fundamental equation: its root is the fundamental linear form in a zero system $\beta$. The factorization therefore gives the product over all zeros. The prime function, or a power of it, is also the norm $N(\mathfrak p)$ of the prime ideal, obtained by regarding $\mathfrak p$ as a module of linear forms in the monomials in the eliminated variables and defining the norm as the product of the elementary divisors. Over a perfect coefficient field this norm is exactly the prime function; over an imperfect field it may be a power.

IV. For the original ideal $\mathfrak a$, one regards it successively as a module of linear forms in the monomials of $x_1,\ldots,x_{i-1}$ for $i=1,\ldots,n$. At each step, with coefficient field $P(x_i,
\ldots,x_n)$, only finitely many elementary divisors differ from one. Their product gives the individual norm, and the norm $N(\mathfrak a)$ is the product of all individual norms. The greatest primary factors of this norm correspond one-to-one to the associated prime ideals $\mathfrak p$ of $\mathfrak a$. Each such factor is a power of the norm of $\mathfrak p$ or, in the inseparable case, of its highest elementary divisor. Thus the norm of $\mathfrak a$ splits according to the associated primes, and the preceding steps give the linear factorization for all zeros of $\mathfrak a$.

The multiplicity attached to the zeros corresponding to an associated prime ideal $\mathfrak p$ is the degree of the corresponding primary factor of the norm. This degree is the number of linearly independent residue classes, over the appropriate parameter field, of the isolated component determined by $\mathfrak p$ and by the prime ideals of higher dimension. In the one-variable case this is exactly the degree of the primary factor and therefore the usual multiplicity.

The case of one indeterminate is thus fully included by passing from a polynomial to its principal ideal. Its factorization may be read either as a least common multiple of primary components or as a decomposition of the norm into norms of prime ideals. The zero-field belongs to the prime ideal defined by $p(x)$, while the norm gives the linear factorization. The multiplicity is the number of residue classes in the relevant quotient, and the lowest fundamental ideal is the unit ideal whereas the first fundamental ideal is already the ideal itself.

In this form the usual elimination theory is incorporated in ideal theory: the norm of the ideal has the meaning of a resultant. The old elimination problem becomes the part of ideal theory in the polynomial domain that treats, alongside ideal decompositions, the decomposition of norms and the relation between the two decompositions.

\newpage

\section{Abstract construction of ideal theory in algebraic number fields}

The following abstract properties were indicated as completely equivalent to the ideal theory in an algebraic number field, and hence as characterizing all rings that possess such an ideal theory. They are domains with identity element and without zero divisors in which the double chain condition holds: every divisor chain of ideals stops after finitely many steps, and so does every multiple chain modulo every ideal different from the zero ideal. In addition the ring is integrally closed in its quotient field. If the last assumption is omitted, one obtains theorems on ideal theory in a finite order, some of which are already found without proof in Dedekind's third edition. A detailed exposition will appear in the \emph{Mathematische Annalen}.

\section{Hilbert numbers in ideal theory}

By Hilbert numbers one generally understands numbers that concern the enumeration of residue classes of ideals. Hilbert himself gave, in his ``characteristic function'', a function which for ideals in a polynomial domain gives, from a certain bound on, the exact number of linearly independent residue classes. For lower degrees Macaulay supplemented this function by a generating function, which Ostrowski could evaluate explicitly.

Macaulay also introduced numbers of another kind. These have proved most essential for general ideal theory because they admit an abstract formulation. In accordance with the general decomposition theorems it suffices to define them for primary ideals $\mathfrak q$. They count the number of linearly independent residue classes with respect to the residue-class field of the associated prime ideal $\mathfrak p$.

These numbers decompose into subsystems given by the numbers of irreducible components of
\[
  \mathfrak q,
  \quad \mathfrak q/\mathfrak p,
  \quad \mathfrak q/\mathfrak p^2,
  \quad \ldots .
\]
The number of the $i$-th subsystem is also characterized as the number of terms in a composition series running from $\mathfrak q/\mathfrak p^i$ to $\mathfrak q/\mathfrak p^{i-1}$. The total composition series and its Hilbert numbers form the substitute for the representation of primary ideals as powers of prime ideals, a representation which is no longer valid in general, for example in finite orders of an algebraic number field.

\section{Group characters and ideal theory}

Let $G$ be a finite group and let $K$ be a splitting field for $G$. The group algebra $K[G]$ decomposes into simple components corresponding to the irreducible characters. The arithmetical questions concern what happens when one replaces $K$ by a subring, or by the ring of algebraic integers in a number field, and considers ideals or modules stable under the group operation.

The character equations give idempotent decompositions after extension to the splitting field. Ideal theory supplies the descent from these decompositions to integral modules. In this way the elementary divisors attached to a group representation can be read in parallel with the elementary divisors of ideals. The irreducible characters describe the simple constituents over the field; the corresponding integral ideal theory describes how these constituents combine when denominators and divisibility are taken into account.

This viewpoint places the arithmetic of group characters beside the arithmetic of polynomial ideals. In both cases a decomposition over a field gives a formal splitting, while ideal theory records the integral and modular data lost in the field decomposition alone. The connection is especially useful when studying representations modulo prime ideals of the coefficient ring: the character-theoretic decomposition may remain valid, refine, or coalesce according to the behavior of the associated ideals after reduction.

\newpage

\section{Abstract construction of ideal theory in algebraic number and function fields: first installment}

In the following an abstract characterization is given of all rings whose ideal theory agrees with the ideal theory of all algebraic integers in an algebraic number field, that is, rings in which every ideal is uniquely expressible as a product of powers of prime ideals. Among the best known examples are also the ring of all integral elements of an algebraic function field in one indeterminate, and more generally, for several indeterminates, the rings obtained by restricting oneself, in Kronecker's sense, to ideals of highest dimension and passing to the corresponding quotient ring, the functional domain.

The axioms completely equivalent to this ideal theory, for the commutative ring under consideration, are the following.

\begin{description}
\item[I. Divisor chain condition.] Every chain of ideals in which each ideal is a proper divisor of the preceding one terminates after finitely many steps; equivalently, in every divisor chain of ideals there is an index from which on all ideals are equal.
\item[II. Multiple chain condition modulo each non-zero ideal.] Every chain of ideals, all of which divide a fixed ideal different from the zero ideal, and in which each ideal is a proper multiple of the preceding one, terminates after finitely many steps.
\item[III.] There is an identity element for multiplication.
\item[IV.] The ring has no zero divisors.
\item[V. Integral closedness in the quotient field.] Every element of the quotient field integral over the ring belongs to the ring.
\end{description}

From these axioms Noether develops step by step an increasingly restricted ideal theory, ending with the desired theory, and conversely shows that all five axioms are actually satisfied in the classical examples.

From the divisor chain condition I follows the representation of an ideal as a least common multiple of finitely many primary ideals belonging to distinct prime ideals. Noether repeats the proof because it is simplified by using the ideal quotient and because a gap in the earlier proof is to be corrected: the old proof used, at one point, the existence of a multiplicative identity, which had not been assumed.

The multiple chain condition II implies that in the residue-class ring modulo every ideal different from the zero ideal, a prime ideal can have no proper divisor other than the unit ideal containing all elements. Hence the primary components of such ideals are uniquely determined except for components belonging to the unit ideal. In a somewhat less sharp form this result already occurs in work of Masazo Sono. Sono's hypothesis of the existence of a composition series in the residue-class ring is equivalent to the double chain condition in that ring. By the double chain condition Noether means the simultaneous validity of the divisor chain condition and the multiple chain condition, without restricting to non-zero ideals.

If axiom III, the existence of the unit element, is added, then the unit ideal does not occur as an associated prime ideal. The primary components are therefore uniquely determined; they are pairwise relatively prime, and their least common multiple is their product. Axioms I, II, and III hold for all finite orders of an algebraic number field. Dedekind had already stated, without a complete proof, the unique representation of an ideal as a product of pairwise relatively prime one-type ideals, i.e. primary components.

Axiom IV, the absence of zero divisors, implies that the different powers of an ideal different from both the zero ideal and the unit ideal are all distinct, and that the quotient field exists. If finally axiom V is added, integral closedness in the quotient field, then the primary ideals, which by IV are uniquely determined, become powers of prime ideals. This gives the usual decomposition theorem. The proof uses a consequence of integral closedness already drawn by Dedekind in the number-field case: a genuinely fractional element can be written as a quotient of integral elements in such a way that even the square of the numerator is not divisible by the denominator.

In later treatments this consequence is often replaced by the less transparent generalized Gauss theorem, or by corresponding module theorems that give somewhat more. Those results require basis elements in their application; the proof here works with ideals themselves.

In the converse direction the axioms other than V follow almost directly from unique factorization of ideals. Axiom V follows from the proof that every genuinely fractional element can be represented so that no power of the numerator is divisible by the denominator. That is even stronger than Dedekind's original consequence of integral closedness. The proof is obtained only after drawing the customary conclusions from the ideal factorization: principal-ideal representation modulo a fixed ideal, and the theory of fractional ideals, i.e. ideal quotients in the field. Dedekind had already known that integral closedness follows from the representation by prime-ideal powers.

The axioms I--V imply that four decompositions which are generally distinct coincide: decomposition into least pairwise relatively prime ideals, into mutually prime ideals, into greatest primary components, and into irreducible ideals. Conversely, this coincidence together with axioms I--IV implies integral closedness. If zero divisors are allowed in the ring, then the remaining axioms, with V interpreted as integral closedness in the total quotient ring, do not necessarily imply this coincidence; the residue-class ring modulo certain ideals of a finite order gives examples.

Section 1 sketches the theory of elements integral over a ring, through Dedekind's consequences of integral closedness. Section 2 proves the transfer of chain conditions from a ring to modules of linear forms over it. Section 3 applies this to passage to finite extension fields of first kind: axioms I--IV are preserved for every order, and axiom V holds for the order consisting of all integral elements. This places both number fields and function fields under the same scheme. In particular, once it is shown directly that the functional domain derived from integral polynomials in several indeterminates satisfies the axioms, Kronecker's ideal theory is fully incorporated. The later ideal theory does not depend on these classificatory sections; sections 4 and 5 give the foundations independent of axioms I--V and show more clearly which parts of the theory do not require finiteness hypotheses.

\subsection{1. Theory of integral elements}

Let $T$ be a commutative ring without zero divisors and with multiplicative identity $e$, and let $R$ be a fixed subring of $T$ containing $e$. The words module, order, and integral are throughout to be understood relative to this fixed ring $R$.

A system $M$ of elements of $T$ is called an $R$-module, or simply a module, if with two elements $\alpha$ and $\beta$ it contains $\alpha-\beta$, and with $\alpha$ it contains $r\alpha$ for arbitrary $r\in R$. If every element of $M$ also belongs to $N$, then $M$ is said to be divisible by $N$, written $M\equiv0\pmod N$; equivalently, $M$ is a multiple and $N$ a divisor. $R$-modules all of whose elements lie in $R$ are ideals of $R$.

The intersection, or least common multiple, of any family of modules is again a module, and so is $T$. Thus the module generated by any system of elements is defined uniquely as the intersection of all modules containing that system. The sum, or greatest common divisor, of a family of modules is the module generated by their union. The product $MN$ of two modules is the module generated by all products $\mu\nu$ with $\mu\in M$ and $\nu\in N$. This product is associative and commutative because the ring multiplication is; distributivity follows from distributivity in $T$.

Since $R$ itself is a module, saying that $R$ is the multiplier domain is expressed by $RM\equiv0\pmod M$. Since $R$ contains the identity element, also $M\equiv0\pmod{RM}$, hence $M=RM$.

A module is finite if it is generated by finitely many elements. If
\[
  M=(\mu_1,\ldots,\mu_s),
\]
then the $\mu_i$ form a module basis. Sums and products of finitely many finite modules are again finite; for products this follows from writing every element of $M$ as an $R$-linear combination of the basis elements and using commutativity of multiplication in $T$.

An extension ring of $R$ lying in $T$, that is, a subring of $T$ containing all elements of $R$, is called an $R$-order, or simply an order. The intersection of arbitrary orders is again an order, and $T$ is an order; hence the order generated by any set of elements is uniquely defined. Every order is also an $R$-module. An order is finite if it is finite as a module. Sums and products of finite orders are finite orders. There is also a transitivity law: if $A$ is a finite $R$-order and $B$ a finite $A$-order, then $B$ is a finite $R$-order, since products of an $R$-basis of $A$ with an $A$-basis of $B$ give an $R$-basis of $B$.

An element $\alpha\in T$ is called integral over $R$ if the order $R[\alpha]$ generated by it is finite. Equivalently, the chain of modules
\[
  R_1=(1),\quad R_2=(1,\alpha),\quad R_3=(1,\alpha,\alpha^2),\ldots
\]
terminates. Equivalently again, there exists a non-zero principal module $C=(\gamma)$ such that $C R[\alpha]$ is finite; that is, the chain $C R_i$ terminates. Finally this is identical with the usual equation formulation: $\alpha$ is integral if it satisfies an equation
\[
  \alpha^m+r_1\alpha^{m-1}+\cdots+r_m=0,
  \qquad r_i\in R.
\]
These formulations are equivalent. If $R[\alpha]$ is finite, a basis can be chosen among the powers of $\alpha$, and termination of the chain gives the monic equation. Conversely such an equation expresses every high power of $\alpha$ in terms of lower powers. The principal-module formulation is equivalent because multiplying by a non-zero principal module in a domain does not destroy the eventual equality of the chains.

The ring property of the integral elements and the transitivity of integral dependence rest on the following lemma.

\textbf{Lemma.} If $\mathcal O$ is a finite order in $T$ and $\alpha\in\mathcal O$, then the order $R[\alpha]$ generated by $\alpha$ is finite. In other words, all elements of a finite order are integral.

Indeed, let $\omega_1,\ldots,\omega_s$ be a module basis of $\mathcal O$. Since $\alpha\omega_i$ again lies in $\mathcal O$, there are relations
\[
  \alpha\omega_i=\sum_j r_{ij}\omega_j.
\]
The determinant of the corresponding system gives a monic equation for $\alpha$. The use of the determinant is legitimate because $T$, and hence $\mathcal O$, has no zero divisors.

The system $\mathcal G$ of all elements of $T$ integral over $R$ is an order. It contains $R$, because the elements of $R$ are integral with basis $e$. It is closed under addition and multiplication: if $\alpha$ and $\beta$ are integral, the order generated by them is the product of two finite orders and is finite, so $\alpha+\beta$ and $\alpha\beta$ are elements of a finite order and hence integral by the lemma.

The transitivity law for integral dependence is this: if $\mathcal O$ is an order consisting of elements integral over $R$, then every element of $T$ integral over $\mathcal O$ is already integral over $R$. For such an element $\alpha$ satisfies a monic equation with coefficients in $\mathcal O$; the order generated over $R$ by those finitely many coefficients is finite, and $\alpha$ is integral over that finite order. By the transitivity of finite orders, $R[\alpha]$ is finite.

The ring $R$ is called integrally closed in $T$ if every element of $T$ integral over $R$ already belongs to $R$. By the second formulation of integrality, an element $\alpha\in T$ belongs to $R$ exactly when there is a non-zero principal module $C$ such that $C R[\alpha]$ is finite.

Dedekind drew two important consequences from integral closedness.

\textbf{Dedekind consequence I.} Suppose $R$ is integrally closed in $T$, and suppose further that the divisor chain condition for ideals of $R$ holds. If $\beta$ is not integral over $R$ and $c\neq0$ is an element of $R$, then there is an exponent $\rho>0$ such that in the sequence
\[
  c,
  \quad c\beta,
  \quad c\beta^2,
  \quad \ldots,
  \quad c\beta^3,\ldots
\]
the terms $c,\ldots,c\beta^{\rho-1}$ are integral and all subsequent terms are not integral.

The point is that if all $c\beta^r$ were integral, then by integral closedness they would all belong to $R$. The ideals generated by the initial segments would form a divisor chain and would therefore stop. That would imply finiteness of $R[\beta]$ after multiplication by $c$, hence integrality of $\beta$, contrary to assumption. Once a non-integral term occurs, every later term remains non-integral; otherwise a power relation would make the earlier term integral. Thus there is a first non-integral term.

Specializing $T$ to the quotient field of $R$ gives the second consequence.

\textbf{Dedekind consequence II.} Suppose $R$ is integrally closed in its quotient field and satisfies the divisor chain condition for ideals. If $\beta$ is a non-integral element of the quotient field, then it has a representation
\[
  \beta=m/n,
  \quad m,n\in R,
\]
such that $m^2/n$ is also not integral.

Indeed, write $\beta=b/c$ with $c\neq0$. In the sequence $c,c\beta=b,c\beta^2,\ldots$ the first two terms are integral, so the exponent $\rho$ of the preceding consequence is greater than one. Put
\[
  m=c\beta^\rho,
  \qquad n=c\beta^{\rho-1}.
\]
Then $m/n=\beta$, while $m^2/n=c\beta^{\rho+1}$ is not integral.

Finally, integral closedness is itself transitive: if $\mathcal G$ is the system of elements of $T$ integral over $R$, then $\mathcal G$ is also the system of elements integral over $\mathcal G$. Any element integral over $\mathcal G$ is integral over $R$ by transitivity, hence belongs to $\mathcal G$.

\subsection{2. Chain conditions in finite module domains}

The module theorem used in the application says that chain conditions pass from the coefficient ring to finite modules. A module domain $M$ over a ring $R$ is an abelian group with scalar multiplication by elements of $R$ satisfying the associative and distributive laws. All module definitions from the preceding subsection that do not use multiplication of module elements remain valid.

A module domain is finite if finitely many elements $\varepsilon_1,\ldots,\varepsilon_k$ generate it; if $R$ has a unit, its elements are all linear forms
\[
  r_1\varepsilon_1+\cdots+r_k\varepsilon_k.
\]

\textbf{Module theorem.} Let $M$ be a finite module domain over a commutative ring $R$ with identity. If the divisor chain condition, respectively the multiple chain condition, holds for ideals of $R$, then the corresponding divisor, respectively multiple, chain condition holds for submodules of $M$.

The proof assigns to each submodule $U$ of $M$ a finite system of ideals $a_1,\ldots,a_k$ of $R$. An element of $M$ is said to have length $i$ if it can be expressed using $\varepsilon_1,
\ldots,\varepsilon_i$ but not using fewer generators. Let $U_i$ be the elements of $U$ of length at most $i$, and let $a_i$ be the set of coefficients of $\varepsilon_i$ occurring in elements of $U_i$. If $V$ is a proper divisor or proper multiple of $U$, then at least one of the ideals assigned to $V$ is a proper divisor or proper multiple of the corresponding $a_i$. Thus an infinite chain of submodules would produce an infinite chain among finitely many ideal chains in $R$, impossible under the corresponding chain condition.

A useful consequence is the following. If in $R$ the multiple chain condition holds modulo every non-zero ideal, then in a finite module domain $M$ the multiple chain condition holds modulo every submodule whose associated ideals $c_1,\ldots,c_k$ are all non-zero. The proof is the same, applied modulo those non-zero ideals.

\subsection{3. Passage to finite extension fields}

The preceding module theorem is used to bring algebraic number fields and function fields under the axioms.

Assume that in $R$ all axioms I--V except possibly the multiple chain condition II hold. Let $K$ be the quotient field of $R$, and let $L$ be a finite extension field of first kind over $K$. Let $\mathcal G$ be the system of elements of $L$ integral over $R$. Then $\mathcal G$ satisfies the same axioms, and every order contained in $\mathcal G$ satisfies all of them except possibly integral closedness.

The axioms III and IV hold because the integral elements form a ring without zero divisors and with the same identity. Integral closedness of $\mathcal G$ in $L$ follows from the transitivity of integral dependence. It remains to see the divisor chain condition. Since the extension is of first kind, the field $L$ is generated by a single element $\alpha$ integral over $R$. Its conjugates in the Galois field are integral, and the square $D$ of the product of their differences is a non-zero element of $K$, hence, by integral closedness of $R$, an element of $R$. The usual discriminant argument shows that the integral closure $\mathcal G$ is contained in a finite $R$-module generated by elements of the form $\alpha^j/D$. Therefore the module theorem gives the divisor chain condition for ideals of $\mathcal G$ and for ideals of every order in $\mathcal G$.

If all axioms I--V hold in $R$, then all hold in $\mathcal G$, and every order in $\mathcal G$ satisfies all but integral closedness. The only remaining point is axiom II. By the consequence of the module theorem it suffices to show that a non-zero ideal of $\mathcal G$, regarded as an $R$-module inside the finite module containing $\mathcal G$, has all associated coefficient ideals non-zero. This follows because multiplication by a suitable non-zero element of $R$ carries each basis element of the finite module into the given ideal. The same argument applies to orders of full rank; for lower-rank orders one passes to their quotient fields, intermediate between $K$ and $L$ and again of first kind.

It remains only to verify the axioms for the basic domains: the rational integers, polynomials in one indeterminate, and the functional domain of polynomials in several indeterminates. For the rational integers and for polynomials in one indeterminate over a field, the axioms follow from unique factorization into irreducible elements and from the fact that every ideal is principal. Integral closedness follows because a non-integral quotient can be written with coprime numerator and denominator, so that no power of the numerator is divisible by the denominator.

For the ring of polynomials in several indeterminates, axioms III, IV, and V again follow from unique factorization of elements. The divisor chain condition is Hilbert's basis theorem. The multiple chain condition need not hold. Passing to the functional domain restores it. One adjoins a new indeterminate $u$ and considers rational functions whose denominators are primitive polynomials in $u$. The elements with primitive numerator and denominator become units, and every element of the functional domain has a unique expression as an element of the original polynomial ring multiplied by a unit. The factorization theory of the original ring therefore transfers to the functional domain. In that domain every ideal is principal: starting with an element of an ideal and another not divisible by it, the greatest common divisor also lies in the ideal and is a proper divisor; after finitely many repetitions one obtains a basis element. Hence unique factorization of ideals and the axioms I and II follow.


\subsection{4. Isomorphism theorems and direct sums}

In the following only the ring or module structure is assumed; no chain axiom is required.

Let $M$ and $\overline M$ be module domains over $R$. We say that $M$ is homomorphic to $\overline M$ if to every element of $M$ there corresponds one and only one element of $\overline M$, so that all of $\overline M$ is reached, and such that difference and multiplication by the same element of $R$ correspond. Thus from
\[
  \beta\sim \bar\beta,
  \qquad \gamma\sim\bar\gamma
\]
one obtains
\[
  \beta-\gamma\sim \bar\beta-\bar\gamma,
  \qquad r\beta\sim r\bar\beta.
\]
Every $R$-module $B$ in $M$ then corresponds to an $R$-module $\bar B$ in $\bar M$. Conversely, if $\bar C$ is a module in $\bar M$, the totality $C^*$ of elements of $M$ to which elements of $\bar C$ correspond is an $R$-module in $M$ uniquely determined by $\bar C$; it is the module associated with $\bar C$. If $A$ is the module corresponding to the zero element of $\bar M$, then $C^*$ is a divisor of $M$ and decomposes into congruence classes modulo $A$, these classes corresponding one-to-one to the elements of $\bar C$. Starting from a module $B$ in $M$, passing to $\bar B$ in $\bar M$, and coming back gives
\[
  B^*=(B,A),
\]
the greatest common divisor of $B$ and $A$. If $B$ divides $M$, then $B^*=B$.

If the correspondence is reversible, the module domains are isomorphic. If $A$ is any submodule of $M$, the residue-class module $M/A$ is obtained by taking congruence modulo $A$ as the equality relation. Equivalently, one replaces equal elements modulo $A$ by one residue class. Every homomorphism is obtained in this way: if $M$ maps homomorphically onto $\bar M$ and $A$ is the kernel, then $\bar M$ is isomorphic to $M/A$.

The two isomorphism theorems follow. First, if $\bar M=M/A$ and $C$ divides $A$, then
\[
  \bar M/C \simeq M/C.
\]
Indeed, equality modulo $A$ is already equality modulo $C$ after the indicated passage to residue classes. Secondly, for two submodules $B$ and $A$ of $M$ one has
\[
  (B,A)/A \simeq B/[B,A],
\]
where $(B,A)$ denotes the greatest common divisor and $[B,A]$ the least common multiple. This is the module form of the familiar second isomorphism theorem.

The same considerations apply to commutative rings when modules are replaced by ideals. If $R$ maps homomorphically onto $\bar R$, difference and product correspond. If $\mathfrak a$ is an ideal, the residue-class ring $R/\mathfrak a$ is formed by using congruence modulo $\mathfrak a$ as equality; products of residue classes are then well-defined. Every ring homomorphism is obtained in this way, and the isomorphism theorems become
\[
  (R/\mathfrak a)/(\mathfrak c/\mathfrak a)\simeq R/\mathfrak c
  \qquad (\mathfrak c\mid \mathfrak a),
\]
and
\[
  (\mathfrak b,\mathfrak a)/\mathfrak a\simeq \mathfrak b/[\mathfrak b,\mathfrak a].
\]

Assume now that the commutative ring $R$ has a multiplicative identity. Two ideals $\mathfrak a,\mathfrak b$ are relatively prime if their greatest common divisor is the unit ideal,
\[
  (\mathfrak a,\mathfrak b)=\mathfrak o.
\]
The basic rules are most simply stated without writing identity elements explicitly. If $(\mathfrak a,\mathfrak b)=\mathfrak o$, then
\[
  (\mathfrak a,\mathfrak b\mathfrak c)=(\mathfrak a,\mathfrak c).
\]
Therefore from $\mathfrak c\mathfrak b\equiv0\pmod{\mathfrak a}$ and $(\mathfrak a,\mathfrak b)=\mathfrak o$ it follows that $\mathfrak c\equiv0\pmod{\mathfrak a}$. If every $\mathfrak a_i$ is relatively prime to every $\mathfrak b_j$, then the product of the $\mathfrak a_i$ is relatively prime to the product of the $\mathfrak b_j$. In particular, if $\mathfrak b_1,\ldots,\mathfrak b_s$ are pairwise relatively prime, their least common multiple is their product.

The Chinese-remainder decomposition is then immediate. If $\mathfrak b_1,
\ldots,\mathfrak b_s$ are pairwise relatively prime and
\[
  \mathfrak v=[\mathfrak b_1,\ldots,\mathfrak b_s]=\mathfrak b_1\cdots\mathfrak b_s,
\]
then the residue-class ring modulo $\mathfrak v$ is the direct sum of the residue-class rings modulo the $\mathfrak b_i$:
\[
  R/\mathfrak v \simeq R/\mathfrak b_1\oplus\cdots\oplus R/\mathfrak b_s.
\]
The idempotents giving the direct summands are obtained by choosing, for each $i$, an element congruent to $1$ modulo $\mathfrak b_i$ and congruent to $0$ modulo the other $\mathfrak b_j$.

Conversely, a decomposition of a ring or module into a direct sum is expressed ideal-theoretically by such idempotent decompositions and by pairwise relatively prime annihilator ideals. This is the form in which Noether uses direct sums later: it is not an additional structure, but the ideal-theoretic expression of pairwise relatively prime components.

\subsection{5. Primary ideals without chain assumptions}

Noether distinguishes weak and strong primary ideals in a general ring before imposing chain conditions. An ideal $\mathfrak q$ is weakly primary if, whenever $ab\equiv0\pmod{\mathfrak q}$ and $a$ is not divisible by $\mathfrak q$, every sufficiently high power of $b$ is divisible by $\mathfrak q$. Equivalently, the zero divisors of the residue-class ring $R/\mathfrak q$ are precisely those residue classes all of whose sufficiently high powers vanish. The associated prime ideal $\mathfrak p$ consists of all elements whose powers are divisible by $\mathfrak q$.

The strong version is obtained by replacing elements by ideals. Thus $\mathfrak q$ is strongly primary if from
\[
  \mathfrak a\mathfrak b\equiv0\pmod{\mathfrak q}
\]
and from the non-divisibility of $\mathfrak a$ by $\mathfrak q$, it follows that some power of $\mathfrak b$ is divisible by $\mathfrak q$. In rings with the divisor chain condition the distinction disappears; without that hypothesis it is real. Noether records an example, due to R. Hölzer, showing that a weak primary ideal need not be strong.

The least common multiple of finitely many weak, respectively strong, primary ideals with the same associated prime ideal is again weak, respectively strong, primary with that same associated prime ideal. But a shortest representation by finitely many primary ideals belonging to different prime ideals is not primary. Indeed, if
\[
  \mathfrak m=[\mathfrak q_1,
\ldots,\mathfrak q_r]
\]
is a shortest representation with distinct associated primes $\mathfrak p_i$, one can choose elements, or ideals in the strong case, which vanish modulo all but one component and not modulo the remaining prime. Multiplying them gives a product divisible by $\mathfrak m$ while no suitable power condition required for primaryness can hold. This proves non-primaryness of such a mixed component.

The quotient of modules is now introduced. Let $T$ be an extension ring of $R$, and let $M,B,C,\ldots$ be $R$-modules in $T$. The quotient
\[
  C=M:B
\]
is the smallest module satisfying
\[
  CB\equiv0\pmod M,
\]
in the sense that if $DB\equiv0\pmod M$, then $D\equiv0\pmod C$. Equivalently, $C$ is the greatest common divisor of all modules $D$ with $DB\equiv0\pmod M$. The zero module always gives existence. From the definition one obtains the formal rules: if $M$ is a multiple of $N$, then $M:B$ is a multiple of $N:B$; if $B$ is a multiple of $C$, then $M:B$ is a divisor of $M:C$; and
\[
  M:BC=(M:B):C=(M:C):B.
\]
If the extension ring is the original ring, the module quotient becomes the ideal quotient $\mathfrak a:\mathfrak b$, with the same rules. Since $\mathfrak a\mathfrak b\equiv0\pmod{\mathfrak a}$, one also has
\[
  \mathfrak a\equiv0\pmod{\mathfrak a:\mathfrak b}.
\]

An ideal $\mathfrak b$ is called prime to $\mathfrak a$ if
\[
  \mathfrak a:\mathfrak b=\mathfrak a.
\]
This says exactly that $d\mathfrak b\equiv0\pmod{\mathfrak a}$ always implies $d\equiv0\pmod{\mathfrak a}$. From the quotient rules it follows that if $\mathfrak b$ and $\mathfrak c$ are prime to $\mathfrak a$, then both their product and their least common multiple are prime to $\mathfrak a$. If $\mathfrak b$ is prime to $\mathfrak a$ and $\mathfrak a$ is prime to $\mathfrak b$, the ideals are mutually prime. Relatively prime ideals in the sense of greatest common divisor one are mutually prime, but the quotient formulation is more flexible and is the one needed for the proofs under chain conditions.

\subsection{6. Ideal theory under the divisor chain condition}

Assume now that the commutative ring $R$ satisfies axiom I, the divisor chain condition. Fix once and for all a well-ordering of all elements of $R$. This also gives a well-ordering of the ideals. Indeed, the divisor chain condition implies that every ideal has a finite basis; ordering finite subsets quasi-lexicographically and assigning to each ideal its first possible basis well-orders the set of ideals.

\textbf{Theorem I.} Under the divisor chain condition every ideal of $R$ can be represented as a least common multiple of finitely many irreducible ideals, i.e. ideals not representable as a least common multiple of two proper divisors.

If the theorem failed for an ideal $\mathfrak m$, then $\mathfrak m$ would be reducible; writing $\mathfrak m=[\mathfrak a,\mathfrak b]$, at least one of $\mathfrak a$ or $\mathfrak b$ would again fail the theorem. Choose the first such proper divisor in the well-ordering. Repeating the construction gives a definite infinite chain of proper divisors, contradicting the divisor chain condition. This is the only part of the argument where the well-ordering is used; the statement is set-theoretic in nature and applies equally to module domains with the corresponding chain condition.

With the divisor chain condition, primary ideals in the weak and strong sense coincide, and one may simply speak of primary ideals. The relation between primary and irreducible is the key.

\textbf{Theorem II.} Under the divisor chain condition every irreducible ideal is primary. Equivalently, every non-primary ideal is reducible.

If $\mathfrak m$ is not primary, then there are elements, or equivalently ideals through the quotient formalism, giving a product divisible by $\mathfrak m$ while neither factor satisfies the required power-divisibility condition. The ideal quotient then furnishes two proper divisors of $\mathfrak m$ whose least common multiple is $\mathfrak m$ itself. Thus $\mathfrak m$ is reducible. Combining this with Theorem I gives the representation of every ideal as a least common multiple of finitely many primary ideals.


Continuing the same proof, once every irreducible ideal is known to be primary, the representation theorem takes the following explicit form. Every ideal $\mathfrak a$ admits a shortest representation
\[
  \mathfrak a=[\mathfrak q_1,\ldots,\mathfrak q_s]
\]
by primary ideals. The associated prime ideals are obtained by taking, for each $\mathfrak q_i$, the totality of elements whose powers are divisible by $\mathfrak q_i$. The primary components may still fail to be unique at this stage. What is fixed by the divisor chain condition alone is the existence of a finite primary representation and the possibility of choosing it shortest.

Noether's reason for introducing the ideal quotient now becomes visible. If $\mathfrak b$ is prime to $\mathfrak a$, the equality $\mathfrak a:\mathfrak b=\mathfrak a$ is an exact algebraic substitute for saying that $\mathfrak b$ does not meet the component defined by $\mathfrak a$. Multiplying by ideals prime to a component cannot introduce new divisibility modulo that component. Thus the quotient calculus supplies the separation of components without choosing elements or bases.

The next step, which belongs to the multiple chain condition, is the uniqueness of isolated components. Passing to the residue-class ring modulo a non-zero ideal, the multiple chain condition prevents a prime ideal from having proper divisors other than the unit ideal. Consequently the primary component corresponding to an isolated associated prime is forced: it is recovered as the quotient of the given ideal by all ideals prime to that associated prime. This is the abstract version of separating one primary component by multiplying away all the others.

When the ring has an identity, the unit ideal no longer appears as an associated prime for non-zero ideals. Hence the ambiguity just described disappears for non-zero ideals once both chain conditions hold. The primary components become unique; they are pairwise relatively prime, and their least common multiple is their product. Thus the theory has reached the form familiar from orders in number fields: an ideal is represented as a product of one-type components, and after integral closedness is added these one-type components become powers of prime ideals.

This order of proof is important. The primary decomposition, its uniqueness, and the passage from primary ideals to prime powers are not bundled together. They are separated according to the axioms that force them: divisor chains give existence, multiple chains give uniqueness of the components, the identity and absence of zero divisors remove exceptional unit phenomena and create a quotient field, and integral closedness turns primary components into prime powers. The classical Dedekind ideal theory is therefore not a single indivisible phenomenon but the conjunction of these four structural facts.


For the continuation of the paper this separation of the axioms is the controlling device. In the next sections Noether proves the converse direction: from the prime-power factorization of ideals one recovers the chain conditions and integral closedness. The proof is again conducted with ideal quotients. Fractional ideals are introduced as quotients in the quotient field, their classes form a group precisely when every non-zero ideal factors into prime powers, and Dedekind's criterion for integral closedness is recovered from the impossibility of writing a genuinely fractional element with all powers of the numerator divisible by the denominator. Thus the abstract axioms and the classical factorization theorem determine one another.

This is the point at which the paper passes from mere primary decomposition to Dedekind ideal theory. The previous sections account for the decomposition of ideals into components; the subsequent ones show why, in integrally closed domains satisfying the double chain condition, those components have the special and much stronger form of powers of prime ideals. The proof is deliberately intrinsic: it does not rely on the existence of bases for ideals in a number field, but only on divisibility, quotients, chain conditions, and integrality.


The distinction is also why the paper can include both algebraic number fields and algebraic function fields without changing the formal proof. In the number-field case the integral closure of the rational integers in a finite extension field is the standard order of all algebraic integers. In the function-field case the functional domain plays the same formal role after the passage that removes the lower-dimensional ideals. Once the axioms are verified for the base domain and shown to survive finite extensions of first kind, the subsequent ideal theory proceeds without referring to the concrete origin of the ring.

The practical mathematical point is that Noether is not merely restating Dedekind's theory abstractly. She is isolating exactly which parts of the classical proof require which finiteness or integrality assumptions. This separation lets later commutative algebra treat primary decomposition, isolated components, ideal quotients, integral closure, and unique factorization of ideals as separate structural layers rather than as one monolithic theorem inherited from algebraic number theory.

\clearpage
\section{Abstract Construction of Ideal Theory in Algebraic Number and
Function
Fields}\label{abstract-construction-of-ideal-theory-in-algebraic-number-and-function-fields}

\subsection{Continuation of §6. Ideal theory under the divisor-chain
condition}\label{continuation-of-6.-ideal-theory-under-the-divisor-chain-condition}

To obtain a representation by greatest primary components it is only
necessary to replace, wherever possible, the representation by
irreducible ideals whose existence is supplied by Theorem I - and which
by Theorem II are primary - by a shortest such representation. If the
primary ideals belonging to the same prime ideal are then collected
together, the desired representation follows from §5, 3. The primary
components may be called greatest components, since the least common
multiple of any several of them is no longer primary.

\subsection{§7. Ideal theory under the double chain
condition}\label{ideal-theory-under-the-double-chain-condition}

The simplifications, compared with the theory developed so far, rest on
the two auxiliary results stated below.

First, if the multiple-chain condition holds in a commutative ring
without zero divisors, then the ring is already a field. It is enough to
show that, for \(a\ne0\), the equation \[
  ax=b
\] always has a solution in the ring. It cannot have more than one
solution, because the ring has no zero divisors. Let \(\mathfrak a\) be
the principal ideal generated by \(a\). By hypothesis the sequence of
powers of \(\mathfrak a\) eventually stops; say \(\mathfrak a^n\) is
equal to all later powers. If \(\mathfrak b\) denotes the principal
ideal generated by \(b\), then \[
  \mathfrak a^n\mathfrak b=\mathfrak a^{n+1}\mathfrak b.
\] In particular, the element \(a^n b\) has a representation of the form
\[
  a^n b = k a^{n+1}b+n a^{n+1}b = a^{n+1}c,
\] where \(k,n\) are ordinary integers and \(c\) is again an element of
the ring. Since zero divisors do not occur, this gives \(b=ac\). Thus
the field property is proved.

The preceding argument immediately gives the following lemma. If the
multiple-chain condition holds in a commutative ring, then a prime ideal
has no proper divisor different from the zero ideal. Likewise, if only
axiom II is assumed - the multiple-chain condition modulo every ideal
different from the zero ideal - then every prime ideal different from
the zero ideal has no proper divisor different from \(0\). Indeed, the
residue class ring modulo a prime ideal is, by definition, a ring
without zero divisors; by homomorphy it too satisfies the multiple-chain
condition, and hence is a field. By the one-to-one correspondence
between divisors of the prime ideal and ideals of the residue class
ring, the prime ideal has no proper divisor different from \(0\).

\textbf{Theorem IV.} If the double chain condition holds in a
commutative ring, then in every shortest representation of an ideal the
primary components which do not belong to the unit ideal are uniquely
determined.

\textbf{Supplement.} Under axioms I and II the same assertion holds for
every ideal different from the zero ideal.

Let \[
  \mathfrak m=[\mathfrak q,\mathfrak q_1,\ldots,\mathfrak q_r]
      =[\bar{\mathfrak q},\bar{\mathfrak q}_1,\ldots,\bar{\mathfrak q}_s]
\] be two shortest representations of \(\mathfrak m\) by greatest
primary components, and let \[
  0,\mathfrak p_1,\ldots,\mathfrak p_r
  \quad\text{and}\quad
  0,\bar{\mathfrak p}_1,\ldots,\bar{\mathfrak p}_s
\] be the corresponding prime ideals. The component \(\mathfrak q\),
respectively \(\bar{\mathfrak q}\), is omitted if no primary ideal
belonging to \(0\) occurs. Since \[
  0\cdot \mathfrak p_1\cdots\mathfrak p_r\equiv0\pmod{\mathfrak m},
\] each \(\bar{\mathfrak p}_j\) must be contained in one of the
\(\mathfrak p_i\), and hence, by the lemma, must be equal to it.
Conversely each \(\mathfrak p_i\) must agree with one of the
\(\bar{\mathfrak p}_j\). Thus the prime ideals different from \(0\)
coincide.

Next, for instance, \[
  \mathfrak q\mathfrak q_2
\cdots \mathfrak q_r\equiv0\pmod{\mathfrak q_1},
\] and therefore \(\bar{\mathfrak q}_1\equiv0\pmod{\mathfrak q_1}\),
since the other components are not divisible by \(\mathfrak p_1\) and
are therefore prime to \(\mathfrak q_1\). The reverse divisibility is
obtained in the same way; hence the primary component belonging to each
non-zero prime ideal is unique. If a component belonging to \(0\)
occurs, then, because the representation is shortest, the corresponding
component occurs in the other representation as well; this gives the
remaining uniqueness. The same proof also shows that every
representation containing no primary component belonging to \(0\) is
itself shortest.

If, in addition to the double chain condition, the existence of a unit
element is assumed, Theorem IV sharpens as follows.

\textbf{Theorem V.} In a commutative ring with unit element satisfying
the double chain condition, every ideal can be represented uniquely as a
product of finitely many pairwise relatively prime primary ideals.

\textbf{Supplement.} Under axioms I, II, III, Theorem V holds for every
ideal different from the zero ideal.

Because the unit element exists, \(0^0=\mathfrak o\); hence every
primary ideal belonging to \(0\) is equal to \(0\). Such a component
cannot occur in a shortest representation of a non-zero ideal. The
primary components are therefore uniquely determined by Theorem IV. At
the same time every representation becomes shortest after omitting the
zero component. Distinct prime ideals are, by the lemma, relatively
prime; hence by the calculation rules of §4, 4, the same is true of the
corresponding primary components. The least common multiple is therefore
equal to their product.

The same hypotheses also give the following lemma. In a commutative ring
with unit element and double chain condition, the prime ideals which are
prime to an ideal \(\mathfrak m\) are precisely all prime ideals, except
the unit ideal, which are not contained in \(\mathfrak m\). The notions
``prime to'' and ``relatively prime'' coincide.

Under axioms I, II, III the ideal \(\mathfrak m\) is to be assumed
different from the zero ideal.

If \(\mathfrak p\) is not contained in \(\mathfrak m\), then it is
different from all prime ideals belonging to \(\mathfrak m\). Hence, by
the lemma above, it is divisible by none of these prime ideals, and so
is prime to each primary component and therefore to \(\mathfrak m\). It
is also relatively prime to each component and hence relatively prime to
\(\mathfrak m\).

If \(\mathfrak p\ne0\) is contained in \(\mathfrak m\), then it must
coincide with one of the prime ideals belonging to \(\mathfrak m\).
Suppose it belongs to the component \(\mathfrak q=\mathfrak q_1\). Then
\(\mathfrak q:\mathfrak p\) is a proper divisor of \(\mathfrak q\),
since \(\mathfrak p^{e-1}\equiv0\pmod{\mathfrak q:\mathfrak p}\) but not
modulo \(\mathfrak q\). At the same time \(\mathfrak q:\mathfrak p\), as
a divisor of \(\mathfrak p^{e-1}\), is divisible only by one non-zero
prime ideal and is therefore primary. By unique factorization the
product \[
  (\mathfrak q:\mathfrak p)\mathfrak q_2\cdots\mathfrak q_r,
\] which is divisible by \(\mathfrak m:\mathfrak p\), is a proper
divisor of \(\mathfrak m\). Thus \(\mathfrak p\) is not prime to
\(\mathfrak m\).

Consequently an ideal \(\mathfrak t\ne\mathfrak p\) is prime to
\(\mathfrak m\) precisely when none of the prime ideals belonging to
\(\mathfrak t\) is contained in \(\mathfrak m\); in that case
\(\mathfrak t\) is also relatively prime to \(\mathfrak m\). Since,
conversely, two relatively prime ideals are always mutually prime (§5,
5), the two notions coincide under the hypotheses above.

\subsection{§8. Ideal theory under integral closedness in the quotient
field}\label{ideal-theory-under-integral-closedness-in-the-quotient-field}

The ideal theory developed up to this point is now to be sharpened to
the usual one by adjoining the last two axioms.

First lemma. Let \(R\) be a commutative ring without zero divisors and
with unit element, and assume the divisor-chain condition in \(R\)
(axioms I, III, IV). Then from \[
  \mathfrak a=\mathfrak a\mathfrak b,
  \qquad \mathfrak a\ne(0),
\] it always follows that \(\mathfrak b=\mathfrak o\). Thus \[
  \mathfrak c\ne\mathfrak c\mathfrak t
\] for every ideal different from the zero and unit ideals.

The proof is the same as in the lemma of §1, 3, since
\(\mathfrak a\mathfrak b\) may be regarded as a finite module with
respect to \(\mathfrak b\). If \(\alpha_1,\ldots,\alpha_n\) is a basis
of the ideal \(\mathfrak a\), which exists by axiom I, then from
\(\mathfrak a=\mathfrak a\mathfrak b\) one obtains a system of equations
\[
  \alpha_i=b_{i1}\alpha_1+\cdots+b_{in}\alpha_n,
  \qquad b_{ij}\in\mathfrak b.
\] Since \(R\) has no zero divisors, the determinant \(|b_{ij}-e_{ij}|\)
is zero, where \(e_{ij}=0\) for \(i\ne j\) and \(e_{ii}=1\). The
equation obtained from this determinant has the form \[
  1+b_1+\cdots+b_n=0,
  \qquad b_i\in\mathfrak b,
\] and hence \(1\in\mathfrak b\), so \(\mathfrak b=\mathfrak o\).

It remains to show that, once integral closedness in the quotient field
is added to axioms I-IV (axiom V), every primary ideal different from
zero is a power of its associated prime ideal. The uniqueness of the
resulting representation \[
  \mathfrak m=\mathfrak p_1^{e_1}\cdots\mathfrak p_r^{e_r}
\] for all ideals different from zero has already been proved by Theorem
V and the preceding lemma; at the same time those prime ideals have no
proper divisors except the unit ideal. The proof will first be given for
exponent two, using Dedekind's second consequence from §1, 4, and then
in general by complete induction.

\textbf{Lemma.} Under axioms I-V there are no primary ideals of exponent
two except \(\mathfrak q=\mathfrak p^2\).

It must be shown that from \[
  \mathfrak p^2\equiv0\pmod{\mathfrak q},\qquad
  \mathfrak q\equiv0\pmod{\mathfrak p},\qquad
  \mathfrak q\not\equiv0\pmod{\mathfrak p^2}
\] one necessarily gets \(\mathfrak q=\mathfrak p^2\). Choose \[
  \mathfrak c\equiv0\pmod{\mathfrak q},\qquad
  \mathfrak c\equiv0\pmod{\mathfrak p},\qquad
  \mathfrak c
ot\equiv0\pmod{\mathfrak p^2}.
\] Then by the lemma of §7, 3, the ideal \(\mathfrak c:\mathfrak p\) is
a proper divisor of \(\mathfrak c\). Thus there is an element \(b\) such
that \[
  b\mathfrak p\equiv0\pmod{\mathfrak c}
      \equiv0\pmod{\mathfrak q},
  \qquad b
ot\equiv0\pmod{\mathfrak c}.
\] Consequently \(\gamma=b/c\) is not integral: it belongs to the
quotient field but not to \(R\).

We shall prove \(b\equiv0\pmod{\mathfrak p}\); then, since
\(\mathfrak q\) is primary and belongs to \(\mathfrak p\), it follows
that \(\mathfrak p\equiv0\pmod{\mathfrak q}\) and hence
\(\mathfrak q=\mathfrak p^2\). By Dedekind's second consequence there
exist elements \(m,n\) of \(R\) with \[
  \gamma=b/c=m/n
\] such that \(m^2/n\) is still not integral. Equivalently, \[
  bn=mc,
  \qquad m\not\equiv0\pmod{nR},
  \qquad m^2
ot\equiv0\pmod{nR}.
\] Multiplying \(b\mathfrak p\equiv0\pmod{\mathfrak c}\) by \(n\) gives
\[
  mb\mathfrak p\equiv0\pmod{n\mathfrak c},
\] and hence \(m\mathfrak p\equiv0\pmod{nR}\), because the ring has no
zero divisors. From \(m^2\equiv0\pmod{\mathfrak p nR}\) and
\(n\not\equiv0\pmod{\mathfrak p}\) it follows that
\(m\equiv0\pmod{\mathfrak p}\); here one again uses that
\(nR:\mathfrak p\) is a proper divisor of \(nR\). The four relations \[
  m\equiv0\pmod{\mathfrak p},
  \quad n\not\equiv0\pmod{\mathfrak p},
  \quad c\equiv0\pmod{\mathfrak p},
  \quad c
ot\equiv0\pmod{\mathfrak p^2},
  \quad bn=mc
\] then force \(b\equiv0\pmod{\mathfrak p}\). Indeed, since
\(\mathfrak p^2\) is primary, first \(mc\equiv0\pmod{\mathfrak p^2}\),
and from \(bn=mc\) one obtains \(b\equiv0\pmod{\mathfrak p}\). This
proves the lemma.

A second lemma removes the restriction on the exponent.

\textbf{Lemma.} Under axioms I-V there are no primary ideals of exponent
\(\rho\) except \(\mathfrak q=\mathfrak p^{\rho}\).

This follows from the preceding lemma without further use of the axioms.
First one proves: \[
  \text{if }\mathfrak c\equiv0\pmod{\mathfrak p},\quad
  \mathfrak c
ot\equiv0\pmod{\mathfrak p^2},
  \quad\text{then}\quad
  \mathfrak p^{\rho}=[\mathfrak c^{\rho},\mathfrak p^{\rho+1}]
\] for every \(\rho\). For \(\rho=1\) this is just the exponent-two
lemma. If \[
  \mathfrak p^{\rho-1}=[\mathfrak c^{\rho-1},\mathfrak p^{\rho}],
\] then multiplication by \(\mathfrak p\) and the distributive law give
\[
  \mathfrak p^{\rho}=
  [\mathfrak c^{\rho},\mathfrak c^{\rho-1}\mathfrak p^{\rho+1},\mathfrak p^{\rho+1}]
  =[
    \mathfrak c^{\rho},\mathfrak p^{\rho+1}].
\]

The assertion can be put in the following equivalent form. If \[
  \mathfrak q\equiv0\pmod{\mathfrak p^{\rho}},
  \qquad \mathfrak q
ot\equiv0\pmod{\mathfrak p^{\rho+1}},
\] and if \[
  \mathfrak p^{\rho+k}\equiv0\pmod{\mathfrak q},\qquad k\ge1,
\] then already \[
  \mathfrak p^{\rho+k-1}\equiv0\pmod{\mathfrak q}.
\] Indeed, for every primary ideal there is an exponent \(\rho>1\) with
\(\mathfrak q\equiv0\pmod{\mathfrak p^{\rho}}\) but not modulo
\(\mathfrak p^{\rho+1}\); and \(\mathfrak q
ot\equiv0\pmod{\mathfrak p^{\rho+1}}\) follows from \(\mathfrak p^{\rho}
e\mathfrak p^{\rho+1}\) by the first lemma of this section. Repeated
application of the second formulation gives
\(\mathfrak p^{\rho}\equiv0\pmod{\mathfrak q}\) and hence
\(\mathfrak q=\mathfrak p^{\rho}\).

From the hypotheses of the second formulation, and from the above
representation of \(\mathfrak p^{\rho}\), one obtains for every element
\(g\) of \(\mathfrak q\) a representation \[
  g=b c^{\rho} \pmod{\mathfrak p^{\rho+1}}
\] with either \(b\equiv0\pmod{\mathfrak p}\) or \[
  b c^{\rho}\equiv0\pmod{[\mathfrak q,\mathfrak p^{\rho+k}]}.
\] Since \([\mathfrak q,\mathfrak p^{\rho+1}]\) is primary, it follows
in the latter case that \[
  c^{\rho}\equiv0\pmod{[\mathfrak q,\mathfrak p^{\rho+1}]},
\] or else \[
  c^{\rho+k-1}\equiv0\pmod{[\mathfrak q,\mathfrak p^{\rho+k}]}.
\] In either case one obtains \[
  \mathfrak p^{\rho+k-1}\equiv0\pmod{\mathfrak q}.
\] This proves the general lemma.

Combining the results gives:

\textbf{Theorem VI.} If axioms I-V hold in a commutative ring, then
every ideal different from the zero and unit ideals is uniquely
representable as a product of finitely many prime ideals different from
the zero and unit ideals. These prime ideals have no proper divisors
except the unit ideal.

\subsection{§9. The axioms as consequences of the assumed
decomposition}\label{the-axioms-as-consequences-of-the-assumed-decomposition}

We now prove conversely that the usual ideal factorization implies the
axioms. It is necessary to assume the factorization in the following
sharpened form.

\textbf{Assumption.} In the commutative ring \(R\), every ideal not
generated by the zero or unit element is uniquely representable as a
product of powers of prime ideals. These prime ideals are simple ideals
not generated by the unit element; that is, they have no proper divisor
different from \(0\). Conversely, all simple ideals are prime ideals.
The uniqueness is to be understood sharply: for every ideal not
generated by the zero or unit element, \[
  \mathfrak a^{\rho}\ne\mathfrak a^{\rho+1}.
\]

The existence of a unit element is not presupposed. If no unit element
exists, the condition ``not generated by the unit element'' imposes no
restriction.

First we prove axiom III, the existence of a unit element. Suppose axiom
III fails. We shall show that \(R^2\) is then a simple ideal without
being a prime ideal, contradicting the assumption. By the uniqueness
assumption \(R\ne R^2\); hence \[
  R\equiv0\pmod{R^2},
  \qquad R^2
ot\equiv0\pmod{R^2},
\] and so \(R^2\) is not a prime ideal. But \(R^2\) is simple. Indeed,
it can be divisible by no non-zero prime ideal, and therefore, by the
assumed product representation, it has no divisors except \(0\) and
\(R^2\). This proves the existence of a unit.

Next comes axiom IV, absence of zero divisors. If \(a\) is a zero
divisor, so that \[
  ab=0,
  \qquad a\ne0,
  \qquad b\ne0,
\] then multiplication of the representations of the principal ideals
generated by \(a\) and \(b\) gives \[
  (0)=\mathfrak p_1^{\rho_1}\cdots\mathfrak p_r^{\rho_r},
\] where the \(\mathfrak p_i\) are different from the zero and unit
ideals. The representation may be assumed shortest in the sense that
omission of any factor gives a proper divisor of zero; this need not be
automatic, since no uniqueness assumption has been made for the zero
ideal. If the representation contains only one prime ideal, then \[
  (0)=\mathfrak p^{\rho}=\mathfrak p^{\rho+1},
\] contrary to the sharpened uniqueness condition. If it contains more
than one prime ideal, then because the unit element exists
\(\mathfrak p_1\) is relatively prime to all the other prime ideals;
hence again the sharpened uniqueness condition is contradicted. Thus
zero divisors do not occur.

The chain conditions, axioms I and II, follow as well. The assumptions
imply that for every non-zero ideal, divisibility is equivalent to the
occurrence of a product representation. More precisely, from \[
  \mathfrak m\equiv0\pmod{\mathfrak a},
  \qquad \mathfrak a
e0,
\] it follows that \[
  \mathfrak m=\mathfrak a\mathfrak b,
  \qquad \mathfrak b\equiv0\pmod{\mathfrak m}.
\] If \[
  \mathfrak m=\mathfrak p_1^{\rho_1}\cdots\mathfrak p_r^{\rho_r},
  \qquad \mathfrak a=\mathfrak q_1^{\sigma_1}\cdots\mathfrak q_s^{\sigma_s},
\] then each \(\mathfrak q_j\) must be identical with one of the
\(\mathfrak p_i\), and the corresponding exponent must not exceed
\(\rho_i\). Hence the divisors of \(\mathfrak m\) are only the finitely
many ideals corresponding to the combinations \[
  0\le \sigma_i\le \rho_i.
\] The residue class ring modulo \(\mathfrak m\) therefore satisfies the
double chain condition. This proves axioms I and II; the divisor-chain
condition holds in \(R\) itself, since every proper divisor of the zero
ideal is different from the zero ideal.

Before proving axiom V, integral closedness in the quotient field, we
derive the usual consequences of ideal factorization.

First, the residue class ring modulo every non-zero ideal is a principal
ideal ring. The ideal may be assumed different from the unit ideal,
since for the residue ring consisting only of the zero element the
assertion is clear. If the ideal is primary, say
\(\mathfrak q=\mathfrak p^{\rho}\), and if \[
  \mathfrak c\equiv0\pmod{\mathfrak p},
  \qquad \mathfrak c\not\equiv0\pmod{\mathfrak p^2},
\] then the argument of §8 gives \[
  \mathfrak p^{\rho}=[\mathfrak c^{\rho},\mathfrak p^{\rho+1}]
\] for every \(\rho\). Thus in the residue class ring modulo a primary
ideal every ideal is principal. In general, let \[
  \mathfrak m=\mathfrak q_1\cdots\mathfrak q_r
\] be the decomposition into pairwise relatively prime primary ideals,
and let \[
  R/\mathfrak m=\mathfrak a_1+\cdots+\mathfrak a_r
\] be the corresponding direct-sum decomposition. Each \(\mathfrak a_i\)
is isomorphic to \(R/\mathfrak q_i\), and every ideal in
\(\mathfrak a_i\) is principal. If \(\mathfrak c\) is an arbitrary ideal
of the residue class ring, then \(\mathfrak a_i\mathfrak c\) is a
principal ideal, say \(\mathfrak a_i c_i\). Setting \[
  c=c_1+\cdots+c_r
\] one obtains \[
  \mathfrak c=Rc.
\] The ideal is therefore principal.

The principal-ideal-ring result may also be stated in the following
form. If \(\mathfrak c\) is any ideal, then it can be transformed into a
principal ideal by multiplication with an ideal relatively prime to a
prescribed ideal \(\mathfrak b\). Indeed, set
\(\mathfrak m=\mathfrak b\mathfrak c\). Then \(\mathfrak c\) becomes
principal modulo \(\mathfrak m\), i.e. \[
  \mathfrak c=(cR,\mathfrak m)=cR+c\mathfrak b=c(R,\mathfrak b),
\] and hence \(\mathfrak c\mathfrak a=cR\) with
\((\mathfrak a,\mathfrak b)=R\).

We now pass to fractional ideals. A fractional ideal is by definition a
finite \(R\)-module in the quotient field \(K\).

\textbf{Theorem.} The non-zero finite \(R\)-modules in \(K\) form an
abelian group under multiplication.

The product of two finite \(R\)-modules is again finite; multiplication
is associative and commutative; and, since \(R=R\cdot R\), the unit
ideal is the identity element. It remains only to show that the equation
\[
  \mathfrak A X=R
\] has a solution in the system of finite \(R\)-modules.

We use the following preliminary remark. If the equation \[
  \mathfrak A T=\mathfrak B
\] has one and only one solution, then \[
  T=\mathfrak B:\mathfrak A.
\] For \(T\equiv0\pmod{\mathfrak B:\mathfrak A}\), and conversely \[
  \mathfrak A(\mathfrak B:\mathfrak A)=\mathfrak B.
\]

If first \(\mathfrak A\) is a principal module,
\(\mathfrak A=(\alpha)\), then \(X=(1/\alpha)\) solves
\(\mathfrak A X=R\); it is again a finite \(R\)-module. It follows that
every finite \(R\)-module is the quotient of two ideals of \(R\), which
justifies the term fractional ideal. Indeed, if \(\tau_1,\ldots,\tau_n\)
is a module basis of \(T\), write \(\tau_i=t_i/\alpha\) and let
\(\mathfrak t\) be the ideal generated by the \(t_i\) in \(R\). Then \[
  \alpha T=\mathfrak t,
\] so by the preceding remark \[
  T=\mathfrak t:(\alpha).
\]

The inverse of a general finite module can now be written down. Suppose
\[
  \mathfrak A=\mathfrak a:\mathfrak c,
\] and choose \(\mathfrak b\) so that \(\mathfrak a\mathfrak b\) is a
principal ideal, say \((\alpha)\). Then \[
  \mathfrak A\mathfrak c=\mathfrak a,
  \qquad \mathfrak A\mathfrak b\mathfrak c=(\alpha),
\] and therefore \[
  \mathfrak A\, (\mathfrak b\mathfrak c:(\alpha))=R.
\] The module \(\mathfrak b\mathfrak c:(\alpha)\) is, as the quotient of
an ideal by a principal ideal, again a finite \(R\)-module. The group
property is proved. It also follows that the quotient of any two ideals
\(\mathfrak a:\mathfrak b\), as the solution of
\(\mathfrak bX=\mathfrak a\), is a finite \(R\)-module.

The group property has the following consequences.

First, cancellation is possible: \[
  T=AM:BM \quad\Longleftrightarrow\quad T=A:B.
\] Indeed, from \(TBM=TAM\) follows \(TB=TA\), and conversely.

Second, every fractional ideal admits a representation \[
  \mathfrak C=\mathfrak a:\mathfrak b
\] where \(\mathfrak a\) and \(\mathfrak b\) are relatively prime; under
this condition \(\mathfrak a\) and \(\mathfrak b\) are uniquely
determined. Existence follows from the previous paragraph and
cancellation. If \[
  \mathfrak C=\mathfrak a:\mathfrak b=\bar{\mathfrak a}:\bar{\mathfrak b},
\] then \(\mathfrak a\bar{\mathfrak b}=\mathfrak b\bar{\mathfrak a}\);
uniqueness follows when both pairs are assumed relatively prime.

Third, every principal module generated by a non-integral element - an
element belonging to the quotient field but not to \(R\) - has a
representation as a quotient of principal ideals such that no power of
the numerator is divisible by the denominator. Let \[
  \eta R=\mathfrak b:\mathfrak c
\] be the reduced representation, so \((\mathfrak b,\mathfrak c)=R\).
Choose, by the principal-ideal statement above, \(\mathfrak a\) such
that \[
  \mathfrak b\mathfrak a=(b),
  \qquad (\mathfrak a,\mathfrak c)=R.
\] Then \[
  \eta R=\mathfrak a\mathfrak b:\mathfrak a\mathfrak c=(b):(c).
\] No power of \(b\) is divisible by \(c\), because
\((\mathfrak a,\mathfrak c)=R\) and \((\mathfrak b,\mathfrak c)=R\).

Now axiom V follows. Every non-integral element \(\eta\) of the quotient
field admits a quotient representation \[
  \eta=a/c
\] such that no power of the numerator is divisible by the denominator.
Indeed, from \(\eta R=(b):(c)\) it follows that \(\eta\) differs from
\(b/c\) only by a unit of \(R\). Such an element can satisfy no equation
characterizing integrality over \(R\); for from \[
  \eta^n+r_1\eta^{n-1}+\cdots+r_n=0
\] one obtains \[
  a^n\equiv0
yal\pmod{cR},
\] contrary to the construction. Thus \(R\) is integrally closed in its
quotient field.

A useful corollary is the following. If axioms I-IV hold in a
commutative ring and every primary ideal is irreducible, then axiom V
also holds. Indeed, by axioms I-III every power \(\mathfrak p^{\rho}\)
of a prime ideal is primary, in particular \(\mathfrak p^2\) is primary.
By assumption \(\mathfrak p^2\) is irreducible. It remains to prove a
representation \[
  \mathfrak p=(cR,\mathfrak p^2),
\] for then §8, 3 shows that every primary ideal is a power of its prime
ideal, and this, together with the remaining hypotheses, implies
integral closedness by the preceding paragraph.

By the divisor-chain condition one can write \[
  \mathfrak p=(c_1R,\ldots,c_kR),
\] where the \(c_i\) are chosen linearly independent over the residue
class field modulo \(\mathfrak p\). If \(k>1\), this independence gives
a decomposition \[
  \mathfrak p^2=[\mathfrak q_1,
\mathfrak q_2],
\] where \[
  \mathfrak q_1=(c_1R,\mathfrak p^2),
  \qquad
  \mathfrak q_2=(c_2R,
\ldots,c_kR,\mathfrak p^2).
\] Both \(\mathfrak q_1\) and \(\mathfrak q_2\) are proper divisors of
\(\mathfrak p^2\), contradicting irreducibility. Therefore \(k=1\) and
the required representation follows.

Conversely, under axioms I-V every primary ideal is irreducible. A power
\(\mathfrak p^{\rho}\) of a prime ideal cannot be written as the least
common multiple of proper divisors of the form \(\mathfrak p^{\alpha}\)
and \(\mathfrak p^{\beta}\).

If zero divisors are allowed in the ring, axioms I-III together with
integral closedness in the quotient ring do not imply that every primary
ideal is irreducible. Let \(T\) be a ring in which all axioms I-IV
except integral closedness hold; such rings exist, for example finite
orders different from the full ring of integral elements in a finite
extension of the quotient field of a ring satisfying I-V. By the
preceding corollary there is in \(T\) a reducible primary ideal
\(\mathfrak q\). Let \(\mathfrak p^{\rho}\) be a proper multiple of
\(\mathfrak q\), and let \[
  R=T/\mathfrak p^{\rho}
\] be the residue class ring. Then the ideal corresponding to
\(\mathfrak q\) becomes in \(R\) a non-zero reducible primary ideal. The
ring \(R\) satisfies axioms I-III and is nevertheless integrally closed
in its quotient ring. For every element of \(T\) not divisible by
\(\mathfrak p\) becomes relatively prime to \(\mathfrak p^{\rho}\), so
in \(R\) every regular element is a unit; hence \(R\) coincides with its
quotient ring.

\subsection{§10. Double chain condition and composition
series}\label{double-chain-condition-and-composition-series}

We shall now show that, for arbitrary well-ordered module domains (§2),
the validity of the double chain condition - every divisor chain and
every multiple chain of modules terminates - is equivalent to the
existence of a composition series. Specializing the module domain to a
commutative ring then gives the corresponding facts for the system of
all ideals of the ring. In that case ``module isomorphism'' below is to
be replaced throughout by ``ring isomorphism''.

A module domain \(G\) is called \textbf{simple} if there are in \(G\) no
\(R\)-modules other than \(G\) itself and the zero module \(0\). A
module \(A\) is called \textbf{simple in} \(G\) if \(G/A\) is simple.

A multiple chain \[
  G,U_1,
\ldots,U_r,0
\] is called a composition series of \(G\) of length \(r\) if all
modules in the chain are distinct and each module is simple in the
preceding one. Thus the residue modules, or quotient groups, \[
  U_{i-1}/U_i
\] are simple module domains, and similarly \(G/U_1\) and \(U_r/0\). By
forming the residue module \(G/A\) the case of a composition series from
\(G\) down to a module \(A\ne0\) is reduced to the absolute case.

Assume first that the double chain condition holds in \(G\). Then there
exists a composition series in \(G\), provided \(G\ne0\). From the
assumed well-ordering of \(G\) and the divisor-chain condition, a
well-ordering of the system of all modules is obtained by the
quasi-lexicographic arrangement used in §6. By this well-ordering and
the divisor-chain condition, at least one module simple in \(G\) exists.
Let \(G_1\) be the first proper divisor of \(G\) in the well-ordering,
and, in general, let \(G_i\) be the first proper divisor of \(G_{i-1}\).
The well-defined chain \[
  G,G_1,G_2,
\ldots
\] terminates by the divisor-chain condition and so necessarily reaches
a module \(A\) simple in \(G\). If \(A\ne0\), the same procedure gives a
module simple in \(A\), and so on. We obtain a well-defined multiple
chain which, by the multiple-chain condition, must terminate. It is
therefore a composition series of \(G\).

Conversely, suppose that a composition series exists in \(G\). Then the
double chain condition holds in \(G\). The proof proceeds by complete
induction, and no well-ordering of \(G\) is needed for this direction.
During the induction one also proves the Jordan-Hölder theorem under the
sole assumption that a composition series exists.

If \(G\) is simple, so that the only composition series is \(G,0\), the
following assertions are immediate:

\begin{enumerate}
\def\labelenumi{(\alph{enumi})}
\item
  Through every module simple in \(G\) one can draw a composition
  series.
\item
  Jordan-Hölder theorem: every composition series has the same length
  and the same system of quotient groups, up to order and isomorphism.
\item
  Through every module of \(G\) different from \(0\) one can draw a
  composition series.
\item
  The multiple-chain condition holds.
\item
  The divisor-chain condition holds.
\end{enumerate}

Assume, then, that the assertions (a)-(e) have been proved for every
module domain possessing a composition series of length \(r<r+1\). We
prove them for a module domain \(G\) with a composition series \[
  G,U,U_1,
\ldots,U_r,0
\] of length \(r+1\).

First, if there is no module simple in \(G\) other than \(U\), there is
nothing to prove. Let \(B\) be a module simple in \(G\) and different
from \(U\). Since \(U\) and \(B\) are both simple in \(G\) and are
distinct, necessarily \[
  (U,B)=G.
\] Put \[
  M=[U,B].
\] The second isomorphism theorem (§4, 2) gives \[
  G/B \simeq U/M,
  \qquad
  G/U \simeq B/M.
\] Thus \(M\) is simple in both \(U\) and \(B\). Since \(U\) has a
composition series of length \(r<r+1\), the induction hypotheses (a) and
(d) give a composition series in \(U\) passing through \(M\), say \[
  U,M,M_1,
\ldots,M_{r-1},0.
\] Then \[
  G,B,M,M_1,
\ldots,M_{r-1},0
\] is a composition series of \(G\) passing through \(B\).

To prove the Jordan-Hölder theorem, let \[
  G,A,A_1,
\ldots,A_s,0
  \quad\text{and}\quad
  G,B,B_1,
\ldots,B_t,0
\] be two composition series of \(G\). If \(A=B\), the result follows
immediately from the induction hypothesis for length \(r<r+1\). If
\(A\ne B\), comparison with the series constructed above, \[
  G,A,M,M_1,
\ldots,M_{r-1},0
  \quad\text{and}\quad
  G,B,M,M_1,
\ldots,M_{r-1},0,
\] shows the following. By the induction hypothesis the systems of
quotient groups of the first given series and of the first constructed
series agree; likewise for the second given series and the second
constructed series, since the latter passes through \(B\) and has length
\(r+1\). Hence the lengths are equal. Moreover, the isomorphism noted
above gives the required agreement between the quotients arising in the
two constructed series. Therefore the two original composition series
have the same quotient groups up to order and isomorphism. This proves
the Jordan-Hölder theorem.

We next need a preliminary observation. Let \(A,B\) be modules in \(G\)
and let \(B\) be simple in \(A\). Then the modules \((B,H)\) and
\((A,H)\) are either equal, or \((B,H)\) is simple in \((A,H)\). Indeed,
by the second isomorphism theorem, \[
  (A,B,H)/(B,H) \simeq A/[A,(B,H)].
\] Since \(B=0(A)\), this may also be written with \[
  C=[A,(B,H)].
\] Now \(C\) lies between \(A\) and \(B\), and hence, by the simplicity
assumption, is equal either to \(A\) or to \(B\). The assertion follows.

Let \[
  G,U,U_1,
\ldots,U_r,0
\] again be a composition series of \(G\), and let \(H\) be any module
of \(G\) different from \(0\). Form the series \[
  G,(U,H),(U_1,H),
\ldots,(U_r,H),H.
\] By the preliminary observation, the distinct modules among these form
a composition series from \(G\) down to \(H\). Therefore \(G/H\) has a
composition series, and hence, by the induction hypothesis, one can draw
a composition series through \(H\) in \(G/H\). Adding \(H\) gives a
composition series of \(G\) through \(H\). Repetition shows moreover
that such a series may be taken as an extension of the one just
constructed from \(G\) down to \(H\).

The multiple-chain condition is now immediate. Suppose again that \(G\)
has a composition series of length \(r+1\), and let \[
  G,B,B_1,
\ldots,B_n,
\ldots
\] be a multiple chain, each module being a proper multiple of the
preceding one. There is no loss in beginning the chain with \(G\), since
\(G\) may always be adjoined as first term. By the preceding paragraph
there is a composition series through \(B\) in \(G\), and therefore
\(B\) has a composition series of shorter length. By the induction
hypothesis the chain \[
  B,B_1,
\ldots,B_n,
\ldots
\] terminates; hence the chain with \(G\) adjoined terminates as well.

Finally, the divisor-chain condition is proved similarly. Let \[
  G,G_1,G_2,
\ldots,G_n,
\ldots
\] be a divisor chain in which every module is a proper divisor of the
preceding one. Again it is no restriction that the chain begins with
\(G\). By the preceding paragraph there is a composition series through
\(G_1\) in \(G\); hence \(G/G_1\) has a composition series of shorter
length. By the induction hypothesis the chain \[
  G/G_1,
  \quad G_2/G_1,
  \quad \ldots
\] terminates in \(G/G_1\). By the one-to-one correspondence supplied by
the first isomorphism theorem (§4, 2), the original chain beginning with
\(G_1\) also terminates; therefore the chain with \(G\) adjoined
terminates. The equivalence of the two assumptions - existence of a
composition series and the double chain condition - is proved.

\newpage

\section{The Discriminant Theorem for the Orders of an Algebraic Number
or Function
Field}\label{the-discriminant-theorem-for-the-orders-of-an-algebraic-number-or-function-field}

The Dedekind discriminant theorem says, as is well known, that a prime
number divides the field discriminant if and only if it is divisible by
at least the square of a prime ideal.

I show here that this theorem remains valid for all finite orders of a
finite algebraic number field - that is, for all rings of algebraic
integers which contain the ring of rational integers - in the following
form: a prime number \(p\) divides the discriminant of an order if and
only if, in the ideal decomposition valid in that order of the principal
ideal \((p)\) into greatest primary components, at least one proper
primary ideal occurs, that is, a primary ideal which is not itself a
prime ideal. Since in the principal order - the system of all algebraic
integers of the field - there are no primary ideals other than powers of
prime ideals, Dedekind's theorem follows as a special case.

The proof is obtained by passing from the order modulo \((p)\) to the
residue class ring. The assertion is thereby reduced to the
decomposition theory of rings of finite rank over a field contained in
the ring. Such a ring is a commutative system of hypercomplex quantities
with principal unit and with coefficients from that field. In the
present case the coefficient field is the residue class field of the
rational integers modulo \(p\), and the decomposition of \((p)\)
corresponds in the residue class ring to the decomposition of the zero
ideal.

Call such a finite-rank ring \textbf{completely reducible} if its zero
ideal can be represented as the least common multiple of finitely many
prime ideals. We are then concerned with criteria for complete
reducibility. First, if the coefficient field is perfect, the ring is
completely reducible if and only if the extension ring obtained by
passing to an algebraically closed coefficient field is completely
reducible. For this extension ring the non-vanishing of the discriminant
is an especially simple necessary and sufficient criterion of complete
reducibility. Since the discriminant of the order modulo every prime
number passes into the discriminant of the residue class ring, the
discriminant theorem follows.

For function fields in place of number fields the order must contain the
ground domain of all integral rational functions; for algebraic
functions in several indeterminates or for integral algebraic functions
it must contain the ground domain of all integral functionals with
integral-rational coefficients. In this situation the imperfect
coefficient-field case may occur in the residue class ring, for example
when, in the integral functional domain case, one passes to the residue
class ring modulo a prime number. With an imperfect coefficient field,
prime ideals of the first and second kind must be distinguished
according as the residue class field modulo the prime ideal is an
extension field of the first or second kind. Correspondingly, complete
reducibility of the first and second kind must be distinguished. The
criteria just stated concern only complete reducibility of the first
kind, which alone occurs in the perfect case.

In its most general form, then, the discriminant theorem says: a prime
element of the ground domain divides the discriminant of an order if and
only if in the ideal decomposition of \((p)\) either at least one proper
primary component or a prime ideal of the second kind occurs.

For the principal order this fact was first observed in examples by
Kronecker, after the earlier Festschrift statement had still contained,
mistakenly, the form valid for the principal order of a number field,
though without a complete proof. For the principal order Ostrowski gave
a proof of this discriminant theorem and further ramification theorems
connected with it, along with a detailed account of the literature. The
discriminant theorem for all orders of relative fields follows directly
by passing to the quotient ring with respect to each prime ideal of the
ground domain. The idea of this passage, known in special cases, has
been emphasized in principle by W. Krull; a systematic theory of
quotient rings in a more general context is given by H. Grell.

The uniform treatment here of all orders, and of both perfect and
imperfect coefficient fields, rests not only on ideal theory but also on
taking from the start the theory of rings of finite rank, and hence the
general concept of finite extension, as fundamental, rather than the
concept of an extension generated by a primitive element. The module
basis of an order, together with the relations among its elements, takes
the place of the powers of a primitive element and its defining
equation, or the place of the powers of a fundamental form and its
fundamental equation.

Such a module basis passes modulo every prime element of the ground
domain into a basis of the residue class ring. This need not be true,
even for the principal order, of a basis consisting of powers of a
primitive element; nor of a basis consisting of powers of the
fundamental form, once integral algebraic functions of several
indeterminates are involved. Equivalently, the residue class ring of the
polynomial ring in one indeterminate, with coefficients in the ground
domain, modulo an arbitrary defining equation - even the fundamental
equation - need not be isomorphic modulo \(p\) to the residue class ring
of the principal order modulo \((p)\). The main difficulties in the
usual representations lie precisely in this failure of isomorphism.
Since the auxiliary polynomial ring does not occur here, these
difficulties are automatically avoided. At the same time the additive
decomposition theory of the residue class ring becomes the equivalent of
the multiplicative decomposition of the equation; when the equation
decomposes into relatively prime factors, a corresponding additive
decomposition occurs in the polynomial residue class ring.

It should be added that the discriminant theorem can only be the first
theorem of a general ramification theory for orders, just as for the
principal order the discriminant theorem is accompanied by the finer
theory of the different or ramification ideal, which there permits a
more precise determination of exponents. This determination of
exponents, resting on the power representation of the primary factors of
equations, may be interpreted as determination of the length of a
composition series, from \(p\) to \(\mathfrak q=p^t\), and will
presumably appear in this form in the general ramification theory.

\subsection{§1. Extension rings of
fields}\label{extension-rings-of-fields}

We begin with a few remarks about arbitrary rings. If \(T\) is a ring
and \(S\) is a system of elements of \(T\), the ``ring derived from
\(S\) in \(T\)'' means the intersection of all rings in \(T\) which
contain \(S\). Correspondingly, one defines the ``ideal derived from
\(S\) in \(T\)'', or the ``\(P\)-module derived from \(S\) in \(T\)''
when \(P\) is a subring of \(T\).

If \(R\) is a subring of \(T\) and \(S\) is a system of elements of
\(T\), the ring derived from \(R\) and \(S\) in \(T\) is denoted by \[
  R[S]
\] and is called the ring obtained by ring-adjoining \(S\) to \(R\). It
consists of all polynomials in the elements of \(S\) with coefficients
in \(R\), where equality and the operations are fixed by the equality
and operation relations already present in \(T\). Because these
relations have already been fixed, in special cases it is enough to form
only a subsystem of all polynomials - for example, all linear forms in
\(S\) with coefficients in \(R\). Such a case occurs, for instance, when
\(R\) is a ring with unit element and \(S\) is a field with the same
unit element. Then all linear combinations \[
  \sum c_i y_i,
\] with \(c_i\in S\) and \(y_i\in R\), form the ring \(R[S]\).

One can also construct the extension ring \(R[S]\) when only the ring
\(R\) is given and, in addition, a system \(S\) of elements not lying in
\(R\) is given with equality relations and perhaps mutual operation
relations. One shows that there is always a ring containing \(R\) and
\(S\), namely the ring formally formed from all polynomials in \(S\)
with coefficients in \(R\), subject to the usual operation rules. Apart
from these operation rules, there remains complete freedom in defining
equality, except that the equality definitions already valid in \(R\)
and \(S\) must be included. Thus one is forced to identify only those
polynomials which, by equality in \(R\) and \(S\) and by the stated
operations, are built in the same way from equal elements of \(R\) and
\(S\). The ring \(T\) so constructed, containing \(R\) and \(S\), is
then identical with the ring derived in \(T\) from \(R\) and \(S\).

From now on, rings \(R\) are assumed to have a unit element and to
contain a field \(P\) as subring; they are thus over-rings of fields.
The unit element \(e\) of \(R\) is also the unit element of \(P\), and
\(R\) is a \(P\)-module.

If \(\bar P\) is an overfield of \(P\), the extension ring \[
  \bar R=R[\bar P]
\] obtained by ring-adjoining \(\bar P\) to \(R\) is to be understood as
follows. Assume that the set-theoretic intersection of \(R\) and
\(\bar P\) is exactly \(P\), and that no operation relations exist
between elements of \(R\) not in \(P\) and elements of \(\bar P\) not in
\(P\). This assumption can always be achieved by replacing \(\bar P\) by
an equivalent extension field of \(P\), or \(R\) by an equivalent
extension ring.

The elements of \(\bar R\) are all formally formed linear combinations
\[
  c_1y_1+\\cdots+c_ny_n,
\] where the \(c_i\) run through the elements of \(\bar P\) and the
\(y_i\) through the elements of \(R\). The definition is to be
unambiguous: if the \(y_i\) are replaced by equal elements of \(R\), and
the \(c_i\) by equal elements of \(\bar P\), the element of \(\bar R\)
is to be the same.

The sum is defined in the evident way, \[
  \sum c_i y_i+
  \sum d_i y_i=
  \sum (c_i+d_i)y_i,
\] where zeros may occur among the \(c_i,d_i\) so that formally the same
\(y_i\) occur. The product is likewise defined by distributivity, \[
  \left(\sum c_i y_i\right)
  \left(\sum d_j y_j\right)
  =\sum c_i d_j y_i y_j.
\] For equality one declares the following. If \[
  y_{i_1},\ldots,y_{i_s}
\] are linearly independent over \(P\), then two elements \[
  \sum c_j y_{i_j},\qquad \sum d_j y_{i_j}
\] are equal if and only if \(c_j=d_j\) for all \(j\) in \(\bar P\). In
other words, elements of \(R\) linearly independent over \(P\) are
declared linearly independent over \(\bar P\).

This gives equality for all elements of \(\bar R\), because every
element can be expressed by linearly independent elements of \(R\).
Indeed, from a relation \[
  u=c_1y_1+
\cdots+c_ny_n
\] in \(\bar R\), multiplication by \(b=be\) gives \[
  bu=bc_1y_1+
\cdots+bc_ny_n,
\] so that all elements are expressible in terms of linearly independent
elements of \(R\). Equality is therefore coefficient equality over
\(\bar P\). For elements lying simultaneously in \(R\) and \(\bar R\),
or in \(P\) and \(\bar P\), this agrees with the old equality. For the
remaining elements of \(\bar R\), the equality in \(R\) and \(\bar P\)
says no more than the required unambiguity of the definition.

Nor does the requirement that sum and product be unambiguous introduce
further relations: equal elements may be assumed coefficientwise equal,
and hence they give formally equal sums and products. One checks at once
that all ring axioms hold. Thus the extension ring \(R[\bar P]\)
generated by ring adjunction always exists.

Now let \(\mathfrak a\) be an ideal of \(R\). The module product \[
  \bar R\mathfrak a
\] is an ideal \(\bar{\mathfrak a}\) of \(\bar R\), the extension ideal
of \(\mathfrak a\). It is also the ideal derived from \(\mathfrak a\) in
\(\bar R\); it consists of all finite linear combinations \[
  \sum c_i a_i,
\] where \(a_i\in\mathfrak a\) and \(c_i\in\bar P\).

If \(\bar{\mathfrak b}\) is an ideal of \(\bar R\), then the
intersection \[
  [\bar{\mathfrak b},R]
\] is an ideal \(\mathfrak b\) of \(R\), the contraction ideal of
\(\bar{\mathfrak b}\). Every ideal \(\mathfrak a\) of \(R\) is a
contraction ideal, namely the contraction of its own extension ideal.
For an element \(a\) of \([R,\bar{\mathfrak a}]\) has the form
\(\sum c_i a_i\) with the \(a_i\) chosen linearly independent over \(P\)
and hence over \(\bar P\). It follows that the elements \(a,a_i\) are
linearly dependent in \(\bar R\), and therefore, by the equality
definition, already in \(R\); hence \(a\) belongs to \(\mathfrak a\).
The converse is immediate.

If \(\bar{\mathfrak p}\) is a prime ideal of \(\bar R\), then \[
  \mathfrak p=[R,\bar{\mathfrak p}]
\] is a prime ideal of \(R\). By the second isomorphism theorem, the
residue ring \(R/\mathfrak p\) is isomorphic to a subring of
\(\bar R/\bar{\mathfrak p}\), and hence has no zero divisors. More
precisely, \[
  (R,\bar{\mathfrak p})/\bar{\mathfrak p}
  \simeq R/[R,\bar{\mathfrak p}],
\] where \((R,\bar{\mathfrak p})\) denotes the sum, the greatest common
divisor, of the modules \(R\) and \(\bar{\mathfrak p}\).

\subsection{§2. Rings of finite rank. Representation as a direct
sum}\label{rings-of-finite-rank.-representation-as-a-direct-sum}

If \(R\) is a finite \(P\)-module, then \(R\) has finite rank over
\(P\), say \(n\): there are \(n\), but no more, linearly independent
elements over \(P\), and every system of \(n\) linearly independent
elements is a module basis. If a \(P\)-module \(B\) of finite rank is
divisible by another \(P\)-module \(A\) of the same finite rank, then
they are equal.

From now on \(R\) is assumed to have finite rank \(n\) over \(P\). Then
every \(P\)-module in \(R\) has finite rank. In every proper divisor or
multiple chain of \(P\)-modules in \(R\), the rank of each module must
therefore be greater, respectively smaller, than the rank of the
preceding module; hence the chains terminate. In particular the double
chain condition holds for the system of all ideals of \(R\).

It follows that a prime ideal \(\mathfrak p\) of \(R\) has no proper
divisor other than the unit ideal. Thus the residue ring
\(R/\mathfrak p\) modulo a prime ideal different from the unit ideal is
a field. If \(\mathfrak q\) is a primary ideal belonging to
\(\mathfrak p\), say \[
  \mathfrak q\equiv0\pmod{\mathfrak p},
  \qquad
  \mathfrak p^{\rho}\equiv0\pmod{\mathfrak q},
\] then \(R/\mathfrak q\) is a primary ring: a power, in fact at latest
the \(\rho\)-th power, of every zero divisor vanishes. Further,
\(R/\mathfrak q\) contains only zero divisors and units. This is
immediate because every ideal not divisible by \(\mathfrak p\) is
relatively prime to \(\mathfrak p^{\rho}\) and hence to \(\mathfrak q\).

From the double chain condition and the existence of the unit element it
follows that every ideal \(\mathfrak a\) of \(R\) is uniquely
representable as a product of finitely many pairwise relatively prime
primary ideals. Because of relative primality, this product is the least
common multiple, and it corresponds to a representation of the residue
class ring \(R/\mathfrak a\) as a direct sum of primary rings:
\(R/\mathfrak a\) is the greatest common divisor of these rings, and the
sum representation of every element of \(R/\mathfrak a\) is unique.

Let the decomposition of the zero ideal of \(R\) be \[
  (0)=\mathfrak q_1\mathfrak q_2\cdots\mathfrak q_r
      =[\mathfrak q_1,\mathfrak q_2,\ldots,\mathfrak q_r],
\] and let the corresponding direct-sum decomposition be \[
  R=R_1+
  \cdots+R_r,
  \qquad R_i R_j=(0)\quad (i\ne j),
  \qquad R_i^2=R_i.
\] Here \(R_i\) is an ideal of \(R\), and \[
  R_i\simeq R/\mathfrak q_i;
\] the element \(y_i\in R_i\) corresponds to the class represented by
\(y_i\) in \(R/\mathfrak q_i\).

If \(\mathfrak a\) is any ideal of \(R\), the distributive law gives \[
  \mathfrak a=R_1\mathfrak a+
  \cdots+R_r\mathfrak a.
\] The components \[
  \mathfrak a_i=R_i\mathfrak a
\] are ideals both in \(R_i\) and in \(R\). As ideals they are finite
\(P\)-modules, and because the sum is direct their ranks add to the rank
of \(\mathfrak a\) or of \(R\).

If the unit element has the decomposition \[
  e=e_1+
\cdots+e_r,
\] where \[
  e_i e_j=0\quad(i\ne j),
  \qquad e_i^2=e_i,
\] then every element \(y\) of \(R\) has the unique decomposition \[
  y=y_1+
\cdots+y_r,
  \qquad y_i=ye_i\in R_i,
\] and \(R_i=Re_i\). The orthogonality relations show that the \(i\)-th
component of a product is \(y_i z_i\), and the \(i\)-th component of a
sum is \(y_i+z_i\).

The ring \(R_i\) contains a subfield \[
  P_i=Pe_i
\] isomorphic to \(P\), and has finite rank over \(P_i\); this rank is
equal to the rank of \(R_i\) as \(P\)-module. For if \(c\in P\), then
\(ce_i=ce\) modulo \(\mathfrak q_i\), so \(ce_i\ne0\) whenever
\(c\ne0\). Hence \(P\to P_i\) is an isomorphism. Linear dependence over
\(P\) among elements of \(R_i\) passes, by multiplication with \(e_i\),
to linear dependence over \(P_i\), and conversely, since \(P_i=Pe_i\),
every relation over \(P_i\) is already a relation over \(P\). Thus the
ranks agree.

If \[
  R=R_1+
\cdots+R_r,
\] then for any extension ring \(\bar R=R[\bar P]\) there is a
corresponding decomposition \[
  \bar R=\bar R_1+
\cdots+\bar R_r,
\] where \(\bar R_i\) denotes the extension ideal of \(R_i\) in
\(\bar R\), and at the same time the ring derived from \(R_i\) in
\(\bar R\). Indeed, from \((0)=\mathfrak q_1\cdots\mathfrak q_r\)
follows by multiplication with \(\bar R\) the product representation \[
  (0)=\bar{\mathfrak q}_1\cdots\bar{\mathfrak q}_r,
\] and the \(\bar{\mathfrak q}_i\) remain pairwise relatively prime. The
direct summands are the corresponding products of the other factors,
hence the extension ideals of the \(R_i\). But this extension ideal is,
by §1, 3, just the extension ring \[
  R_i[P_i]
\] with \(P_i=Pe_i\). Thus in extending one may extend the individual
components separately. In general \(\bar R_i\) need no longer be a
primary ring.

We shall also need a distinction between prime ideals of first and
second kind. Modulo every prime ideal \(\mathfrak p\) different from the
unit ideal, the residue class ring \(R/\mathfrak p\) is a field. It
contains a subfield isomorphic to \(P\), namely the classes represented
by elements of \(P\). This follows from the second isomorphism theorem:
\[
  (P,\mathfrak p)/\mathfrak p \simeq P/[P,\mathfrak p]\simeq P,
\] because \(\mathfrak p\) contains no non-zero element of \(P\). The
field \(R/\mathfrak p\) has finite rank over this copy of \(P\), and is
therefore a finite algebraic extension field. According as this
extension is of the first or second kind, \(\mathfrak p\) is called a
prime ideal of the first or second kind.

\subsection{§3. Completely reducible
rings}\label{completely-reducible-rings}

The ring \(R\) is called \textbf{completely reducible} if in the product
decomposition of its zero ideal all primary components are prime ideals.
Equivalently, \(R\) can be represented as a direct sum of finitely many
fields. The completely reducible ring is called completely reducible of
the first kind if all its prime ideals are of the first kind; otherwise
it is completely reducible of the second kind.

If \(P\) is a perfect field, only complete reducibility of the first
kind can occur; this is in particular always the case when \(P\) is
algebraically closed. In what follows \(\bar P\) denotes the algebraic
closure of \(P\), and \[
  \bar R=R[\bar P].
\]

\textbf{Theorem.} The ring \(R\) is completely reducible of the first
kind if and only if the extension ring \[
  \bar R=R[\bar P]
\] with algebraically closed \(\bar P\) is completely reducible.

The proof has three parts.

First, complete reducibility of an extension ring implies complete
reducibility of \(R\). More generally, if the zero ideal of an extension
ring has a representation as a least common multiple of prime ideals, \[
  (0)=[\bar{\mathfrak p}_1,
\ldots,\bar{\mathfrak p}_r],
\] then by intersecting with \(R\) one obtains the representation \[
  (0)=[R,\bar{\mathfrak p}_1],
\ldots,[R,\bar{\mathfrak p}_r]
\] of the zero ideal in \(R\); and the contracted ideals are prime by
§1, 3.

Second, complete reducibility of the first kind of \(R\) implies
complete reducibility of the first kind of every extension ring, and
more generally of every intermediate ring between \(R\) and the
extension ring. Since by §2, 3 the direct-sum decomposition of \(R\)
induces one of any extension ring, it is enough to consider one
component. Thus we may assume without loss of generality that \(R\) is a
field, and indeed a finite extension field of the first kind of its
subfield \(P\). Then \(R\) is obtained from \(P\) by adjoining one
primitive element of the first kind. Equivalently, \[
  R\simeq P[x]/(f(x)),
\] where \(f(x)\) is a prime polynomial of the first kind in the
polynomial ring \(P[x]\). If \(Q\) is any extension field of \(P\), then
\[
  Q[x]/(f(x))\simeq R[Q],
\] for the rank of the residue ring remains equal to the degree of
\(f(x)\) when passing from \(P[x]\) to \(Q[x]\), and hence this is
exactly the construction of the extension ring described in §1, 2. In
\(Q[x]\) the polynomial \(f(x)\) decomposes into pairwise relatively
prime prime functions of the first kind; each generates a prime ideal in
\(Q[x]\). Therefore \(Q[x]/(f(x))\), and hence \(R[Q]\), is completely
reducible of the first kind. Taking \(Q=\bar P\) proves the second step.

Third, if \(R\) is completely reducible of the second kind, then
\(\bar R\) is not completely reducible. As before it suffices to take
\(R\) to be a field, now a finite extension field of the second kind of
\(P\). We must show that in the direct-sum decomposition of \(\bar R\)
into primary rings at least one proper primary component occurs. It is
enough to exhibit a non-zero element of \(\bar R\) whose power is zero:
then some component is non-zero but nilpotent, and the corresponding
component ring is not a field.

Let \(\gamma\) be an element of the second kind in \(R\), of exponent
\(\ell>1\), and let \[
  F(\gamma)=\gamma^{kp^{\ell}}+a_1\gamma^{kp^{\ell}-1}+
\cdots+a_{kp^{\ell}}=0
\] be a lowest-degree equation satisfied by \(\gamma\) over \(P\), where
\(p\) is the characteristic of \(P\). The elements \[
  1,\gamma,
\gamma^2,
\ldots,
\gamma^{kp^{\ell}-1}
\] are linearly independent over \(P\) and therefore remain linearly
independent over \(\bar P\) by the equality definition. If \[
  G(\gamma)=\gamma^{kp^{\ell-1}}+
\cdots
\] is chosen so that \[
  G(\gamma)^p=F(\gamma),
\] then \(G(\gamma)\) is non-zero in \(\bar R\) but its \(p\)-th power
vanishes. This proves the third step.

The first and third steps show that complete reducibility of \(R\) of
the first kind follows from complete reducibility of \(\bar R\); the
second step gives the converse.

As a consequence, if \(R\) is completely reducible of the first kind,
then \(\bar R\) is a direct sum of rings of rank one, hence of fields
isomorphic to \(\bar P\). Such a decomposition already exists in some
ring \(R[Q]\), where \(Q\) is a finite extension field of \(P\).

Indeed, by the theorem, \(\bar R\) is a direct sum of fields which, as
finite extensions of subfields isomorphic to \(\bar P\), are identical
with those subfields. Thus \[
  \bar R=\bar P e_1+
\cdots+\bar P e_n,
  \qquad e=e_1+
\cdots+e_n,
\] and the components \(e_i\) of the identity form a linearly
independent module basis of \(\bar R\). Expressing the \(e_i\) as linear
combinations of elements of \(R\) with coefficients in \(\bar P\), the
set of all coefficients generates a finite extension field \(Q\) of
\(P\). The \(e_i\) therefore already belong to \(R[Q]\), and the same
decomposition exists there. In every intermediate ring between \(R[Q]\)
and \(R\) the indecomposable components arise as extension rings of the
\(Qe_i\).

\subsection{§4. Matrix representation, traces, and discriminants in
rings of finite
rank}\label{matrix-representation-traces-and-discriminants-in-rings-of-finite-rank}

In this and the next section the hypotheses on the base ring are
slightly more general than before, in order also to cover orders in
number and function fields. The treatment is a general and uniform
formulation of essentially known facts.

Assume that \(R\) is an extension ring of an integral domain \(Z\), and
that the unit element of \(R\) already lies in \(Z\). For every
\(Z\)-module in \(R\), in particular for \(R\) itself, at least one
finite module basis exists which consists of elements linearly
independent over \(Z\). Two such bases \(a_1,
\ldots,a_n\) and \(b_1,
\ldots,b_n\) are bases of the same \(Z\)-module if and only if the
transformation determinant is a unit of \(Z\). This hypothesis is
plainly fulfilled for the finite-rank rings considered above, where
\(Z\) is the field \(P\), and it is also fulfilled for orders in number
and function fields, where \(Z\) is the ring of rational integers,
respectively the functional domain.

Let \(a_1,
\ldots,a_n\) be a linearly independent \(Z\)-module basis of an ideal
\(\mathfrak a\), and let \(y\) be an arbitrary element of \(R\). Since
\(ya_i\) again belongs to \(\mathfrak a\), equations with coefficients
in \(Z\) hold: \[
  ya_i=c_{i1}a_1+
\cdots+c_{in}a_n.
\] In matrix form, if \((a_1,
\ldots,a_n)\) denotes a row matrix, \[
  y(a_1,
\ldots,a_n)=(a_1,
\ldots,a_n)C.
\] As \(y\) ranges over \(R\), the matrices \(C\) form a ring
\(R_{\mathfrak a}\) homomorphic to \(R\). More precisely, \[
  R_{\mathfrak a}\simeq R/((0):\mathfrak a),
\] where \((0):\mathfrak a\) is the ideal quotient of the zero ideal by
\(\mathfrak a\). Differences and products of elements correspond to
differences and products of matrices; scalar elements \(c\in Z\)
correspond to diagonal matrices \(cE\). Exactly those elements \(y\) for
which \(y\mathfrak a=(0)\) correspond to the zero matrix, which proves
the stated isomorphism.

A different basis of \(\mathfrak a\) gives an equivalent matrix ring. If
\[
  (\bar a_1,
\ldots,\bar a_n)=(a_1,
\ldots,a_n)P
\] with determinant of \(P\) a unit of \(Z\), then the corresponding
matrices satisfy \[
  \bar C=P^{-1}CP.
\]

Two ideals \(\mathfrak a\) and \(\mathfrak b\) of \(R\) belong to the
same ideal class if they are isomorphic as \(R\)-modules; that is, their
elements correspond one-to-one in such a way that difference and
multiplication by the same element of \(R\) correspond. Two matrix rings
obtained as above belong to the same representation class if they are
equivalent, i.e.~if \[
  R_2=P^{-1}R_1P
\] for a unimodular matrix \(P\). Ideal classes and representation
classes correspond one-to-one.

Indeed, different bases of the same ideal give equivalent matrix rings,
as just shown. Isomorphic ideals also give the same representation
class: if \(b_1,
\ldots,b_n\) correspond to a basis \(a_1,
\ldots,a_n\), then \(ya_i\) and \(yb_i\) correspond and the same
matrices occur. Conversely, if the matrix rings are equivalent, replace
the basis of \(\mathfrak b\) by the transformed basis. Then the matrices
coincide, and the correspondence sending \[
  c_1a_1+
\cdots+c_na_n
  \quad\text{to}\quad
  c_1b_1+
\cdots+c_nb_n
\] is an \(R\)-module isomorphism. Thus the ideals are in the same
class. The class of the unit ideal is called the principal class; every
basis of \(R\) gives the principal representation class.

Let \(C\) be the matrix corresponding to the element \(y\) in the
representation class determined by \(\mathfrak a\). The polynomial \[
  |tE-C|
\] is invariant under equivalence of representations. Hence its
coefficients are invariants of the ideal class. In particular, the
coefficient of \(-t^{n-1}\) is called the trace of \(y\) with respect to
the class \(\mathfrak a\), written \[
  S_{\mathfrak a}(y)=c_{11}+c_{22}+
\cdots+c_{nn}.
\] The constant coefficient, up to sign, is the norm of \(y\) with
respect to the class.

For a fixed class \((\mathfrak a)\), the traces \(S_{\mathfrak a}(y)\)
form a \(Z\)-module to which \(R\) is homomorphic as a \(Z\)-module.
This follows from the matrix homomorphism: \[
  S_{\mathfrak a}(x-y)=S_{\mathfrak a}(x)-S_{\mathfrak a}(y),
  \qquad
  S_{\mathfrak a}(cy)=cS_{\mathfrak a}(y).
\] If \(R\) has no zero divisors, the trace and norm of an element are
independent of the class. For every non-zero ideal has the same rank as
\(R\), and a basis of an ideal is at the same time a basis of the unit
ideal in the quotient field. Thus all traces and norms are generated in
the quotient field by the unit ideal and are independent of the class.

Now let \(\omega_1,
\ldots,
\omega_n\) be a \(Z\)-module basis of an ideal \(\mathfrak r\). The
determinant \[
  |S_{\mathfrak a}(\omega_i\omega_j)|
\] is well defined only up to multiplication by the square of a unit of
\(Z\), but the principal ideal it generates is an invariant of the ideal
\(\mathfrak r\) and the class \((\mathfrak a)\). A change of basis
transforms the matrix cogrediently; the determinant therefore changes by
the square of the transformation determinant. This principal ideal is
called the discriminant ideal of \(\mathfrak r\) with respect to the
class \((\mathfrak a)\). If \(R\) has no zero divisors, the discriminant
ideal of an ideal is independent of the chosen class, because each entry
in the determinant is already independent of the class.

\subsection{§5. Direct sums of classes}\label{direct-sums-of-classes}

If the ideal \(\mathfrak a\) is a direct sum, \[
  \mathfrak a=\mathfrak b+\mathfrak c,
\] then every ideal in the class of \(\mathfrak a\) is likewise a direct
sum, with components isomorphic to \(\mathfrak b\) and \(\mathfrak c\)
respectively. Thus one may speak of the direct sum of the ideal class
and of its component classes.

Similarly, if a ring \(R_{\mathfrak a}\) in a representation class is a
direct sum of subrings which are at the same time ideals in
\(R_{\mathfrak a}\), then the same is true for every equivalent
representative, and the components are equivalent rings. Thus one may
speak of the direct sum of the representation class and of its component
classes. The direct sum of ideal classes corresponds to the direct sum
of representation classes; in this sense we simply speak of the direct
sum of a class.

Let \(b_1,
\ldots,b_s\) be a basis of \(\mathfrak b\) and \(c_1,
\ldots,c_t\) a basis of \(\mathfrak c\). Since the sum is direct, the
two bases together form a basis of \(\mathfrak a\). With respect to this
basis, the matrix associated with an arbitrary element of \(R\) has
block-diagonal form \[
  \begin{pmatrix}D^{(b)}&0\\0&D^{(c)}\end{pmatrix}.
\] The systems of matrices with only one block non-zero form subrings,
and in fact ideals, whose direct sum is the whole representation ring.
They are isomorphic to the representation rings obtained from
\(\mathfrak b\) and \(\mathfrak c\) themselves.

The trace of an element with respect to a class is the sum of the traces
with respect to the component classes. If \(\beta\in\mathfrak b\), then
the trace of \(\beta\) with respect to \((\mathfrak a)\) is equal to the
trace of \(\beta\) with respect to \((\mathfrak b)\). Hence the
discriminant ideal of the component ideal \(\mathfrak b\), taken with
respect to \((\mathfrak a)\), is equal to its discriminant ideal taken
with respect to the component class \((\mathfrak b)\).

Finally, the discriminant ideal of \(\mathfrak a\) with respect to
\((\mathfrak a)\) is the product of the discriminant ideals of the
components with respect to their component classes. This follows by
taking the basis corresponding to the direct sum: the discriminant
determinant becomes block diagonal, and its determinant is the product
of the component determinants.

Now let \(X\) be an extension ring without zero divisors over \(Z\) such
that the intersection of \(R\) and \(X\) is \(Z\), and let \[
  \bar R=R[X]
\] be the extension ring of the same rank as in §1, 2. If
\(\mathfrak a\) and \(\mathfrak r\) are ideals of \(R\), and
\(\bar{\mathfrak a}\) and \(\bar{\mathfrak r}\) their extensions in
\(\bar R\), then the discriminant ideal of \(\bar{\mathfrak r}\) with
respect to \((\bar{\mathfrak a})\) is the extension ideal in \(X\) of
the discriminant ideal of \(\mathfrak r\) with respect to
\((\mathfrak a)\). One can use the same bases in the discriminant
construction. In particular, the discriminant ideal of \(\bar R\) with
respect to the principal class is the extension of the discriminant
ideal of \(R\) with respect to the principal class. This ideal, in \(X\)
or \(Z\), is called simply the discriminant ideal of the ring.

\subsection{§6. Discriminant criterion for complete reducibility of the
first
kind}\label{discriminant-criterion-for-complete-reducibility-of-the-first-kind}

From now on \(R\) is again an extension ring of a field \(P\), and
\(\bar P\) is an extension field of \(P\). The non-vanishing of the
discriminant will be a necessary and sufficient condition for complete
reducibility of the first kind. By §5, 3, the discriminant ideal of
\(R\) may be determined after passing to the extension ring
\(\bar R=R[\bar P]\). By §5, 2, it is enough to restrict attention to
the primary rings out of which \(\bar R\) is additively composed. We
therefore begin with the discriminant ideals of primary rings.

If \(\mathfrak a\) and \(\mathfrak r\) are ideals of \(R\) and
\(\mathfrak a\) is divisible by \(\mathfrak r\), then every basis of
\(\mathfrak a\) can be completed by adjoining elements to a linearly
independent module basis of \(\mathfrak r\). This is an immediate
consequence of finite rank over the field \(P\).

Assume \(R\) is a primary ring; let \(\mathfrak p\) be the prime ideal
belonging to the zero ideal, and let \(\rho\) be its exponent, so \[
  \mathfrak p^{\rho}=(0),
  \qquad \mathfrak p^{\rho-1}\ne(0).
\] A special \(P\)-basis of \(R\) is obtained by successively completing
a basis of \(\mathfrak p^{\rho-1}\) to a basis of
\(\mathfrak p^{\rho-2}\), then a basis of \(\mathfrak p^{\rho-2}\) to
one of \(\mathfrak p^{\rho-3}\), and so on, counting \(R\) as
\(\mathfrak p^0\). With such a basis one proves:

If \(R\) is a primary ring, the trace, with respect to the principal
class, of every element divisible by the associated prime ideal
\(\mathfrak p\) is zero.

Indeed, with the special basis the basis elements are arranged in
layers, where elements in the layer \(\mathfrak p^i\) are not all in
\(\mathfrak p^{i+1}\). Multiplication by an element of \(\mathfrak p\)
sends each layer into the next; the corresponding matrix has zeros on
the diagonal. Its trace therefore vanishes.

If the zero ideal of \(R\) is already a prime ideal and \(P\) is
algebraically closed, then \(R\) has rank one and is identical with
\(P\). The unit element \(e\) may be used as basis, and \[
  S(e^2)=S(e)=e.
\] The discriminant ideal is therefore the unit ideal.

On the other hand, the discriminant ideal of a proper primary ring with
algebraically closed subfield \(P\) is the zero ideal. Use the special
basis just described, or simply complete any basis of \(\mathfrak p\) to
a basis of \(R\). Since the residue class ring \(R/\mathfrak p\) has
rank one, the first basis element may be taken to be the unit element
\(e\). All basis products other than \(e^2\) are then divisible by
\(\mathfrak p\), and their traces vanish. The discriminant determinant,
having at least two rows because \(R\) is properly primary, therefore
has all entries except possibly \(S(e^2)\) equal to zero and
consequently vanishes.

\textbf{Theorem.} A ring \(R\) of finite rank over a subfield \(P\) is
completely reducible of the first kind if and only if its discriminant
ideal is the unit ideal in \(P\), equivalently if its discriminant,
determined up to the square of a unit, is non-zero.

Complete reducibility of the first kind is, by §3, 2, equivalent to
complete reducibility of \[
  \bar R=R[\bar P]
\] with algebraically closed \(\bar P\). By §5, 3 the discriminant
ideals of \(R\) and \(\bar R\) are simultaneously zero or unit ideals.
By §5, 2 the discriminant ideal of \(\bar R\) is the product of the
discriminant ideals of its direct-sum components, each taken with
respect to its component class. But those component discriminants are
exactly of the type just considered. If \(\bar R\) is completely
reducible, all components are fields isomorphic to \(\bar P\) and all
component discriminant ideals are unit ideals. Their product is a unit
ideal. If \(\bar R\) is not completely reducible, at least one component
is properly primary and its discriminant ideal is zero; the product is
then zero. The theorem follows.

When complete reducibility of the first kind holds, trace, norm, and
discriminant admit the usual representation by conjugate elements. Let
\[
  R[Q]=Qe_1+
\cdots+Qe_n
\] be a decomposition as a direct sum of rank-one rings over a suitable
finite extension field \(Q\) of \(P\), chosen minimally. The components
of \(R\) are represented in subfields \(K_i\) of \(Q\), and \(R\) maps
homomorphically to each \(K_i\). With respect to the basis \(e_1,
\ldots,e_n\), an element \[
  y=c_1e_1+
\cdots+c_ne_n
\] has diagonal matrix with entries \(c_i\). Hence \[
  S(y)=c_1+
\cdots+c_n,
  \qquad
  N(y)=c_1c_2\cdots c_n.
\] If \(a_1,
\ldots,a_n\) is a module basis of \(R\) and \[
  a_j=a_{1j}e_1+
\cdots+a_{nj}e_n,
\] then the discriminant is \[
  |S(a_i a_j)|=
  \left|
  \begin{matrix}
  a_{11}&\cdots&a_{1n}\\
  \vdots& &\vdots\\
  a_{n1}&\cdots&a_{nn}
  \end{matrix}
  \right|^2.
\] If \(R\) itself is a field, i.e.~an extension field of the first kind
of \(P\), then the homomorphisms \(R\to K_i\) are isomorphisms fixing
\(P\). Thus the \(K_i\) are the \(n\) conjugate fields in a common
Galois extension field, each isomorphic to \(R\). The formulas for
trace, norm, and discriminant become the familiar formulas in terms of
conjugates. It should again be stressed that \(R\) itself is assumed
only equivalent to the \(K_i\) and need not be a subfield of the Galois
extension field \(Q\).

\subsection{§7. The discriminant theorem for orders of an algebraic
number or function
field}\label{the-discriminant-theorem-for-orders-of-an-algebraic-number-or-function-field}

The number and function fields in question - including algebraic
functions of more than one indeterminate and integral algebraic
functions, where one is extending a functional domain - all fit the
following type of field with a prescribed notion of integral element.

Let \(\mathfrak S\) be a principal ideal ring without zero divisors and
with unit element: for example, the ring of rational integers, a
polynomial ring in one indeterminate over a field, or a functional
domain. Let \(K\) be its quotient field. Let \(L\) be a finite extension
of the first kind of \(K\), and let \(\mathfrak G\) be the ring of all
elements of \(L\) integral with respect to \(\mathfrak S\). Every
subring \(T\) of \(\mathfrak G\) which contains \(\mathfrak S\) is
called an order; \(\mathfrak G\) itself is the principal order.

Every \(\mathfrak S\)-module in \(\mathfrak G\), in particular every
ideal of an order and the order itself, has a linearly independent
module basis over \(\mathfrak S\). If \(L\) has degree \(n\) over \(K\),
it is enough to consider orders of rank \(n\) over \(\mathfrak S\): an
order of smaller rank would, by forming quotients, generate an
intermediate field of that smaller degree, and all hypotheses would be
preserved there.

Thus every order \(T\) is a ring satisfying the more general hypotheses
of §§4 and 5, and is in fact a ring without zero divisors. Consequently
the discriminant \(D_T\) of the order is defined uniquely up to
multiplication by the square of a unit of \(\mathfrak S\). This
discriminant agrees, up to a unit of \(\mathfrak S\), with the
discriminant of the field \(L\) considered as a \(K\)-module, and is
therefore non-zero because \(L/K\) is an extension of the first kind. It
also admits the representation in §6, 4 as the square of a determinant
of conjugate basis elements.

The ideal theory in \(T\) is given by the following theorem. Every ideal
of \(T\) different from the zero ideal can be represented uniquely as a
product of finitely many pairwise relatively prime primary ideals, whose
associated prime ideals have no proper divisors except the unit ideal.

Let \(p\) be a prime element of \(\mathfrak S\), so that the principal
ideal \(\mathfrak S p\) is a prime ideal different from the zero and
unit ideals. Let \(Tp\) be the corresponding principal ideal of \(T\).
The residue class ring \[
  R=T/Tp
\] is a ring of rank \(n\) over a subfield \(P\) isomorphic to
\(\mathfrak S/\mathfrak S p\). The natural homomorphism \(T\to T/Tp\) is
always defined. The fact that \(R\) contains a subfield isomorphic to
\(\mathfrak S/\mathfrak S p\) follows from the second isomorphism
theorem: \[
  (\mathfrak S,Tp)/Tp\simeq
  \mathfrak S/[\mathfrak S,Tp]
  =\mathfrak S/\mathfrak S p.
\] Moreover \(R\) has rank \(n\) over \(P\). If \[
  \alpha_1,
\ldots,
\alpha_n
\] is a linearly independent module basis of \(T\) over \(\mathfrak S\),
and \((\alpha_i)\) denotes the residue class of \(\alpha_i\) modulo
\(Tp\), then the \((\alpha_i)\) are linearly independent over \(P\). For
a linear relation over \(P\) would lift to a congruence \[
  c_1\alpha_1+
\cdots+c_n\alpha_n\equiv0\pmod{Tp},
\] i.e.~to \[
  c_1\alpha_1+
\cdots+c_n\alpha_n=p\gamma,
\] with \(\gamma\in T\). Expressing \(\gamma\) in the basis \(\alpha_i\)
and using uniqueness of the representation shows that every \(c_i\) is
divisible by \(p\), hence every class \((c_i)\) vanishes in \(P\).

Under the homomorphism the discriminant \(D_T\) passes to the
discriminant of \(R\), and the ideal decomposition of \(Tp\) passes to
the decomposition of the zero ideal of \(R\). The trace used here is the
one defined by matrix representation, not the representation by
conjugate elements, which is valid for \(T\) but not necessarily for
\(R\). The trace and hence the discriminant are therefore built from
ring operations and ring elements and pass through the homomorphism.

If \[
  Tp=\mathfrak q_1\cdots\mathfrak q_r
       =[\mathfrak q_1,
\ldots,
\mathfrak q_r]
\] is the decomposition of \(Tp\) in \(T\), then it passes to \[
  (0)=[(\mathfrak q_1),
\ldots,(
\mathfrak q_r)]
\] in \(R\). By the first isomorphism theorem, \[
  R/(\mathfrak q_i)\simeq T/\mathfrak q_i,
\] so that \(\mathfrak q_i\) and \((\mathfrak q_i)\) are simultaneously
proper primary ideals or prime ideals. The components remain pairwise
relatively prime under the homomorphism.

\textbf{Discriminant theorem.} A prime element \(p\) of \(\mathfrak S\)
divides the discriminant \(D_T\) of an order \(T\) if and only if, in
the decomposition of the principal ideal \(Tp\) into pairwise relatively
prime primary components in \(T\), at least one proper primary component
or at least one prime ideal of the second kind occurs.

If the residue class field \(\mathfrak S/\mathfrak S p\) is perfect,
then precisely those prime elements \(p\) which divide \(D_T\) have at
least one proper primary component. If \(\mathfrak S/\mathfrak S p\) is
perfect and \(T\) is the principal order, then \(p\) divides the
discriminant if and only if \(Tp\) is divisible at least by the square
of a prime ideal of the principal order.

Indeed, by the correspondence between the ideal decomposition of \(Tp\)
in \(T\) and the zero-ideal decomposition in \[
  R=T/Tp,
\] the occurrence of a proper primary component or of a prime ideal of
the second kind means exactly that \(R\) is not completely reducible of
the first kind. By §6, 3 this is the case if and only if the
discriminant of \(R\) vanishes. By the homomorphic passage just
described, this is exactly the divisibility of \(D_T\) by \(p\).

\subsection{§8. The discriminant theorem for orders of relative
fields}\label{the-discriminant-theorem-for-orders-of-relative-fields}

If, instead of the principal ideal ring \(\mathfrak S\), one starts with
the principal order of a number or function field - more generally, with
a multiplication ring, i.e.~a ring in which the ideal theory of the
principal order holds - then the discriminant \(D_T\) is replaced by the
discriminant ideal of \(T\) with respect to \(\mathfrak S\). The
discriminant theorem transfers completely by the following
considerations.

Let \(\mathfrak S\) be a multiplication ring: a ring without zero
divisors and with unit element, in which every ideal different from the
zero and unit ideals is uniquely representable as a product of powers of
prime ideals, and where these prime ideals have no proper divisors
except the unit ideal. Let \(K,L,\mathfrak G,T\) be defined as above.
Then the ideal theory stated in §7, 1 holds in \(T\).

In particular, if a multiplication ring \(\mathfrak S\) has only one
prime ideal different from the zero and unit ideals, then
\(\mathfrak S\) is a principal ideal ring. For let this prime ideal be
\(\mathfrak p\), and choose \(p\in\mathfrak p\) with
\(p\not\equiv0\pmod{\mathfrak p^2}\). Since every ideal different from
zero and unit is a power of \(\mathfrak p\), the principal ideal
\(\mathfrak S p\) must be equal to \(\mathfrak p\). Hence every power of
\(\mathfrak p\) is principal, and the zero and unit ideals are principal
as well.

The trace of an element \(\gamma\in T\) is defined as the trace of
\(\gamma\) from \(L\), where \(L\) is considered as a module over \(K\).
Any system of \(n\) elements of \(L\) linearly independent over \(K\)
may be used as a basis to generate a matrix representation defining the
trace. If \(\alpha_1,
\ldots,
\alpha_n\) is such a system, then because \(L/K\) is an extension of the
first kind, \[
  |S(\alpha_i\alpha_j)|
\] is a non-zero element of \(K\), namely the square of the determinant
of the conjugate basis elements. From integral closedness of
\(\mathfrak S\) in \(K\) it follows that this determinant lies in
\(\mathfrak S\) whenever all \(\alpha_i\) lie in \(\mathfrak G\).

The discriminant ideal is defined as follows. Let \(\tau_1,
\ldots,
\tau_n\) run through all systems of \(n\) elements of \(T\), and form
the determinant \[
  |S(\tau_i\tau_j)|.
\] This determinant is an element of \(\mathfrak S\), and is non-zero
when the \(\tau_i\) are linearly independent. The ideal of
\(\mathfrak S\) generated by all these determinants is called the
discriminant ideal of the order \(T\) with respect to \(\mathfrak S\);
it is different from the zero ideal.

If \[
  \beta_1,
\ldots,
\beta_m,
  \qquad m\ge n,
\] is a module basis of \(T\) over \(\mathfrak S\), then the finitely
many determinants corresponding to all choices of \(n\) of the
\(\beta_i\) form a basis of the discriminant ideal. This follows from
the module homomorphy of \(T\) to the system of traces from \(T\): the
determinants \(|S(\tau_i\tau_j)|\) are linear combinations, with
coefficients in \(\mathfrak S\), of the finitely many determinants
\(|S(\beta_i\beta_j)|\) obtained by choosing \(n\) basis elements. If
the order has a linearly independent module basis \(\alpha_1,
\ldots,
\alpha_n\), then \(|S(\alpha_i\alpha_j)|\) itself may be taken as a
basis of the discriminant ideal.

We recall the passage to the quotient ring. If \(\mathfrak a\) is an
ideal of \(\mathfrak S\), the quotient ring
\(\mathfrak S_{\mathfrak a}\) is the set of elements of the quotient
field of \(\mathfrak S\) having in the denominator an element prime to
\(\mathfrak a\). This quotient ring has no prime ideals other than zero,
unit, and the extension ideals of the prime ideals belonging to
\(\mathfrak a\). There is a one-to-one correspondence, together with
isomorphism of residue class rings, between the ideals of
\(\mathfrak S_{\mathfrak a}\) different from zero and unit and those
ideals of \(\mathfrak S\) whose associated prime ideals are among those
belonging to \(\mathfrak a\); zero and unit ideals correspond as well.

The same statements hold for a multiplication ring. If \(\mathfrak p\)
is a prime ideal, then \(\mathfrak S_{\mathfrak p}\) has only one prime
ideal different from zero and unit; hence, by the preceding paragraph,
\(\mathfrak S_{\mathfrak p}\) is a principal ideal ring. A generator
\(p\) of the prime ideal derived from \(\mathfrak p\) in
\(\mathfrak S_{\mathfrak p}\) is a prime element. The extension ideal of
an arbitrary ideal \(\mathfrak c\) of \(\mathfrak S\) is generated by
the primary component of \(\mathfrak c\) belonging to \(\mathfrak p\);
hence it is \((p)^{\rho}\) or the unit ideal according as
\(\mathfrak p\) divides \(\mathfrak c\) to exponent \(\rho\) or does not
divide it. Therefore if two ideals have the same extensions in all
quotient rings \(\mathfrak S_{\mathfrak p}\), they are equal: they are
divisible by the same prime ideals to the same powers.

Now let \(\mathfrak p\) be a prime ideal of \(\mathfrak S\), and write
\[
  B=T\mathfrak p
\] for its extension in \(T\). The quotient ring \(T_{B}\) has a module
basis of linearly independent elements over the principal ideal ring
\(\mathfrak S_{\mathfrak p}\); it is at the same time a finite
\(\mathfrak S_{\mathfrak p}\)-order in the extension field \(L\) of the
quotient field \(K\) of \(\mathfrak S_{\mathfrak p}\). Therefore the
discriminant theorem of §7, 3 applies to \(T_B\) over
\(\mathfrak S_{\mathfrak p}\).

The discriminant \(D_{T_B}\) generates the extension in
\(\mathfrak S_{\mathfrak p}\) of the discriminant ideal of \(T\) over
\(\mathfrak S\). Indeed, one may choose a module basis of \(T_B\) over
\(\mathfrak S_{\mathfrak p}\) consisting of elements \(\alpha_1,
\ldots,
\alpha_n\) of \(T\). Then \(|S(\alpha_i\alpha_j)|\) belongs to the
discriminant ideal, and every determinant \(|S(\tau_i\tau_j)|\) in
\(\mathfrak S_{\mathfrak p}\) is divisible by it.

We finally obtain the general discriminant theorem.

\textbf{General discriminant theorem.} A prime ideal \(\mathfrak p\) of
a multiplication ring \(\mathfrak S\) divides the discriminant ideal of
an order \(T\) over \(\mathfrak S\) if and only if, in the decomposition
of \(T\mathfrak p\) into pairwise relatively prime primary components in
\(T\), at least one proper primary component or at least one prime ideal
of the second kind occurs.

By the preceding paragraphs, \(\mathfrak p\) divides the discriminant
ideal of \(T\) if and only if the corresponding prime element \(p\) of
\(\mathfrak S_{\mathfrak p}\) divides the discriminant of \(T_B\). By
§7, this is equivalent to the failure of complete reducibility of the
first kind of the residue class ring \[
  T_B/T_Bp.
\] Since this residue class ring is isomorphic to \[
  T/T\mathfrak p,
\] the general discriminant theorem follows.

As a specialization to relative discriminants of number fields one
obtains: a prime ideal \(\mathfrak p\) of the principal order of a
number field divides the relative discriminant of an extension field if
and only if, in the principal order of the extension field,
\(\mathfrak p\) is divisible by the square of a prime ideal.

Göttingen, 30 March 1926.

\newpage

\section{With R. Brauer: On Minimal Splitting Fields of Irreducible
Representations}\label{with-r.-brauer-on-minimal-splitting-fields-of-irreducible-representations}

The question of the number fields of smallest degree over a ground field
\(P\) in which an irreducible representation of a finite group over
\(P\) decomposes into absolutely irreducible constituents was first
treated by I. Schur. Suppose, for example, that \(P\) contains the
character of such a constituent. Schur showed that the degree of these
fields over \(P\) - the \textbf{index} - agrees with the number of those
absolutely irreducible constituents, all of which belong to the same
class. Consequently the same field is already determined by the
requirement that one absolutely irreducible constituent be split off. He
also showed that the degree over \(P\) of every field in which such a
splitting takes place is a multiple of the index.

All such fields will be called \textbf{splitting fields}, even when
\(P\) does not contain the character. Schur showed by example that
splitting fields of smallest degree need not be isomorphic. If \(P\)
contains no associated simple character, the splitting fields must
contain the field of such a character and its conjugates over \(P\). The
absolutely irreducible representations belonging to different conjugate
characters are themselves conjugate, but they plainly belong to
different classes. Even in that case, in order to split off a system of
absolutely irreducible representations belonging to one class, it
suffices to take the field of the corresponding character together with
one of the corresponding splitting fields.

We now continue the study of splitting fields. Splitting fields of
smallest degree can be characterized by associated division algebras;
general splitting fields are characterized by two-sided simple rings
invariantly connected with these division algebras. In addition, the
example of the quaternion algebra shows that the degrees of minimal
splitting fields are not bounded. Here a splitting field is called
\textbf{minimal} if none of its proper subfields is a splitting field.

This unboundedness rests essentially on a number-theoretic existence
theorem proved by H. Hasse in the immediately following note. The
example is also treated elementarily and computationally. In this way,
without using the general theory or the number-theoretic existence
theorem, one still obtains the unboundedness of the degree numbers from
a weaker but sufficient existence theorem. It should also be noted that
this example again concerns splitting fields of a finite group, namely
the quaternion group, which also underlies Schur's example. The
representations written there are simultaneously isomorphic
representations of the quaternion algebra.

\subsection{1. Characterization of splitting
fields}\label{characterization-of-splitting-fields}

We first recall the basic concepts. Let \(\mathfrak S\) be a
hypercomplex system over a field \(P_0\). A representation \(\Pi\) of
\(\mathfrak S\) over \(P\) means a system of matrices with entries in
\(P\) such that \(\Pi\) is a homomorphic image of \(\mathfrak S\), and
such that elements \(p_0\) of \(P_0\) correspond to diagonal matrices
\(p_0E\). Thus \(P\) must contain \(P_0\). A representation \(\Pi\),
together with all its transforms \[
  P^{-1}\Pi P,
\] forms a representation class.

The representations of finite groups are included in this definition.
The associated hypercomplex system \(\mathfrak S\), the group ring,
consists of all group numbers, that is, all linear combinations of the
group elements with coefficients in \(P_0\), the original group elements
being considered linearly independent and multiplication being defined
by group multiplication and the laws of arithmetic. The terms reducible,
irreducible, completely reducible, and absolutely irreducible are used
in the usual sense; they will be applied also to representation classes.
A system \(\mathfrak S\) is called completely reducible over \(P\) if
all its representation classes over \(P\) are completely reducible.

Assume now that \(\mathfrak S\) is completely reducible over \(P_0\),
and that \(P_0\) is a perfect field. Then the same is true for every
algebraic extension \(P\) of \(P_0\), and the representation classes
remain completely reducible after every such extension. The center of
\(\mathfrak S\) is a direct sum of finitely many fields. Every extension
\(P\) of \(P_0\) isomorphic to one of these fields is the field of a
simple character.

If \(P_0\) is extended to such a field \(P\), a two-sided simple ring
splits off from \(\mathfrak S\); to it belongs the absolutely
irreducible representation class of the character. By Wedderburn's
theorem this ring is the direct product of a matrix ring over \(P\) with
a division algebra \(A\), determined uniquely up to isomorphism. The
center of \(A\) is \(P\), and \(A\) has degree \(m^2\) over \(P\), where
\(m\) is Schur's index attached to the character.

Every splitting field of \(\mathfrak S\), with respect to the
representation belonging to that character, is at the same time a
splitting field of \(A\) with respect to the representations over \(P\),
and conversely. The conjugate representations of \(\mathfrak S\) are
obtained by starting with the corresponding two-sided simple ring over
\(P_0\); the associated division algebra is then isomorphic to \(A\) and
has the conjugate field as its center. Thus the problem of splitting
fields is completely reduced to the associated division algebra \(A\)
and to its splittings.

In particular, the following theorem holds.

All maximal commutative subfields of \(A\) have the same degree \(m\)
over \(P\), where \(m\) is Schur's index. All these maximal commutative
subfields are splitting fields of smallest degree, and every splitting
field of smallest degree is contained in one of them. The degree of
every splitting field is a multiple of \(m\). The algebra \(A\) has only
one absolutely irreducible representation class, and likewise only one
irreducible representation class over \(P\).

To characterize general splitting fields, let \(A_r\) denote the direct
product of \(A\) with the full matrix ring of degree \(r\) over \(P\).
Then \(A_r\) is two-sided simple and has \(A\) as its associated
division algebra. Conversely, every two-sided simple hypercomplex ring
over \(P\) to which \(A\) is associated is obtained in this way.
Equivalently, \(A_r\) may be characterized as the ring of all \(r\) by
\(r\) matrices with entries in \(A\).

Every splitting field of \(A\) of degree \(mr\), if such a field exists,
is isomorphic to a maximal commutative subfield of \(A_r\). Conversely,
all maximal commutative subfields of \(A_r\) are splitting fields, each
of degree \(ms\), where \(s\) divides \(r\). In the regular case all
maximal commutative subfields of \(A_r\) have degree \(mr\). The regular
case means that the ground field \(P\) has the following property: for
every commutative field \(Z\) over \(P\), finite over \(P\), there exist
commutative extensions of \(Z\) of arbitrarily prescribed degree.

Whether, for \(r>1\), minimal fields occur among these splitting fields
is not decided merely by the embedding in \(A_r\). The following example
shows that this can in fact happen for infinitely many \(r\).

\subsection{2. Unboundedness of the degrees of minimal splitting
fields}\label{unboundedness-of-the-degrees-of-minimal-splitting-fields}

For this unboundedness we use the quaternion algebra as example,
regarded as a hypercomplex system over the rational field \(P\). It has
basis elements \(i,j,k\), besides the unit \(1\), with the familiar
relations \[
  i^2=j^2=k^2=-1,
  \qquad ij=k,
  \qquad jk=i,
  \qquad ki=j,
\] with the reversed products changing sign.

We shall show that the splitting fields are exactly those algebraic
number fields in which, after adjoining the field, there exists an
idempotent element \(r\), i.e. \[
  r^2=r,
\] with \(r\) neither zero nor one. These number fields can also be
characterized as those in which \(-1\) is representable as a sum of
three squares, and consequently as a sum of two squares. Indeed, write
\[
  r=a+\frac{\beta}{2}i+
       \frac{\gamma}{2}j+
       \frac{\delta}{2}k.
\] Then \[
  r^2=r
\] is equivalent to the system \[
  4a^2-4a-\beta^2-\gamma^2-\delta^2=0,
  \qquad
  2a\beta-\beta=0,
\] \[
  2a\gamma-\gamma=0,
  \qquad
  2a\delta-\delta=0.
\] Since \(\beta,\gamma,
\delta\) are not all to vanish, this is equivalent to \[
  -1=\beta^2+
\gamma^2+\delta^2.
\] The fact that representability of \(-1\) as a sum of three squares
implies representability as a sum of two squares follows from the
identity \[
  (c^2+d^2)(a^2+b^2+c^2+d^2)
  =(ac+bd)^2+(ad-bc)^2+(c^2+d^2)^2,
\] which is also a consequence of the product theorem for quaternion
norms.

Every such idempotent element produces a splitting, and the existence of
such an idempotent is necessary for splitting. This follows from
standard facts about completely reducible systems. Every irreducible
representation of a completely reducible system is generated by a
one-sided simple ideal, more precisely by the corresponding ideal class,
and every such ideal class generates the representation. These one-sided
simple ideals all occur in one-sided direct decompositions, and those
decompositions determine the primitive idempotent elements, the
components of the identity. Conversely, every idempotent element gives a
decomposition \[
  \mathfrak S=r\mathfrak S+(1-r)\mathfrak S;
\] primitive idempotents split off simple ideals.

A division algebra, regarded as a hypercomplex system over its center
and then extended to a splitting field, becomes a direct sum of \(m\)
absolutely simple one-sided ideals, where \(m\) is Schur's index. For
the quaternion algebra \(m=2\). Hence every idempotent different from
\(0\) and \(1\) already gives the decomposition into absolutely simple
ideals, and this corresponds to the splitting into absolutely
irreducible representation classes. Thus the fields described above are
exactly the splitting fields.

It follows further that all cyclic fields \(K_n\) of degree \(2^n\) over
the rational field in which \(-1\) is representable as a sum of two
squares are minimal splitting fields of the quaternion algebra. Such a
splitting field cannot be real, because \(-1\) is a sum of squares in
it. On the other hand, the subfield of degree \(2^{n-1}\) of \(K_n\),
and hence every proper subfield of \(K_n\), is necessarily real: it is
the intersection of \(K_n\) with the field of all real algebraic
numbers. Thus \(K_n\) is Galois of degree two over this intersection,
which must be the subfield of degree \(2^{n-1}\).

By Hasse's existence theorem there are such fields \(K_n\) for every
\(n\); hence the degrees of the minimal splitting fields are not
bounded. In fact every power of the index occurs here as the degree of a
minimal splitting field.

\subsection{3. Elementary treatment of the splitting fields of the
quaternion
algebra}\label{elementary-treatment-of-the-splitting-fields-of-the-quaternion-algebra}

Let \(\mathfrak S\) denote the representation of the quaternion group
generated by the usual two by two matrices \(I,J,K\) over a field. The
group ring generated by \(\mathfrak S\) is just a representation of the
quaternion algebra as a hypercomplex system. We give an elementary proof
that \(\mathfrak S\) is representable over a field \(K\) if and only if
\(-1\) is representable in \(K\) as a sum of two squares.

Suppose first that \[
  \mathfrak S^*=T^{-1}\mathfrak S T
\] is a representation rational over \(K\). Write \[
  I^*=T^{-1}IT,
  \qquad
  J^*=T^{-1}JT.
\] Since \(J\) and \(J^*\) are similar matrices rational over \(K\),
there is a transformation rational over \(K\) carrying \(J^*\) to \(J\).
Replacing \(\mathfrak S^*\) by an equivalent \(K\)-rational
representation, we may therefore assume \(J^*=J\). Let \[
  I^*=\begin{pmatrix}a&b\\ c&d\end{pmatrix}
\] with \(a,b,c,d\in K\). From \[
  IJ=-JI
\] one obtains \(I^*J=-JI^*\), hence \(a=-d\) and \(b=c\). From
\(I^2=-E\) follows \[
  (I^*)^2=-E,
\] which gives \[
  a^2+b^2=-1.
\] Thus \(-1\) is indeed a sum of two squares in \(K\).

Conversely, suppose that in \(K\) one has \[
  -1=\alpha^2+\beta^2.
\] Set \[
  I^*=\begin{pmatrix}\alpha&\beta\\ \beta&-\alpha\end{pmatrix},
  \qquad
  J^*=\begin{pmatrix}0&-1\\1&0\end{pmatrix},
\] and put \(K^*=I^*J^*\). A direct calculation shows that \[
  \pm E,
  \quad \pm I^*,
  \quad \pm J^*,
  \quad \pm K^*
\] form a group rational over \(K\) similar to the quaternion
representation. This proves the characterization.

It remains to show, elementarily, that for infinitely many \(n\) there
are cyclic fields of degree \(2^n\) which are minimal splitting fields
of the quaternion group. Let \(p\) be a rational prime for which, for
some \(r>0\), \[
  2^r+1\equiv0\pmod p,
\] and let \(2^n\) be the highest power of \(2\) dividing \(p-1\). There
are primes \(p\) belonging to arbitrarily large values of \(n\). For if
\(k\) is any positive integer and \(p\) is a prime divisor of
\(2^{2^k}+1\), then the congruence just required holds with \(r=2^k\).
Moreover \(2\) modulo \(p\) has exponent \(2^{k+1}\), and hence the
highest power \(2^n\) dividing \(p-1\) has \(n>k\).

Let \(\varepsilon\) be a primitive \(p\)-th root of unity. We show that
in \(P(\varepsilon)\), where \(P\) is the rational field, one can
represent \(-1\) as a sum of two squares: \[
  -1=a^2+b^2=(a+ib)(a-ib),
  \qquad a,b\in P(\varepsilon).
\] This is equivalent to saying that in the field of \(4p\)-th roots of
unity, \[
  P(\varepsilon,i)=P(\varepsilon_{4p}),
\] there is an element whose relative norm over \(P(\varepsilon)\) is
exactly \(-1\).

The element \(1+i\varepsilon^{\nu}\) has relative norm \[
  (1+i\varepsilon^{\nu})(1-i\varepsilon^{\nu})
  =1+
\varepsilon^{2\nu}.
\] Thus all numbers \(1+
\varepsilon^{2\nu}\) are relative norms from \(P(\varepsilon,i)\) over
\(P(\varepsilon)\). Hence so is their product \[
  \prod_{
u=0}^{r-1}(1+
\varepsilon^{2\nu}).
\] Using the congruence \(2^r\equiv-1\pmod p\), this product becomes
\(-1\). Also \(\varepsilon\) itself is a relative norm. Hence \(-1\) is
a relative norm, and so \(-1\) is representable in \(P(\varepsilon)\) as
a sum of two squares. Thus \(P(\varepsilon)\) is a splitting field for
the quaternion group.

Let \(K\) be the cyclic subfield of \(P(\varepsilon)\) of degree
\(2^n\). We claim that \(K\) is itself a splitting field. Let \(m\) be
the index of the quaternion algebra with respect to \(K\) as ground
field. Then \(m=1\) or \(m=2\). But \(P(\varepsilon)\) is a splitting
field of odd relative degree over \(K\), so the case \(m=2\) is
impossible. Therefore \(m=1\), which means that \(K\) is a splitting
field.

Thus for infinitely many \(n\) there exist cyclic fields of degree
\(2^n\) which are minimal splitting fields. \# Hypercomplex quantities
and representation theory in arithmetical conception

\subsection{Introduction}\label{introduction}

The most important general theorems on hypercomplex systems go back to
Molien. Almost independently, Frobenius soon afterwards developed the
theory of hypercomplex systems and their representations, especially the
representation theory of finite groups, in a unified way. The basis is
Dedekind's notion of the group determinant, more generally the
determinant of an arbitrary hypercomplex system.

Frobenius shows that the different irreducible factors of this
determinant, with the complex numbers as coefficient field, correspond
to the different irreducible representation classes, and that all
irreducible representation classes are exhausted in this way. Such a
representation class is completely characterized by its character
system. For finite groups, this character system arises by decomposing
the determinant of the commutative hypercomplex system obtained from the
conjugacy classes of the group. That determinant decomposes into linear
factors whose coefficients are the characters. This is a direct
generalization of Dedekind's result that the group determinant of a
finite abelian group decomposes into linear factors whose coefficients
are the several characters of the abelian group.

Two questions are separated. The first is the reduction and
decomposition of a given representation of a given ring over a given,
not necessarily commutative, representation field. The second is the
determination of all possible representations, or even only of all
irreducible representations, of the ring in that field.

The first question is the simpler one. No special assumptions about the
ring to be represented, or about the possibly non-commutative
representation field, are needed. The fact that one deals with matrices
of a fixed degree already supplies the necessary finiteness. The
question can be treated completely by means of the theory of groups with
operators.

A group \(M\) is said to have an operator domain if, together with the
group, there is given a system of symbols \[
  O,H,\ldots
\] so that \(O(a),H(a),\ldots\) are defined uniquely for every
\(a\in M\), again as elements of \(M\), and so that the distributive law
\[
  O(ab)=O(a)O(b)
\] holds. Two groups with the same operator domain are
operator-isomorphic if they are isomorphic as ordinary groups and if,
whenever \(a\) corresponds to \(\bar a\), the elements \(O(a)\) and
\(O(\bar a)\), \(H(a)\) and \(H(\bar a)\), and so on, correspond as
well.

Only subgroups which themselves admit the operator domain are allowed.
Thus a group is simple if it has no proper admissible normal subgroup
except the identity. With this convention the Jordan-Hoelder theorem
remains valid: if a group with operators has a composition series at
all, the system of factor groups is uniquely determined, up to order, in
the sense of operator-isomorphism. Under the same hypothesis there is
also a direct-product theorem: the decomposition into directly
indecomposable factors is unique in the same sense.

The connection with representation theory is that ideals and modules are
examples of groups with operators. They are abelian groups with respect
to addition, restricted by admitting specified multiplications by
elements of a ring. Every representation is generated by a
representation module, that is, by a module with respect to the ring and
the representation field. More precisely, every module class, meaning a
class of mutually operator-isomorphic representation modules, generates
the same representation.

The uniqueness of the reduction of a representation is therefore
answered by the composition-series theorem and the direct-product
theorem. Applied to the representation module, a composition series
gives, for every matrix, a block reduction in which the diagonal blocks
corresponding to the factor groups generate irreducible representations,
in the sense that they allow no further such reduction. Since the
classes of those factor groups are uniquely determined, the diagonal
representations are uniquely determined up to equivalence. Similarly,
the direct-product theorem gives the uniqueness of the splitting into
indecomposable constituents; each constituent can then be reduced
uniquely by its own composition series.

The second question is only sketched here for hypercomplex systems, and
more generally for rings satisfying the double chain condition. The
result is that every irreducible, or simple, module class is an ideal
class. Thus composition series of one-sided ideals, through their factor
groups, already supply all irreducible representations. More exactly, it
is enough to take composition series in the residue class ring modulo
the radical; this quotient becomes a completely reducible ring.

The representation problem is thereby reduced to the structural study of
completely reducible rings. It also becomes natural to start with
representations in the automorphism division rings. All simple right
ideals of a two-sided simple component of such a ring belong to one
class and therefore have the same automorphism ring. Because the ideals
are simple, this automorphism ring is a division ring, and it is at the
same time the automorphism ring of the simple left ideals. The simple
ring itself becomes isomorphic to the full matrix ring over this
automorphism division ring.

All representations over a subfield, and in particular over the
coefficient domain of the hypercomplex system, arise from this matrix
representation by replacing the elements of the automorphism division
ring by the corresponding matrices. In this way Wedderburn's structure
theorem for hypercomplex systems is reinterpreted through the general
group-theoretic notion of the automorphism division ring, and is at the
same time recognized as the source of the representation theory of
hypercomplex systems.

The new questions which appear are those of a Galois theory for
non-commutative division rings, closely connected with the question of
splitting fields. At the same time one obtains a prospect of structure
theorems for more general, not completely reducible, rings through the
automorphism rings of indecomposable ideals. Group theory, extended to
the most general groups with operators, again appears as the organizing
principle.

\section{Hypercomplex quantities and representation
theory}\label{hypercomplex-quantities-and-representation-theory}

\subsection{Introduction}\label{introduction-1}

The general theory just indicated is now developed in a strictly
arithmetical form. Hypercomplex quantities and representation theory are
not treated as separate subjects, but as parts of one theory: the theory
of module classes and ideal classes over non-commutative rings
satisfying suitable finiteness conditions.

The central assertion is that the irreducible module classes are already
exhausted by the corresponding ideal classes. In particular, for rings
without radical, every module class decomposes into irreducible classes.
This is the arithmetical counterpart of Frobenius's theorem that the
irreducible factors of the regular group determinant, or more generally
of the system determinant, already give all irreducible representation
classes, and that for systems without radical complete reducibility
holds.

The determinant and its factorization may be regarded as a passage to
the norm. Here, instead, the direct-sum decomposition of ideals and
their composition series are treated directly. Representation theory
becomes a theory of module classes by means of the representation
module.

The introduction of module and ideal classes is group-theoretic. Modules
and ideals are abelian groups under addition, with additional operators
supplied by multiplication by ring elements. They are therefore groups
with operators. For such groups the ordinary isomorphism is replaced by
operator-isomorphism; groups which are operator-isomorphic are collected
into a class. In the case of ideals of a number field this class notion
coincides with the usual one.

The theory of groups with operators supplies the group-theoretic
foundation. Composition series, direct products or direct sums, and
completely reducible groups all retain their usual uniqueness
properties, interpreted in the sense of operator-isomorphism. These
results will be used first to obtain the general uniqueness theorem for
the irreducible diagonal constituents of arbitrary representations.
Since representation classes correspond one-to-one to the classes of
representation modules, and these are classes of groups with operators,
the assertion follows formally from the group theory.

The further question, the determination of all representations, requires
a closer analysis of the structure of the rings to be represented. In
the case of hypercomplex systems this means an arithmetical
reconstruction and extension of the Wedderburn results, guided by the
study of one-sided ideals. It turns out that the multiple-chain
condition for right ideals, or equivalently the minimal condition for
right ideals, is already sufficient as the finiteness hypothesis.

Rings without radical satisfying this minimal condition are exactly the
completely reducible right rings with identity. In such rings all simple
right ideals of a two-sided indecomposable ring belong to one class;
their automorphism ring is a division ring; and the ring itself is a
full matrix ring over that division ring. From this, the representation
theory and the questions about splitting fields follow in the same
conceptual framework.

\subsection{Chapter I. Group-theoretic
foundations}\label{chapter-i.-group-theoretic-foundations}

\subsection{§1. Groups with operators}\label{groups-with-operators}

Let \(G\) be a group, finite or infinite. An operator domain \(\Omega\)
for \(G\) is a set of new symbols \[
  O,H,\ldots
\] such that for each \(a\in G\) and each \(O\in\Omega\) an element
\(Oa\in G\) is uniquely defined and the distributive law \[
  O(ab)=Oa\cdot Ob
\] is satisfied. Thus every operator defines a homomorphism of the group
into itself, possibly onto a proper subgroup. The identity is carried to
the identity, and inverses to inverses.

An operator domain is called absolute if distinct operators define
distinct homomorphisms. From any operator domain one obtains an absolute
one by identifying all operators which define the same homomorphism. An
absolute operator domain is then in one-to-one correspondence with a
subset of the set of all endomorphisms of the group.

A subgroup \(A\) of \(G\) is admissible, with respect to the fixed
operator domain \(\Omega\), if it is stable under the operators: for
every \(a\in A\) and every \(O\in\Omega\) one has \(Oa\in A\).

As examples, take first the inner automorphisms \[
  a\mapsto c^{-1}ac.
\] The admissible subgroups are precisely the normal subgroups. If all
automorphisms are taken as operators, then the admissible subgroups are
the characteristic subgroups, invariant under every automorphism.

A ring is also an example. A ring is here an abelian group under
addition equipped with multiplication satisfying \[
  r(a+b)=ra+rb,
  \\qquad
  (a+b)r=ar+br,
  \\qquad
  (ab)c=a(bc).
\] Each element \(r\) gives two operators, left multiplication
\(x\mapsto rx\) and right multiplication \(x\mapsto xr\). The admissible
subgroups are then the ideals: left ideals if they are stable under all
left multiplications, right ideals if they are stable under all right
multiplications, and two-sided ideals if they are stable under both.

The ideals become only the trivial ideals when the ring is a division
ring, that is, when the non-zero elements form a group under
multiplication. No commutativity of multiplication is assumed here,
either for rings or for division rings.

Let \(o\) be a ring, and let \(M\) be an abelian group, written
additively. Suppose that a multiplication of elements of \(o\) with
elements of \(M\) is given on the left, and that \[
  r(a+b)=ra+rb,
  \qquad
  (rs)a=r(sa),
  \qquad
  (r+s)a=ra+sa
\] for \(r,s\in o\) and \(a,b\in M\). Then \(M\) is a left \(o\)-module.
The ring \(o\) is at the same time an operator domain. If different
elements of \(o\) induce different operations, the multiplier domain is
absolute. As before, one passes from an arbitrary multiplier domain to
the absolute one by identifying elements which induce the same
operation; the absolute multiplier domain is again a ring. The
admissible subgroups of a module are its submodules. Taking \(M=o\)
gives the left ideals.

Right modules are defined analogously by the rules \[
  (a+b)r=ar+br,
  \qquad
  a(rs)=(ar)s,
  \qquad
  a(r+s)=ar+as.
\] If \(M\) is simultaneously a left module over \(o\) and a right
module over \(o'\), and if \[
  r(ar')=(ra)r'
\] for \(r\in o\), \(a\in M\), \(r'\in o'\), then \(M\) is a bimodule.

Finally, for an abelian group \(M\) the endomorphisms of \(M\) form a
ring. Addition and multiplication are defined by \[
  (H+O)a=Ha+Oa,
  \qquad
  (HO)a=H(Oa).
\] With respect to this ring, the abelian group itself is a module. From
now on, when subgroups are mentioned in the context of groups with
operators, only admissible subgroups are meant. The intersection of two
admissible subgroups is admissible. The product of two such subgroups is
admissible whenever at least one of them is normal.

\subsection{§2. The isomorphism
theorems}\label{the-isomorphism-theorems}

Let two groups have the same operator domain. A map from \(G\) onto
\(\bar G\) is an operator-homomorphism if it is a homomorphism in the
ordinary sense and, moreover, if \(a\) maps to \(\bar a\) then \(Oa\)
maps to \(O\bar a\) for every operator \(O\). If the correspondence is
one-to-one, it is an operator-isomorphism.

The homomorphism theorem has the same form as for ordinary groups. If
\(\bar G\) is an operator-homomorphic image of \(G\), then \(\bar G\) is
operator-isomorphic to a factor group \(G/N\), where \(N\) is the
admissible normal subgroup consisting of the elements of \(G\) which map
to the identity of \(\bar G\). Conversely, every admissible normal
subgroup \(N\) defines a factor group \(G/N\) which admits the same
operator domain and is an operator-homomorphic image of \(G\).

A group is simple if it has no normal subgroups except itself and the
identity group. Every homomorphism of a simple group is either an
isomorphism or sends every element to the identity.

The first isomorphism theorem says: if \(\bar G\) is a homomorphic image
of \(G\), if \(A\) is a normal subgroup of \(\bar G\), and if \(B\) is
the set of elements of \(G\) which map into \(A\), then \(B\) is a
normal subgroup and \[
  \bar G/A\simeq G/B.
\] In the notation of the homomorphism theorem, if \(\bar G=G/N\), then
\(A\) necessarily contains the image of \(N\), and one recovers \(A\)
from its inverse image. The formula can be written as \[
  (G/N)/(B/N)\simeq G/B.
\]

The second isomorphism theorem says: if \(A\) is a subgroup and \(B\) is
a normal subgroup of \(G\), then \(A\cap B\) is normal in \(A\), and \[
  AB/B\simeq A/(A\cap B).
\] Indeed, under the homomorphism \(G\to G/B\), the elements of \(A\) go
to the cosets making up \(AB/B\), and the kernel is \(A\cap B\).

An operator-homomorphic map of a group into itself, or into a subgroup
of itself, will be called an endomorphism of the group. If the group is
abelian and written additively, the operator-endomorphisms form a ring
under the addition and multiplication just described; this ring is the
automorphism ring in Noether's terminology.

The automorphism ring of a simple abelian group with operators is a
division ring. For every endomorphism either maps the group onto itself
or maps it to the zero group. The non-zero endomorphisms are
isomorphisms, and the isomorphisms of the group onto itself form a group
under multiplication. Hence the ring with the zero operator omitted is a
group, which is exactly the assertion that the ring is a division ring.

If \(G\) is a left \(o\)-module and the operator-endomorphisms are
written as right operators, then \(G\) becomes a bimodule over \(o\) on
the left and over its automorphism ring on the right. The identities \[
  (a+b)T=aT+bT,
  \qquad
  (ra)T=r(aT)
\] state precisely that \(T\) is a homomorphism with respect to the left
operators. Conversely, in any bimodule the same formulas say that each
right multiplier induces an operator-endomorphism relative to the left
multipliers.

\subsection{§3. Composition series}\label{composition-series}

A finite chain of admissible subgroups \[
  G=A_0\supset A_1\supset\cdots\supset A_r=E
\] is a composition series if each \(A_i\) is normal in the preceding
subgroup and no further subgroup with the same property can be inserted
between two consecutive terms. The factor groups \[
  A_0/A_1,
  A_1/A_2,
  \ldots,
  A_{r-1}/A_r
\] are the composition factors. They are simple groups.

If all inner automorphisms are included among the operators, such a
series consists entirely of normal subgroups of \(G\) and is called a
principal series. If all automorphisms are included, it is called a
characteristic series. In the examples above one obtains composition
series of ideals, modules, or bimodules.

The Jordan-Hoelder theorem remains valid in this setting. If a group
with operators has two composition series, then the two series have the
same length, and their composition factors are pairwise
operator-isomorphic after a suitable rearrangement. Moreover, if a
composition series exists, then one can draw such a series through every
admissible normal subgroup.

The existence of a composition series is automatic for finite groups,
and more generally for groups satisfying both the maximal and minimal
conditions. The maximal condition says that every set of admissible
subgroups contains a maximal member, equivalently that every ascending
chain \[
  A_1\subset A_2\subset A_3\subset\cdots
\] terminates. The minimal condition says that every set of admissible
subgroups contains a minimal member, equivalently that every descending
chain terminates.

When only normal subgroups are admissible, as for abelian groups, the
existence of a composition series conversely implies both chain
conditions. Each normal subgroup has a composition series, and its
length is the length of the normal subgroup. If \(A\subset B\), the
length of \(A\) is smaller than the length of \(B\), since a composition
series for \(B\) can be drawn through \(A\). Hence a subgroup of
shortest length is minimal, and one of greatest length is maximal.

\subsection{§4. Direct products and direct
sums}\label{direct-products-and-direct-sums}

A group \(G\) is the direct product of two factors, written \[
  G=A\times B,
\] if \(A\) and \(B\) are normal subgroups, \(AB=G\), and \(A\cap B=E\).
Equivalently, every element \(g\in G\) has a unique expression \[
  g=ab,
  \qquad a\in A,\ b\in B,
\] and the elements of \(A\) commute with the elements of \(B\).

For finitely many factors one writes \[
  G=A_1\times\cdots\times A_n.
\] This means that, for each \(i\), \(G\) is the direct product of
\(A_i\) and the product of all remaining factors. Equivalently, every
element has a unique representation \[
  g=a_1a_2\cdots a_n,
  \qquad a_i\in A_i,
\] and the elements of different factors commute with one another.

A group is directly indecomposable if it cannot be written as a direct
product of two proper factors. Every group satisfying the minimal
condition is a direct product of finitely many directly indecomposable
groups. In additive notation, products become sums and direct products
become direct sums.

The direct-product theorem used later is a uniqueness theorem in the
sense of operator-isomorphism. Under the chain hypotheses stated above,
the decomposition into directly indecomposable factors is unique up to
order and operator-isomorphism. This is the additive form that will be
applied to ideals and modules.

\clearpage
\section{34. Hypercomplex Quantities and Representation
Theory}\label{hypercomplex-quantities-and-representation-theory}

\subsection{Continuation of §4. Direct products and direct
intersections}\label{continuation-of-4.-direct-products-and-direct-intersections}

A group \(D\) is called the \textbf{direct intersection} of two groups
\(R\) and \(S\) with respect to \(G\) if

\begin{enumerate}
\def\labelenumi{\alph{enumi})}
\item
  \(R\) and \(S\) are normal subgroups of \(G\),
\item
  \(R\cap S=D\),
\item
  \(RS=G\).
\end{enumerate}

Similarly, \(D\) is called the direct intersection of \(n\) groups
\(S_1,\ldots,S_n\) if, after putting \[
  R_i=S_1\cap\cdots\cap S_{i-1}\cap S_{i+1}\cap\cdots\cap S_n,
\] one has \[
  D=R_i\cap S_i
\] directly for each \(i\).

It follows at once from the definition that \(D\) must be a normal
subgroup of \(G\).

If \[
  D=S_1\cap\cdots\cap S_n
\] is direct and, under the homomorphism \(G\to G/D\), the groups
\(G,S_i,D\) go over into \(\overline G,\overline S_i,E\), then \[
  E=\overline S_1\cap\cdots\cap\overline S_n
\] is direct. Conversely, from every such representation \[
  E=\overline S_1\cap\cdots\cap\overline S_n
\] one returns to a direct representation \[
  D=S_1\cap\cdots\cap S_n.
\] This correspondence reduces theorems on direct intersections with
arbitrary \(D\) to the special case \(D=E\).

For \(D=E\) and two factors, the conditions for direct product and
direct intersection coincide. For \(n\) factors there is the following
one-to-one relation between direct product and direct intersection.

\textbf{I.} If \[
  G=A_1\times\cdots\times A_n=A_i\times B_i,
\] then \[
  E=B_1\cap\cdots\cap B_n=B_i\cap C_i
\] directly, and \(C_i=A_i\).

\textbf{II.} If \[
  E=S_1\cap\cdots\cap S_n=R_i\cap C_i
\] directly, then \[
  G=R_1\times\cdots\times R_n=R_i\times T_i,
\] and \(T_i=S_i\).

For I, write \[
  B=B_1\cap\cdots\cap B_n=B_i\cap C_i.
\] Then \(C_i\supseteq A_i\); we show that \(C_i\subseteq A_i\). If, for
example, \(c\in C_1\), then \(c\) belongs to \(B_2,\ldots,B_n\), and its
component representation with respect to the factors \(A_i\) gives \[
  c=a_1a_2\cdots a_n
   =a_1ea_3\cdots a_n
   =a_1a_2e\cdots a_n
   =\cdots=a_1\cdots a_{n-1}e.
\] Since the product of all the \(A_i\) is direct, these representations
agree. At every position except the first there appears once the
identity element. Hence \(c=a_1\), and so \(C_1\subseteq A_1\). Thus
\(C_i=A_i\) for all \(i\). It follows that \[
  B=B_i\cap A_i=E,
  \qquad B_iA_i=G,
\] so that the intersection is direct.

For II, write \[
  H=R_1\cdots R_n=R_iT_i.
\] Then \(T_i\subseteq S_i\); we show that \(S_i\subseteq T_i\). Let,
for instance, \(s\in S_1\). Since \(G=S_1\times R_1\), one may write \[
  s=s_1e=s_2r_2=\cdots=s_nr_n.
\] Put \[
  t_1=er_2\cdots r_n=t_ir_i.
\] Using the elementwise commutativity of \(R_i\) and \(S_i\) one
obtains \[
  s t_1^{-1}=s_it_i^{-1},
\] an element of all the \(S_i\), hence the identity. Thus
\(S_1\subseteq T_1\), and so \(T_1=S_1\); similarly \(T_i=S_i\) for
every \(i\). Consequently \[
  H=R_iT_i=R_iS_i=G,
  \qquad R_i\cap T_i=R_i\cap S_i=E,
\] and therefore \(G\) is the direct product of the \(R_i\).

\subsection{§5. Completely reducible
groups}\label{completely-reducible-groups}

A group is called \textbf{completely reducible} if it is the direct
product of finitely many simple groups: \[
  G=A_1\times\cdots\times A_n.
\] In this case the groups \[
  G,
  A_2\times\cdots\times A_n,
  A_3\times\cdots\times A_n,
  \ldots,
  A_n,
  E
\] form a composition series. Its length is \(n\), and its composition
factors are \[
  A_1,A_2,\ldots,A_n.
\]

If \(G\) is completely reducible, then every normal subgroup is a direct
factor. The complementary factor can be chosen as the product of simple
groups occurring in a prescribed product decomposition of \(G\). The
normal subgroup itself is completely reducible.

Indeed, suppose \[
  G=H A_1A_2\cdots A_n.
\] Put \(H_0=H\) and \(H_i=H_{i-1}A_i\). Either
\(A_i\subseteq H_{i-1}\), in which case the factor \(A_i\) may be
omitted, or \(H_{i-1}\cap A_i\) is a proper normal subgroup of \(A_i\)
and hence equals \(E\); in that case \[
  H_i=H_{i-1}\times A_i.
\] After omitting the superfluous factors the decomposition becomes \[
  G=H\times A_{i_1}\times\cdots\times A_{i_s},
\] so \(H\) is a direct factor. Moreover \[
  H\simeq G/(A_{i_1}\times\cdots\times A_{i_s}),
\] which exhibits \(H\) as completely reducible.

It follows that every decomposition of a completely reducible group into
directly indecomposable factors is already a decomposition into simple
factors. It therefore gives a composition series, and the factors are
uniquely determined, up to isomorphism and order.

\subsection{§6. Modules relative to a field. Hypercomplex
systems}\label{modules-relative-to-a-field.-hypercomplex-systems}

An almost trivial example of completely reducible groups with operators
is furnished by modules with finite basis relative to a field, not
necessarily commutative.

Let \(G\) be a right \(K\)-module, and assume that the unit element
\(e\) of \(K\) is at the same time the identity operator: \(ae=a\) for
\(a\in G\). Every submodule \(aK\) generated by a non-zero element \(a\)
is simple; indeed it is equal to the module generated by any one of its
non-zero elements. Thus, if \(H\) is the submodule generated by certain
elements and if \(a\notin H\), then \[
  H\cap aK=0,
\] and \(H+aK\) is direct. Starting with a basis element \(a_1\) one may
adjoin further basis elements \(a_2,a_3,\ldots\), obtaining \[
  G=a_1K+a_2K+\cdots+a_nK
\] as a direct sum. Hence \(G\) is completely reducible.

The number \(n\), the length of the composition series, is called the
\textbf{rank} of \(G\) relative to \(K\). Because the representation of
the elements is unique, and because \(e\) is assumed to be the identity
operator, the basis \(a_1,\ldots,a_n\) is linearly independent.

Every submodule \(H\) is a direct summand: \[
  G=H+R=(a_1K+\cdots+a_rK)+(a_{r+1}K+\cdots+a_nK).
\] Equivalently, every linearly independent basis of \(H\) can be
extended to a linearly independent basis of \(G\). The complementary
basis elements may even be chosen from the original basis elements.

If \((a_1,\ldots,a_n)\) and \((c_1,\ldots,c_n)\) are two linearly
independent bases, then \[
  c_j=\sum_i a_i p_{ij},
\] or, in matrix notation, \[
  (c_1,\ldots,c_n)=(a_1,\ldots,a_n)P.
\] Conversely, \[
  (a_1,\ldots,a_n)=(c_1,\ldots,c_n)Q.
\] It follows that \[
  (c_1,\ldots,c_n)=(c_1,\ldots,c_n)QP,
\] and, since the \(c_i\) are linearly independent, \[
  QP=E.
\] The corresponding relation gives \(PQ=E\) as well.

Special \(K\)-modules which are at the same time rings are the
\textbf{hypercomplex systems}.

A ring \(\mathfrak{o}\) which is also a right module relative to a
commutative field \(K\) is called a hypercomplex system relative to
\(K\) (also called, in the literature, an algebra over \(K\)) if the
following hold.

\begin{enumerate}
\def\labelenumi{\arabic{enumi}.}
\item
  The rank is finite; let \(u_1,\ldots,u_n\) be a linearly independent
  basis.
\item
  Elements of \(K\) commute with multiplication in \(\mathfrak{o}\): \[
    (ab)x=a(bx)=(ax)b.
  \] One says that \(K\) is connected commutatively with
  \(\mathfrak{o}\).
\item
  The unit element \(e\) of \(K\) is the identity operator: \[
    ae=a \qquad (a\in\mathfrak{o}).
  \]
\end{enumerate}

Because \[
  \left(\sum_i u_i\alpha_i\right)
  \left(\sum_j u_j\beta_j\right)
  =\sum_{i,j} (u_i u_j)(\alpha_i\beta_j),
\] the system is completely determined, besides \(K\), by the
multiplication table \[
  u_i u_j=\sum_k u_k\gamma_{ij}^{k}.
\] The coefficients \(\gamma_{ij}^{k}\) have to satisfy the familiar
relations imposed by associativity.

If, as will often be assumed, \(\mathfrak{o}\) has a unit element \(1\)
with \(1a=a1=a\) for all \(a\), then the elements \(x\in K\) may be
identified with \(1x\). Thus \(K\) may be regarded as a subfield of
\(\mathfrak{o}\). In that case the hypercomplex system may be described
as a ring of finite rank over a field lying in the center.

One example is the group ring of a finite group: take the group elements
as basis elements \(u_i\), and define multiplication by the group
multiplication. The field \(K\) is arbitrary.

\subsection{§7. Matrices}\label{matrices}

The square matrices \((x_{ij})\) with entries in a ring \(\mathfrak{o}\)
form a ring under the usual matrix multiplication and addition.

A matrix with entries in a field \(K\) is called \textbf{regular} if it
has both a right inverse and a left inverse.

In §6 we saw that the transition matrix between two linearly independent
bases of a \(K\)-module is regular. For regularity it is enough to have
a right inverse. To see this, form a right \(K\)-module with linearly
independent basis \((a_1,\ldots,a_n)\), and put \[
  c_j=\sum_i a_i p_{ij}.
\] If \(P\) has a right inverse, then the \(a_i\) may be expressed
through the \(c_j\); hence the \(c_j\) are again a basis, and in
particular linearly independent. Thus \(P\) is regular. The right
inverse is then also the left inverse. The same argument applies if a
left inverse is given first. If \(P\) has a left inverse, \(P\) is not a
left zero divisor: from \(PA=0\) one obtains \(A=P^{-1}PA=0\).

Let \(P_{ij}\) denote, as usual, the matrix whose \((i,j)\) entry is
\(x\), all other entries being zero; and let \(E_{ij}\) denote the
matrix which has the unit element at position \((i,j)\) and zeros
elsewhere. Then \[
  (x_{ij})=\sum E_{ij}x_{ij}.
\] Thus the matrices over \(K\) form a \(K\)-module of rank \(n^2\). The
matrix units satisfy \[
  E_{ij}E_{jk}=E_{ik},
  \qquad
  E_{ij}E_{hk}=0\quad (j\ne h).
\] Consequently the matrices over \(K\) form a finite \(K\)-module and,
when \(K\) is commutative, a hypercomplex system.

\subsection{Chapter II. Non-commutative ideal
theory}\label{chapter-ii.-non-commutative-ideal-theory}

\subsection{§8. The homomorphism theorem for
rings}\label{the-homomorphism-theorem-for-rings}

If \(\mathfrak a\) is a two-sided ideal of \(\mathfrak{o}\), then the
residue class domain \(\mathfrak{o}/\mathfrak a\) is not only an
\(\mathfrak{o}\)-module but also a ring. Every homomorphic image of
\(\mathfrak{o}\) is again a ring and is isomorphic to a residue class
ring \(\mathfrak{o}/\mathfrak a\), where \(\mathfrak a\) is a two-sided
ideal. The proof is the same as for groups.

In particular, if \(\mathfrak{o}\) is a division ring, then there are no
ideals other than \((0)\) and \(\mathfrak{o}\); hence a homomorphism is
either an isomorphism or sends every element to zero.

If the ring \(\mathfrak{o}\) is a right \(K\)-module and \(K\) is a ring
which commutes elementwise with \(\mathfrak{o}\), \[
  (ab)x=a(bx)=(ax)b,
\] for example if \(\mathfrak{o}\) is a hypercomplex system relative to
a commutative field \(K\), then only those right, left, or two-sided
ideals are admitted which are at the same time \(K\)-modules. Besides
ring homomorphisms, one demands operator homomorphisms with respect to
multiplication by \(K\). The admissible ideals are bimodules relative to
\(\mathfrak{o}\) and \(K\). In hypercomplex systems, the word ``ideal''
by itself always means such an ideal, admissible relative to the
underlying coefficient field \(K\). With this enlargement, the
homomorphism theorem just stated remains valid.

The passage from a multiplier domain \(\mathfrak{o}\) to the absolute
multiplier domain (§1, Example 4) is an example of ring homomorphism.
Thus the absolute multiplier domain is isomorphic to a residue class
ring \(\mathfrak{o}/\mathfrak d\), where \(\mathfrak d\) is the
two-sided ideal of all elements of \(\mathfrak{o}\) which annihilate the
module \(M\).

From the homomorphism theorem one obtains, as in §2, the first and
second isomorphism theorems for ring isomorphism and
operator-isomorphism.

All products \(AB\) of subsets of \(\mathfrak{o}\) are to be understood
as module products: \(AB\) is the set of all sums of products \(ab\)
with \(a\in A\), \(b\in B\); and \(aB\) is the set of all sums \(ab\)
with \(b\in B\). If \(B\) is a right module with respect to any ring as
operator domain, then \(AB\), respectively \(aB\), is again a right
module. In particular it may be a right ideal, an admissible ideal
relative to a coefficient domain, and so on. If \((A,B)\) denotes the
sum of the modules \(A\) and \(B\), the calculation rules are \[
  A(B,C)=(AB,AC),
  \qquad
  (A,B)C=(AC,BC).
\] The direct sum is denoted by \(A+B\).

In what follows we often consider rings satisfying a maximal or a
minimal condition for right ideals. These include, in particular, the
hypercomplex systems, since by convention only \(K\)-modules are
admitted as ideals, and their ranks are therefore bounded.

\subsection{§9. Idempotent elements. Direct-sum decomposition into right
ideals}\label{idempotent-elements.-direct-sum-decomposition-into-right-ideals}

Let \(\mathfrak{o}\) be a ring with unit element.

If \(\mathfrak r\) is a right ideal, then
\(\mathfrak r\mathfrak{o}\subseteq\mathfrak r\); because there is a unit
element, in fact \(\mathfrak r\mathfrak{o}=\mathfrak r\). Similarly for
left ideals, \(\mathfrak{o}\mathfrak l=\mathfrak l\).

An element \(e\) is called \textbf{idempotent} if \(e^2=e\).

Let \[
  \mathfrak{o}=\mathfrak r_1+\cdots+\mathfrak r_s
\] be a decomposition into right ideals, and let \[
  1=e_1+\cdots+e_s,
  \qquad e_i\in\mathfrak r_i.
\] Then \[
  e_i e_k=0\quad(i\ne k),
  \qquad
  \mathfrak r_i=e_i\mathfrak{o}.
\] These are the orthogonality relations.

For if \(r\in\mathfrak r_i\), then \[
  r=re=r e_1+\cdots+r e_s,
\] while also \[
  r=0+\cdots+0+r+0+\cdots+0.
\] Hence \(re_k=0\) for \(k\ne i\), and \(re_i=r\). The first formula
gives \(\mathfrak r_i\subseteq e_i\mathfrak{o}\), while the reverse
inclusion is clear, so \(\mathfrak r_i=e_i\mathfrak{o}\). Taking
\(r=e_i\) gives \[
  e_i^2=e_i,
  \qquad e_i e_k=0\quad(i\ne k).
\]

Conversely, if \[
  1=e_1+\cdots+e_s,
  \qquad e_i^2=e_i,
  \qquad e_i e_k=0\ (i\ne k),
\] and if \(\mathfrak r_i=e_i\mathfrak{o}\), then \[
  \mathfrak{o}=\mathfrak r_1+\cdots+\mathfrak r_s
\] directly. Every element \(r\) of \(\mathfrak{o}\) is \[
  r=e_1r+\cdots+e_sr.
\] The representation is unique, since from \[
  0=e_1a_1+\cdots+e_sa_s
\] multiplication by \(e_i\) gives \(e_i a_i=0\).

From a right decomposition \[
  \mathfrak{o}=e_1\mathfrak{o}+\cdots+e_s\mathfrak{o}
\] one obtains a left decomposition \[
  \mathfrak{o}=\mathfrak{o}e_1+\cdots+\mathfrak{o}e_s.
\] If the right ideals \(\mathfrak r_i\) are directly indecomposable,
then so are the corresponding left ideals
\(\mathfrak l_i=\mathfrak{o}e_i\); otherwise a decomposition of some
\(\mathfrak l_i\) would give a corresponding decomposition of
\(\mathfrak r_i\) by the same orthogonal idempotent calculation.

There is also the following converse form. If \(e_i e_k=0\) for
\(i\ne k\), if \(e_i^2=e_i\), and if \(e=e_1+\cdots+e_s\), then \(e\) is
a left identity for the ring \[
  e_1\mathfrak{o}+\cdots+e_s\mathfrak{o}.
\] The directness of the sum is proved as above.

\subsection{§10. Decomposition into two-sided
ideals}\label{decomposition-into-two-sided-ideals}

Assume again that \(\mathfrak{o}\) is a ring with unit element, and that
\[
  \mathfrak{o}=\mathfrak a_1+\cdots+\mathfrak a_s
\] is a decomposition of \(\mathfrak{o}\) into two-sided ideals. Then
the orthogonality relations hold also for the ideals: \[
  \mathfrak a_i\mathfrak a_k=0\quad(i\ne k),
  \qquad
  \mathfrak a_i^2=\mathfrak a_i.
\] Indeed, if \(\mathfrak b_i\) is the sum of all \(\mathfrak a_j\) with
\(j\ne i\), then \(\mathfrak{o}=\mathfrak a_i+\mathfrak b_i\) and \[
  \mathfrak a_i\mathfrak b_i\subseteq\mathfrak a_i\cap\mathfrak b_i=0.
\] A fortiori \(\mathfrak a_i\mathfrak a_k=0\) for \(i\ne k\), and \[
  \mathfrak a_i\mathfrak{o}=\mathfrak a_i(\mathfrak a_i+\mathfrak b_i)=\mathfrak a_i^2.
\] Since \(\mathfrak a_i\mathfrak{o}=\mathfrak a_i\), this gives
\(\mathfrak a_i^2=\mathfrak a_i\).

Conversely, if a direct sum decomposition satisfies these orthogonality
relations, then each summand \(\mathfrak a_i\) is a two-sided ideal.
Right ideals in the ring \(\mathfrak a_i\) are right ideals in
\(\mathfrak{o}\).

If \(\mathfrak r\) is a right ideal of \(\mathfrak{o}\), then
\(\mathfrak r\) is the direct sum of the right ideals \[
  \mathfrak r\mathfrak a_1,\ldots,\mathfrak r\mathfrak a_s.
\] For \[
  \mathfrak r=\mathfrak r\mathfrak{o}=(\mathfrak r\mathfrak a_1,\ldots,\mathfrak r\mathfrak a_s),
\] and the orthogonality relations make the sum direct. If
\(\mathfrak r\) is directly indecomposable, it can have only one
component, and therefore must be contained in one of the
\(\mathfrak a_i\).

These facts reduce the ideal theory of \(\mathfrak{o}\) to the ideal
theory of the individual \(\mathfrak a_i\). For this reason we shall
usually restrict ourselves to rings which are two-sided indecomposable.

Two right ideals contained in different \(\mathfrak a_i\) can never be
operator-isomorphic. An ideal contained in \(\mathfrak a_i\) is
annihilated by all the other \(\mathfrak a_j\), but not by
\(\mathfrak a_i\); the opposite holds for an ideal contained in
\(\mathfrak a_j\).

If a ring can be decomposed as \[
  \mathfrak{o}=\mathfrak a_1+\cdots+\mathfrak a_s
\] with the \(\mathfrak a_i\) two-sided indecomposable, and if each
\(\mathfrak a_i\) is decomposed into one-sided indecomposable right
ideals, then operator-isomorphic right ideals are to be looked for only
among those belonging to the same \(\mathfrak a_i\). It may nevertheless
happen that there are several classes of non-isomorphic right ideals
within the same \(\mathfrak a_i\).

For example, let \(K\) be the field of rational numbers and let \[
  \mathfrak{o}=e_1K+e_2K+uK
\] be a hypercomplex system in which \(e_1,e_2,u\) commute with the
elements of \(K\), and multiplication is given by \[
\begin{array}{c|ccc}
       & e_1&e_2&u\\ \hline
 e_1   & e_1&0&u\\
 e_2   & 0&e_2&0\\
 u     & u&0&0
\end{array}
\] The unit element is \(e=e_1+e_2\). The \(e_i\) satisfy the
orthogonality relations. Thus the right decomposition is \[
  \mathfrak{o}=e_1\mathfrak{o}+e_2\mathfrak{o}=(e_1,u)+(e_2),
\] and the left decomposition is \[
  \mathfrak{o}=\mathfrak{o}e_1+\mathfrak{o}e_2=(e_1)+(e_2,u).
\] Since \(e_1\mathfrak{o}\) and \(\mathfrak{o}e_1\) are visibly
indecomposable, the corresponding opposite one-sided ideals are
indecomposable as well. A two-sided decomposition is impossible, since
one component would have to contain \(e_1\mathfrak{o}\) and the other
\(e_2\mathfrak{o}\); then the first would contain \(u\), but the second
would contain \(ue_1=u\) as well. Yet \(e_1\mathfrak{o}\) and
\(e_2\mathfrak{o}\) are not operator-isomorphic, since they have
different ranks over \(K\).

Now suppose \[
  \mathfrak{o}=\mathfrak a_1+\cdots+\mathfrak a_s
\] with the \(\mathfrak a_i\) two-sided indecomposable. Then the
\(\mathfrak a_i\) are uniquely determined. Indeed, suppose also \[
  \mathfrak{o}=\mathfrak c_1+\cdots+\mathfrak c_t.
\] Then \[
  \mathfrak a_i=\mathfrak a_i\mathfrak{o}=(\mathfrak a_i\mathfrak c_1,\ldots,\mathfrak a_i\mathfrak c_t),
\] and this sum is direct, since
\(\mathfrak a_i\mathfrak c_j\subseteq \mathfrak c_j\). Since
\(\mathfrak a_i\) is indecomposable, all these summands vanish except
one. Thus \(\mathfrak a_i=\mathfrak a_i\mathfrak c_j\). Symmetrically
each \(\mathfrak c_j\) contains precisely one of the \(\mathfrak a_i\),
and the two decompositions coincide up to order.

\subsection{§11. The center}\label{the-center}

The \textbf{center} \(Z\) of a ring \(\mathfrak{o}\) is the set of all
elements \(z\) commuting with every ring element: \[
  za=az\qquad(a\in\mathfrak{o}).
\] It is a commutative ring.

Every one-sided ideal \(\mathfrak a\) of \(\mathfrak{o}\) has a
contracted ideal \(\mathfrak a\cap Z\) in \(Z\), since both
\(\mathfrak a\) and \(Z\) are \(Z\)-modules.

Every ideal \(A\) in \(Z\) has an extended ideal: the ideal generated by
\(A\) in \(\mathfrak{o}\), \[
  \mathfrak a=(A,\mathfrak{o}A).
\] This ideal is two-sided, since \[
  \mathfrak a\mathfrak{o}=(\mathfrak{o}A,A\mathfrak{o})\mathfrak{o}
      = (\mathfrak{o}A\mathfrak{o},A\mathfrak{o}^2)
      \subseteq (\mathfrak{o}A,A\mathfrak{o})=\mathfrak a.
\] When \(\mathfrak{o}\) has a unit element one simply writes
\(\mathfrak a=\mathfrak{o}A\).

In that case the following theorem holds. Every two-sided decomposition
of \(\mathfrak{o}\) corresponds one-to-one to a decomposition of the
center. From \[
  \mathfrak{o}=\mathfrak a_1+\cdots+\mathfrak a_s,
  \qquad
  A_i=\mathfrak a_i\cap Z
\] one obtains \[
  Z=A_1+\cdots+A_s,
  \qquad
  \mathfrak a_i=\mathfrak{o}A_i,
\] and conversely.

For let \(z\in Z\) and write its decomposition with respect to the
two-sided decomposition of \(\mathfrak{o}\): \[
  z=z_1+\cdots+z_s.
\] Then the \(z_i\) are again central. Indeed, for all \(a\), \[
  za=z_1a+\cdots+z_sa=az=az_1+\cdots+az_s,
\] and comparison of components gives \(z_i a=a z_i\), since the
\(\mathfrak a_i\) are two-sided. Hence
\(z_i\in\mathfrak a_i\cap Z=A_i\), and the above is the asserted
direct-sum decomposition of \(Z\). In particular, if \[
  1=e_1+\cdots+e_s
\] is the decomposition of the unit, then the \(e_i\) are central
elements. Finally \[
  \mathfrak{o}A_i=\mathfrak{o}Ze_i=\mathfrak{o}e_i=\mathfrak a_i.
\]

Conversely, let \[
  Z=A_1+\cdots+A_s
\] be a decomposition of the center. By §9 the components of the
identity satisfy \[
  1=e_1+\cdots+e_s,
  \qquad e_i^2=e_i,
  \qquad e_ie_j=0\ (i\ne j).
\] The orthogonality relations give \[
  \mathfrak{o}=\mathfrak{o}e_1+\cdots+\mathfrak{o}e_s.
\] Putting \(\mathfrak a_i=\mathfrak{o}A_i\) gives, by the first part of
the proof, \[
  A_i=\mathfrak a_i\cap Z.
\] It follows that \(\mathfrak a_i=\mathfrak{o}e_i\), and the
decomposition is the desired two-sided decomposition.

Consequently \(\mathfrak{o}\) and \(Z\) are at the same time two-sided
directly decomposable, or not.

\subsection{§12. Nilpotent ideals}\label{nilpotent-ideals}

An ideal \(\mathfrak c\) is called \textbf{nilpotent} if
\(\mathfrak c^r=0\) for some \(r\).

The ideal \((u)\) in the example of §10 is nilpotent.

If \(\mathfrak{o}\) contains a nilpotent right ideal \(\mathfrak r\),
then it contains a nilpotent left ideal, indeed a two-sided nilpotent
ideal. If \(\mathfrak r^r=0\), then \(\mathfrak{o}\mathfrak r\) is a
left ideal; if \(\mathfrak{o}\mathfrak r\) could be zero because
\(\mathfrak{o}\) has no unit element, take instead
\((\mathfrak{o}\mathfrak r,\mathfrak r)\). This ideal is nilpotent.

The sum of two nilpotent right ideals is again nilpotent. If
\(\mathfrak c^r=0\) and \(\mathfrak d^s=0\), then in \[
  (\mathfrak c,\mathfrak d)^{r+s-1}
\] every term contains either at least \(r\) factors from
\(\mathfrak c\) or at least \(s\) factors from \(\mathfrak d\), and so
vanishes. Hence the whole power is zero.

If the maximal condition for right ideals holds in \(\mathfrak{o}\),
there is therefore a maximal nilpotent right ideal. It contains all
other nilpotent right ideals, for otherwise one could form the sum with
such an ideal. Call it \(\mathfrak c\). Since
\(\mathfrak{o}\mathfrak c\) is again a nilpotent right ideal,
\(\mathfrak{o}\mathfrak c\subseteq\mathfrak c\), and therefore
\(\mathfrak c\) is a left ideal as well. Every nilpotent left ideal
\(\mathfrak d\) is contained in \(\mathfrak c\), because
\((\mathfrak d,\mathfrak d\mathfrak{o})\) is a nilpotent right ideal.
Thus there is a maximal two-sided nilpotent ideal \(\mathfrak c\), the
\textbf{radical}, which contains all nilpotent right and left ideals.

In the example of §10, \((u)\) is the maximal nilpotent ideal.

If the radical is the zero ideal, one speaks of a \textbf{ring without
radical}: equivalently, a semisimple ring, or a Dedekind system. The
residue class ring modulo the radical is always a ring without radical.

Let \(Z\) be the center of \(\mathfrak{o}\) and \(\mathfrak c\) the
radical of \(\mathfrak{o}\). Then \[
  Z\cap\mathfrak c
\] is the radical of \(Z\). It is plainly nilpotent. If there were a
larger nilpotent ideal \(C\) in \(Z\), it would extend to a nilpotent
ideal of \(\mathfrak{o}\) not wholly contained in \(\mathfrak c\): \[
  (C\mathfrak{o})^r=C^r\mathfrak{o}^r=0.
\] This is impossible by maximality of \(\mathfrak c\).

It follows that if \(\mathfrak{o}\) is without radical, then so is its
center. The converse is false, as the example in §10 shows. The center
of \(\mathfrak{o}/\mathfrak c\) contains a ring isomorphic to
\(Z/(Z\cap\mathfrak c)\), but this ring may be a proper subring.

\subsection{§13. Completely reducible
rings}\label{completely-reducible-rings}

By §5, a ring is called \textbf{right completely reducible} if it is the
direct sum of finitely many simple right ideals.

A right completely reducible ring with unit element has no non-zero
nilpotent ideal, hence no radical. Indeed every right ideal
\(\mathfrak r\) is a direct summand: \[
  \mathfrak{o}=\mathfrak r+\mathfrak s=e_1\mathfrak{o}+e_2\mathfrak{o}.
\] Since \(e_1^2=e_1\), one also has \(e_1=e_1^r\). If \(\mathfrak r\)
were nilpotent, \(e_1^r=0\) and so \(e_1=0\), whence \(\mathfrak r=0\).

We now prove the two converses.

First: A ring without radical satisfying the minimal condition for right
ideals is completely reducible relative to right ideals.

Let \(\mathfrak t\ne0\) be a minimal right ideal. We first show that
\(\mathfrak t\) contains an idempotent. Since
\(\mathfrak t^2\subseteq\mathfrak t\) and \(\mathfrak t^2\ne0\), we have
\(\mathfrak t^2=\mathfrak t\). Thus there is an \(a\in\mathfrak t\) with
\(a\mathfrak t\ne0\), and therefore \(a\mathfrak t=\mathfrak t\). The
set of all \(b\in\mathfrak t\) annihilated by \(a\), so that \(ab=0\),
is a right ideal. It is not all of \(\mathfrak t\), hence it is zero.
Therefore \(ab=0\) implies \(b=0\).

Because \(\mathfrak t=a\mathfrak t\), one may write \(a=ac\) with
\(c\ne0\). It follows that \[
  a(c^2-c)=0,
\] and therefore \(c^2=c\).

We next show that \(\mathfrak t\) is a direct summand. The right ideal
\(c\mathfrak{o}\) is contained in \(\mathfrak t\) and is non-zero
because \(c^2=c\) lies in it; hence \(c\mathfrak{o}=\mathfrak t\). Every
element \(r\in\mathfrak{o}\) has the Peirce decomposition \[
  r=cr+(r-cr).
\] The elements \(cr\) form \(\mathfrak t\). The elements \(r-cr\) form
another right ideal, annihilated by \(c\): \[
  c(r-cr)=0.
\] Thus \[
  \mathfrak{o}=\mathfrak t+\mathfrak r.
\] Because the elements of \(\mathfrak t\) are not annihilated by \(c\),
the sum is direct.

Now decompose \(\mathfrak{o}\), using the minimal condition, as a direct
sum of directly indecomposable ideals, \[
  \mathfrak{o}=\mathfrak r_1+\cdots+\mathfrak r_s.
\] The \(\mathfrak r_i\) must be simple, equivalently minimal. If, for
instance, \(\mathfrak r_1\) were not simple and \(\mathfrak t\) were a
minimal ideal in \(\mathfrak r_1\), the direct-summand result above
would give a decomposition of \(\mathfrak r_1\), contradiction. Thus
\(\mathfrak{o}\) is completely reducible.

Second: A ring without radical satisfying the minimal condition has a
unit element.

To establish first the existence of a left unit, it suffices to prove
the following. If a non-zero right ideal \(\mathfrak a=e_1\mathfrak{o}\)
has a left unit \(e_1\), then there is a right ideal \(\mathfrak b\) of
larger length which also has a left unit \(e\). Since complete
reducibility, hence boundedness of all lengths, has already been proved,
finitely many repetitions yield a left unit for \(\mathfrak{o}\) itself.

Assume \(\mathfrak a=e_1\mathfrak{o}\) and \(e_1^2=e_1\). The formula \[
  r=e_1r+(r-e_1r)
\] gives a Peirce decomposition
\(\mathfrak{o}=e_1\mathfrak{o}+\mathfrak r\). Let \(e_2\) be an
idempotent in \(\mathfrak r\) as above. Then \(e_1e_2=0\), because
\(e_1\) annihilates all elements of \(\mathfrak r\). Put \[
  e_2'=e_2-e_2e_1.
\] Then \[
  (e_2')^2=e_2',
  \qquad
  e_2'e_1=0,
  \qquad
  e_1e_2'=0,
\] so \(e_2'\mathfrak{o}\) is non-zero; and \(e=e_1+e_2'\) is a left
unit for the ring \(e_1\mathfrak{o}+e_2'\mathfrak{o}\), which has larger
right-ideal length than \(\mathfrak a\).

To show that the left unit thus constructed is also a right unit, make
the Peirce decomposition into left ideals: \[
  \mathfrak{o}=\mathfrak{o}e+\mathfrak l.
\] We have \(\mathfrak l e=0\) and \(\mathfrak l=e\mathfrak l\), so \[
  \mathfrak l^2=\mathfrak l e\mathfrak l=0.
\] Since by hypothesis no nilpotent left ideal can exist,
\(\mathfrak l=0\). Thus \(e\) is a right unit as well, and hence a unit
element.

\textbf{Theorem.} From right complete reducibility and the existence of
a unit element follows two-sided complete reducibility: \[
  \mathfrak{o}=\mathfrak d_1+\cdots+\mathfrak d_s,
\] where the \(\mathfrak d_i\) are two-sided simple.

As before, it is enough to show that every minimal two-sided ideal
\(\mathfrak a\) is a direct summand. As a right ideal, \[
  \mathfrak{o}=\mathfrak a+\mathfrak r=e_1\mathfrak{o}+e_2\mathfrak{o}.
\] The corresponding left decomposition is \[
  \mathfrak{o}=\mathfrak{o}e_1+\mathfrak{o}e_2=\mathfrak a+\mathfrak l.
\] One obtains \[
  \mathfrak a\mathfrak l=0,
  \qquad
  (\mathfrak a\mathfrak l)^2=0,
\] and hence \(\mathfrak a\mathfrak l=0\); similarly
\(\mathfrak l\mathfrak a=0\). It follows that
\(\mathfrak r=\mathfrak l\) and is two-sided, so \(\mathfrak a\) is a
two-sided direct summand. The rest of the proof is the same as for
one-sided complete reducibility.

The \(\mathfrak d_i\) are two-sided simple also as rings, because all
ideals in \(\mathfrak d_i\) are ideals in \(\mathfrak{o}\). We now
examine their structure.

\subsection{§14. Two-sided simple completely reducible rings with unit
element}\label{two-sided-simple-completely-reducible-rings-with-unit-element}

Let \(\mathfrak{o}\) be such a ring, and let \[
  \mathfrak{o}=\mathfrak r_1+\cdots+\mathfrak r_n=e_1\mathfrak{o}+\cdots+e_n\mathfrak{o}
\] be a decomposition into simple right ideals. Then \[
  \mathfrak{o}=\mathfrak l_1+\cdots+\mathfrak l_n
      =\mathfrak{o}e_1+\cdots+\mathfrak{o}e_n,
\] where the \(\mathfrak l_i\) are directly indecomposable left ideals.

The product \(\mathfrak l_i\mathfrak r_j\) is a non-zero two-sided ideal
whenever \(i,j\) are fixed and the product is non-zero; since
\(\mathfrak{o}\) is two-sided simple, one has, in fact, enough non-zero
components to recover all of \(\mathfrak{o}\). Put \[
  A_{ij}=\mathfrak l_i\mathfrak r_j.
\] Then \[
  \mathfrak{o}=\sum_{i,j} A_{ij}
\] directly: from \[
  \sum a_{ij}=0,
  \qquad a_{ij}\in A_{ij},
\] one first uses the directness of the right decomposition and then the
directness of the left decomposition to conclude each \(a_{ij}=0\).

Moreover \[
  A_{ij}\mathfrak r_j=A_{ij}\ne0,
\] and for any non-zero \(a_{ij}\in A_{ij}\) one has \[
  a_{ij}\mathfrak r_j\subseteq\mathfrak r_i,
  \qquad
  a_{ij}\mathfrak r_j\ne0.
\] Since \(\mathfrak r_j\) is simple, multiplication by \(a_{ij}\) gives
an operator homomorphism \[
  \mathfrak r_j\to\mathfrak r_i.
\] The kernel is an ideal in \(\mathfrak r_j\), hence zero. Thus it is
an isomorphism. Therefore any two simple right ideals occurring in a
decomposition of \(\mathfrak{o}\) are operator-isomorphic; the
isomorphisms are mediated by elements of \(A_{ij}\). Since any two
simple right ideals can be made to occur together in some decomposition
of \(\mathfrak{o}\), all simple right ideals are mutually isomorphic.

All homomorphisms from \(\mathfrak r_j\) to \(\mathfrak r_i\) are
mediated by elements of \(A_{ij}\). In particular, endomorphisms of
\(\mathfrak r_i\) are mediated by elements of \(A_{ii}\). The product
and sum of elements of \(A_{ii}\) correspond to product and sum of
endomorphisms. Distinct elements give distinct endomorphisms, since
\(e_i\) goes to \(a_{ii}e_i=a_{ii}\). Hence the ring \(A_{ii}\) is
isomorphic to the endomorphism ring of \(\mathfrak r_i\).

But the endomorphism ring of a simple ideal is a division ring.
Therefore \(A_{ii}\) is a division ring. Since all \(\mathfrak r_i\) are
operator-isomorphic, their endomorphism rings are ring-isomorphic. Thus
\(A_{ii}\) is determined by \(\mathfrak{o}\) uniquely up to ring
isomorphism. The various possible \(A_{ii}\) are merely different
concrete realizations of the abstract endomorphism division ring of the
simple right ideals.

\textbf{Main theorem.} Every completely reducible two-sided simple ring
\(\mathfrak{o}\) with unit element is isomorphic to the ring of all
\(n\) by \(n\) matrices over a division ring \(K\). The division ring
\(K\) is isomorphic to the automorphism division ring of the right
ideals of \(\mathfrak{o}\).

We construct matrix units. Let \(I_i\) be the identity automorphism of
\(\mathfrak r_i\), and let \(I_{ij}\) be an arbitrary isomorphism
\(\mathfrak r_j\to\mathfrak r_i\). Put \[
  I_{ik}=I_{ij}I_{jk}.
\] If \(I_{ij}\) is mediated by an element \(c_{ij}\in A_{ij}\), then \[
  c_{ij}=e_i c_{ij}=c_{ij}e_j,
\] and \[
  c_{ij}c_{jk}=c_{ik},
  \qquad
  c_{ij}c_{hk}=0\quad(j\ne h).
\] The \(c_{ij}\) are therefore matrix units.

Now define \(K\). To an element \(a_{11}\in A_{11}\) assign its
conjugates \[
  a_{ii}=c_{i1}a_{11}c_{1i}\in A_{ii},
\] and set \[
  \alpha=a_{11}+a_{22}+\cdots+a_{nn}.
\] The correspondence \(a_{11}\mapsto\alpha\) is a ring isomorphism from
\(A_{11}\) onto the set of all such \(\alpha\); call this set \(K\).
Every \(\alpha\in K\) commutes with all matrix units \(c_{ij}\).
Moreover \[
  \mathfrak{o}=\sum_{i,j} c_{ij}K.
\] Thus every element of \(\mathfrak{o}\) can be written uniquely in the
form \[
  \sum_{i,j} c_{ij}\alpha_{ij},
\] and assigning to it the matrix \((\alpha_{ij})\) gives an isomorphism
of \(\mathfrak{o}\) onto the full matrix ring over \(K\). This proves
the theorem.

Conversely, the ring of all \(n\) by \(n\) matrices over a division ring
\(A\) is two-sided simple and completely reducible; \(A\) is isomorphic
to the automorphism ring of the right ideals, and the construction above
recovers the subring \(K=A\).

Let \(c_{ij}\) be the usual matrix units, and put
\(\mathfrak r_i=c_{ii}\mathfrak{o}\). Then \[
  \mathfrak{o}=\mathfrak r_1+\cdots+\mathfrak r_n.
\] The \(\mathfrak r_i\) are simple right ideals, since every non-zero
element \(\sum_j c_{ij}\alpha_j\) generates the whole \(\mathfrak r_i\).
Indeed, if \(\alpha_k\ne0\), then multiplication by
\(\alpha_k^{-1}c_{ki}\) produces \(c_{ii}\), and \(c_{ii}\beta\) runs
through all elements of \(\mathfrak r_i\).

The right ideals \(\mathfrak r_i\) are operator-isomorphic by
\(\mathfrak r_i\mapsto c_{ji}\mathfrak r_i\). It follows that the ring
is two-sided indecomposable; by the theorem of §13, since it is already
right completely reducible and has a unit, it is two-sided simple.
Equivalently, every non-zero element generates the full two-sided ideal.

The center of \(\mathfrak{o}\) is the center of \(K\), and hence is a
field. Indeed, a central element in the matrix ring must commute with
all matrix units, so it has the form \[
  \alpha c_{11}+\cdots+\alpha c_{nn},
\] where \(\alpha\) lies in the center of \(K\). Conversely, every such
scalar diagonal matrix is central.

Since a matrix ring is also left completely reducible, and every right
completely reducible ring with unit is a direct sum of matrix rings, it
follows that every right completely reducible ring with unit is also
left completely reducible, and conversely. The center of a completely
reducible ring with unit is a direct sum of commutative fields
corresponding to the two-sided simple matrix-ring components.

We shall also use this fact: any two decompositions \[
  \mathfrak{o}=\sum c_{ij}K,
  \qquad
  \mathfrak{o}=\sum c'_{ij}K'
\] are carried into each other by an inner automorphism, \[
  x\mapsto y^{-1}xy,
  \qquad
  c'_{ij}=y^{-1}c_{ij}y.
\] The proof is by choosing elements mediating reciprocal isomorphisms
between corresponding right ideals and forming the usual conjugating
element.

\subsection{Chapter III. Module and representation
theory}\label{chapter-iii.-module-and-representation-theory}

\subsection{§15. Representations and representation
modules}\label{representations-and-representation-modules}

Let \(\mathfrak{o}\) be a ring and \(K\) a ring with unit element. In
the later applications \(K\) will always be a field.

A representation of degree \(n\) of \(\mathfrak{o}\) in \(K\) is a
homomorphism \[
  \mathfrak{o}\to D,
\] where \(D\) is a ring of \(n\) by \(n\) matrices with entries in
\(K\).

A \textbf{representation module} of \(\mathfrak{o}\) relative to \(K\)
is a bimodule \(M\) which is a left \(\mathfrak{o}\)-module and a right
\(K\)-module: \[
  \mathfrak{o}M\subseteq M,
  \qquad
  MK\subseteq M,
\] which is a direct sum of finitely many one-generator \(K\)-modules,
\[
  M=x_1K+\cdots+x_nK,
\] and in which the unit element of \(K\) is the identity operator.

Every representation module gives a representation. If
\(c\in\mathfrak{o}\) and \[
  cx_i=\sum_j x_j\gamma_{ji},
\] then, in matrix notation, \[
  c(x_1,
x_n)=(x_1,\ldots,x_n)C,
  \qquad C=(\gamma_{ji}).
\] The matrices \(C\) form a representation, for \[
  (b+c)x_i=bx_i+cx_i,
  \qquad
  (bc)x_i=b(cx_i),
\] and these are exactly the matrix rules \(B+C\) and \(BC\).

Conversely, every representation \(D\) of \(\mathfrak{o}\) belongs to a
representation module, and indeed to a particular basis of that module.
Let \(M\) be the set of formal linear forms \[
  y=x_1\alpha_1+\cdots+x_n\alpha_n,
  \qquad \alpha_i\in K.
\] Then \(M\) is a right \(K\)-module. If the element \(c\) is
represented by the matrix \(C=(\gamma_{ji})\), define \[
  cx_i=\sum_j x_j\gamma_{ji},
\] and extend by linearity. The homomorphism relations for the matrices
give \[
  (c+d)y=cy+dy,
  \qquad
  (cd)y=c(dy),
\] and the remaining bimodule laws \[
  c(y+z)=cy+cz,
  \qquad
  c(y\alpha)=(cy)\alpha
\] are immediate. Thus \(M\) is a representation module, and, relative
to the chosen basis, it gives exactly the representation
\(\mathfrak{o}\to D\).

If \((y_1,\ldots,y_n)\) is another basis of the same representation
module and \[
  (y_1,\ldots,y_n)=(x_1,\ldots,x_n)P,
\] then the matrices obtained from this new basis are \[
  P^{-1}CP.
\] Representations related by \[
  C\mapsto P^{-1}CP
\] are called equivalent and are counted in the same representation
class. Since every regular matrix \(P\) transforms one basis of \(M\)
into another, every representation module determines uniquely a
representation class.

Equivalently, one may formulate the correspondence invariantly, without
a special choice of basis. A representation module assigns to the
multiplier domain \(\mathfrak{o}\) a homomorphic system of linear
transformations of the module. Choosing a basis only writes these
transformations as matrices. Conversely, any homomorphic system of
linear transformations of a finite linear-form module gives a
representation module. Thus the relation between modules and
representation classes is independent of coordinates.

It is clear that operator-isomorphic representation modules determine
the same representation class. Conversely, if two representation modules
with bases \((x_i)\) and \((z_i)\) produce the same representation, then
\[
  x_i\mapsto z_i
\] extends to an operator-isomorphism. The multiplication law is
completely determined by the matrices \(C\).

In particular, if two bases of the same module produce the same
representation, then they differ by an operator-isomorphism of the
module onto itself. Equivalently, if \[
  C=P^{-1}CP
\] for every \(C\) and one fixed \(P\), then the transformation of the
module induced by \(P\) is an automorphism. Conversely, every
automorphism of the bimodule is mediated by a matrix commuting with all
matrices of the representation.

The notion of representation can also be extended to rings which are
commutatively connected with a commutative multiplier domain \(P\). In
that case the representing rings \(K\) must contain \(P\) in their
center, and a representation is required to be not only a ring
homomorphism but also an operator homomorphism: from \(a\mapsto A\) one
must have \[
  a\rho\mapsto A\rho\qquad(\rho\in P).
\] For representation modules this gives the additional rule \[
  am\cdot\rho=a(m\rho)=a\rho\cdot m.
\] One says that the module and the ring are commutatively connected
with \(P\).

\subsection{§16. Reducible
representations}\label{reducible-representations}

Let \(A\) be a submodule of the representation module \(M\), and suppose
it is possible to choose a basis for \(M\) consisting of a basis
\(z_1,\ldots,z_t\) of \(A\), completed by elements \(y_1,\ldots,y_r\):
\[
  M=y_1K+\cdots+y_rK+z_1K+\cdots+z_tK.
\] Then the representing matrices have block triangular form \[
  C=\begin{pmatrix} R&S\\0&T\end{pmatrix}.
\] The matrices \(T\) alone form a representation of degree \(t\),
mediated by \(A\), and the matrices \(R\) form a representation of
degree \(r\), mediated by \(M/A\).

For the products \(cz_i\) are again elements of \(A\), so they are
expressible in terms of the \(z_i\) alone; hence the lower-left block is
zero in the chosen order. The remaining blocks give the two induced
representations. Conversely, if a representation is given in such a
block form, then the subspace generated by the last \(t\) basis elements
is a submodule.

Let \[
  M=A_0\supset A_1\supset\cdots\supset A_s=0
\] be a composition series for \(M\), and suppose that each \(A_i\) is a
direct summand of the preceding module as a \(K\)-module. Then one can
choose a basis adapted to the series, and the matrices have the form \[
\begin{pmatrix}
R_{11}&*&*&\cdots&*\\
0&R_{22}&*&\cdots&*\\
0&0&R_{33}&\cdots&*\\
\vdots&\vdots&\vdots&\ddots&\vdots\\
0&0&0&\cdots&R_{ss}
\end{pmatrix}.
\] The diagonal blocks \(R_{ii}\) are representations mediated by the
composition factors \[
  A_{i-1}/A_i.
\] These representations are irreducible, since the composition factors
are simple.

If \(K\) is a field, every submodule is a direct summand, and conversely
every representation reduced as far as possible in the above sense gives
a composition series. By the Jordan-Hölder theorem, the composition
factors are uniquely determined up to operator-isomorphism; hence the
diagonal representations \(R_{ii}\) are uniquely determined up to
equivalence and order.

If \[
  M=A+B
\] is a direct sum of representation modules, then with a basis adapted
to the two summands the representation has block diagonal form \[
  \begin{pmatrix} R&0\\0&S\end{pmatrix},
\] where \(R\) is mediated by \(A\) and \(S\) by \(B\). Conversely, such
a block diagonal form corresponds to a direct-sum decomposition of the
representation module.

If one first decomposes \(M\) as a direct sum of directly indecomposable
constituents and then draws composition series in these constituents,
the representation can be put into a block form in which the large
diagonal boxes correspond to the directly indecomposable summands, and
inside each box the diagonal pieces correspond to the composition
factors. If \(K\) is a field, every representation reduced as far as
possible in this way conversely gives such a module decomposition. By
the theorem of Krull and Otto Schmidt, the classes of the directly
indecomposable constituents are uniquely determined up to order.

If the module \(M\) is completely reducible, all boxes consist of single
irreducible representations. The corresponding representation is called
\textbf{completely reducible}.

\subsection{§17. Direct-sum decomposition of representation modules for
rings with unit
element}\label{direct-sum-decomposition-of-representation-modules-for-rings-with-unit-element}

Let \(M\) be an \(\mathfrak{o}\)-module, with \(\mathfrak{o}\) a ring
with unit element. Then \[
  M=M_e+M_0,
\] where the unit is the identity operator on \(M_e\) and the zero
operator on \(M_0\). If \(M\) is also a right \(K\)-module, then \(M_e\)
and \(M_0\) are right \(K\)-modules as well.

Every \(m\in M\) is decomposed as \[
  m=em+(m-em).
\] Let \(M_e\) be the set of all \(em\), and \(M_0\) the set of all
\(m-em\). These are additive groups, and right \(K\)-modules if \(M\) is
one. Moreover \[
  r(em)=er m,
\] so \(M_e\) is an \(\mathfrak{o}\)-module; also \[
  r(m-em)=rm-erm=rm-rm=0,
\] so \(M_0\) is an \(\mathfrak{o}\)-module. The unit \(e\) acts as the
identity on the elements \(em\) and as zero on the elements \(m-em\).
Hence only the zero element belongs to both \(M_e\) and \(M_0\), and the
sum is direct.

For representation modules \(M_0\) gives the trivial representation, in
which every element is assigned the zero matrix. After splitting off
\(M_0\) there remains the representation mediated by \(M_e\), in which
the unit element is represented by the identity matrix. From now on we
may and shall restrict ourselves to such representations.

Let \[
  \mathfrak{o}=\mathfrak a_1+\cdots+\mathfrak a_s
\] be a direct sum of two-sided ideals, let \[
  e=e_1+\cdots+e_s
\] be the corresponding decomposition of the unit, and let \(M\) be an
\(\mathfrak{o}\)-module on which \(e\) is the identity operator. Then \[
  M=\mathfrak a_1M+\cdots+\mathfrak a_sM
\] directly.

Indeed \(M=\mathfrak{o}M=(\mathfrak a_1M,\ldots,\mathfrak a_sM)\). If
\(\mathfrak b_i\) denotes the sum of all components except
\(\mathfrak a_i\), then \[
  M=(\mathfrak a_iM,\mathfrak b_iM),
\] and this sum is direct because \(e_i\) acts as identity on
\(\mathfrak a_iM\) and as zero on \(\mathfrak b_iM\).

We shall therefore usually restrict ourselves to representations of
two-sided indecomposable rings.

\subsection{§18. Module and representation theory of completely
reducible
rings}\label{module-and-representation-theory-of-completely-reducible-rings}

Let \(\mathfrak{o}\) be completely reducible and two-sided simple, and
let \(M\) be a finite \(\mathfrak{o}\)-module on which the unit of
\(\mathfrak{o}\) is the identity operator. Then \(M\) is completely
reducible, and its simple constituents are operator-isomorphic to the
simple left ideals of \(\mathfrak{o}\).

For let \[
  \mathfrak{o}=\mathfrak l_1+\cdots+\mathfrak l_n
\] be the decomposition into simple left ideals, and suppose \[
  M=(\mathfrak{o}m_1,\ldots,\mathfrak{o}m_r).
\] Then \[
  \mathfrak{o}m_j=(\mathfrak l_1m_j,\ldots,\mathfrak l_nm_j).
\] The non-zero modules \(\mathfrak l_i m_j\) are isomorphic to
\(\mathfrak l_i\), since the map \(a\mapsto am_j\) is an operator
homomorphism whose kernel is a left ideal. Omitting those summands
already contained in previous ones, one obtains a direct sum of simple
modules. Thus \(M\) is completely reducible. If \(M\) is directly
indecomposable, it is simple and isomorphic to one of the
\(\mathfrak l_i\).

Let \(A\) be the automorphism division ring of such a simple module
\(M\), and let \(A_0\) be that of \(\mathfrak l_1\). Since \(M\) and
\(\mathfrak l_1\) are operator-isomorphic, the two division rings are
ring-isomorphic. By §1, both \(M\) and \(\mathfrak l_1\) may be regarded
as bimodules with these automorphism division rings on the right. They
are representation modules: \(\mathfrak l_1\), and therefore \(M\), has
finite rank over its automorphism division ring. The representations
mediated by them correspond under the ring isomorphism \(A\simeq A_0\).
Thus the study of this representation class may be reduced to the
representation mediated by \(\mathfrak l_1\) in its own automorphism
division ring.

Choose for \(\mathfrak l_1\) a basis consisting of matrix units \[
  \mathfrak l_1=(c_{11},\ldots,c_{n1})K,
\] where \(K\) denotes the concrete realization, from §14, of the
automorphism division ring. Write every element \(a\in\mathfrak{o}\) as
\[
  a=\sum_{i,j} c_{ij}\alpha_{ij}.
\] Then the representation mediated by \(\mathfrak l_1\) over \(K\) is
\[
  a\mapsto (\alpha_{ij}).
\] Indeed, multiplication of \(a\) on the basis column
\((c_{11},\ldots,c_{n1})\) gives the corresponding matrix of
coefficients.

Now let \(\Gamma\) be a subfield of \(K\) such that \(K\) has finite
rank over \(\Gamma\): \[
  K=x_1\Gamma+\cdots+x_t\Gamma.
\] Then \(K\) itself is a representation module over \(\Gamma\). Each
element \(\beta\in K\) is represented over \(\Gamma\) by a matrix
\(B=(\beta_{ij})\) determined by \[
  \beta x_i=\sum_j x_j\beta_{ji}.
\] Regarding \(\mathfrak l_1\) as representation module relative to
\(\Gamma\), the representation of \(\mathfrak{o}\) over \(\Gamma\) is
obtained from \[
  a\mapsto (\alpha_{ij})
\] by replacing each entry \(\alpha_{ij}\) by the matrix \(A_{ij}\)
which represents that element of \(K\) over \(\Gamma\).

The representations just studied - those in the automorphism division
ring of the left ideals and in its finite subfields - are, up to
isomorphism, the only ones that can be mediated by simple
\(\mathfrak{o}\)-modules, or by left ideals. For if one regards an
\(\mathfrak{o}\)-module \(M\) as a bimodule relative to some field
\(K\), then the elements of \(K\) induce module homomorphisms. Hence
\(K\) must be mapped isomorphically into a subfield of the automorphism
division ring of \(M\). The full automorphism division ring must have
finite degree over that subfield, otherwise the representation module
would have infinite rank.

These irreducible representations are not yet all irreducible
representations in general. There may be representation modules which
are reducible as \(\mathfrak{o}\)-modules but irreducible as bimodules,
and which therefore still yield irreducible representations.

\subsection{§19. The simple composition factors in modules and
representation
modules}\label{the-simple-composition-factors-in-modules-and-representation-modules}

Let \(\mathfrak{o}\) be a ring satisfying the maximal and minimal
conditions for left ideals, and let \(\mathfrak c\) be its maximal
nilpotent ideal. Then every simple \(\mathfrak{o}\)-module \(M\) is
either annihilated by \(\mathfrak{o}\), or else it is isomorphic to a
simple left ideal \(\mathfrak l\) of the radical-free ring
\(\mathfrak{o}/\mathfrak c\). In the latter case the module is
annihilated by precisely those two-sided ideals of
\(\mathfrak{o}/\mathfrak c\) which do not contain \(\mathfrak l\), and
its absolute multiplier domain is isomorphic to the two-sided simple
ring containing \(\mathfrak l\).

If \(\mathfrak c^r=0\), then necessarily \(\mathfrak cM=0\); otherwise
\[
  M=\mathfrak cM=\mathfrak c^2M=\cdots=0,
\] a contradiction. Thus \(M\) may be regarded as an
\(\mathfrak{o}/\mathfrak c\)-module. The ring
\(\mathfrak{o}/\mathfrak c\) has a unit element and no radical. Hence
\(M\) decomposes into one part annihilated by
\(\mathfrak{o}/\mathfrak c\) and one part on which the unit is the
identity; since \(M\) is simple, only one part can occur. If the unit is
the identity, the argument of §18 gives an isomorphism to a simple left
ideal. The remaining assertions are clear for such left ideals and pass
immediately to all modules isomorphic to them.

Consequently, if \(M\) is an \(\mathfrak{o}\)-module with a composition
series, its simple composition factors are either annihilated by
\(\mathfrak{o}\) or are isomorphic to simple left ideals of
\(\mathfrak{o}/\mathfrak c\).

If \(P\) is a ring commutatively connected with \(\mathfrak{o}\), all
results remain valid after replacing modules by bimodules relative to
\(P\) on the right and \(\mathfrak{o}\) on the left, subject to \[
  am\cdot\rho=a(m\rho)=a\rho\cdot m.
\] The unit element of \(P\) is to be the identity operator. Submodules
are always required to admit multiplication by \(P\). The submodules
constructed in the proofs, such as
\(\mathfrak cM,\mathfrak c^2M,\ldots\), have this property.

In particular, if \(P\) is a field and \(M\) is a finite-rank
\(P\)-module, then \(M\) is a representation module. The representation
mediated by \(M\) is a ring homomorphism and also an operator
homomorphism with respect to \(P\). Conversely, all such representations
over \(P\) arise from such modules. Irreducible representations
correspond to simple modules, and conversely. Therefore the preceding
theorems may be read as theorems about representations of
\(\mathfrak{o}\) over \(P\): apart from the zero representations, all
irreducible representations are mediated by the simple left ideals of
\(\mathfrak{o}/\mathfrak c\). If an arbitrary representation is reduced
by means of a composition series as in §16, then the only diagonal
representations which occur, besides zeros, are equivalent to those
mediated by simple left ideals of \(\mathfrak{o}/\mathfrak c\).

For these statements, no finiteness condition on \(\mathfrak{o}\) is
needed beyond the hypotheses already placed on the representation over
\(P\). Every representation is an isomorphic representation of its
absolute multiplier domain. This absolute multiplier domain is, as the
inverse image of a finite matrix ring, itself a finite-rank
\(P\)-module. Since the admissible ideals are by definition
\(P\)-modules, the maximal and minimal conditions hold for it.

If \(P\) is a commutative field, this absolute multiplier domain is a
hypercomplex system over \(P\). This assumption is automatically
fulfilled whenever non-zero representations occur on the diagonal: in
that case \(P\) is isomorphic to a subfield of the automorphism division
ring of a simple left ideal, and, in the realization by the subfield
\(K\) of \(\mathfrak{o}/\mathfrak c\), the commutative connection forces
it to lie in the center of \(K\).

\subsection{Chapter IV. Representations of groups and hypercomplex
systems}\label{chapter-iv.-representations-of-groups-and-hypercomplex-systems}

\subsection{§20. Inclusion of hypercomplex
systems}\label{inclusion-of-hypercomplex-systems}

We consider hypercomplex systems \(\mathfrak{o}\) relative to a
commutative field \(P\). By convention, only \(P\)-modules are counted
as ideals in \(\mathfrak{o}\). Hence the maximal and minimal conditions
hold for left and right ideals, and the entire ring theory of Chapter II
applies.

In what follows, a representation of the hypercomplex system
\(\mathfrak{o}\) always means a representation in the field \(P\), and
more precisely one satisfying: from \(a\mapsto A\) follows \[
  a\rho\mapsto A\rho\qquad(\rho\in P),
\] that is, module homomorphy relative to \(P\). For representation
modules this says \[
  a\rho\cdot m=a(m\rho).
\] The results of §19 therefore apply. All irreducible representations
are mediated by simple left ideals of \[
  \mathfrak{o}_0=\mathfrak{o}/\mathfrak c,
\] where \(\mathfrak c\) is the radical. Thus there are as many
inequivalent irreducible representations as there are two-sided simple
summands \(\mathfrak a\) in a decomposition \[
  \mathfrak{o}_0=\mathfrak a_1+\cdots+\mathfrak a_s.
\]

To determine explicitly the representation defined by such a left ideal,
first pass from the elements of \(\mathfrak{o}\) to their residue
classes in \(\mathfrak{o}_0\). Multiplication by an element of \(P\) is
an isomorphism for all ideals of \(\mathfrak{o}_0\); hence \(P\) is
isomorphic to a subfield \(e^{(\nu)}P\) of the automorphism division
ring \(K_\nu\) of every simple left ideal. The representation over \(P\)
mediated by a simple left ideal \(\mathfrak l_\nu\) is obtained by first
studying the representation over this subfield \(e^{(\nu)}P\), and then
transporting it to \(P\) through the isomorphism \(e^{(\nu)}P\simeq P\).
The division ring \(K_\nu\) has finite rank over \(P\); thus we are
dealing with representation in a finite-degree subfield of the
automorphism division ring of \(\mathfrak l_\nu\).

To find the representation, write each element \(c\) of
\(\mathfrak{o}_0\) according to its two-sided components, \[
  c=c_1+\cdots+c_s.
\] Only the matrix of \(c_\nu\) is then relevant, since the other
\(c_i\) annihilate \(\mathfrak l_\nu\) and are represented by zero
matrices. Write \(c_\nu\) with respect to the matrix units
\(c_{ij}^{(\nu)}\) of \(\mathfrak a_\nu\): \[
  c_\nu=\sum_{i,j} c_{ij}^{(\nu)}\alpha_{ij}^{(\nu)}.
\] Then replace each entry \(\alpha_{ij}^{(\nu)}\) by the matrix which
represents it in the representation of \(K_\nu\) over \(P\).

The division ring \(K_\nu\) has finite rank over \(P\). Every element
\(x\in K_\nu\) satisfies an equation \(f(x)=0\) with coefficients in
\(e^{(\nu)}P\), since its powers are linearly dependent. If \(P\) is
algebraically closed, this equation splits into linear factors; because
\(K_\nu\) is a division ring, \(x\) is already a root of one of those
linear factors, and so \(x\in e^{(\nu)}P\). Thus \(K_\nu=e^{(\nu)}P\).
In this case the matrices \((\alpha_{ij}^{(\nu)})\) themselves are the
irreducible representations.

Combining this observation with the end of §19 gives Burnside's theorem
in the following form.

Let \(P\) be algebraically closed and commutatively connected with a
ring \(\mathfrak{o}\). Then every irreducible representation of degree
\(n\) over \(P\) contains exactly \(n^2\) linearly independent matrices.

This immediately implies the usual formulation: a system of \(n\) by
\(n\) matrices over \(P\), closed under multiplication, is irreducible
if and only if there is no non-trivial homogeneous linear relation \[
  \sum \alpha_{ij}x_{ij}=0
\] with coefficients in \(P\) satisfied by the entries \(x_{ij}\) of
every matrix in the system. For by adjoining all linear combinations one
obtains a ring and a \(P\)-module, and may regard this ring as its own
representation.

Indeed, the absolute multiplier domain is isomorphic to a two-sided
simple completely reducible ring with unit, hence to a full matrix ring
over the automorphism division ring of its left ideals. Algebraic
closedness makes that division ring equal to \(P\). An irreducible
representation is generated by a simple left ideal. If this left ideal
has rank \(m\), and the representation has degree \(n\), then \(m=n\),
and there are exactly \(n^2\) linearly independent elements in the full
matrix ring.

Similarly one obtains the generalized Burnside theorem. Under the same
hypotheses, if \(D\) is a completely reducible representation which
decomposes into pairwise inequivalent constituents of degrees \[
  n_1,n_2,\ldots,n_s,
\] then \(D\) has rank \[
  n_1^2+n_2^2+\cdots+n_s^2
\] over \(P\).

If \(\mathfrak{o}\) is a hypercomplex system without radical over an
arbitrary field, then \(\mathfrak{o}\) is completely reducible and has a
unit element. It follows from §§17 and 18 that every representation of
\(\mathfrak{o}\) is completely reducible.

Conversely, every completely reducible representation of a ring
\(\mathfrak{o}\) commutatively connected with \(P\) has as absolute
multiplier domain a hypercomplex system without radical. Indeed, the
absolute multiplier domain is isomorphic to a ring and \(P\)-module of
matrices over \(P\), hence is a hypercomplex system. The elements of its
radical act as zero in every irreducible constituent and hence in the
whole representation; therefore the radical is zero.

The \textbf{regular representation} of a hypercomplex system
\(\mathfrak{o}\) is the representation mediated by the unit ideal
\(\mathfrak{o}\) itself. For a group ring one chooses specifically the
basis consisting of the group elements. Since all left ideals of
\(\mathfrak{o}/\mathfrak c\) occur as composition factors in
\(\mathfrak{o}\), all irreducible representations already occur on the
diagonal of the regular representation, after reducing the regular
representation by a composition series as in §16.

A representation is known once the representing matrices of the basis
elements \(u_1,\ldots,u_h\) are known. If the system is a group ring, so
that the \(u_i\) are group elements, every homomorphic matrix
representation of the group induces a representation of the group ring.
Thus the representation problem for groups is a special case of the
representation problem for hypercomplex systems.

\subsection{§21. Extension of the ground field. Representations of the
center}\label{extension-of-the-ground-field.-representations-of-the-center}

Let \[
  \mathfrak{o}=a_1P+\cdots+a_hP
\] be a hypercomplex system, and let \(\Omega\) be an extension field of
\(P\). One forms \[
  \mathfrak{o}\Omega=a_1\Omega+\cdots+a_h\Omega
\] with the old multiplication rules for the basis elements. This is
again hypercomplex over \(\Omega\).

Every representation of \(\mathfrak{o}\) gives a representation of
\(\mathfrak{o}\Omega\), since the representation of all elements is
known as soon as the representations of the basis elements are known.

An ideal, a representation module, or the corresponding representation
is called \textbf{absolutely irreducible} if it remains irreducible
after passage to the algebraic closure.

If \(\mathfrak{o}\Omega\) has no radical, then neither has
\(\mathfrak{o}\): a nilpotent ideal of \(\mathfrak{o}\) would give a
nilpotent extended ideal in \(\mathfrak{o}\Omega\). The converse is
false, as will appear below.

The center of \(\mathfrak{o}\Omega\) is the extended center \(Z\Omega\).
If \(\mathfrak{o}\) is completely reducible, then \[
  Z=Z_1+\cdots+Z_s
\] is a direct sum of fields. If, after passage to \(\Omega\), the
algebra \(\mathfrak{o}\Omega\) remains completely reducible, so that no
nilpotent ideal appears, then the new center \[
  Z\Omega=Z_1\Omega+\cdots+Z_s\Omega
\] is again a direct sum of fields. If \(\Omega\) is algebraically
closed, these fields have rank one over \(\Omega\).

For the representations of the center the following statements hold.

Every irreducible representation of a commutative hypercomplex system
\(Z\) over an algebraically closed field \(\Omega\) has degree one.
Equivalently, irreducible representations are the same as homomorphisms
\[
  Z\to\Omega.
\] Indeed, every irreducible representation of \(Z\) gives one of
\(Z\Omega\), and hence of \(Z\Omega/\mathfrak c\), where \(\mathfrak c\)
is the radical. This quotient is commutative and without radical, hence
a direct sum of fields. Over an algebraically closed field these fields
are of first degree, and they generate all irreducible representations.
Since equivalent one-dimensional representations are necessarily equal,
the number of distinct homomorphisms \(Z\to\Omega\) is the rank of
\(Z\Omega/\mathfrak c\).

For the special case in which \(Z\) is a field over \(P\), this says:
\(Z\) is completely reducible, i.e.~has no radical after scalar
extension, if and only if \(Z\) is an extension of the first kind over
\(P\). For a field the homomorphisms are isomorphisms. Their number is
the rank of \(Z\Omega/\mathfrak c\), and it equals the field degree
precisely when \(\mathfrak c=0\).

If \[
  Z=Z_1+\cdots+Z_s
\] is completely reducible, then under every irreducible representation
exactly one of the fields \(Z_i\) is mapped isomorphically, and the
others are represented by zero.

Applying these statements to hypercomplex systems gives first, for
systems without radical: if \(\mathfrak{o}\Omega\), and therefore
\(\mathfrak{o}\), is without radical, then in every absolutely
irreducible representation of \(\mathfrak{o}\) the central elements
\(z\) are represented by scalar diagonal matrices \[
  \xi E,
\] and \(z\mapsto \xi\) is a one-dimensional representation of the
center. The correspondence between absolutely irreducible representation
classes of \(\mathfrak{o}\) and those of \(Z\) is one-to-one. Hence the
number of such classes equals the rank of the center.

For the decomposition into two-sided indecomposable ideals \[
  \mathfrak{o}\Omega=\mathfrak a_1+\cdots+\mathfrak a_s
\] corresponds to a decomposition of the center into fields \[
  Z\Omega=Z_1+\cdots+Z_s.
\] In an irreducible representation one \(\mathfrak a_i\) is represented
isomorphically, the others by zero; this two-sided summand determines
the representation class. The summand \(\mathfrak a_i\) is represented
by a full matrix ring, and the center of such a ring consists of scalar
diagonal matrices. Thus the representation of the center is a
homomorphism onto the corresponding scalar field, and this proves the
assertion.

The following consequence is useful. If the center has a component
\(Z_i\) which is an extension of the second kind over \(P\), then
\(\mathfrak{o}\Omega\) certainly has a radical, since \(Z_i\Omega\)
already has one.

\subsection{§22. Application to abelian
groups}\label{application-to-abelian-groups}

Let \[
  G=G_1\times G_2\times\cdots\times G_s
\] be a finite abelian group, decomposed into cyclic groups \(G_i\) of
orders \(h_i\). Let \[
  G_i=\{1,a_i,a_i^2,\ldots,a_i^{h_i-1}\}.
\] Let \(P\) be a field whose characteristic does not divide the group
order \[
  h=h_1h_2\cdots h_s.
\] The group ring \(Z\) consists of all sums \[
  \sum c_{i_1\ldots i_s}a_1^{i_1}\cdots a_s^{i_s},
  \qquad c_{i_1\ldots i_s}\in P.
\] It is a homomorphic image of the polynomial ring
\(P[x_1,\ldots,x_s]\) by the substitution \(x_i\mapsto a_i\). Hence \[
  Z\simeq P[x_1,\ldots,x_s]/\mathfrak m,
\] where \[
  \mathfrak m=(x_1^{h_1}-1,\ldots,x_s^{h_s}-1).
\] In an extension field \(\Omega\) in which all these polynomials
split, each \(x_i^{h_i}-1\) decomposes into distinct linear factors,
since the characteristic does not divide \(h_i\). Thus \(\mathfrak m\)
decomposes as the intersection of distinct prime ideals \[
  (x_1-\xi_1,
    \ldots,
    x_s-\xi_s),
\] where each \(\xi_i\) runs through the \(h_i\)-th roots of unity. The
direct-sum decomposition corresponding to this intersection gives one
field for each system \((\xi_1,
\ldots,\xi_s)\). The associated representations are \[
  a_i\mapsto \xi_i,
  \qquad
  a_1^{i_1}\cdots a_s^{i_s}
     \mapsto \xi_1^{i_1}\cdots\xi_s^{i_s}.
\] These are the characters. Their number is \[
  h_1h_2\cdots h_s=h,
\] as the general theory requires.

\subsection{§23. Determinant of a hypercomplex
system}\label{determinant-of-a-hypercomplex-system}

Let \[
  \mathfrak{o}=a_1P+
\cdots+a_hP
\] be a hypercomplex system. Adjoin to \(P\) independent variables
\(x_1,\ldots,x_h\) and form \[
  \mathfrak{o}^*=a_1P(x_1,\ldots,x_h)+\cdots+a_hP(x_1,\ldots,x_h)
\] with the old multiplication rules; the variables commute with the
\(a_i\). Thus the rules in \(\mathfrak{o}^*\) are already determined.

In \(\mathfrak{o}^*\) lies the \textbf{general element} of
\(\mathfrak{o}\), \[
  w=a_1x_1+
\cdots+a_hx_h.
\] If \(a_i\mapsto A_i\) in a representation, then \(w\) is represented
by \[
  W=A_1x_1+\cdots+A_hx_h.
\] The matrix \(W\) is the system matrix attached to the representation;
when the \(a_i\) are group elements and \(\mathfrak{o}\) is the group
ring, it is the group matrix. For the regular representation one has the
regular system matrix.

The entries of \(W\) are linear forms in the \(x_i\). The determinant
\(|W|\) is therefore a homogeneous form of degree \(n\) when the
representation has degree \(n\). In particular, the regular system
determinant has degree \(h\).

The system determinant does not change under passage to an equivalent
representation, because \[
  |PWP^{-1}|=|P|\,|W|\,|P^{-1}|=|W|.
\] If one changes from the basis \((a_1,\ldots,a_h)\) to a new basis
\((b_1,\ldots,b_h)\) and writes \[
  w=\sum_i b_i y_i,
\] then the new determinant is obtained from the old one by an
invertible linear substitution \[
  x_i=\sum_j \gamma_{ij}y_j.
\]

If a composition series of the representation module is given, then, for
a suitable choice of basis, the matrix \(W\) is upper triangular in
blocks: \[
\begin{pmatrix}
W_1&*&\cdots&*\\
0&W_2&\cdots&*\\
\vdots&\vdots&\ddots&\vdots\\
0&0&\cdots&W_r
\end{pmatrix}.
\] The determinants \(|W_i|\) which occur in an arbitrary representation
are among those occurring in the regular representation of
\(\mathfrak{o}\), indeed already in that of
\(\mathfrak{o}/\mathfrak c\), where \(\mathfrak c\) is the radical.

Assume now that \(P\) is algebraically closed. Then the determinant
belonging to an irreducible representation is an irreducible polynomial
in the \(x_i\), and inequivalent representations give different
irreducible factors.

It is enough to work in the radical-free ring
\(\mathfrak{o}/\mathfrak c\). There we use matrix units
\(c_{ij}^{(\nu)}\) as a basis; the general element has the form \[
  w=\sum_{\nu,i,j} c_{ij}^{(\nu)}x_{ij}^{(\nu)}.
\] The matrices of the irreducible representations are precisely \[
  W_\nu=(x_{ij}^{(\nu)}).
\] The determinants \[
  |W_\nu|=|x_{ij}^{(\nu)}|
\] are the familiar irreducible determinant forms and are plainly
distinct.

To compute the \(|W_\nu|\), one may factor the regular system matrix of
\(\mathfrak{o}/\mathfrak c\). Each irreducible factor appears with
multiplicity equal to the degree of the corresponding irreducible
representation, since the corresponding left ideal occurs that many
times in a composition series.

One may also start from the regular system matrix of \(\mathfrak{o}\)
itself. Then each irreducible factor appears with a higher multiplicity,
namely as often as the corresponding simple left ideal occurs as
composition factor in the left ideals of \(\mathfrak{o}\). If
\(\mathfrak{o}\) is instead regarded as right ideal, one gets a second
regular system matrix, Frobenius's antistrophic matrix. It contains the
same irreducible factors - the system determinants of all irreducible
representations - possibly with different exponents.

The system determinant of a commutative system splits into linear
factors, because all irreducible representations have degree one. These
linear factors are themselves the irreducible representations and, for
abelian groups, the characters. This was Dedekind's starting point in
the study of the group determinant of non-abelian groups.

\subsection{§24. Traces and characters}\label{traces-and-characters}

If, in a representation of a hypercomplex system \(\mathfrak{o}\), the
element \(a\) is assigned the matrix \(A\), set \[
  \operatorname{Sp}(a)=\operatorname{Sp} A.
\] Traces are linear functions: \[
  \operatorname{Sp}(c+d)=\operatorname{Sp}(c)+\operatorname{Sp}(d),
  \qquad
  \operatorname{Sp}(\alpha c)=\alpha\operatorname{Sp}(c).
\] Equivalent representations have the same traces.

The trace in a reducible representation is the sum of the traces in the
representations mediated by the composition factors, because the
representing matrix is block triangular and its trace is the sum of the
traces of the diagonal blocks.

The \textbf{principal trace} is the trace in the regular representation.
The \textbf{reduced trace} is the sum of the traces in the distinct
irreducible representations.

If \(c\) belongs to the maximal nilpotent ideal \(\mathfrak c\), then \[
  \operatorname{Sp}(c)=0
\] for every representation. It suffices to consider the representations
mediated by the simple left ideals of \(\mathfrak{o}/\mathfrak c\),
since in every other representation only these composition factors occur
and trace is additive. But \(\mathfrak c\) annihilates all elements of
\(\mathfrak{o}/\mathfrak c\); hence every \(c\in\mathfrak c\) is
represented by the zero matrix.

Let \(\mathfrak{o}\) be two-sided simple and let \(P\) be algebraically
closed; thus \(\mathfrak{o}\) is a matrix ring \[
  \mathfrak{o}=\sum c_{ij}P.
\] In the irreducible representation, the element \[
  a=\sum c_{ij}\alpha_{ij}
\] is assigned the matrix \((\alpha_{ij})\). Hence \[
  \operatorname{Sp}_{\mathrm{red}}(a)=\sum_i\alpha_{ii},
\] while the principal trace is \[
  \operatorname{Sp}_{\mathrm{pr}}(a)=n\sum_i\alpha_{ii}.
\] In particular, \[
  \operatorname{Sp}_{\mathrm{red}}(c_{ij})=0\quad(i\ne j),
  \qquad
  \operatorname{Sp}_{\mathrm{red}}(c_{ii})=1,
\] and \[
  \operatorname{Sp}_{\mathrm{pr}}(c_{ii})=n.
\] Passing from \(P\) to an algebraically closed field \(\Omega\) does
not change the principal trace, since it may be computed from the same
basis. This is not true for the reduced trace.

The traces of the elements \(a\) of the radical-free system
\(\mathfrak{o}\) in the absolutely irreducible representations are
called \textbf{characters} and are denoted by \[
  \chi(a),
\] or, if the representation is to be specified, by \[
  \chi^{(\nu)}(a).
\]

In an irreducible representation of degree \(n_\nu\), the central
elements, by §21, are represented by scalar matrices \[
  \xi^{(\nu)}E,
\] where \(z\mapsto \xi^{(\nu)}\) is a one-dimensional representation of
the center, i.e.~a homomorphism of the center into \(\Omega\). Hence the
trace has value \(n_\nu\xi^{(\nu)}\). Thus the homomorphisms \(\Theta\)
of the center and the characters \(\chi\) are connected by \[
  \chi(z)=n_\nu\Theta(z).
\] In the commutative case \(n_\nu=1\), and the characters themselves
are the homomorphisms.

If the algebraically closed field \(\Omega\) has characteristic zero, as
will be assumed below, one may divide by \(n_\nu\): \[
  \Theta(z)=\frac{\chi(z)}{n_\nu}.
\] The homomorphism property of \(\Theta\) becomes \[
  \frac{\chi(z)}{n_\nu}+\frac{\chi(z')}{n_\nu}
    =\frac{\chi(z+z')}{n_\nu},
\] and \[
  \frac{\chi(z)}{n_\nu}\frac{\chi(z')}{n_\nu}
    =\frac{\chi(zz')}{n_\nu}.
\]

A representation class is already uniquely determined by the traces of
its matrices. To know all traces it suffices, of course, to know the
traces of the basis elements in the given representation. If the
representation module is \(R\), it is determined by the multiplicities
with which the irreducible modules \(M_\nu\) occur as direct summands.
If this multiplicity is \(p_\nu\), then the trace of the idempotent
\(e^{(\nu)}\) in the representation equals \[
  p_\nu n_\nu.
\] Thus \[
  p_\nu=\frac{\operatorname{Sp}(e^{(\nu)})}{n_\nu}.
\]

\subsection{§25. Discriminants}\label{discriminants}

Let \[
  \mathfrak{o}=a_1P+
\cdots+a_hP
\] be a hypercomplex system.

The \textbf{discriminant matrix} is the matrix whose entries are the
principal traces \[
  \operatorname{Sp}_{\mathrm{pr}}(a_i a_j).
\] The \textbf{reduced discriminant matrix} is the analogous matrix
formed from reduced traces after passing to the algebraically closed
field. The determinants of these matrices are the discriminant and
reduced discriminant.

The discriminant remains the same under extension of the ground field.
Under passage to another basis it is multiplied by the square of the
transformation determinant.

Indeed, suppose \[
  b_j=\sum_i a_i p_{ij}.
\] Then \[
  \bigl(\operatorname{Sp}(a_i a_j)\bigr)
    =\bigl(\operatorname{Sp}(a_i b_j)\bigr)P,
\] and similarly, using the transposed matrix, \[
  \bigl(\operatorname{Sp}(a_i b_j)\bigr)
    =P^t\bigl(\operatorname{Sp}(b_i b_j)\bigr).
\] Thus \[
  \bigl(\operatorname{Sp}(a_i a_j)\bigr)
    =P^t\bigl(\operatorname{Sp}(b_i b_j)\bigr)P,
\] and taking determinants proves the assertion. The determinant is
therefore determined only up to a square factor from \(P\), but its
vanishing or non-vanishing is invariant. The same calculation also shows
that, up to a non-zero factor, one may compute the discriminant from two
different bases as \[
  \left|\operatorname{Sp}(a_i b_j)\right|.
\]

The discriminant vanishes if \(\mathfrak{o}\) has a nilpotent ideal
\(\mathfrak c\). Choose a basis \[
  c_1,\ldots,c_r,d_1,
\ldots,d_s
\] so that \(c_1,\ldots,c_r\) is a basis of \(\mathfrak c\). Since
products \(c_i c_j\) are nilpotent, and the trace of nilpotent elements
is zero in every representation, the corresponding block of the
discriminant matrix is zero, forcing the determinant to vanish. The same
holds for the reduced discriminant.

Consequently, the discriminant also vanishes if a nilpotent ideal
appears only after passing to the algebraic closure.

If \[
  \mathfrak{o}=\mathfrak a_1+
\cdots+
\mathfrak a_s
\] is a direct two-sided decomposition, and if \(M,M_1,\ldots,M_s\) are
the discriminant matrices of \(\mathfrak{o},\mathfrak a_1,
\ldots,\mathfrak a_s\), then \(M\) is block diagonal with blocks
\(M_i\), and hence \[
  |M|=|M_1|\cdots|M_s|.
\] Indeed, choose a basis of \(\mathfrak{o}\) composed of bases of the
\(\mathfrak a_i\). If \(a\in\mathfrak a_i\), then its principal trace
computed in \(\mathfrak{o}\) is the same as its principal trace computed
in \(\mathfrak a_i\). Products from different components vanish, so the
off-diagonal trace blocks are zero. The same reasoning applies to the
reduced discriminant.

For the full matrix ring \[
  \sum c_{ij}P
\] one can compute the discriminant by taking two bases consisting of
the matrix units, once ordered by first indices and once by second
indices. The product table shows that \[
  \operatorname{Sp}(c_{ij}c_{kl})
\] vanishes except in the matching cases. The resulting determinant is
\[
  D_{\mathrm{red}}=1
\] for reduced traces, and \[
  D=n^{n^2}
\] for principal traces.

Thus, if \(\mathfrak{o}\) decomposes over the algebraic closure into
matrix rings of degrees \[
  n_1,\ldots,n_s,
\] then \[
  D=(n_1^{n_1^2}\cdots n_s^{n_s^2})\,e
\] up to a non-zero square factor, whereas the reduced discriminant is
non-zero. In particular, the reduced discriminant is non-zero for
radical-free systems, and the principal discriminant is non-zero
precisely when none of the \(n_i\) is divisible by the characteristic of
the field.

Hence non-vanishing of the reduced discriminant is necessary and
sufficient for the absence of radical after algebraic closure of the
field \(P\). In characteristic zero, non-vanishing of the principal
discriminant is likewise necessary and sufficient for the absence of
radical after algebraic closure.

\subsection{§26. The group ring}\label{the-group-ring}

Let \(a_1,\ldots,a_h\) be the elements of a finite group, and form the
group ring over a field whose characteristic does not divide \(h\). Then
the discriminant \(D\) is non-zero; hence the group ring is
radical-free.

For the principal trace in the regular representation satisfies \[
  \operatorname{Sp}(1)=h,
\] and \[
  \operatorname{Sp}(a_i)=0
\] for \(a_i\ne1\). In the regular representation the identity is
represented by a diagonal matrix with \(h\) diagonal entries equal to
\(1\). If \(a_i\ne1\), then multiplication by \(a_i\) permutes the group
elements without fixed points; the corresponding matrix therefore has
zeros on the main diagonal, so its trace is zero.

Take the two bases \[
  a_1,\ldots,a_h
\] and \[
  a_1^{-1},\ldots,a_h^{-1}.
\] The matrix \[
  \bigl(\operatorname{Sp}(a_i a_j^{-1})\bigr)
\] is diagonal with diagonal entries \(h\). Therefore \[
  D=h^h\ne0.
\]

The \textbf{class} of a group element is the sum, formed in the group
ring, \[
  K_a=\sum_s s^{-1}as,
\] where the sum ranges over the distinct conjugates. The class sums
\(K_a\) are central, since they commute with all group elements. They
generate the center: if an element \(\sum \alpha_i a_i\) of the group
ring commutes with all \(s\), then \[
  \sum_i \alpha_i a_i=s^{-1}\left(\sum_i \alpha_i a_i\right)s
\] for every \(s\), so the coefficients are constant on conjugacy
classes.

Thus the rank of the center is the number of conjugacy classes.
Consequently the number of absolutely irreducible representations is
equal to the number of classes of conjugate group elements.

The matrices \(A\) and \(S^{-1}AS\) have the same trace; hence conjugate
group elements have the same character. Taking the trace of the class
sum gives \[
  \chi(K_a)=h_a\chi(a),
\] where \(h_a\) is the number of elements in the conjugacy class of
\(a\). In words: the character of a group element is the character of
its class sum divided by the number of elements of that class.

\newpage

\newpage

\section{35. On Maximal Domains of Integral
Functions}\label{on-maximal-domains-of-integral-functions}

The following investigations are closely connected with Hilbert's
problem on relatively integral functions. They go back to an occasional
question of E. Hecke, who needed a special case as an auxiliary lemma
and handled it by direct calculation.

The question is this. By a maximal domain of integral functions in the
variables \(x_1,\ldots,x_n\) - polynomials, power series, or quotients
of such - inside a given domain \(\mathfrak B\), I mean a domain which
can no longer be enlarged by adjoining integral denominators. Thus,
whenever \(a\,g(x)\) belongs to it, with \(g(x)\in\mathfrak B\) and
\(a\) a non-zero rational integer, then \(g(x)\) itself must already
belong to it.

Every integral domain \(\mathfrak A\) in \(\mathfrak B\) generates
uniquely such a maximal domain \(M(\mathfrak A)\). It consists of all
functions \(g(x)\) in \(\mathfrak B\) for which there is at least one
non-zero integer \(a\) such that \[
  a g(x)\in\mathfrak A.
\] The question is what can be said about the denominators \(a\) of this
maximal domain when the original integral domain \(\mathfrak A\) was
finite, that is, when it consists of all integral polynomials in
finitely many functions \[
  f_1(x),\ldots,f_r(x)
\] from \(\mathfrak B\).

For \(\mathfrak B\) one may take all integral polynomials or power
series in \(x_1,\ldots,x_n\), or else those quotients which contain no
denominators except those already occurring among the \(f_i(x)\), and
their powers. I show - under the additional hypothesis, in the case of
power series or quotients of power series, that only algebraic relations
occur among the \(f_i(x)\) - that the integers \(a\) may always be
chosen so that only finitely many different prime numbers occur in them.
The powers of those primes, however, need not be bounded.

This last point is immediate. If \(a g(x)\) belongs to \(\mathfrak A\)
but \(a^\rho g(x)\) does not belong to \(\mathfrak A\) for any fixed
bound on \(\rho\), then all powers of \(a\) may occur. Thus the
reasonable boundedness question is different: can one give, in general,
an infinite system of generators of \(M(\mathfrak A)\) in which the
exponents of the primes in the corresponding denominators remain
bounded? Since only finitely many primes are involved, this is
equivalent, by known arguments, to the finiteness of \(M(\mathfrak A)\)
itself. It is therefore Hilbert's problem on relatively integral
functions in its integral form. The methods used here do not decide that
question.

The proof that only finitely many different primes \(p\) occur in the
denominators rests on the following idea. To every such prime there
corresponds at least one relation modulo \(p\) among the functions
\(f_i(x)\) generating \(\mathfrak A\), a relation with integral
coefficients which is not identically satisfied in the variables \(x\)
before reduction. In another language, let \(\mathfrak o\) be the
polynomial ring over the integers in indeterminates \(u_1,\ldots,u_r\).
The assignment \[
  u_i\longmapsto f_i(x)
\] makes \(\mathfrak A\) a homomorphic image of \(\mathfrak o\); the
homomorphism is defined by a prime ideal \(\mathfrak a\) of
\(\mathfrak o\), so that \[
  \mathfrak A\simeq\mathfrak o/\mathfrak a.
\] After reduction modulo \(p\), \(\mathfrak B_p\) is likewise a
homomorphic image of the reduced polynomial ring \(\mathfrak o_p\)
except for the finitely many primes already appearing in the
denominators of the \(f_i(x)\). The new homomorphism is defined by a
prime ideal \(\mathfrak c_p\) of \(\mathfrak o_p\) containing
\(\mathfrak a_p\).

A prime \(p\) occurs in the denominator of \(M(\mathfrak A)\) if and
only if \(\mathfrak c_p\) is a proper over-ideal of \(\mathfrak a_p\).
This can happen only if either the transcendence degree of
\(\mathfrak B_p\) drops in comparison with that of \(\mathfrak A\), or
if \(\mathfrak a_p\) ceases to be a prime ideal. Both events occur only
for finitely many primes. The first is controlled by functional
determinants modulo \(p\); the second by the fact that \(\mathfrak a\)
is an absolute prime ideal and can lose primality only modulo finitely
many primes.

All arguments remain valid, with minor modifications, if the rational
integers are replaced by the integers of a finite algebraic number field
and the rational primes by the prime ideals of that field.

\subsection{\texorpdfstring{§1. Functional determinants over coefficient
fields of characteristic
\(p\)}{§1. Functional determinants over coefficient fields of characteristic p}}\label{functional-determinants-over-coefficient-fields-of-characteristic-p}

Let \[
  f_1(x_1,\ldots,x_n),\ldots,f_r(x_1,\ldots,x_n)
\] be rational functions, or more generally formally defined power
series or quotients of such, with coefficients in a field \(P\) of
characteristic \(p\). Assume that the derivatives
\(\partial f_i/\partial x_j\) are formally defined and satisfy the usual
rules of calculation. Let \(\Omega\) be an algebraically closed
extension field of \(P\).

Then the following statement holds as in characteristic zero. If there
is an algebraic dependence among the \(f_i(x)\) with coefficients in
\(\Omega\), then the rank of the functional matrix \[
  \left(\frac{\partial f_i}{\partial x_j}\right)
\] is smaller than the number of algebraically independent functions
under consideration, unless the dependence is purely inseparable. More
precisely, if the functions are algebraically dependent and \(G(u_1,
\ldots,u_r)\) is a non-zero polynomial relation of least degree, then \[
  G(f_1,
\ldots,f_r)=0.
\] Differentiation gives \[
  \sum_i \frac{\partial G}{\partial u_i}(f)
      \frac{\partial f_i}{\partial x_j}=0
\] for every \(j\). If the functional matrix had full rank, all partial
derivatives \(\partial G/\partial u_i\) would vanish. In characteristic
\(p\), this means that \(G\) is a polynomial in the \(p\)-th powers, \[
  G(u)=H(u_1^p,
\ldots,u_r^p).
\] Since \(\Omega\) is algebraically closed, this contradicts the choice
of \(G\) as an irreducible polynomial of least degree. Thus the rank
must drop.

The consequence needed below is the following. Let \[
  F_1(x),\ldots,F_t(x)
\] be integral functions - polynomials, power series, or quotients of
such - whose functional matrix has rank exactly \(t\). Choose a prime
\(p\) which occurs neither in the denominators of the \(F_i(x)\) nor in
all \(t\)-rowed determinants of the functional matrix. Then the reduced
functions \[
  f_1(x),\ldots,f_t(x)
\] modulo \(p\) remain algebraically independent.

Indeed, their functional matrix is obtained by reducing the entries of
the original functional matrix modulo \(p\). By the choice of \(p\), the
rank remains \(t\). Algebraic dependence would force the rank to be
smaller, by the preceding argument.

\subsection{§2. Formulation of the
theorem}\label{formulation-of-the-theorem}

Let \(\mathfrak A\) be an integral domain of integral functions in
\(x_1,\ldots,x_n\): polynomials, power series, or quotients of such,
with rational integral coefficients. In the case of power series one may
understand formal power series or convergent power series; only the fact
that the functions form an integral domain is used.

Let \(\mathfrak Q\) be the quotient field of \(\mathfrak A\), and let
\(\mathfrak B\) be an integral domain between \(\mathfrak A\) and
\(\mathfrak Q\). The domain \(\mathfrak M\) is called maximal in
\(\mathfrak B\) if the following condition holds: whenever \(a\ne0\) is
an integer, \(g(x)\in\mathfrak B\), and \(a g(x)\in\mathfrak M\), then
\(g(x)\in\mathfrak M\).

The maximal domain generated by \(\mathfrak A\) in \(\mathfrak B\)
exists uniquely. It is the intersection of all maximal domains in
\(\mathfrak B\) containing \(\mathfrak A\), and it consists exactly of
those \(g(x)\in\mathfrak B\) for which \(a g(x)\in\mathfrak A\) for at
least one non-zero integer \(a\).

Assume now that \(\mathfrak A\) is finite and has a unit element. Let \[
  f_1(x),\ldots,f_r(x)
\] be an integral basis, so that \(\mathfrak A\) consists of all
integral polynomials in these functions. Let \(\mathfrak B\) consist of
the allowed polynomials, power series, or quotients, with no
denominators except those appearing among the finitely many \(f_i(x)\).

In the case of power series or quotients one adds the following
condition. If \(t\) is the rank of the functional matrix of the
\(f_i(x)\), and if, after renumbering, \(f_1,
\ldots,f_t\) already have a functional matrix of rank \(t\), then the
quotient field \(\mathfrak Q\) is to be a finite algebraic extension of
\[
  \mathbb Q(f_1,
\ldots,f_t).
\] Equivalently, the remaining \(f_{t+1},
\ldots,f_r\) are algebraic over the field generated by \(f_1,
\ldots,f_t\).

Let \(\mathfrak o=\mathbb Z[u_1,
\ldots,u_r]\). Then \(\mathfrak A\) is a homomorphic image of
\(\mathfrak o\) under \(u_i\mapsto f_i(x)\), hence \[
  \mathfrak A\simeq \mathfrak o/\mathfrak a
\] for a prime ideal \(\mathfrak a\). The theorem says that, except for
finitely many rational primes \(p\), reduction modulo \(p\) gives no new
denominator obstruction: if \(a\) is not divisible by any of those
exceptional primes and \(a g(x)\) lies in \(\mathfrak A\), then \(g(x)\)
already lies in \(\mathfrak A\).

Equivalently, only finitely many primes can occur in the denominators
required to pass from \(\mathfrak A\) to its maximal domain.

\subsection{§3. Proof of the theorem}\label{proof-of-the-theorem}

The first auxiliary assertion is that the prime ideal \(\mathfrak a\)
which defines the homomorphism from \(\mathfrak o\) to \(\mathfrak A\)
is absolutely prime, under the stated hypotheses. In other words, after
extending the coefficient field to an algebraic closure, the
corresponding ideal remains prime. This is a standard consequence of the
irreducibility of the system \(f_1,
\ldots,f_t\) and of the algebraic dependence of the remaining functions
over it.

Therefore \(\mathfrak a\) can cease to be prime after reduction modulo
\(p\) for at most finitely many primes. This is one part of the
exceptional set. A second finite exceptional set is formed by the primes
in the denominators of the \(f_i(x)\). A third is formed by the primes
at which all suitable \(t\)-rowed functional determinants vanish modulo
\(p\). Outside these finitely many primes, three facts hold
simultaneously:

\begin{enumerate}
\def\labelenumi{\arabic{enumi}.}
\tightlist
\item
  the reduction of the functions \(f_i(x)\) is defined;
\item
  the reduced ideal \(\mathfrak a_p\) is still prime;
\item
  the reduced functions \(f_1,
  \ldots,f_t\) remain algebraically independent.
\end{enumerate}

Choose \(p\) outside the exceptional set. The homomorphism from the
reduced polynomial ring \(\mathfrak o_p\) to \(\mathfrak B_p\) is
defined by a prime ideal \(\mathfrak c_p\) containing \(\mathfrak a_p\).
It remains to show that in fact \[
  \mathfrak c_p=\mathfrak a_p.
\] For this it is enough to compare transcendence degrees. The residue
field of \(\mathfrak c_p\) is the quotient field of the reduced image
domain, so its transcendence degree is at least \(t\). The residue field
of \(\mathfrak a_p\) also has transcendence degree exactly \(t\),
because the classes of \(u_1,
\ldots,u_t\) remain algebraically independent and the remaining classes
are algebraic over them.

To see the latter point explicitly, use the algebraic relations in
characteristic zero. The ideal \(\mathfrak a\) contains primitive
integral polynomials \[
  G_{t+1}(u_1,
\ldots,u_t,u_{t+1}),
  \ldots,
  G_r(u_1,
\ldots,u_t,u_r)
\] expressing the algebraic dependence of the remaining generators. For
all but finitely many primes these polynomials do not vanish after
reduction and still involve the corresponding variable \(u_j\)
genuinely. Hence the classes of \(u_{t+1},
\ldots,u_r\) are algebraic over the classes of \(u_1,
\ldots,u_t\) modulo \(\mathfrak a_p\).

Thus \(\mathfrak a_p\) and \(\mathfrak c_p\) are prime ideals of the
same transcendence degree, and \(\mathfrak c_p\) contains
\(\mathfrak a_p\). A prime ideal has no proper prime over-ideal of the
same transcendence degree. Therefore \(\mathfrak c_p=\mathfrak a_p\).

This proves that no such prime \(p\) can occur in a denominator. The
theorem follows.

\subsection{§4. Algebraic integer coefficient
domains}\label{algebraic-integer-coefficient-domains}

The same reasoning applies if the rational integers are replaced by the
ring of integers of a finite algebraic number field \(K\), and rational
primes by prime ideals of \(K\).

The determinant argument and the formulation of the theorem are
unchanged. In the proof, however, one must add to the exceptional prime
ideals those which divide the finitely many algebraic integers occurring
as contents of the polynomials expressing the algebraic dependences;
these polynomials can no longer all be assumed primitive over
\(\mathbb Z\) in the same naive sense.

The final formulation is as follows. Let \(a\) be an integer of \(K\)
not divisible by any of the finitely many exceptional prime ideals. If
\(a g(x)\) lies in the original domain \(\mathfrak A\) and
\(g(x)\in\mathfrak B\), then \(g(x)\) lies in the maximal domain
generated by \(\mathfrak A\) without needing any of the non-exceptional
prime ideals as denominators.

Equivalently, after localizing away from the exceptional prime ideals,
the same argument as above shows that the relevant reduced homomorphisms
have the same prime ideal kernel. If \[
  (a)=\mathfrak p_1^{e_1}\cdots\mathfrak p_s^{e_s}
\] is the prime-ideal factorization and none of the \(\mathfrak p_i\) is
exceptional, the localization makes these prime ideals principal and
leaves the ideal decomposition unchanged. The preceding proof then
applies word for word.

If the exceptional prime ideals are generated in the localized ring by
elements \[
  \lambda_1,
\ldots,
\lambda_q,
\] then every needed denominator may be chosen as a product of powers of
these finitely many \(\lambda_i\). Returning from the localized ring to
the original ring of integers gives the asserted finite-prime statement.

\newpage

\section{36. Ideal Differentiation and the
Different}\label{ideal-differentiation-and-the-different}

It is shown how the different of an algebraic number field may be
understood as a differential quotient of a defining ideal. The agreement
of this definition with the usual one follows from structural theorems
which are interesting in themselves and which, in an analogous sense,
generalize Lagrange's interpolation formula. A full account is to appear
in the \emph{Mathematische Annalen}.

\newpage

\section{37. Normal Basis for Fields Without Higher
Ramification}\label{normal-basis-for-fields-without-higher-ramification}

By a normal basis of a Galois number field \(K/k\) one means a basis of
the maximal order \(\mathfrak O\) of \(K\) over the maximal order
\(\mathfrak o\) of \(k\) which consists of conjugate elements. It is
known that a necessary condition for the existence of such a basis is
the absence of higher ramification groups in \(K/k\). The same remains
true locally: after passing to the \(\mathfrak p\)-adic extension
\(k_{\mathfrak p}\) of \(k\) and the corresponding extension
\(K_{\mathfrak p}\) of \(K\).

I prove the converse in the following form. For every prime
\(\mathfrak p\) whose residue characteristic does not divide the degree
of \(K/k\), there exists a normal basis. In particular, if \(K/k\) has
no higher ramification, then at every ramified prime there is a normal
basis; the discriminant at such a prime becomes the square of a group
determinant.

This last assertion must be supplemented. Contrary to a conjecture of E.
Artin, the passage from this factorization to Artin conductors still
requires further ideal-theoretic considerations. The decomposition of
the group determinant according to irreducible characters gives
generalized root numbers. A suitable collection of these gives a
factorization in the ground field enlarged by the character values. In
the simplest cases these new factors can be identified with Artin's
conductors by means of the ideal-theoretic factorization of the root
numbers.

Even in the general case, where no normal basis exists, a decomposition
into generalized root numbers is always possible by means of a
composition series as Galois modules. This again gives a factorization
in the enlarged ground field, and again ideal-theoretic questions enter.

The proof of the existence of the normal basis under the hypotheses
stated above rests on this point: the module
\(\mathfrak O/\mathfrak o\), regarded as a Galois module - that is, as
an \(\mathfrak o\)-module with the substitutions of the Galois group as
operator domain - becomes operator-isomorphic to a one-sided ideal of
the integral group ring enlarged by \(\mathfrak o\). At all primes with
only ordinary ramification this group ring is a maximal order. Ideals in
maximal orders of \(p\)-adic semisimple systems are principal. Hence the
corresponding Galois module is generated by one element under the
substitutions of the group, which is exactly the assertion that there is
a normal basis.

At primes with higher ramification the integral group ring passes to a
non-maximal order, and the ideal attached to \(\mathfrak O/\mathfrak o\)
passes to a non-principal ideal. The same reasoning applies to function
fields of one variable, provided that the characteristic of the constant
field does not divide the degree of \(K/k\).

\subsection{\texorpdfstring{§1. The \(p\)-adically extended integral
group
ring}{§1. The p-adically extended integral group ring}}\label{the-p-adically-extended-integral-group-ring}

Let \(\mathfrak o\) and \(\mathfrak o_{\mathfrak p}\) denote the maximal
orders of an algebraic number field \(k\) and of its
\(\mathfrak p\)-adic extension \(k_{\mathfrak p}\), respectively. Let
\(G\) be a finite group of order \(n\). Write \((G)\) for the rational
group ring and \([G]\) for the integral group ring. Then \((G)_k\) and
\((G)_{k_{\mathfrak p}}\) denote the corresponding group rings with
coefficients in \(k\) and \(k_{\mathfrak p}\); similarly
\([G]_{\mathfrak o}\) and \([G]_{\mathfrak o_{\mathfrak p}}\) denote the
integral group rings after extension of the coefficient ring.

Thus \((G)_k\) and \((G)_{k_{\mathfrak p}}\) are semisimple systems, and
\([G]_{\mathfrak o}\) and \([G]_{\mathfrak o_{\mathfrak p}}\) are orders
in them.

\textbf{Theorem 1.} If the prime ideal \(\mathfrak p\) does not divide
the group order \(n\), then \([G]_{\mathfrak o_{\mathfrak p}}\) is a
maximal order in \((G)_{k_{\mathfrak p}}\). If \(\mathfrak p\) divides
\(n\), it is not maximal.

The discriminant of the integral group ring, and hence of the extended
group rings, is a power of the group order \(n\). If \(\mathfrak p\)
does not divide \(n\), this discriminant is a unit in
\(\mathfrak o_{\mathfrak p}\). Thus \([G]_{\mathfrak o_{\mathfrak p}}\)
cannot be enlarged to a strictly larger order while keeping the
coefficient ring fixed: such an enlargement would split off a non-unit
square factor from the discriminant. Since \(\mathfrak o_{\mathfrak p}\)
is already the maximal order of the local field, this proves the first
part of the theorem.

If \(\mathfrak p\) divides \(n\), the direct summand corresponding to
the trivial representation gives a genuine enlargement. The idempotent
attached to the trivial character has denominator \(n\), and the
corresponding order is strictly larger than the integral group ring.
Thus the group ring is then not maximal.

\textbf{Theorem 2.} If \(\mathfrak p\) does not divide \(n\), then every
integral or fractional ideal with respect to
\([G]_{\mathfrak o_{\mathfrak p}}\) is principal.

This follows because \([G]_{\mathfrak o_{\mathfrak p}}\) is then a
maximal order in a \(p\)-adic semisimple system, and ideals in such
maximal orders are principal. The assertion is first proved for simple
components and then follows for semisimple systems by decomposing into
components.

\subsection{§2. Galois modules, operator-isomorphism, normal
bases}\label{galois-modules-operator-isomorphism-normal-bases}

Let \(K/k\) be Galois with group \(G\). If \(S\in G\), write \(z^S\) for
the conjugate of \(z\in K\) under \(S\). The substitutions of \(G\) act
as operator automorphisms on \(K/k\) considered as a \(k\)-module. Hence
\(K/k\) becomes a module over the rational group ring \((G)_k\) by
putting \[
  z\left(\sum_i c_iS_i\right)=\sum_i c_i z^{S_i}.
\] A module over \((G)_k\) contained in \(K/k\) will be called a
rational Galois module. Similarly, \([G]_{\mathfrak o}\)-modules
contained in \(K/k\) are called integral Galois modules; in particular
\(\mathfrak O/\mathfrak o\) is an integral Galois module.

\textbf{Theorem 3.} As a rational Galois module, \(K/k\) is
operator-isomorphic to the group ring \((G)_k\).

This is the usual existence of a field normal basis: there is an element
\(z\in K\) whose conjugates \(z^S\) form a \(k\)-basis of \(K\). The map
\[
  \sum_i c_iS_i\longmapsto \sum_i c_i z^{S_i}
\] is an operator homomorphism, and because the two sides have the same
rank over \(k\), it is an isomorphism.

\textbf{Theorem 4.} As an integral Galois module,
\(\mathfrak O/\mathfrak o\) is operator-isomorphic to a
\([G]_{\mathfrak o}\)-module of highest rank \(n\) inside \((G)_k\).
Locally, \(\mathfrak O_{\mathfrak p}/\mathfrak o_{\mathfrak p}\) is
operator-isomorphic to a \([G]_{\mathfrak o_{\mathfrak p}}\)-module of
rank \(n\) inside \((G)_{k_{\mathfrak p}}\).

Indeed, the rational isomorphism of Theorem 3 carries the integral
module \(\mathfrak O/\mathfrak o\) to a module over the integral group
ring. Extending the coefficients to \(k_{\mathfrak p}\) gives the
corresponding local statement.

\textbf{Theorem 5.} At every prime \(\mathfrak p\) not dividing
\(n=[K:k]\), the local module
\(\mathfrak O_{\mathfrak p}/\mathfrak o_{\mathfrak p}\) has a normal
basis.

By Theorem 1, \([G]_{\mathfrak o_{\mathfrak p}}\) is a maximal order at
such a prime. By Theorem 2, the module corresponding to
\(\mathfrak O_{\mathfrak p}/\mathfrak o_{\mathfrak p}\) is a principal
ideal: \[
  W [G]_{\mathfrak o_{\mathfrak p}}.
\] Its basis is \(WS\), where \(S\) runs through \(G\). Passing back
through the operator-isomorphism, if \(w\) is the element corresponding
to \(W\), then the conjugates \(w^S\) form an
\(\mathfrak o_{\mathfrak p}\)-basis of \(\mathfrak O_{\mathfrak p}\).
This is a normal basis.

If \(\mathfrak p\) divides \(n\), the associated module cannot be
principal when no normal basis exists. This is precisely what happens at
primes with higher ramification.

The operator-isomorphism of Theorem 3 also gives a concise form of
Speiser's results on Klein's problem of forms. Every representation of
the Galois group irreducible over \(k\) is generated by an irreducible
rational Galois module contained in \(K/k\); every absolutely
irreducible one is generated after passing to a splitting field \(Z\),
using the extended system \(K_Z/Z\). This is exactly the formulation in
which a basis \(v_1,
\ldots,v_m\) of such a Galois module satisfies \[
  v_i^S=\sum_j s_{ji}v_j,
\] where \(S\mapsto(s_{ji})\) is the corresponding representation.

\subsection{§3. The discriminant as a group
determinant}\label{the-discriminant-as-a-group-determinant}

To decompose the discriminant it is enough to work at the individual
ramified primes. In fields without higher ramification, these are
precisely the primes at which the normal basis just constructed is
available.

\textbf{Theorem 6.} If \(K/k\) is a Galois field without higher
ramification, then at every ramified prime its discriminant is the
square of the group determinant of the Galois group, with the
indeterminates replaced by the elements of a normal basis. Corresponding
to the irreducible representations, the group determinant decomposes
into generalized root numbers; the discriminant decomposes into ideals
in the ground field enlarged by the character values, with
conjugate-complex characters giving the same ideals.

Let \(w^S\) be a normal basis. The discriminant is the square of \[
  D=|w^{ST}|_{S,T\in G},
\] which is the group determinant with the \(w^S\) in place of the
indeterminates. The decomposition of the group determinant according to
the absolutely irreducible representations gives factors \(D_1,
\ldots,D_r\), the generalized root numbers. In the cyclic case these are
simply the Lagrange resolvents \[
  \sum_S w^S \chi(S).
\] In general, \(D_i\) lies in the coefficient field obtained by
adjoining the values of the corresponding character. Pairing an
irreducible representation with its adjoint gives real factors \[
  A_i=D_i\overline{D_i}
\] in the real subfield of the character field. The discriminant ideal
is thereby expressed as a product of such factors.

If one can show that the ideals generated by the \(A_i\) actually
descend to the ground field and are equal for conjugate characters, then
they agree with Artin's conductors, provided compound characters are
assigned the corresponding products. The functional equation and the
behavior for trivial representations of subgroups give the same result
for rationally irreducible characters.

\textbf{Theorem 7.} Let \(k=\mathbb Q\), and let \(K/k\) be cyclic of
prime degree \(\ell\), with \(\ell\) prime to the discriminant. Then for
the non-trivial representations the ideals generated by the factors
\(A_i\) are all equal to one another and agree with the conductor of
\(K/k\).

This follows from the classical factorization of Lagrange root numbers.
At a ramified prime \(p\), after adjoining the \(\ell\)-th roots of
unity, the Lagrange resolvents factor in such a way that \[
  (A_i)=(D_i\overline{D_i})
\] becomes the same prime-ideal factor for every non-trivial character.
Taking norms down to the rational field gives exactly the conductor. The
existence of a normal basis, already supplied above, replaces the usual
appeal to the realization of absolutely cyclic fields by cyclotomic
fields.

\newpage

\section{38. With R. Brauer and H. Hasse: Proof of a Main Theorem in the
Theory of
Algebras}\label{with-r.-brauer-and-h.-hasse-proof-of-a-main-theorem-in-the-theory-of-algebras}

We have finally succeeded, by joint effort, in proving the following
theorem, which is fundamental for the structure theory of algebras over
algebraic number fields and beyond.

\textbf{Main theorem.} Every normal division algebra over an algebraic
number field is cyclic, or, in the usual terminology, of Dickson type.

It is a special pleasure to present this result, essentially an
achievement of the \(p\)-adic method, to Kurt Hensel, the founder of
that method, on his seventieth birthday.

The proof consists of three reductions, contributed by the three
authors.

\subsection{1. Hasse's first reduction}\label{hasses-first-reduction}

The first reduction is due to H. Hasse and is based on his theory of
cyclic algebras over algebraic number fields.

\textbf{Reduction 1.} The main theorem follows if the following
assertion is proved:

\textbf{I.} Every algebra over \(Q\) which is split everywhere is
already split over \(Q\).

Here, as in what follows, an ``algebra over \(Q\)'' means a normal
simple algebra over the algebraic number field \(Q\), that is, a simple
hypercomplex system with center \(Q\). We write \[
  A\sim Q
\] when \(A\) is a full matrix algebra over \(Q\), i.e.~belongs to the
special similarity class determined by \(Q\) itself as division algebra.
An algebra is said to be split everywhere if, for every prime place
\(\mathfrak p\) of \(Q\), the \(\mathfrak p\)-adic extension \[
  A_{\mathfrak p}=A\otimes_Q Q_{\mathfrak p}
\] is split over \(Q_{\mathfrak p}\).

Let \(D\) be a division algebra over \(Q\). Choose a cyclic field
\(Z/Q\) such that for every prime place \(\mathfrak p\) of \(Q\) the
local degree of \(Z\) at \(\mathfrak p\) is a multiple of the local
index of \(D\) at \(\mathfrak p\). Then at every prime of \(Z\) over
\(\mathfrak p\), the local field \(Z_{\mathfrak P}\) is a splitting
field for \(D_{\mathfrak p}\). Hence the algebra \(D_Z\) obtained by
extending the center from \(Q\) to \(Z\) is split everywhere. If
assertion I is known for every ground field, then \[
  D_Z\sim Z.
\] Thus \(D\) has a cyclic splitting field. Hasse's theory then gives a
cyclic splitting field whose degree is the degree of \(D\) itself; hence
\(D\) is cyclic. Therefore the main theorem follows from I.

\subsection{2. Brauer's second
reduction}\label{brauers-second-reduction}

R. Brauer supplied the decisive idea for the second reduction. He had
observed, in the related question of determining the exponent of a
division algebra, that Sylow's theorem reduces the problem to the case
of a solvable splitting field. The same idea applies to cyclicity.

\textbf{Reduction 2.} Assertion I follows if the following assertion is
proved:

\textbf{II.} Every algebra over \(Q\) which is split everywhere and has
a solvable splitting field is already split over \(Q\).

Let \(A\) be split everywhere over \(Q\), and let \(K/Q\) be a Galois
splitting field for \(A\). Fix a rational prime \(p\), and let \(L\) be
the fixed field of a Sylow \(p\)-subgroup of the Galois group. Then
\(A_L\) is again split everywhere and has a solvable splitting field,
namely \(K/L\). If II is known, then \[
  A_L\sim L.
\] Thus \(A\) has a splitting field of degree prime to \(p\). The index
of \(A\), which divides the degree of any splitting field, is therefore
prime to \(p\). Since this holds for every prime \(p\), the index is
\(1\), and \(A\) is split. Hence I follows from II.

\subsection{3. Noether's third
reduction}\label{noethers-third-reduction}

The third reduction, which together with the first two leads to the
proof, was found by E. Noether after a communication from Hasse. Hasse
had proved the main theorem in the special case of an abelian splitting
field by using Reduction 1 and a formula reduction of the corresponding
factor system; Noether extracted the essential reduction from that
argument. Independently, Brauer had also considered this reduction.

\textbf{Reduction 3.} Assertion II follows if the following assertion is
proved:

\textbf{III.} Every algebra over \(Q\) which is split everywhere and has
a cyclic splitting field of prime degree is already split over \(Q\).

Let \(A\) be split everywhere and have a solvable splitting field
\(K/Q\). Choose a chain of fields \[
  K=A_0\supset A_1\supset\cdots\supset A_r=Q
\] such that each \(A_i/A_{i+1}\) is cyclic of prime degree. Then
\(A_{A_1}\) is split everywhere over \(A_1\) and has the cyclic
splitting field \(K=A_0\) of prime degree. If III is known, \(A_{A_1}\)
is split. Hence \(A_1\) is a splitting field for \(A\). Repeating the
argument down the chain gives finally \(A_Q\sim Q\). Thus II follows
from III.

But III is known from Hasse's results. Therefore, proceeding backwards
through Reductions 3, 2, and 1, the main theorem is proved.

Noether adds that the truth of III, more generally for cyclic splitting
fields of arbitrary degree, is equivalent to Hasse's norm theorem. For
Reduction 3 only the prime-degree case is required, hence only the
original Hilbert-Furtwängler norm theorem. Thus the reduction gives a
new simple proof of Hasse's norm theorem.

Indeed, let \(Z/Q\) be cyclic and let \(\alpha\in Q^*\) have trivial
norm residue symbol at every prime place of \(Q\). Then the cyclic
algebra \[
  A=(\alpha,Z)
\] splits everywhere. By Reduction 3, using only the Hilbert-Furtwängler
norm theorem, \(A\) is split over \(Q\). But this is equivalent to
saying that \(\alpha\) is the norm of an element of \(Z\).

The proper generalization of the norm theorem to higher, even
non-abelian, cases is therefore not the literal norm statement but
assertion I: an algebra which splits everywhere splits globally.

\subsection{Consequences}\label{consequences}

The immediate consequence of the main theorem is that Hasse's theory of
cyclic algebras over algebraic number fields becomes a general structure
and invariant theory of normal simple algebras, in particular normal
division algebras, over algebraic number fields.

\textbf{Theorem 1.} The exponent of a normal simple algebra over an
algebraic number field is equal to its index.

\textbf{Theorem 2.} The ground ideal of a genuine normal division
algebra over \(Q\) is divisible by at least one prime place of \(Q\),
and in fact by at least two.

Here the ground ideal may be defined as the reduced norm of the reduced
different; an infinite prime is included exactly when the algebra
reduces there to the quaternion algebra rather than to the real or
complex field. Hasse's local theory says that a prime place appears in
the ground ideal exactly when the local index is not \(1\). If all local
indices were \(1\), the algebra would be split everywhere and hence
would not be a genuine division algebra. A single exceptional prime is
also impossible, because the product formula for the norm residue
symbols forces the product of all local invariants to be trivial.

The theorem also gives the arithmetic characterization of splitting
fields.

\textbf{Theorem 3.} For an algebra \(A\) over \(Q\), an algebraic number
field \(K/Q\) is a splitting field for \(A\) if and only if, for every
prime place \(\mathfrak p\) of \(Q\), the local degree of \(K\) at
\(\mathfrak p\) is divisible by the \(\mathfrak p\)-index of \(A\).

Necessity is local. For sufficiency, extend \(A\) to \(K\). At each
prime of \(K\), the local algebra is split by the divisibility
condition. Hence \(A_K\) is split everywhere, and by assertion I it is
split globally.

Further, the main theorem yields a decomposition theorem for prime
divisors in splitting fields and an ordering theorem for splitting
fields. If \(K\) and \(K'\) are splitting fields and \(K\subset K'\),
then the local degree conditions for \(K\) imply those for \(K'\);
conversely the partial order among splitting fields is controlled by
divisibility of the corresponding local degrees. In particular, minimal
splitting fields may be studied arithmetically through local indices.

Finally, the main theorem applies to representations of finite groups.
All absolutely irreducible representations of a finite group are
realizable over algebraic number fields with the expected splitting
behavior, because the simple components of the rational group algebra
are normal simple algebras over number fields and hence cyclic algebras.
Schur's representation-theoretic indices are therefore exactly the
indices of cyclic algebras and are determined by the local invariants.

\newpage

\section{39. Hypercomplex Systems in Their Relations to Commutative
Algebra and Number
Theory}\label{hypercomplex-systems-in-their-relations-to-commutative-algebra-and-number-theory}

The theory of hypercomplex systems, or algebras, has developed rapidly
in recent years. Only very recently, however, has the significance of
this theory for commutative questions become clear. I will describe that
significance through two classical problems going back to Gauss: the
principal genus theorem and the closely related norm theorem.

These problems have repeatedly changed their form. In Gauss they appear
at the conclusion of his theory of binary quadratic forms. Later they
play an essential role in characterizing relatively cyclic and abelian
number fields through class field theory. Finally, they can be stated as
theorems about automorphisms and splitting of algebras. This last
formulation at the same time extends the theorems to arbitrary
relatively Galois number fields.

The guiding principle is the use of non-commutative theory for
commutative problems. One seeks, through the theory of algebras,
invariant and simple formulations of known facts, and then those
formulations lead to generalizations.

\subsection{1. Crossed products}\label{crossed-products}

Let \(K/k\) be a Galois extension of degree \(n\) with group \(G\). A
crossed product is a simultaneous embedding of \(K\) and the group \(G\)
into an algebra \(A\) in such a way that the automorphisms of \(K\)
become inner automorphisms in \(A\).

Let \(u_S\) be symbols corresponding to the group elements \(S\in G\).
First put \[
  A=\sum_{S\in G} u_S K,
\] so that \(A\) is a module of rank \(n\) over \(K\). The condition
that \(u_S\) implement the automorphism \(S\) is \[
  z u_S=u_S z^S
\] for every \(z\in K\). Multiplication among the \(u_S\) is prescribed
by \[
  u_Su_T=u_{ST}a_{S,T},
  \qquad a_{S,T}\in K^*.
\] Associativity is equivalent to the factor-system law \[
  a_{S,T}a_{ST,U}=a_{S,TU}^{U}a_{T,U}.
\] The resulting algebra is the crossed product of \(K\) with \(G\) for
the factor system \(a_{S,T}\).

Changing \(u_S\) to \[
  v_S=u_Sc_S,
  \qquad c_S\in K^*,
\] changes the factor system to an associated one, \[
  a'_{S,T}=a_{S,T}c_S^T c_T c_{ST}^{-1}.
\] Associated factor systems form one class. Likewise, similar algebras
are collected into one algebra class. The classes of crossed products
with fixed splitting field \(K\) form an abelian group under direct
product; this is Brauer's group of algebra classes, expressed in the
language of factor systems.

\subsection{2. Cyclic fields and norms}\label{cyclic-fields-and-norms}

If \(K/k\) is cyclic and \(S\) is a generator of its group, the crossed
product becomes particularly simple. One may choose one element \(u\)
corresponding to \(S\), and then \[
  uz=z^S u,
  \qquad
  u^n=a\in k^*.
\] The factor system is determined by the single element \(a\), and
changing \(u\) to \(uc\) replaces \(a\) by \[
  a\,N_{K/k}(c).
\] Thus cyclic algebra classes with splitting field \(K\) are described
by the quotient group \[
  k^*/N_{K/k}(K^*).
\] A cyclic algebra \((a,K)\) splits if and only if \(a\) is the norm of
an element of \(K\).

This is the bridge from the norm theorem to the theorem on split
algebras. The norm theorem for cyclic fields says, in one formulation,
that if \(a\) is a local norm at every finite and infinite place, then
\(a\) is a global norm. In algebraic language, the cyclic algebra
\((a,K)\) splits at every place and therefore splits globally.

The general theorem on algebras says: if an algebra over an algebraic
number field splits at every place, then it splits globally. The cyclic
norm theorem is exactly its cyclic special case. The proof of the
general theorem from the cyclic case proceeds by purely
algebraic-arithmetic reductions.

One important consequence, drawn by Hasse, is that every normal simple
algebra over an algebraic number field is cyclic. This long-conjectured
statement was the problem whose solution led to the general formulation.

\subsection{3. Principal genus theorem}\label{principal-genus-theorem}

A second consequence of the theorem on split algebras is the principal
genus theorem. Its invariant formulation rests on the fact that the
defining relations of the crossed product are purely multiplicative.
They continue to make sense if \(K^*\) is replaced by an abelian group
\(I\) on which the Galois group operates.

Taking \(I\) to be the group of ideals of \(K\), one obtains factor
systems of ideals. A class division of ideals induces a class division
of the factor systems, usually finer than the original ideal class
division. The identity class of factor systems is defined by those
\(a_{S,T}\) which generate split algebras at all ramified finite and
infinite places.

Let \(G^*\) denote the extension of the Galois group obtained in this
way. The invariant form of the principal genus theorem is this:

If, under the class division just defined, the substitution \[
  v_S=u_S(c_S)
\] induces an automorphism of \(G^*\), then this automorphism is inner
and is generated by an ideal class \((b)\).

Equivalently, if the transformation factors \[
  (c_S)^T(c_T)(c_{ST})^{-1}
\] belong to the identity class of factor systems, then there is an
ideal class \((b)\) such that \[
  (c_S)=\frac{(b)}{(b)^S}
\] for all \(S\in G\). In the cyclic case this says that if the norm of
an ideal class lies in the identity class of factor systems, then the
class is a symbolic \((1-S)\)-th power.

\subsection{4. Relation with the classical
form}\label{relation-with-the-classical-form}

If one restricts to ideals prime to the ramified places, the content
remains the same. For quadratic fields the principal genus just defined
agrees with Gauss's principal genus. Ideal classes correspond to
quadratic forms, and norms of ideal classes correspond to integers
represented by those forms. The condition that the identity class
generate split algebras at the ramified places says exactly that those
represented integers are quadratic residues at the ramified places. Thus
the forms have the total character of the principal form.

For cyclic fields the theorem becomes the usual statement: the principal
genus consists of those ideal classes whose norms are norm residues at
the finite and infinite ramified places. This is equivalent to being a
norm residue modulo the conductor.

In the general Galois case one can likewise introduce a conductor built
only from the ramified primes, normalized so that the factor systems at
each ramified place contain only elements of the identity class. The
relation with Artin's conductors, which involve the same prime ideals,
leads back to the theory of Galois modules and hence to the second
hypercomplex method. How far these two methods reach remains a question
for future work.

\section{Non-commutative Algebras}\label{non-commutative-algebras}

It is well known that the main theorems of commutative algebra are
contained in Galois theory, preceded by the theory of detachment fields
and splitting fields: fields large enough, respectively, to split off
one linear factor of a given polynomial or to decompose it completely
into linear factors.

The corresponding parts of algebra are developed here in the
non-commutative setting, and especially for hypercomplex systems. The
method is deliberately non-commutative: it works with representations in
non-commutative division rings. At the end I show how the corresponding
theorems of commutative algebra may be proved in a completely parallel
manner by representation in commutative fields.

The representation theory used here is a further development of the
theory based on representation modules. It is preceded by a short
automorphism theory. Alongside direct representations I consider
reciprocal representations. Each can be reduced to the other, but the
reciprocal representation is generated by the reciprocal representation
module. The advantage is that this reciprocal representation module is a
bimodule and may also be regarded as a one-sided module over an
extension ring. Thus representation in non-commutative division rings is
reduced to ideal theory in the extension ring: the irreducible
representation classes correspond to the irreducible ideal classes of
that extension ring, in exact analogy with the facts for representations
of hypercomplex systems over their commutative coefficient domain.

From this one obtains structure theorems for matrix rings over
non-commutative division rings. The basic observation is that every
subring defines a representation by the division ring, or reciprocally
by the oppositely isomorphic division ring. This oppositely isomorphic
division ring is the first analogue of a minimal detachment or splitting
field in the commutative case: it mediates all reciprocal
representations of first degree and, when the division ring has finite
rank over its center, also gives a complete splitting into direct
summands of rank one.

This yields the Galois theory for division rings and at the same time
expresses the fact that reciprocally isomorphic division algebras
generate inverse classes in Brauer's group of algebra classes. A second
proof of the Galois theory, applying more generally to simple systems,
is almost an immediate consequence of a theorem on mutually commuting
subrings. That theorem itself is tied directly to the preliminary
automorphism considerations.

Up to this point the arguments are purely non-commutative. But the
problem of commutative splitting fields is also handled
non-commutatively, by representing the splitting field in the division
algebra. The final sections transfer the same proof method to the
commutative case and develop splitting fields for arbitrary systems,
thereby giving the connection with the ordinary commutative
representation theory.

\subsection{§1. Automorphisms, modules, and
bimodules}\label{automorphisms-modules-and-bimodules}

The representation theory based on representation modules rests on the
theory of the automorphism ring of abelian groups with or without
operators. The elementary arguments concerning mappings and laws of
calculation will be collected here to avoid repetition.

\subsubsection{1. Multiplicative mapping and
associativity}\label{multiplicative-mapping-and-associativity}

Let \(G\) first be a group without operators, and let \(A\) be its
absolute endomorphism domain, the set of all homomorphisms of \(G\) into
itself. This domain is multiplicatively closed. If \(s,t\in A\), the
product \(st\) is defined by \[
  g(st)=(gs)t,
\] and associativity is immediate.

Now let \(G\) be a group with operators. Thus there is a set \(B\) of
symbols \(O,H,\ldots\) such that \(gO,gH,\ldots\) are uniquely defined
elements of \(G\) and each symbol induces an endomorphism of \(G\).
Hence there is a definite, though not necessarily one-to-one, map from
\(B\) into the absolute endomorphism domain \(A\).

If the operator domain \(B\) is itself multiplicatively closed, then the
map from \(B\) to \(A\) is multiplicatively homomorphic if and only if
the corresponding associativity relation \[
  g(OH)=(gO)H
\] holds for all \(g\in G\) and \(O,H\in B\). For left operators the
same assertion gives a reciprocal homomorphism, since left
multiplication must be read in the opposite order whereas operators are
read from left to right.

\subsubsection{2. Operator-homomorphic
mapping}\label{operator-homomorphic-mapping}

If \(G\) is a group with operator domain \(B\), a mapping \(G\to G'\) is
operator-homomorphic when it is a group homomorphism and respects the
operators. When the operator domains themselves are also mapped, this
includes ring homomorphisms as a special case.

The essential principle is this. Whenever a group is made into a group
with operators by a multiplicatively closed operator domain, the
associativity law is exactly the statement that the operator domain maps
homomorphically into the endomorphism ring. Conversely, a homomorphic
image of an operator domain in endomorphisms supplies the required
associative action.

\subsubsection{3. Modules and bimodules over
rings}\label{modules-and-bimodules-over-rings}

Let \(R\) be a ring and \(M\) an abelian group. A right \(R\)-module
structure on \(M\) is an operator structure satisfying \[
  (m_1+m_2)r=m_1r+m_2r,
\] \[
  m(r_1+r_2)=mr_1+mr_2,
\] \[
  m(r_1r_2)=(mr_1)r_2.
\] Left modules are defined similarly, with the order in the last
formula reversed. If \(M\) is a right module over \(R\) and a left
module over \(S\), and \[
  s(mr)=(sm)r,
\] then \(M\) is a bimodule.

The automorphism ring of a module is the ring of all module
endomorphisms. If these endomorphisms are written on the side opposite
the original ring action, a module becomes a bimodule over the original
ring and its automorphism ring. Conversely, every bimodule determines
homomorphisms from one of its multiplier rings into the automorphism
ring relative to the other.

\subsubsection{4. Passage from a right bimodule to a one-sided
module}\label{passage-from-a-right-bimodule-to-a-one-sided-module}

Suppose \(M\) is a right module over a ring \(R\), and suppose that
\(R\) contains two subrings which commute elementwise. Then \(M\) can be
regarded as a one-sided module over the extension ring generated by
those two subrings. Conversely, a one-sided module over such an
extension ring gives a bimodule by restriction to the two commuting
subrings.

This elementary passage is the key technical device. It permits one to
replace reciprocal representation modules, which are naturally
bimodules, by one-sided modules over an extension ring. Ideal theory for
the extension ring can then be applied.

\subsection{§2. Representation modules}\label{representation-modules}

\subsubsection{1. Modules of linear
forms}\label{modules-of-linear-forms}

The transition from the preceding automorphism theory to representation
theory proceeds through modules of linear forms. Let a hypercomplex
system \(S\) act on a module \(M\) by right multiplication. Choosing a
basis of \(M\) over the coefficient division ring gives matrices for the
action of every element of \(S\). The entries of those matrices depend
on the chosen basis, but the module class does not.

The automorphism ring of a module of linear forms commutes with the
representing matrices. Conversely, matrices commuting with a given
representation arise as endomorphisms of the corresponding
representation module. This is the module-theoretic form of Schur's
lemma and its generalizations.

\textbf{Theorem.} If \(M\) is a module of linear forms for a system
\(S\), then its \(S\)-automorphism ring is the complete commutant of the
matrices representing \(S\) on \(M\).

In particular, when \(M\) is irreducible, this automorphism ring is a
division ring.

\subsubsection{2. Representation and representation
module}\label{representation-and-representation-module}

A direct representation of a system \(S\) assigns to every element
\(s\in S\) a matrix in such a way that sums and products are respected.
The representation module is the module whose basis vectors are the rows
or columns on which these matrices act. Equivalent representations yield
operator-isomorphic modules, and operator-isomorphic representation
modules yield equivalent representations.

A reciprocal representation is defined by reversing the multiplication
order. It is generated by the reciprocal representation module. The
reciprocal module is naturally a bimodule: the represented system
operates on one side, and the coefficient division ring operates on the
other.

\textbf{Theorem.} Every reciprocal or direct representation module
generates a representation, and equivalent modules generate equivalent
representations. Conversely, a representation determines its module
class.

The usefulness of the reciprocal module is that it may be viewed as a
one-sided module over a suitable extension ring. Thus questions about
irreducible reciprocal representations become questions about simple
one-sided ideals in that extension ring.

\subsubsection{3. Passage from reciprocal modules to ordinary
modules}\label{passage-from-reciprocal-modules-to-ordinary-modules}

Let \(A\) be the coefficient division ring and let \(S\) be represented
reciprocally in \(A_n\), the full matrix ring over \(A\). The reciprocal
representation module is a right module over an extension ring generated
by \(S\) and \(A\). By the preceding passage, it may be treated as a
one-sided module.

The consequence is the following fundamental connection. The classes of
irreducible reciprocal representations of \(S\) over \(A\) correspond
one-to-one to the irreducible ideal classes of the extension ring. Under
complete reducibility, direct-sum decompositions of the module
correspond to decompositions of representations into irreducible
components.

This is the non-commutative analogue of the familiar fact that, for a
hypercomplex system represented over its own commutative coefficient
field, irreducible representations are obtained from one-sided ideals.

\subsubsection{4. The theorem on commuting
matrices}\label{the-theorem-on-commuting-matrices}

The module formulation gives the theorem on commuting matrices. If a
system of matrices over a division ring is irreducible, its commutant is
a division ring. If the representation is completely reducible, the
commutant is a full matrix ring over the division ring attached to the
irreducible class, with one full matrix factor for each multiplicity.

Conversely, if two systems of matrices are full mutual commutants in a
full matrix ring, each is completely reducible relative to the other.
This double-centralizer idea is the mechanism behind the later theorem
on mutually commuting subrings.

\subsection{§3. Extension modules}\label{extension-modules}

\subsubsection{1. Normal bases of
submodules}\label{normal-bases-of-submodules}

Let \(A\) be a division ring and let \(P\) be a subfield of its center
or, more generally, a central subfield. If \(M\) is a module over \(A\),
a \(P\)-submodule generated by a basis may be extended by multiplication
with \(A\) to the whole module. Linear independence over \(P\) and over
\(A\) are related exactly as in the commutative case once left and right
conventions are fixed.

This gives a normal-basis type statement for modules. A suitable basis
of a submodule may be completed to a basis of the larger module after
scalar extension. Such bases are used to compare ranks in the
double-centralizer calculations.

\subsubsection{2. Extension modules}\label{extension-modules-1}

Let \(A\) be a division ring and \(P\) a subfield. An \(A\)-module \(M\)
of rank \(n\) is called an extension module of \(P\)-rank \(r\) when it
is generated over \(A\) by a \(P\)-module of the corresponding rank and
when the two scalar structures are compatible. The extension module
records precisely how a representation over a smaller field becomes a
representation over the larger division ring.

If \(z_1,
\ldots,z_m\) are linearly independent over \(P\), then their images
remain independent under suitable scalar extension. Every \(P\)-module
\[
  T=z_1P+
\cdots+z_mP
\] inside \(M\) generates an extension module by multiplication with
\(A\).

\subsubsection{3. Invariant modules}\label{invariant-modules}

Let \(M=M_1+\cdots+M_r\) be a direct decomposition into simple
components. A submodule invariant under a second commuting system
corresponds to a sum of certain components. The precise theorem is that
extension modules of invariant submodules and invariant submodules of
extension modules are in one-to-one correspondence.

The proof is a straightforward use of the uniqueness of direct-sum
decompositions into simple modules. If a submodule is invariant, its
projection onto each simple summand is either zero or the full summand.
Conversely, any such sum is invariant.

This will be used to identify two-sided ideals with invariant modules.

\subsection{§4. Simple systems}\label{simple-systems}

\subsubsection{1. Simple systems and
modules}\label{simple-systems-and-modules}

Let \(S\) be a simple hypercomplex system over \(P\). A representation
over a division ring \(A\) is irreducible if its representation module
is simple. The corresponding extension ring is then a simple ring.
Conversely, a simple ideal class of the extension ring gives an
irreducible representation.

If \(S_A\) denotes the extension of \(S\) to \(A\), then a decomposition
of \(S_A\) as a module reflects the decomposition of the representation
after extension. When \(A\) is large enough to split the system, the
simple components have rank one in the relevant sense, and the
representation is absolutely irreducible.

\subsubsection{2. Representation classes}\label{representation-classes}

A representation class is a class of operator-isomorphic representation
modules. For a simple system with center \(P\), representation classes
over a division ring \(A\) correspond to simple ideal classes in the
extension ring. The number of irreducible classes, their dimensions, and
their automorphism division rings are obtained from the decomposition of
this extension ring.

\textbf{Theorem on representation classes.} Let \(S\) be hypercomplex
over \(P\) and let \(A\) be a division ring containing \(P\) in its
center in the appropriate sense. The irreducible reciprocal
representation classes of \(S\) in \(A\) are in one-to-one
correspondence with the irreducible right ideal classes of the extended
system. The direct-sum decomposition of the extended system gives the
decomposition of any completely reducible representation.

\subsubsection{3. Rank relation}\label{rank-relation}

When \(S\) is simple, the rank of an irreducible representation, the
rank of the corresponding ideal, and the rank of the associated division
algebra satisfy the usual product formula. If a simple algebra \(S\) is
embedded irreducibly in a full matrix ring over \(A\), and if \(B\) is
the commuting division ring, then the product of the ranks of \(S\) and
\(B\) equals the rank of the full matrix ring. This is the numerical
form of the double-centralizer theorem.

\subsubsection{4. Sharpened rank
relation}\label{sharpened-rank-relation}

If \(A\) has finite rank over its center, the rank relation can be
sharpened. The division ring reciprocal to \(A\) appears as the full
commutant of an irreducible embedding. The two central simple algebras
obtained in this way have reciprocal classes in the Brauer group. Their
indices multiply in the expected manner, and splitting of one
corresponds to complete reduction of the other.

\subsection{§5. Commuting subrings}\label{commuting-subrings}

\subsubsection{1. Reducible and irreducible
embeddings}\label{reducible-and-irreducible-embeddings}

Let \(A_n\) be a full matrix ring over a division ring \(A\), and let
\(S\) be a subring. An embedding of \(S\) is irreducible if the natural
module has no proper invariant submodule. It is completely reducible if
the module is a direct sum of irreducible invariant submodules.

If the embedding is reducible, the matrices of \(S\) may be
simultaneously put into block triangular form. If it is completely
reducible, they may be put into block diagonal form. These statements
are merely the module-theoretic reduction theorems translated into
matrices.

\subsubsection{2. Inner automorphisms}\label{inner-automorphisms}

If two simple subrings of \(A_n\) both contain the center \(P\) and are
isomorphic over \(P\), then under suitable irreducibility hypotheses an
isomorphism between them is induced by an inner automorphism of \(A_n\).
That is, there exists an invertible element \(u\in A_n\) such that \[
  \varphi(s)=u^{-1}su.
\] This is the non-commutative analogue of the extension of field
isomorphisms.

The proof uses the representation modules attached to the two
embeddings. Since the embeddings are irreducible and isomorphic, the
modules belong to the same class. An operator-isomorphism of the modules
is given by multiplication with an invertible matrix, and this matrix
implements the desired inner automorphism.

\subsubsection{3. The theorem on mutually commuting
subrings}\label{the-theorem-on-mutually-commuting-subrings}

The central theorem is the following.

\textbf{Theorem.} Let \(S\) and \(T\) be simple subrings of a full
matrix ring over a division ring, each containing the center \(P\), and
suppose that every element of \(S\) commutes with every element of
\(T\). Then the full commutant of \(S\) is generated by \(T\) together
with the center in the appropriate way precisely when the corresponding
representation is irreducible; in general the two full commutants are
obtained by passing to the completely reducible decomposition.

In the finite-rank case this becomes the usual double-centralizer
theorem: if \(S\) is a central simple subalgebra of a central simple
algebra \(A\), then the centralizer \(C_A(S)\) is central simple over
the same center, and \[
  [S:P]
  [C_A(S):P]
  =[A:P].
\] Moreover, \(S\) and its centralizer are mutual centralizers.

\subsubsection{4. Sharpened commutation
theorem}\label{sharpened-commutation-theorem}

If \(S\) is a simple subring and \(A\) is a division algebra finite over
its center, then the centralizer of \(S\) in a full matrix algebra over
\(A\) has class inverse to the class represented by the embedding of
\(S\). In particular, if one passes from \(S\) to its reciprocal system,
the associated algebra class is inverted in the Brauer group.

This sharpened theorem is the bridge to Galois theory. It says that
intermediate division rings and closed subgroups of the inner
automorphism group determine one another through centralizers.

\subsection{§6. Galois theory of simple
systems}\label{galois-theory-of-simple-systems}

The sharpened commutation theorem expresses the Galois theory of simple
systems relative to the center. We first consider division rings, then
simple systems in general.

\subsubsection{1. Galois theory of division rings finite over the
center}\label{galois-theory-of-division-rings-finite-over-the-center}

Let \(A\) be a division ring with center \(P\). Its Galois group is the
group of automorphisms of \(A\) which extend the identity on \(P\). By
the Skolem-Noether theorem in this setting, these are the inner
automorphisms. Thus the group is isomorphic to \[
  A^*/P^*,
\] where \(A^*\) and \(P^*\) denote the multiplicative groups of
non-zero elements.

Subgroups of this group correspond to subgroups \(H^*\) of \(A^*\)
containing \(P^*\). A subgroup is called closed if, after adjoining
zero, \(H^*\) becomes a division subring \(H\) of \(A\).

The statement that a division subring \(C\) is fixed by a subgroup is
equivalent to saying that \(C\) commutes elementwise with the division
ring \(H\) attached to the subgroup. Therefore the commutation theorem
gives the main theorem of non-commutative Galois theory:

\textbf{Main theorem of the Galois theory of division rings.} The
division rings \(C\) between \(P\) and \(A\) and the closed subgroups
\(H\) of the Galois group correspond one-to-one, in such a way that
\(C\) is the full invariant division ring of \(H\) and \(H\) is the full
invariant group of \(C\).

\subsubsection{2. Galois theory of simple
systems}\label{galois-theory-of-simple-systems-1}

Let \(A\) now be a simple system with center \(P\). Its Galois group is
again the group of inner automorphisms, isomorphic to \(A^*/P^*\), where
now \(A^*\) denotes the group of regular elements of \(A\).

A subgroup is called simply closed if the subring generated by the
corresponding regular elements is a simple system and if the group
consists of all regular elements of that simple system, modulo \(P^*\).
The key lemma is that every simple system containing \(P\) is generated
by its regular elements. This is clear when \(P\) is infinite, since a
general element with indeterminate coefficients is regular and may be
specialized to regular basis elements. Shoda proved it in general.

Using this lemma, commutativity with a simple system is equivalent to
invariance under the automorphisms induced by its regular elements.
Hence the commutation theorem gives:

\textbf{Main theorem of the Galois theory of simple systems.} The simple
systems \(S\) in \(A\) containing \(P\) and the simply closed subgroups
of the Galois group correspond one-to-one, in such a way that \(S\) is
the full invariant domain of the subgroup and the subgroup is the full
invariant group of \(S\).

\subsubsection{3. Proof by the extension
principle}\label{proof-by-the-extension-principle}

For division rings there is another proof of the main theorem, based on
extending isomorphisms, i.e.~representations of first degree. It avoids
rank counting and therefore requires fewer finiteness hypotheses. It
also shows in what direction an extension to division rings of infinite
rank might be possible.

Assume that a simple system \(S\) over \(P\) admits a reciprocal
representation of first degree in \(A\), so that \(S\) is itself a
division ring. Let \[
  S_A=e_1S_A+
\cdots+e_nS_A
\] be a decomposition into simple right ideals. The representation
generated by \(e_i\) is defined by \(e_iw=e_iw\) in \(A\).

\textbf{Lemma 1.} The \(n\) reciprocal representations generated by
\(e_1,
\ldots,e_n\) are distinct. Thus \(S\) has at least as many distinct
representations as its rank.

If two of them were the same, then \(e_iw=e_jw\) for every \(w\in S_A\).
But this is impossible because \(e_ie_j=0\) for \(i\ne j\).

\textbf{Lemma 2.} Let \(T\) be an intermediate division ring between
\(P\) and \(S\), and let \(s=[S:T]\). Then every reciprocal isomorphism
of \(T\) into \(A\) has at least \(s\) distinct extensions to \(S\).

Indeed, decompose the representation of \(T\) by idempotents.
Multiplication by \(S\) splits each component into exactly \(s\)
components, and Lemma 1 gives their distinctness.

Define the class division of the reciprocal isomorphisms of \(S\)
induced by \(T\): two isomorphisms are equivalent if they induce the
same isomorphism on \(T\). Then:

\textbf{Main theorem in this form.} If \(\mathfrak I_T\) is the class
division induced by \(T\), then \(T\) is a maximal intermediate division
ring with respect to this class division.

Passing from the set of isomorphisms to the automorphism group by
multiplication with the inverse of one fixed isomorphism gives the usual
subgroup formulation. The assertion that closed subgroups are full
invariant groups is included by taking the division ring \(H\) attached
to such a subgroup as \(T\).

\subsection{§7. Similarity classes of algebras and splitting
fields}\label{similarity-classes-of-algebras-and-splitting-fields}

From now on \(A\), \(B\), and their matrix rings are assumed finite over
their centers \(P\). Thus they are normal simple algebras over \(P\);
\(A\) and \(B\) may be division algebras.

All algebras \(A_n\) with the same associated division algebra \(A\) are
collected into one similarity class. Isomorphic division algebras are
identified.

\subsubsection{1. The group of algebra
classes}\label{the-group-of-algebra-classes}

If \(A\) and \(B\) are algebras over \(P\), then their direct product
over \(P\) is again an algebra over \(P\). At the level of associated
division algebras this product determines a class, independent of matrix
factors. Thus the classes form a multiplicatively closed commutative
system.

Brauer's theorem states that these similarity classes form an abelian
group under direct product. The identity is the class of full matrix
algebras over \(P\). The inverse of the class of \(A\) is represented by
the algebra \(\bar A\) reciprocally isomorphic to \(A\). The relation \[
  A\otimes_P \bar A\sim P
\] follows from the sharpened commutation theorem.

This is the Brauer group of \(P\).

\subsubsection{2. Splitting fields}\label{splitting-fields}

An extension field \(Z/P\) is a splitting field of \(A\) if \[
  A_Z\sim Z.
\] Equivalently, \(Z\) contains a maximal commutative subfield of the
division algebra associated with \(A\), or at least is large enough to
make the associated division algebra into a matrix ring.

A field is a detachment field if it splits off one simple component of
the representation; in the central simple case detachment and splitting
agree for maximal commutative subfields. In general, every maximal
commutative subfield of a division algebra is a splitting field, and
every splitting field has degree divisible by the index.

\subsubsection{3. Rank relations and
index}\label{rank-relations-and-index}

Let \(A\) be central simple over \(P\), and let \(Z\) be a maximal
commutative subfield. If \([A:P]=m^2\), then \([Z:P]=m\). The integer
\(m\) is the Schur index. It is the degree of any maximal commutative
subfield embedded in the division algebra, and it is also the minimal
possible degree of a splitting field in the regular case.

For a general splitting field of degree \(n\), one has \[
  n=mr
\] for some integer \(r\). This follows either from rank comparison or
from the fact that the index divides every splitting-field degree.

If \(A\) is embedded in a larger division algebra \(D\) and \(L\) is the
centralizer, then the refined relation reads \[
  [A:P][L:P]=[D:P]
\] with the appropriate square-root conventions for indices. These are
exactly the double-centralizer rank formulas.

\subsubsection{4. Existence of separable splitting
fields}\label{existence-of-separable-splitting-fields}

Every algebra class has separable splitting fields, even such as can be
embedded in the algebra itself. Let the ground field have characteristic
\(p\), and write the index as \[
  m=s p^r,
\] with \(s\) prime to \(p\). Suppose first that a splitting field has
inseparable part of degree \(p^r\). One shows that this inseparable part
can be reduced. The reduced discriminant of the division algebra after
extension to an algebraic closure is non-zero, so there exists an
element with non-zero reduced trace. Such an element cannot lie in the
purely inseparable ground part, since traces there vanish by
multiplication by \(p\). Adjoining it gives a proper separable extension
and reduces the remaining index. Repetition yields a separable splitting
field of degree equal to the index.

\subsection{§8. Splitting fields, detachment fields, and Galois theory
in the commutative
case}\label{splitting-fields-detachment-fields-and-galois-theory-in-the-commutative-case}

We now pass from splitting fields of central simple systems to those of
arbitrary simple systems by developing the splitting theory of the
center. Everything in this paragraph is commutative.

\subsubsection{1. Commutative simple
systems}\label{commutative-simple-systems}

A commutative simple system \(Z\) over \(P\) is a field extension. Let
\(\Omega\) be an algebraically closed extension field of \(P\). Then
\(Z_\Omega\) has no radical if and only if \(Z/P\) is separable.
Equivalently, \(Z\) has as many distinct embeddings into \(\Omega\) as
its degree.

An extension field \(A\) of \(P\) is called a splitting field of \(Z\)
if \(Z_A\) decomposes into composition factors of first degree. It is
called a detachment field if at least one composition factor of first
degree splits off. Thus \(A\) is a detachment field if it contains a
subfield isomorphic to \(Z\), and is a splitting field if it contains a
field isomorphic to the Galois closure of \(Z\).

A minimal detachment field is isomorphic to \(Z\). It is also a
splitting field if and only if \(Z/P\) is normal.

\subsubsection{2. Galois theory for separable fields by the hypercomplex
method}\label{galois-theory-for-separable-fields-by-the-hypercomplex-method}

Assume \(Z/P\) is finite separable, and let \(\Omega\) be a splitting
field, finite or infinite. Then \[
  Z_\Omega=e_1Z_\Omega+
\cdots+e_nZ_\Omega
\] decomposes into a direct sum of fields corresponding to the distinct
embeddings of \(Z\) into \(\Omega\). The idempotents \(e_i\) determine
the embeddings.

If \(T\) is an intermediate field between \(P\) and \(Z\), and
\(s=[Z:T]\), then every embedding of \(T\) into \(\Omega\) has exactly
\(s\) extensions to \(Z\). This is proved exactly as in the
non-commutative extension-principle proof: the decomposition of
\(T_\Omega\) refines to that of \(Z_\Omega\), and the number of
idempotents above each one is \(s\).

The class division of the embeddings of \(Z\) induced by \(T\) has \(T\)
as its maximal intermediate field. When \(Z/P\) is Galois, composition
with the inverse of a fixed embedding identifies the embedding set with
the automorphism group, and one obtains the usual Galois correspondence.

\subsubsection{\texorpdfstring{3. Idempotents of
\(Z_\Omega\)}{3. Idempotents of Z\_\textbackslash Omega}}\label{idempotents-of-z_omega}

Let \(S\) run through the embeddings which carry \(Z\) into its
conjugate copies inside the splitting field. The idempotents \(e^S\) of
\(Z_\Omega\) are conjugate: \[
  e^S=eS,
\] with \(e=e^1\). Applying \(S\) to the defining relations of \(e\)
carries them to the defining relations of \(e^S\). Since the direct-sum
decomposition is unique, the idempotents are uniquely determined and
hence conjugate.

If \(Z/P\) is Galois, this immediately gives the subgroup part of the
main theorem of Galois theory. For if \(H\) is a subgroup, the sum \[
  \sum_{S\in H} e^S
\] is fixed by all substitutions in \(H\) and by no larger subgroup.
Hence the invariant field attached to \(H\) is recovered.

\subsubsection{4. Complementary bases}\label{complementary-bases}

Let \(a_1,
\ldots,a_n\) be a basis of \(Z/P\). Write the idempotent in the form \[
  e=a_1\beta_1+
\cdots+a_n\beta_n.
\] Then, for every embedding \(S\), the element \[
  e^S=a_1\beta_1^S+
\cdots+a_n\beta_n^S
\] expresses the image of the basis complementary to \(a_1,
\ldots,a_n\) in the embedding defined by \(e^S\). The relations \[
  e^Se^R=0\quad(S\ne R),
  \qquad (e^S)^2=e^S
\] translate precisely into the defining equations for complementary
bases. Thus the contragredience of complementary bases follows from the
invariance of the idempotents, independently of the chosen basis.

This recovers Dedekind's relation between complementary bases and
idempotents in the language of hypercomplex decompositions.

\subsection{§9. Splitting and detachment fields of arbitrary
systems}\label{splitting-and-detachment-fields-of-arbitrary-systems}

\subsubsection{1. Definitions and reduction to the simple
case}\label{definitions-and-reduction-to-the-simple-case}

Let \(S\) be hypercomplex over \(P\). A commutative extension field
\(A/P\) is a splitting field of \(S\) if, in \(S_A\), every composition
series by one-sided ideals has absolutely simple composition factors.
Equivalently, all irreducible representations of \(S\) over \(A\) are
absolutely irreducible. An extension field \(T/P\) is a detachment field
if \(S_T\) splits off at least one absolutely simple composition factor,
equivalently if at least one irreducible representation over \(T\) is
already absolutely irreducible.

These definitions include the earlier ones: for commutative simple
systems they reduce to the definitions of §8, and for central simple
algebras they reduce to the definitions of §7. In the latter case
detachment and splitting fields coincide, because after a commutative
extension of the center the algebra remains completely reducible, and
all simple components are operator-isomorphic. Splitting off one
absolutely simple component forces complete splitting.

The splitting fields of a general system are obtained by taking the
compositum of splitting fields of the irreducible representations over
\(P\). A detachment field of the system is a detachment field of at
least one irreducible representation. Because irreducible
representations correspond one-to-one to the simple components of the
radical quotient, it is enough to consider simple systems.

\subsubsection{2. Simple systems with separable
center}\label{simple-systems-with-separable-center}

Let \(A\) be a simple system over \(P\) with center \(Z\), and assume
\(Z/P\) is separable. Let \(\Omega\) be a splitting field of \(Z\). The
direct-sum decomposition of \(Z_\Omega\) induces a two-sided
decomposition of \(A_\Omega\): \[
  A_\Omega=e_1A_\Omega+\cdots+e_nA_\Omega.
\] The components are conjugate and are isomorphic to \(A\) in the
appropriate sense; each is a normal simple algebra over its own center
\(e_iZ_\Omega\).

Therefore \(A_Q\) has no radical for every extension field \(Q\) of
\(P\); it is completely reducible. The components remain simple after
arbitrary coefficient extension, because the center has already
separated. The detachment and splitting fields are then characterized by
combining §7 with §8.

If \(A\) is a division algebra with separable center \(Z\), the
irreducible embeddings of detachment fields are precisely all maximal
commutative subfields of \(A\) which contain a copy of \(Z\). Splitting
fields are precisely the composita of a detachment field with a
splitting field of the center, in particular with the Galois closure of
\(Z\).

A minimal detachment field can be a splitting field even if the center
is not Galois. This is one of the differences from the commutative case:
for example, a quaternion algebra with center isomorphic to a quadratic
field may have the Galois closure of the center already as a minimal
detachment and splitting field.

\subsubsection{3. Inseparable center}\label{inseparable-center}

If the center is inseparable over \(P\), then \(Z_\Omega\) has a
radical, and consequently \(A_\Omega\) has a radical. Let
\(\mathfrak C\) be the radical of \(Z_\Omega\), and let \(\mathfrak c\)
be the ideal it generates in \(A_\Omega\). The quotient
\(Z_\Omega/\mathfrak C\) decomposes into conjugate components, exactly
as in the separable case, and this induces a direct-sum decomposition of
\[
  A_\Omega/\mathfrak c.
\] The components are again isomorphic to the original simple system in
the corresponding sense and have centers equal to the corresponding
components of \(Z_\Omega/\mathfrak C\).

Rank counting shows that the composition factors belonging to a
composition series of \(\mathfrak c\) are operator-isomorphic to those
of \(A_\Omega/\mathfrak c\), not merely homomorphic. Hence the
representation-theoretic conclusions remain valid after replacing
direct-sum components by composition factors.

\subsubsection{4. Final representation-theoretic
formulation}\label{final-representation-theoretic-formulation}

For irreducible representations the result may be summarized as follows.

Let a representation of a hypercomplex system be irreducible over \(P\).
It becomes absolutely completely reducible after extension of the
coefficient field if and only if the corresponding center is separable
over \(P\). In every case, splitting fields and detachment fields of the
representation can be characterized by embeddings in the associated
simple system as above.

In particular, the lowest possible degree of a detachment field is the
product of the degree of the center over \(P\) and the index of the
corresponding algebra class over that center. The degree of every
detachment field is a multiple of this number.

The representation splits, in the sense of composition series, into as
many distinct conjugate absolutely irreducible representation classes as
the degree of the largest separable subfield contained in the center.
Each class occurs with multiplicity equal to the product of the index
and the degree of the center over that separable subfield.

When the ground field is perfect, so that all finite extensions are
separable, this recovers the familiar results of I. Schur, with the
additional embedding characterization already found in Brauer's work.

\clearpage
\section{41. The Principal Genus Theorem for Relatively Galois Number Fields}

\textit{Original publication: Math. Ann. 108 (1933), pp. 411--419.}

Recent work has shown that non-commutative methods, and especially the theory of algebras, make it possible to formulate and prove known theorems on relatively cyclic and relatively abelian number fields for arbitrary relatively Galois fields. The chief example is the norm theorem, which in its general form is a theorem on split algebras. I shall show here that the principal genus theorem can be formulated in the same hypercomplex language: it is a theorem on inner automorphisms, or equivalently on crossed representations.

The proof is obtained from the theorem on split algebras by purely algebraic and arithmetic considerations. The transcendental core of that theorem is already the norm theorem in the cyclic case, and in fact it is enough there to treat prime degree. To make the formulation intelligible I first put before it a purely algebraic analogue, not used later in the proof, which I call the principal genus theorem in the minimal sense. In the cyclic special case it becomes the familiar theorem 90 of Hilbert's Zahlbericht: the norm of an element is one if and only if the element is of the form
\[
  a=b^{1-S}=b/b^S .
\]
I shall state this already as a theorem on inner automorphisms and on crossed representations, and prove it from the general facts on inner automorphisms and crossed product representations of algebras.

In the actual principal genus theorem the crossed product of a Galois group with the multiplicative group of a field is replaced by the crossed product of the ideal-class group and the Galois group. The factor systems are then divided into a finer induced ideal-class classification. This is the analogue of the ray-class division in class-field theory. Since the corresponding general theory of automorphisms and representations is not yet available in this form, the proof below reduces the statement to the theorem on split algebras. More concretely, it reduces the theorem for ideal classes to the corresponding theorem for ideals. At the ideal level Brandt's theorems on maximal orders in an algebra may be regarded as the substitute for the required automorphism theorems. They yield a proof parallel in spirit to the minimal proof, although complicated by the need to pass to the separate places. I give instead a nearly trivial proof communicated by Artin, which is an even simpler analogue of Speiser's argument.

\subsection{The principal genus theorem in the minimal sense}

Let \(k\) be an arbitrary commutative field, not necessarily a number field. Let \(K/k\) be a separable Galois extension of degree \(n\), and let \(G\) denote its Galois group. Recall first the standard facts on crossed products.

A crossed product is obtained by embedding simultaneously the field \(K\) and the group \(G\) into an algebra \(A\) in such a way that the automorphisms of \(K\) become inner automorphisms. Let
\[
  G=\{S,T,R,\ldots\}.
\]
Choose symbols \(u_S\), one for each element \(S\) of \(G\), and put first
\[
  A=\sum_{S\in G} u_S K
\]
as a right \(K\)-linear module. The requirement that \(u_S\) induce the automorphism \(S\) is expressed by
\[
  z u_S = u_S z^S\qquad(z\in K).
\]
Multiplication among the new basis symbols is prescribed by
\[
  u_S u_T = u_{ST} a_{S,T},\qquad a_{S,T}\in K^{\times} .
\]
Associativity is equivalent to the usual cocycle relation
\[
  a_{S,T}^{R} a_{ST,R}=a_{T,R}a_{S,TR}.
\]
With multiplication extended bilinearly, \(A\) is the crossed product
\[
  A=(a_{S,T},K,G).
\]
It is a central simple algebra over \(k\); \(K\) is a maximal commutative subfield and hence a splitting field. Conversely, for a given division algebra one obtains matrix algebras over it which can be represented in this way as crossed products.

Replacing \(u_S\) by
\[
  v_S=u_S c_S,\qquad c_S\in K^{\times},
\]
changes the factor system into the associated factor system
\[
  a'_{S,T}=a_{S,T}c_S^T c_T c_{ST}^{-1}.
\]
In particular, factor systems of the form
\[
  a_{S,T}=c_S^T c_T c_{ST}^{-1}
\]
are associated to the unit factor system. The quantities \(c_S^T c_T c_{ST}^{-1}\) are called transformation quantities. Associated factor systems are collected into a class; similarly, similar algebras are collected into a Brauer class. The two class notions correspond bijectively. With fixed splitting field \(K\), the classes form an abelian group under direct product; this is the Brauer group of algebra classes, and it corresponds to componentwise multiplication of cocycle classes. Its identity is the class of split algebras, or equivalently the class of all transformation systems.

For later comparison note the cyclic special case. Suppose \(K/k\) is cyclic and \(S\) is a generator of the Galois group. One can choose one symbol \(u\), with its powers corresponding to the powers of \(S\), and the crossed product is described by
\[
  A=K+uK+\cdots+u^{n-1}K,
\]
\[
  zu=uz^S,
\]
\[
  u^n=a,
\]
where \(a\in k^{\times}\). If \(v=uc\), then
\[
  a'=a\,N_{K/k}(c).
\]
Thus the group of cyclic algebra classes is represented by the factor group
\[
  k^{\times}/N_{K/k}(K^{\times}).
\]

The formulation of the minimal principal genus theorem uses the fact that the defining relations are purely multiplicative. They therefore also define a group extension
\[
  G^*=\{u_S K^{\times}: S\in G\}
\]
of \(G\) by \(K^{\times}\). The group \(K^{\times}\) is an abelian normal subgroup of \(G^*\), and
\[
  G^*/K^{\times}\cong G.
\]
The elements of \(G^*\) are exactly the regular elements of the algebra which transform \(K\) into itself. In other words they are exactly those \(g^*\) for which
\[
  (g^*)^{-1}Kg^*=K.
\]
They generate precisely the ring automorphisms of \(K\). Indeed, every such element gives an automorphism of \(K\), and conversely any such automorphism is represented by some \(u_S w\). If the induced automorphism on \(K\) is trivial, then \(w\) commutes with all elements of \(K\). Since \(K\) is a maximal commutative subfield, \(w\in K\), and because the element is regular, \(w\in K^{\times}\).

The minimal principal genus theorem can now be stated in three equivalent forms.

\textbf{First form.} Every automorphism of the group \(G^*\) which extends the identity automorphism of \(K^{\times}\) is inner, and is induced by an element of \(K^{\times}\).

\textbf{Second form.} If all transformation quantities
\[
  c_S^T c_T c_{ST}^{-1}
\]
have the value one, then the \(c_S\) are symbolic \((1-S)\)-st powers. Thus there exists \(b\in K^{\times}\) such that
\[
  c_S=b^{1-S}=b/b^S
\]
for every \(S\in G\).

\textbf{Third form.} The group \(G\) has only one crossed representation class of degree one in \(K^{\times}\) belonging to the unit factor system.

Here a representation \(u_S\mapsto c_S\) is called crossed with factor system \(a_{S,T}\) if
\[
  c_S^T c_T=c_{ST}a_{S,T}.
\]
Two such representations belong to the same class if
\[
  c_S=b^{-S}d_S b
\]
for an element \(b\in K^{\times}\). In degree one this is simply the usual equivalence by multiplication with a coboundary.

It remains to prove the first form and to see that the other two follow. Let an automorphism of \(G^*\), fixing \(K^{\times}\) elementwise, send \(u_S\) to \(v_S\). Since all the relations defining the crossed product are preserved, the direct correspondence
\[
  \sum_S u_S b_S\longmapsto \sum_S v_S b_S
\]
defines a ring automorphism of \(A\) which fixes every element of \(K\). Such an automorphism of a central simple algebra is inner. By the characterization of the normalizer of \(K\) just recalled, the implementing element must lie in \(K^{\times}\). This proves the first form.

For the second form, the relations force each \(v_S\) to have the shape \(u_S c_S\). Preservation of the factor system is exactly the assertion that the transformation quantities are one. Conversely, if this condition holds, then \(v_S=u_Sc_S\) defines an automorphism of \(G^*\) of the kind just considered. The first form gives
\[
  v_S=b^{-1}u_Sb=u_S(b/b^S),
\]
so \(c_S=b^{1-S}\). In the cyclic case the condition becomes \(N(c)=1\), and the conclusion \(c=b^{1-S}\) is Hilbert's theorem 90. The third form is merely the same statement expressed in representation language: a degree-one crossed representation with unit factor system is equivalent to the identity representation.

The third form is only a special case of a broader statement. Given any factor system \(a_{S,T}\), there is only one irreducible crossed representation class attached to it. Its degree is the Schur index of the corresponding algebra, that is, the degree of the associated division algebra. The present degree-one case is the part of this general assertion needed here.

\subsection{The principal genus theorem}

We now pass from the minimal theorem to the arithmetic theorem. In the cyclic form of the principal genus theorem of class-field theory two classifications occur at once: absolute ideal classes in the upper field and ray classes in the lower field. In the general Galois case this is replaced by the following mechanism. A classification of the ideals of the upper field induces a finer classification of the factor systems.

Let \(\mathcal I\) denote the group of all ideals of \(K\). Since \(G\) acts on ideals, one can form the same extension by the symbolic elements \(u_S\). Thus one has complexes \(u_S\mathcal I\) and multiplication rules parallel to the crossed-product rules. Passing from ideals to absolute ideal classes, write \(I\) for the ideal-class group. The complexes \(u_S I\) are still meaningful, since the principal class is stable under the action of \(G\). Hence the equivalence relation which passes from ideals to classes transfers from \(\{u_S\mathcal I\}\) to \(\{u_S I\}\). If \(\mathfrak a\) and \(\mathfrak b\) are equivalent ideals, then \(u_S\mathfrak a\) and \(u_S\mathfrak b\) are equivalent as well.

Consequently the symbols \(u_S\) are equivalent to all products \(u_S(c_S)\), where \((c_S)\) ranges over principal ideals of \(K\). If one requires the product relations to be single-valued in the induced class sense, one is requiring that the principal-ideal transformation systems
\[
  (c_S)^T(c_T)(c_{ST})^{-1}
\]
belong to the identity factor-system class. The identity class for factor systems is not merely the class of trivial systems; it consists of those principal-ideal factor systems \((a_{S,T})\) which contain representatives such that the associated algebras
\[
  (a_{S,T},K,G)
\]
split at all ramified places of \(K/k\). These principal ideals form a group under multiplication of factor systems. This group contains all transformation systems arising from principal ideals, because those give split algebras everywhere.

The principal genus theorem can now be stated exactly parallel to the minimal theorem.

\textbf{First form.} Suppose the induced ideal-class classification of factor systems is used. If the substitutions
\[
  v_S=u_S c_S
\]
define an automorphism of the complexes \(u_S I\), in the sense that the relations among the \(v_S\) agree with the original relations under the induced class division, then this automorphism is inner. It is generated by a single ideal class \(\mathfrak b\).

\textbf{Second form.} If the transformation classes
\[
  c_S^T c_T c_{ST}^{-1}
\]
all lie in the identity class of the induced factor-system classification, then the vector of ideal classes \(\{c_S\}\) belongs to the principal genus and is a symbolic \((1-S)\)-st power. Thus there exists an ideal class \(\mathfrak b\) such that
\[
  c_S=\mathfrak b^{1-S}=\mathfrak b/\mathfrak b^S
\]
for every \(S\in G\).

\textbf{Third form.} Relative to the induced ideal-class division of the factor systems, the group \(G\) has only one crossed representation class of degree one in the ideal-class group belonging to the identity class of factor systems.

The equivalence of these three formulations is the same as in the minimal case. The passage from the first to the second is even simpler here, because the automorphism has already been assumed to be generated by substitutions of the form \(v_S=u_Sc_S\). The assumption is necessary: the ideal-class group need not be a maximal commutative subgroup; for instance all classes may be ambiguous. In the third formulation the representation matrices are necessarily monomial matrices, since only multiplication is defined in the class group. For representations of degree one this makes no difference.

\subsection{Proof}

I prove the second form. The proof rests on two lemmas.

\textbf{Lemma 1 (the ideal form of the principal genus theorem).} If the ideals \(\mathfrak c_S\) have transformation systems equal to the unit ideal,
\[
  \mathfrak c_S^T\mathfrak c_T\mathfrak c_{ST}^{-1}=\mathcal O_K,
\]
then the \(\mathfrak c_S\) are symbolic \((1-S)\)-st powers. There is an ideal \(\mathfrak b\) such that
\[
  \mathfrak c_S=\mathfrak b^{1-S}
\]
for every \(S\in G\).

Equivalently, the first and third formulations hold at the ideal level, with the same proviso as above that automorphisms are represented by substitutions \(v_S=u_S\mathfrak c_S\).

The proof is immediate. Put
\[
  \mathfrak b=\sum_S \mathfrak c_S,
\]
where the sum is the module sum of ideals, equivalently the greatest common divisor in the sense of ideal theory. The displayed transformation equations imply
\[
  \mathfrak b^T=\sum_S \mathfrak c_S^T=\sum_S \mathfrak c_{ST}\mathfrak c_T^{-1}=\mathfrak b\mathfrak c_T^{-1}.
\]
Hence \(\mathfrak c_T=\mathfrak b/\mathfrak b^T=\mathfrak b^{1-T}\).

\textbf{Lemma 2 (the split-algebra theorem in the required form).} Let
\[
  A=(a_{S,T},K,G)
\]
be a crossed-product algebra. Suppose \(A\) splits at all finite and infinite ramified places of \(K/k\), and suppose the principal ideals \((a_{S,T})\) are transformation quantities of ideals:
\[
  (a_{S,T})=\mathfrak c_S^T\mathfrak c_T\mathfrak c_{ST}^{-1}.
\]
Then \(A\) splits everywhere, hence splits absolutely. Consequently the ideals \((a_{S,T})\) are transformation quantities of principal ideals:
\[
  (a_{S,T})=(d_S)^T(d_T)(d_{ST})^{-1}.
\]
Conversely, every everywhere split algebra satisfies the stated local condition.

This is just the familiar behavior of crossed products under passage to local fields. Let \(k_{\mathfrak p}\) be the completion of \(k\), and let \(K_{\mathfrak P}\) be the corresponding completion of \(K\) at a prime \(\mathfrak P\) over \(\mathfrak p\). Let \(Z\) be the decomposition group of \(\mathfrak P\), and let \(a_Z\) be the part of the factor system indexed by elements of \(Z\). Then the local algebra is
\[
  A_{\mathfrak p}\sim (a_Z,K_{\mathfrak P}/k_{\mathfrak p},Z).
\]
If \(\mathfrak p\) is unramified, the decomposition group is cyclic and the local algebra is the distinguished cyclic presentation attached to the unramified inertia field. Under the hypothesis of the lemma, the restricted factor system is associated to one consisting of local units. Passing to the normalized cyclic presentation, the parameter is a unit. In an unramified local extension all units that occur here are norms; hence the local cyclic algebra splits. At the infinite places splitting is part of the hypothesis. Thus \(A\) splits at every place. The global theorem on split algebras now gives global splitting.

The final step is immediate. Choose ideals \(\mathfrak c_S\) representing the ideal classes \(c_S\). The assumption that the class transformation systems belong to the identity class of the induced factor-system division says precisely that the ideal transformation systems
\[
  \mathfrak c_S^T\mathfrak c_T\mathfrak c_{ST}^{-1}
\]
are principal ideals \((a_{S,T})\) satisfying the conditions of Lemma 2. Hence there are principal ideals \((d_S)\) such that, after replacing
\[
  \mathfrak c_S\quad\text{by}\quad \mathfrak d_S=\mathfrak c_S/(d_S),
\]
their transformation systems become the unit ideal. Lemma 1 then gives an ideal \(\mathfrak b\) with
\[
  \mathfrak d_S=\mathfrak b^{1-S}.
\]
Passing to classes, the original classes \(c_S\) are also symbolic \((1-S)\)-st powers of the class of \(\mathfrak b\). This proves the theorem.

In the cyclic case the induced identity class becomes, after normalization, the group of principal ideals \((a)\) for which \(a\) is a local norm at all ramified places. Thus the theorem becomes the classical principal genus theorem.

\clearpage
\section{42. Split Crossed Products and Their Maximal Orders}

\textit{Original publication: Actualités scientifiques et industrielles 148 (1934), pp. 5--15.}

In the simplest case of split crossed products, that is, those whose factor systems are associated to one and hence give split algebras, the situation can be followed quite explicitly. The same remains true when the splitting field, or maximal commutative subfield, is not Galois. The elementary algebraic facts needed for this are developed first. I then give, for an algebraic number field as ground field, explicit descriptions of all maximal orders and of their ideals by means of modules and complementary modules of the splitting field. I also introduce a division into regions. Maximal orders are put into the same region when they have the same intersection with the chosen splitting field. The regions thus correspond bijectively to the orders in that field; the principal order corresponds to the principal region. The maximal orders in a region transform into one another by ideals built from modules of the corresponding order. In the principal region this is simply the extension of ideals from the splitting field, or of modules over its principal order.

Finally I pass to arbitrary simple algebras. By refining the division into regions - requiring not only the same intersection, but also equality of the local components at the ramified places - the theorems obtained in the split case transfer to the general case. In particular, the maximal orders of a principal region still transform into one another by extension ideals coming from the chosen maximal subfield.

\subsection{Matrix units of split crossed products and complementary bases}

Let \(\Omega\) be an arbitrary ground field. Let \(k/\Omega\) be a separable Galois extension of degree \(n\), with Galois group \(G\). Consider the crossed product with factor system one:
\[
  K=\sum_{S\in G} u_S k,
\]
with
\[
  z u_S=u_S z^S\qquad(z\in k),
\]
and
\[
  u_Su_T=u_{ST}.
\]
It is a split algebra. Put
\[
  E_1=\sum_{S\in G} u_S .
\]
Then \(E_1\) generates the identity representation of the group and satisfies the elementary relations
\[
  zE_1=E_1 z,
\]
for invariant \(z\), and more generally the trace relation
\[
  E_1 z E_1=E_1\operatorname{Tr}_{k/\Omega}(z).
\]
Here \(\operatorname{Tr}\) is the trace of the element \(z\) of \(k\) over \(\Omega\).

Let \(a_1,\ldots,a_n\) be a basis of \(k/\Omega\), and let \(a_1^*,\ldots,a_n^*\) be the complementary trace basis, so that
\[
  \operatorname{Tr}_{k/\Omega}(a_i a_j^*)=\delta_{ij}.
\]
Then
\[
  e_{ij}=a_i E_1 a_j^*
\]
form a complete system of matrix units in \(K\). Indeed,
\[
  e_{ij}e_{rs}=a_iE_1(a_j^*a_r)E_1a_s^*
  =a_iE_1\operatorname{Tr}(a_j^*a_r)a_s^*
  =\delta_{jr}e_{is}.
\]
Thus the split crossed product is explicitly the full matrix algebra, and the basis-complement relation is the device which supplies the matrix units.

The construction survives scalar extension. If \(Z\) is any commutative extension field of \(\Omega\), the same formulas hold after replacing \(k\) by \(kZ\). When \(kZ\) becomes a direct sum of fields, the group substitutions are defined to act trivially on \(Z\), in accord with the fact that the relevant elements commute elementwise with \(Z\). Thus the normal form is not a peculiarity of field extensions with one component; it is a statement about the separable algebra obtained by extension of scalars.

Assume now that \(\Omega\) has characteristic zero. The displayed normal form also holds when \(k\) is only a non-Galois splitting field. Let \(L\) be the Galois closure of \(k/\Omega\), with Galois group \(G\), and let \(H\) be the subgroup corresponding to \(k\). For the split crossed product built from \(L\) one forms the idempotent belonging to the identity representation of \(H\). Compressing by that idempotent gives back the algebra split by \(k\). The same trace-basis calculation then gives the matrix units. This proves, in particular, that a full matrix algebra over \(\Omega\) contains every degree-\(n\) field as a maximal subfield in a way compatible with the normal form.

The same passage from the Galois closure to the non-Galois splitting field works even when the algebra is not split. If \(L\) is a Galois splitting field and \(k\) corresponds to \(H\), the idempotent attached to \(H\) cuts out an algebra isomorphic to the original algebra with maximal subfield \(k\). It is the endomorphism ring of a suitable left ideal, and its rank is the right one because the left ideal decomposes into \(h=|H|\) mutually isomorphic summands after passing to the Galois closure. In the split case this reduces to the preceding construction.

\subsection{Maximal orders in the split case}

Let the ground field \(\Omega\) now be an algebraic number field. Let \(o\) be its maximal order, and let \(k/\Omega\) be a splitting field with maximal order \(\mathfrak o\). If \(\mathfrak a\) is a full \(o\)-module in \(k\), let \(\mathfrak a'\) denote the complementary module with respect to the trace form:
\[
  \mathfrak a'=\{x\in k: \operatorname{Tr}_{k/\Omega}(x\mathfrak a)\subseteq o\}.
\]
In the split algebra write again \(E_1\) for the idempotent belonging to the identity representation. The maximal orders and their ideals are described by the following formulas.

For a full module \(\mathfrak a\) put
\[
  \mathcal O_{\mathfrak a}=\mathfrak a E_1\mathfrak a'.
\]
Then \(\mathcal O_{\mathfrak a}\) is a maximal order. Conversely, every maximal order of the split algebra is obtained in this way. The left ideals between such maximal orders are exactly the modules
\[
  \mathfrak aE_1\mathfrak b,
\]
where \(\mathfrak a\) and \(\mathfrak b\) are full modules in \(k\). The reciprocal right ideal is
\[
  \mathfrak b'E_1\mathfrak a'.
\]
The intersection of \(\mathcal O_{\mathfrak a}\) with \(k\) is the coefficient order of \(\mathfrak a\). Hence the maximal orders whose intersection with \(k\) is a fixed order \(\mathfrak t\) form the region attached to \(\mathfrak t\).

The proof is local in nature. A module in the algebra is determined by its components at all places. If a module is contained in another at every completion, then it is contained globally; if it is properly enlarged, it is properly enlarged at some place. Hence maximality can be checked componentwise. At a local principal ideal ring choose complementary bases for \(\mathfrak a\) and \(\mathfrak a'\). The products
\[
  a_iE_1a_j^*
\]
are then matrix units, so the local order is the full matrix order and therefore maximal. The trace formula also gives
\[
  \operatorname{Tr}(\mathfrak a\mathfrak a')=o,
\]
which is the same assertion written in the language of complementary modules.

Let \(\mathcal O=\mathfrak aE_1\mathfrak a'\). A left ideal \(\mathcal L\) of \(\mathcal O\) is generated, as a module, by components of the form \(E_1\lambda\). Writing each \(\lambda\) in the splitting field gives
\[
  \mathcal L=\mathfrak aE_1\mathfrak b
\]
for some module \(\mathfrak b\). Regularity of the ideal is exactly the condition that \(\mathfrak b\) have full rank. Composition of such ideals corresponds to multiplication of the right-hand modules. In particular, two maximal orders in the same region transform into one another by ideals made from the modules belonging to the common coefficient order.

In the principal region, where the intersection with \(k\) is the maximal order \(\mathfrak o\), the statement becomes especially simple. The maximal orders in the principal region transform into one another by extension of ideals of \(k\). Equivalently, those left ideals of a maximal order in the principal region which are prime to the different and whose right order contains \(\mathfrak o\) are precisely the extended ideals from \(k\). This gives the arithmetic content of the explicit formulas above.

\subsection{Maximal orders of arbitrary crossed products}

The results just obtained remain true for an arbitrary simple algebra once the division into regions is refined. Let \(A\) be a central simple algebra over \(\Omega\), and let \(k\) be a maximal commutative subfield. Two maximal orders of \(A\) are assigned to the same region relative to \(k\) if they have the same intersection with \(k\) and, in addition, have the same local components at the ramified places of \(A\). If the common intersection is the maximal order of \(k\), the region is a principal region.

For principal regions one has the following theorem. The maximal orders of a principal region transform into one another by extension ideals of ideals from \(k\). In other words, the left ideals of maximal orders in a principal region which are prime to the different and whose right order contains the maximal order of \(k\) are precisely the extended ideals from \(k\).

The reason is that two maximal orders in a principal region can differ only at finitely many places, and by definition not at the ramified ones. At the remaining places the algebra is split. There it is of the form treated above, even if the splitting field is not Galois, by passage to the Galois closure and compression by the suitable idempotent. Thus the local ideals are extension ideals. Gluing the finitely many local data gives the global assertion.

If the algebra has prime degree, there is only one principal region. Indeed, locally at a ramified place it remains a division algebra, so no further ramified-place distinction can occur within the principal intersection condition. Hence the principal region is stable under extension ideals from \(k\). The most general transformation is obtained by composing these extension ideals with the ideals which transform a fixed maximal order into itself, namely the two-sided ideals. These differ from center ideals only at the ramified places; thus the statement about ideals of a fixed maximal order follows.

For arbitrary regions the corresponding assertion is local. The maximal orders of a region transform into one another by ideals prime to the different whose local components are built, at every place, from modules of the associated order in exactly the manner described in the split case. For a fixed ideal, the components differ from the order at only finitely many places. This is a direct consequence of the split calculation after replacing the operators \(u_S\) at the separate places by equivalent local ones.

There remains a finer question. Does every order in the maximal subfield occur as the intersection of a maximal order with that field? This is a question at the ramified places. Principal regions always exist, since a crossed-product presentation, or the generalized version of such a presentation, gives an order containing the maximal order of the subfield after the factor systems have been chosen integral.

One can also ask whether the arithmetic division into regions characterizes the algebraic splitting fields. Do distinct splitting fields always give distinct region divisions? Do they already give different principal regions? Do the principal regions belonging to all splitting fields exhaust all maximal orders? A single maximal order cannot provide such a characterization. Already in the quaternion algebra one finds examples where one maximal order contains both the principal order of a quadratic subfield and the principal order of the field of third roots of unity.

\clearpage
\section{43. Ideal Differentiation and the Different}

\textit{Original publication: J. reine angew. Math. 188 (1950), pp. 1--21. The paper was left by Noether and written in the winter of 1927/28; it was published posthumously.}

The theory of the different is usually developed in parallel with the differential calculus of algebraic equations. If an integral element \(x\) satisfies a defining equation \(f(x)=0\), the derivative \(f'(x)\) enters the different. This suggests that the different is not merely analogous to differentiation but is itself an ideal-theoretic form of differentiation. The following pages make this precise. The method begins with purely algebraic constructions for rings, then specializes to commutative rings and orders. The usual arithmetic different appears as the ideal derivative of a defining ideal.

The guiding idea is this. To extend the coefficient ring of a ring \(D\) from a base ring \(o\) to a larger ring \(h\), one has to construct a direct product of \(D\) with \(h\) over \(o\). If \(D\) is presented by generators and relations over \(o\), the same relations after coefficient extension define the product, provided the compatibility conditions are satisfied. The ideal which records these relations is the defining ideal. Differentiating this defining ideal, in the ordinary formal sense with respect to the generators, produces the different.

\subsection{Extension of the coefficient domain of a ring}

Let \(D\) and \(h\) be rings with common subring \(o\). A ring \(D_h\) is called the direct product of \(D\) with \(h\) over \(o\) if it is generated by isomorphic images of \(D\) and \(h\), if the two images have the same copy of \(o\), and if every element of \(D_h\) is a finite sum of products
\[
  x y\qquad(x\in D,\\ y\in h).
\]
The multiplication has to be bilinear over \(o\), and the two factors must commute with each other in the sense required by the construction. Thus, in the commutative case, one may think of
\[
  D_h=D\otimes_o h,
\]
but the definition is set up so as to apply before commutativity is imposed.

A necessary and sufficient condition for the direct product to exist is that all equality and multiplication relations among the chosen generators of \(D\) remain valid after the coefficients have been extended to \(h\), and similarly for the relations in \(h\). If \(D\) is given by generators \(Z_1,\ldots,Z_m\) over \(o\), let
\[
  o[Z_1,\ldots,Z_m]
\]
be the corresponding polynomial ring or non-commutative polynomial ring, as appropriate. The kernel of the natural map to \(D\) is the defining ideal \(M\). After coefficient extension the candidate direct product is
\[
  h[Z_1,
  \ldots,Z_m]/hM.
\]
It is the true direct product precisely when no new relation collapses the image of \(h\) or of \(D\) beyond those forced by \(M\).

Equivalently, the direct product has the following universal property. Given any ring \(R\) containing compatible images of \(D\) and \(h\) over \(o\), there is a unique homomorphism
\[
  D_h\longrightarrow R
\]
extending the two maps. In this form the construction is independent of the chosen generators. The defining ideal, however, is most convenient for calculation.

If \(D\) and \(h\) are commutative and \(D\) is finite over \(o\), the coefficient extension can be checked by bases. Let \(x_1,\ldots,x_n\) be elements of \(D\) which form an independent \(o\)-module basis. Then the products
\[
  x_ix_j=\sum_k c_{ij}^{k}x_k,
\]
with \(c_{ij}^k\in o\), determine the multiplication after extension to \(h\) by the same structure constants. This gives a direct product with \(h\)-basis \(x_1,
\ldots,x_n\), unless the extension of scalars introduces new linear dependence. Thus independence of the basis after extension is the essential condition.

For ideals and modules one uses the same language. If \(B\) is an \(o\)-module in \(D\), its extension is
\[
  B_h=B\otimes_o h
\]
inside the coefficient-extended ring. If \(C\) is an \(h\)-module in \(D_h\), its contraction is its intersection with \(D\). When \(B\) is an ideal, its extension is an ideal of \(D_h\); conversely, contraction of ideals is compatible with multiplication when the direct product is exact enough. This extension-contraction formalism is the module-theoretic counterpart of passing from an equation to its equations after scalar extension.

\subsection{Rings with an independent module basis}

Assume now that \(D\) has an independent basis \(x_1,
\ldots,x_n\) over \(o\). The direct product with \(h\) is constructed by taking all formal sums
\[
  \sum_i x_i a_i,
  \qquad a_i\in h,
\]
and multiplying them by the same structure constants as in \(D\). Associativity and distributivity follow from the corresponding identities in \(D\), because the basis is independent and the structure constants lie in \(o\). Thus the direct product exists and has \(h\)-basis \(x_1,
\ldots,x_n\).

Let \(B\) be an \(o\)-module in \(D\). Suppose \(B\) has an independent \(o\)-basis which can be enlarged to an independent basis of \(D\). Then extension and contraction behave without loss:
\[
  (B_h)\cap D=B.
\]
If \(B\) is an ideal, then \(B_h\) is an ideal, and contraction of \(B_h\) is again \(B\). This is the basic technical fact needed later. It is the reason that the different can be calculated after passing to a convenient coefficient extension.

A sufficient criterion for the existence of the direct product is therefore the following. If the ring and the modules under consideration possess independent bases over the coefficient ring, and if the multiplication tables have coefficients in that coefficient ring, then scalar extension gives the direct product. In arithmetic language this applies to orders in separable extensions and to the modules which occur as ideals of those orders.

\subsection{The different of a ring over a subring}

From here on the rings are assumed commutative unless explicitly stated otherwise. Let \(D\) be a commutative ring over \(o\), and suppose the direct product
\[
  D\otimes_o D
\]
exists in the sense described above. Multiplication in \(D\) gives a homomorphism
\[
  \mu:D\otimes_o D\longrightarrow D,
  \qquad x\otimes y\mapsto xy.
\]
The kernel of \(\mu\) is generated by the formal differences
\[
  x\otimes1-1\otimes x.
\]
The different is the ideal of \(D\) determined by the annihilator of this kernel. In modern notation this is the same as the zeroth Fitting ideal of the module of differentials, but Noether's construction is phrased entirely in terms of ideals and direct products.

Let \(D\) be given over \(o\) by generators \(Z_1,
\ldots,Z_m\) and defining ideal \(M\). For a polynomial \(F\in M\), form the formal differential
\[
  dF=\sum_i \frac{\partial F}{\partial Z_i}
  dZ_i.
\]
The ideal generated in \(D\) by the appropriate minors of the matrix of the partial derivatives of a system of generators of \(M\) is the ideal derivative of \(M\). This is the differential quotient of the defining ideal. The theorem is that this ideal derivative is exactly the different just defined by means of the direct product.

In the simplest case one generator is enough. Suppose
\[
  D=o[x],\qquad f(x)=0,
\]
and \(f(X)\) is a defining polynomial. Then the direct product is described by two elements \(x\) and \(x'\), with
\[
  f(x)=f(x')=0.
\]
The kernel of multiplication is generated by \(x'-x\). Since
\[
  f(x')-f(x)=(x'-x)\,q(x,x'),
\]
putting \(x'=x\) gives
\[
  q(x,x)=f'(x).
\]
Thus the different is generated by
\[
  f'(x),
\]
the ordinary derivative of the defining equation. The general case is the same calculation with several generators and the Jacobian matrix of the defining relations.

This proves that the classical derivative appearing in the different is not an accident of notation. It is the formal derivative of the defining ideal, transported through the direct-product description of coefficient extension.

\subsection{The different of a direct sum}

If a ring decomposes as a direct sum of components,
\[
  D=D_1\oplus\cdots\oplus D_r,
\]
then the different decomposes correspondingly. The direct product of \(D\) over \(o\) decomposes into the direct sum of the pairwise products \(D_i\otimes_o D_j\). The multiplication map is zero on the mixed components and is the ordinary multiplication on each diagonal component. Hence the kernel and its annihilator split by components. Consequently
\[
  \mathfrak D(D/o)=\mathfrak D(D_1/o)\oplus\cdots\oplus\mathfrak D(D_r/o),
\]
with the evident interpretation of the direct-sum ideal.

This observation is important because coefficient extension of a field often turns a field into a direct sum of conjugate fields. The different can then be computed componentwise. In particular, after a separable scalar extension to a splitting field, the structure of the direct product becomes diagonal, and the different is read off from the complementary bases of the components.

\subsection{Galois extension rings}

Let \(K/k\) be a finite Galois extension, and let \(D\) be an order in \(K\) over an order \(o\) of \(k\). After extending coefficients from \(o\) to \(D\), the direct product
\[
  D\otimes_o D
\]
decomposes into components corresponding to the automorphisms \(S\) of the Galois group. In the field case this is the familiar decomposition
\[
  K\otimes_k K \cong \prod_{S\in G} K,
\]
where the component indexed by \(S\) is obtained by the map
\[
  x\otimes y\longmapsto x\,y^S.
\]
For orders, one obtains modules inside the corresponding components.

The kernel of multiplication is the part which vanishes on the identity component. The other components measure the differences between an element and its conjugates. Thus the different is controlled by the product of the ideals generated by
\[
  x-x^S,
  \qquad S\ne 1,
\]
as \(x\) ranges through the order. This is the ideal-theoretic form of the discriminant calculation.

Equivalently, choose complementary bases \(\omega_1,
\ldots,
\omega_n\) and \(\omega_1^*,
\ldots,\omega_n^*\) for the trace pairing. Then the elements
\[
  \sum_i \omega_i\otimes \omega_i^{*S}
\]
are the component idempotents of the extended ring. The denominator ideals needed to make these idempotents integral are exactly the complementary module, and their inverse is the different. In this way the component representation of the Galois extension ring gives the standard formula
\[
  \mathfrak D^{-1}=\{x\in K: \operatorname{Tr}_{K/k}(xD)\subseteq o\}
\]
for the inverse different, and hence identifies the different itself with the inverse of the complementary module.

This is the structural core of the theory. The complementary module is not introduced as an auxiliary analytic object; it is the module of denominators needed to separate the components of the direct product. The different is its inverse.

\subsection{The different of an order}

Let \(D\) now be an order in a separable field extension \(K/k\), with coefficient order \(o\). The different \(\mathfrak D_{D/o}\) is the inverse of the complementary module
\[
  D^*=
  \{x\in K: \operatorname{Tr}_{K/k}(xD)\subseteq o\}.
\]
Thus
\[
  \mathfrak D_{D/o}=(D^*)^{-1}.
\]
If \(D\) has an independent \(o\)-basis \(\omega_1,
\ldots,
\omega_n\), the discriminant is
\[
  \Delta(\omega_1,
\ldots,
\omega_n)=
  \det(\operatorname{Tr}(\omega_i\omega_j))
\]
and the product of the basis ideal with the different gives the discriminant ideal. More precisely, changing the basis multiplies the determinant by the square of the determinant of the change-of-basis matrix, exactly as in the usual discriminant theory; the ideal formulation absorbs this dependence and gives an invariant ideal.

If the order is monogenic,
\[
  D=o[\theta],
\]
with fundamental equation \(f(\theta)=0\), then the different is generated by
\[
  f'(\theta)
\]
provided the basis \(1,\theta,
\ldots,
\theta^{n-1}\) is an integral basis of the order. In the non-monogenic case one replaces the one derivative by the ideal generated by the Jacobians of a system of defining relations. This is precisely the differential quotient of the defining ideal.

The formula is compatible with localization. At a prime ideal \(\mathfrak p\), localizing the order and the trace dual localizes the different:
\[
  (\mathfrak D_{D/o})_{\mathfrak p}
  =\mathfrak D_{D_{\mathfrak p}/o_{\mathfrak p}}.
\]
Therefore all assertions can be checked locally. Locally at a Dedekind prime, the different exponent is the valuation of the derivative in a primitive-element presentation, or equivalently the length of the module of differentials.

\subsection{Ramification theory of the principal order modulo powers of a prime}

The last part of the paper relates the ideal-theoretic different to ramification. Let \(\mathfrak p\) be a prime of \(o\), and let
\[
  \mathfrak pD=\prod_i \mathfrak P_i^{e_i}
\]
be its decomposition in the order or maximal order of the extension. The different exponent at \(\mathfrak P_i\) measures the failure of the residue extension and ramification to be unramified. In the tame case it is simply
\[
  e_i-1.
\]
In the wild case additional terms occur, represented by the higher ramification groups. Noether's formulation obtains these terms from the ideal derivative: the higher powers of \(\mathfrak P_i\) which divide all differences \(x-x^S\) are exactly the powers which enter the local different.

In a Galois extension the local different can therefore be written in terms of the ramification groups. If \(G_j\) denotes the usual decreasing filtration of the decomposition group, then the exponent of the different is
\[
  \sum_{j\ge0} (|G_j|-1),
\]
which is the Dedekind-Hensel formula. In the present framework this formula is not an external ramification theorem but the local expression of the same ideal-differentiation construction. The derivative of the defining ideal records, component by component in the direct product, the orders of contact of conjugate components; these orders of contact are exactly the ramification jumps.

Thus the different is simultaneously three objects: the annihilator attached to multiplication in the direct product, the ideal derivative of the defining relations, and the inverse of the trace-complementary module. The equality of these descriptions is the content of ideal differentiation.


\clearpage
\section{44. Algebra of Hypercomplex Quantities}

\textit{Lecture by Emmy Noether, winter semester 1929/30; written up by Max Deuring. This batch gives the lecture text through the applications of the cyclic special case.}

\subsection{Introduction}

These lectures treat the general representation theory of rings. This theory grew, on one side, out of the representation theory of finite groups and, on the other side, out of the algebraic and arithmetic theory of hypercomplex number systems. The more general point of view has an advantage over the older theory of Frobenius and Schur. One does not represent merely a group; one represents the group ring, and more generally a hypercomplex system or an arbitrary ring. Thus one may use the theory of rings and ideals. This frees representation theory from many unwieldy computations.

The use of general ring and ideal theory does more than simplify the subject. It also carries it further. It clarifies the relation of a group to its subgroups, and it opens new fields of application. In particular, general representation theory gives a new and elegant foundation of Galois theory; still more importantly, it gives a Galois theory for non-commutative fields.

The basic concepts of group, ring, and ideal theory, and the general theorems of representation theory, are those developed in Noether's paper \textit{Hypercomplex Quantities and Representation Theory}. The present lecture notes are to be read as a supplement to that work. The notions of representation and representation module are recalled briefly, chiefly in order to introduce reciprocal representation modules and reciprocal representations.

\subsection{Chapter I. Representations}

Let \(o\) be a ring, the ring to be represented, and let \(T\) be another ring, the ring in which the representation is to take place.

A direct representation of degree \(n\) of \(o\) in \(T\) is a subring \(O\) of the full matrix ring \(T_n\), together with a ring homomorphism from \(o\) onto \(O\). Symbolically,
\[
  o\twoheadrightarrow O\subset T_n.
\]
A reciprocal representation of degree \(n\) of \(o\) in \(T^*\) is defined similarly, but the map from \(o\) is a reciprocal ring homomorphism. Thus if \(c,d\in o\) correspond to \(c^*,d^*\), then \(c+d\) corresponds to \(c^*+d^*\), while \(cd\) corresponds to
\[
  d^*c^*.
\]
As an example, in a commutative coefficient field, transposition of matrices is a reciprocal ring isomorphism. In general, for every ring \(T\) one constructs a reciprocal ring \(T^*\) by assigning to each \(c\in T\) a symbol \(c^*\), putting
\[
  c^*+d^*=(c+d)^*,\qquad c^*d^*=(dc)^*.
\]
The symbols \(c^*\) then form a ring reciprocally isomorphic to \(T\).

Representations are collected into classes. A direct representation \(O\subset T_n\) is transformed into an equivalent representation by a regular matrix \(P\): instead of the matrix \(C\) attached to an element \(c\in o\) one uses
\[
  P^{-1}CP.
\]
All representations obtained from one another in this way form a direct representation class. The definition for reciprocal representations is the same.

A direct representation module of \(o\) with respect to \(T\) is an \(o,T\)-bimodule \(M\): it is a left \(o\)-module, a right \(T\)-module, and it satisfies the associativity law
\[
  (cm)t=c(mt)
\]
for \(c\in o\), \(m\in M\), and \(t\in T\). It is also required to be a finite direct sum of one-generator right \(T\)-modules,
\[
  M=x_1T+\cdots+x_nT,
\]
with the usual independence condition. A reciprocal representation module is a left \(o\)-module and a left \(T^*\)-module \(M^*\), satisfying
\[
  c(t^*m)=t^*(cm),
\]
and again decomposing as a finite direct sum of simple one-generator \(T^*\)-modules.

Every representation module determines a representation class. For a direct module \(M=x_1T+
\cdots+x_nT\), multiplication by an element \(c\in o\) is a \(T\)-linear transformation of \(M\) into itself, because
\[
  c(mt)=(cm)t.
\]
With respect to the basis \(x_1,
\ldots,x_n\) it is therefore given by a matrix \(C=(c_{ij})\), determined by
\[
  c x_j=\sum_i x_i c_{ij}.
\]
The matrices belonging to the elements of \(o\) form a homomorphic image of \(o\) in \(T_n\). The compatibility with addition and multiplication is immediate:
\[
  (c+d)x_j=cx_j+dx_j,
\]
\[
  (cd)x_j=c(dx_j).
\]
Thus a representation of degree \(n\) is obtained.

Changing the basis changes the matrices by similarity. If
\[
  (y_1,
\ldots,y_n)=(x_1,
\ldots,x_n)P
\]
with \(P\) regular, then the matrix \(C\) is replaced by
\[
  P^{-1}CP.
\]
Hence a module gives not a single representation but a whole representation class. Conversely, every representation class comes from a module. Starting with a representation \(c\mapsto C\), introduce a free right \(T\)-module
\[
  M=x_1T+
\cdots+x_nT
\]
and define left multiplication by \(o\) through
\[
  c(x_1,
\ldots,x_n)=(x_1,
\ldots,x_n)C.
\]
The homomorphism relations for the matrices are exactly the module axioms. Thus representation classes and representation modules are equivalent descriptions. Reciprocal representations and reciprocal modules behave in the same way, with the order of multiplication reversed.

\subsection{Chapter II. Galois theory of commutative fields}

The representation theory just developed is already sufficient to found Galois theory. Its advantage over the usual exposition is that it does not depend on special generating systems or on theorems about polynomials in indeterminates. Instead, it rests on coefficient extension of hypercomplex systems.

Let
\[
  A=c_1P+
\cdots+c_nP
\]
be a hypercomplex system over a commutative field \(P\), and let \(Z\) be an extension ring of \(P\) whose center contains \(P\). The extended system \(A_Z\) consists of the formal sums
\[
  \sum_i c_i z_i,
  \qquad z_i\in Z,
\]
with addition and scalar multiplication defined componentwise, and with multiplication determined by the multiplication table of \(A\). Thus if
\[
  c_ic_j=\sum_k a_{ij}^k c_k,
  \qquad a_{ij}^k\in P,
\]
then
\[
  \left(\sum_i c_i z_i\right)
  \left(\sum_j c_j w_j\right)
  =\sum_{i,j,k} c_k a_{ij}^k z_iw_j.
\]
This is the ordinary extension of coefficients.

Let now \(K/P\) be a commutative field extension of degree \(n\). If \(K\) is separable, then after a suitable extension of coefficients the algebra \(K\) decomposes as a direct sum of fields, one for each embedding. In representation-theoretic language, the regular representation decomposes into one-dimensional constituents. The automorphisms of the field are the isomorphisms of these constituents compatible with the multiplication.

For a field \(K\), an isomorphism over \(P\) is an operation preserving addition, multiplication, and the elements of \(P\). If \(K\) is of the first kind over \(P\), meaning that its scalar extension is completely reducible and has no radical, then the number of distinct \(P\)-isomorphisms into a normal closure equals the degree exactly when \(K/P\) is Galois. In that case the isomorphisms form a group.

The main theorem of Galois theory may then be stated without choosing a primitive element. Let \(K/P\) be a Galois extension with group \(G\). For a subfield \(T\) of \(K\), its invariant group is
\[
  G_T=\{S\in G: x^S=x\text{ for all }x\in T\}.
\]
For a subgroup \(H\subset G\), its invariant field is
\[
  K^H=\{x\in K: x^S=x\text{ for all }S\in H\}.
\]
Then:
\begin{enumerate}[label=\arabic*.]
\item If \(H\) is the invariant group of \(T\), then \(T\) is the invariant field of \(H\).
\item If \(T\) is the invariant field of \(H\), then \(H\) is the invariant group of \(T\).
\end{enumerate}
The first assertion follows from the extension theorem for isomorphisms: an isomorphism of a subfield can be extended to an isomorphism of the whole normal field. If an element of \(K\) is fixed by every automorphism fixing \(T\), then the one-dimensional component it defines in the extended regular representation is already determined by \(T\), and hence the element belongs to \(T\). The second assertion is reduced to the first by applying it to the invariant field of the subgroup.

This proof is deliberately structural. The usual arguments with equations are replaced by the decomposition of the extended system and by the behavior of isomorphisms under restriction and extension. It is precisely this form of the proof that can later be transferred to non-commutative fields.

\subsection{Chapter III. Abelian groups}

The next application concerns abelian groups and their group rings. Let \(G\) be a finite abelian group and let \(P\) be a field containing enough roots of unity, or else pass to an extension field which does. The group ring
\[
  P[G]
\]
is then a commutative semisimple algebra. Its irreducible representations are one-dimensional and are given by characters
\[
  \chi:G\longrightarrow P^{\times}.
\]
The algebra decomposes into a direct sum of fields corresponding to these characters.

The character relations are obtained from the matrix-unit or idempotent decomposition. For a character \(\chi\) define
\[
  e_\chi=\frac1{|G|}\sum_{g\in G}\chi(g^{-1})g.
\]
Then, when the characteristic does not divide \(|G|\), the \(e_\chi\) are primitive orthogonal idempotents and
\[
  1=\sum_\chi e_\chi.
\]
The orthogonality relations follow immediately:
\[
  \sum_{g\in G}\chi(g)\psi(g^{-1})=
  \begin{cases}
  |G|,&\chi=\psi,\\
  0,&\chi\ne\psi.
  \end{cases}
\]
No separate character calculation is needed; it is simply the formula for the decomposition of the group ring into primitive components.

The Galois theory of abelian groups is then the Galois theory of the commutative algebra \(P[G]\). Subgroups correspond to quotient character groups, and quotient groups correspond to subalgebras generated by the corresponding idempotent sums. In this form the duality between subgroups and character groups is just the duality between invariant rings and invariant groups already established.

If the field of coefficients is not large enough to contain the values of all characters, one extends coefficients to a splitting field and then descends. The Galois group of the coefficient extension permutes the primitive idempotents. The rational irreducible constituents over the original field are the sums over the orbits of this action. This is again the same method as before: extend to a completely reducible algebra, decompose, and then collect conjugate components.

\subsection{Chapter IV. Two-sided simple rings}

Let \(A\) be a two-sided simple hypercomplex system over a field \(P\). The basic lemma used throughout is the density statement for simple modules: if \(M\) is a simple representation module, then the ring of endomorphisms commuting with the action is a division ring, and the action of \(A\) on \(M\) is the full ring of endomorphisms over that division ring. Equivalently, a two-sided simple algebra is a full matrix algebra over a division algebra:
\[
  A\cong M_r(D).
\]
The center \(Z(D)\) is a field containing \(P\), and \(A\) is central simple over its center after replacing \(P\) by that center.

Representations of a two-sided simple system are controlled by its simple modules. Every representation module is a direct sum of isomorphic simple modules once the coefficient field is large enough to split the division algebra. Over the original field, the endomorphism division algebra measures the failure of the representation to be absolutely irreducible. The degree of this division algebra is the Schur index.

The passage from an arbitrary coefficient field to a splitting field has two effects. First, the center may decompose into conjugate components if it is not already a field after scalar extension. Second, each central simple component may become a full matrix algebra. The irreducible representations over the splitting field are then the natural column modules of these matrix algebras. Their conjugacy under the coefficient-field Galois group describes the irreducible representation classes over the original field.

This yields the non-commutative analogue of the field-theoretic Galois statements. In the commutative case a field is recovered as its own maximal commutative subfield. In the non-commutative case one must work with maximal commutative subfields and with their normalizers. The automorphism groups are replaced by groups of inner automorphisms, or by factor systems measuring the obstruction to representing automorphisms by elements satisfying the group law strictly.

\subsection{The Galois theory of non-commutative fields}

Let \(D\) be a division algebra over a field \(P\), and let \(H\) be a group of automorphisms of \(D\). The invariant field is
\[
  D^H=\{x\in D:x^S=x\text{ for all }S\in H\}.
\]
Conversely, for a subfield or subdivision ring \(E\subset D\), the invariant group consists of automorphisms fixing \(E\) elementwise. The aim is to obtain the same two assertions as in the commutative case:
\begin{enumerate}[label=\arabic*.]
\item the invariant group of \(E\) has invariant field \(E\);
\item the invariant field of \(H\) has invariant group \(H\), under the appropriate closedness condition.
\end{enumerate}
The proofs run parallel to the commutative proofs after replacing embeddings by simple components of representation modules and replacing ordinary automorphisms by inner automorphisms in a suitable simple algebra.

The normalizer of a maximal subfield plays the role of the Galois group. If \(K\) is a maximal commutative subfield of a central simple algebra \(A\), then an element \(u\in A^{\times}\) which normalizes \(K\),
\[
  u^{-1}Ku=K,
\]
induces an automorphism of \(K\). The obstruction to choosing such elements multiplicatively is measured by a factor system:
\[
  u_Su_T=u_{ST}a_{S,T},\qquad a_{S,T}\in K^{\times}.
\]
Thus non-commutative Galois theory naturally produces crossed products.

\subsection{Chapter V. Factor systems}

Let \(G\) be a group acting on a field \(K\). A factor system is a set of elements
\[
  a_{S,T}\in K^{\times}\qquad(S,T\in G)
\]
satisfying
\[
  a_{S,T}^{R}a_{ST,R}=a_{T,R}a_{S,TR}.
\]
This is exactly the associativity condition for symbols \(u_S\) with
\[
  u_Su_T=u_{ST}a_{S,T},
  \qquad z u_S=u_Sz^S.
\]
The resulting crossed product is
\[
  A=(a_{S,T},K,G)=\sum_{S\in G}u_SK.
\]

Changing the representatives \(u_S\) by
\[
  v_S=u_Sc_S
\]
changes the factor system to
\[
  a'_{S,T}=a_{S,T}c_S^Tc_Tc_{ST}^{-1}.
\]
Associated factor systems determine similar algebras. The classes of associated factor systems form an abelian group under componentwise multiplication. The identity consists of all transformation systems, i.e. all systems associated to the unit system. This group is the subgroup of the Brauer group consisting of classes split by \(K\).

A crossed representation with factor system \(a_{S,T}\) consists of matrices \(C_S\) over a coefficient ring satisfying
\[
  C_S^T C_T=C_{ST}a_{S,T}
\]
with the appropriate interpretation of the action of \(G\) on coefficients. Two such representations are equivalent if they differ by a change of basis. In this language ordinary representations are the special case \(a_{S,T}=1\). The representation theory of factor systems is therefore ordinary representation theory with multiplication twisted by a two-cocycle.

Multiplication of factor systems corresponds to tensor product of crossed representations. The inverse class is obtained by taking the reciprocal system. Powers of a factor-system class have a direct algebraic meaning: the \(m\)-th power corresponds to the tensor product of \(m\) copies of the associated crossed product. The index of the class is the degree of the division algebra in its Brauer class; it is also the degree of the irreducible crossed representations belonging to the system.

\subsection{Normal form with Galois maximal subfield}

Suppose a simple algebra \(A\) has a Galois maximal subfield \(K\) over its center \(P\). Then \(A\) can be written as a crossed product. For each automorphism \(S\) of \(K/P\) choose an element \(u_S\in A^{\times}\) inducing it by conjugation:
\[
  u_S^{-1}zu_S=z^S.
\]
Then every element of \(A\) has a unique expression
\[
  \sum_{S
\in G} u_S z_S,
  \qquad z_S\in K.
\]
The products of the \(u_S\) yield the factor system
\[
  u_Su_T=u_{ST}a_{S,T}.
\]
Associativity gives the cocycle relation. Conversely, any such crossed product is a central simple algebra with maximal subfield \(K\). Thus crossed products are not an artificial construction; they are the normal form of central simple algebras with Galois maximal subfield.

The normal form also clarifies splitting. The algebra is split if and only if the factor system is associated to one. Equivalently, there are elements \(c_S\in K^{\times}\) such that
\[
  a_{S,T}=c_S^Tc_Tc_{ST}^{-1}.
\]
Then the substitution \(v_S=u_Sc_S^{-1}\) makes
\[
  v_Sv_T=v_{ST},
\]
and the algebra is the endomorphism algebra of \(K\) over \(P\), hence a full matrix algebra.

\subsection{Chapter VI. Theory of crossed products}

Let \(K/P\) be Galois with group \(G\). The crossed product attached to a factor system \(a_{S,T}\) is the direct sum
\[
  A=\bigoplus_{S\in G}u_SK.
\]
It satisfies
\[
  z u_S=u_S z^S,
\]
\[
  u_Su_T=u_{ST}a_{S,T}.
\]
The center of \(A\) is \(P\), and \(K\) is a maximal commutative subfield. The algebra is simple. To see this, let \(I\) be a nonzero two-sided ideal and choose an element of \(I\) with the smallest possible number of nonzero components. Commuting it with elements of \(K\) eliminates all but one component unless the element already lies in \(K\). Since \(I\cap K\ne0\), and \(K\) is a field, \(1\in I\). Thus \(I=A\).

The product theorem for crossed products says that multiplying factor-system classes corresponds to taking the tensor product of the corresponding algebras over the common center and then passing to the component split by the diagonal copy of \(K\). In class notation,
\[
  [(a_{S,T})]+[(b_{S,T})]=[(a_{S,T}b_{S,T})].
\]
This gives the abelian group law on Brauer classes with fixed splitting field.

The automorphism ring of a suitable left ideal supplies the reciprocal algebra. If \(L\) is a left ideal of absolute length equal to the degree of \(K/P\), its endomorphism ring is a crossed product with inverse factor system. This gives the inverse in the Brauer group and explains why factor systems and their reciprocal systems cancel under multiplication.

If a factor system has index \(m\), its \(m\)-th power is split in the appropriate sense. Conversely the least positive exponent for which a power becomes split is tied to the exponent of the Brauer class, while the index is the degree of the underlying division algebra. In cyclic cases the two numbers can be read from norm groups; in general they are governed by the structure of the factor-system group.

\subsection{Applications of the cyclic special case}

The cyclic crossed product is especially explicit. Let \(K/P\) be cyclic of degree \(n\), generated by \(S\), and let
\[
  A=(a,K,S)
\]
be given by
\[
  A=K+uK+
\cdots+u^{n-1}K,
\]
\[
  zu=uz^S,
\]
\[
  u^n=a\in P^{\times}.
\]
Changing \(u\) to \(uc\) changes \(a\) to
\[
  aN_{K/P}(c).
\]
Thus the cyclic classes split by \(K\) are represented by
\[
  P^{\times}/N_{K/P}(K^{\times}).
\]

Several classical consequences follow. First, every finite field is commutative. Indeed, the multiplicative group of a finite field is cyclic; the center is a finite field; and the cyclic norm theorem forces the relevant crossed products to split. A finite division algebra over its center must therefore be commutative.

Second, finite Galois fields are cyclic over each of their subfields. The Frobenius automorphism generates the Galois group. In the crossed-product language this is the statement that the automorphism group controlling the extension is generated by one substitution, and all norm obstructions vanish because the multiplicative group is finite cyclic.

Third, for a \(p\)-adic ground field and a cyclic splitting field \(K/P\) of degree \(n\), the cyclic algebras are classified by the residue of \(a\) modulo the norm group. The local norm theorem determines exactly when
\[
  (a,K,S)
\]
splits. Hence every cyclic splitting field of degree \(n\) over a \(p\)-adic field gives the full set of division algebras whose invariant has denominator dividing \(n\). This is the local arithmetic form of the theory of crossed products.

The point of the lecture is that these results, which historically came from group characters, equations, and special computations, are all consequences of the same ring-theoretic mechanism: extend coefficients, decompose into simple components, identify the obstruction to strict multiplication by a factor system, and interpret the resulting class as an algebra class. In this form the theory is ready for the arithmetic applications to maximal orders, ideals, norm theorems, and the principal genus theorem.



\subsection{Further details from the lecture: two-sided simple systems}

We record the auxiliary lemma used for two-sided simple systems in a form suited to later applications. Let \(A\) be a simple system over \(P\), and let \(M\) be a simple left module. If elements \(m_1,
\ldots,m_r\) of \(M\) are linearly independent over the commuting division ring \(D=\operatorname{End}_A(M)\), and if \(n_1,
\ldots,n_r\) are arbitrary elements of \(M\), then there is an element \(a\in A\) with
\[
  a m_i=n_i\qquad(i=1,
\ldots,r).
\]
This is the familiar density theorem. In Noether's terminology it says that the representation system contains all linear substitutions allowed by the operator ring. The proof is the usual minimal-annihilator proof: one chooses a relation of least length among the conditions that fail, multiplies by suitable elements of \(A\), and contradicts the simplicity of the module.

From this lemma one obtains the structure theorem. If \(A\) is two-sided simple, then there is a division ring \(D\) and an integer \(r\) such that
\[
  A\cong D_{r}=M_r(D).
\]
The simple left modules are column modules of length \(r\) over \(D\), and the right operator ring of such a module is \(D\) itself. Conversely, the full matrix ring over a division ring is two-sided simple. This establishes the equivalence between simple systems and full matrix systems over skew fields.

Representations of a two-sided simple system are therefore governed by the decomposition of the coefficient extension of its center. Let \(Z\) be the center of \(A\). If \(Z/P\) is separable and \(\Omega\) is a splitting field, then
\[
  Z_\Omega=Z\otimes_P\Omega
\]
becomes a direct sum of conjugate fields. Consequently
\[
  A_\Omega=A\otimes_P\Omega
\]
becomes the direct sum of conjugate simple algebras, and over a sufficiently large \(\Omega\) each of these becomes a full matrix algebra. The irreducible constituents of the extended representation correspond to these components. If the center is inseparable, the extended center has a radical; passing to the quotient by the radical gives the same component decomposition for the semisimple quotient, and the radical contributes only composition factors operator-isomorphic to those already obtained.

Thus a representation irreducible over \(P\) is absolutely completely reducible after extension of scalars if and only if its center is separable over \(P\). If it is not separable, one must replace the direct-sum decomposition by a composition series; nevertheless the same conjugacy and multiplicity statements survive at the level of composition factors.

\subsection{Non-commutative fields and invariant theory}

A skew field \(D\) over a ground field \(P\) is treated as a simple system. Its subfields and automorphism groups are studied through their action on the associated representation modules. Let \(H\) be a group of automorphisms of \(D\). The invariant division ring is
\[
  D^H=\{x\in D: h(x)=x\text{ for every }h\in H\}.
\]
If \(E\subset D\) is a subdivision ring, its invariant group is
\[
  H_E=\{h\in\operatorname{Aut}(D):h(x)=x\text{ for every }x\in E\}.
\]
The analogy with ordinary Galois theory is exact only after adding a closedness condition. The invariant group of a subdivision ring is closed in the automorphism group. Conversely, for closed groups one recovers the group from its invariant ring.

The proof differs from the commutative proof only in language. One embeds the skew field in a full matrix algebra over a maximal commutative subfield and uses the theorem on extension of isomorphisms. The essential point is that any isomorphism of a suitable substructure can be extended to the whole simple system; in the matrix algebra this is an inner automorphism. Therefore the invariant operations are again described by normalizers. The formula
\[
  u^{-1}Eu=E
\]
for a normalizing element \(u\) replaces the equation \(S(E)=E\) of the commutative case.

If \(E\) is the invariant ring of \(H\), and \(u\) is an element which induces an automorphism fixing \(E\), then the extended-isomorphism theorem supplies an element of the group generated by \(H\) with the same action. Hence the automorphism lies in \(H\). This proves the non-commutative analogue of the Galois correspondence in the cases considered in the lecture.

\subsection{Factor systems and crossed representations in detail}

The lecture next develops factor systems as the exact algebraic measure of the failure of a choice of representatives to be multiplicative. Let \(K/P\) be Galois with group \(G\), and suppose that for every \(S\in G\) a representative \(u_S\) has been chosen. The representatives implement the automorphisms:
\[
  u_S^{-1} z u_S=z^S,
  \qquad z\in K.
\]
Since \(u_Su_T\) and \(u_{ST}\) implement the same automorphism, they differ by an element of \(K^{\times}\):
\[
  u_Su_T=u_{ST}a_{S,T}.
\]
Associativity of multiplication gives
\[
  (u_Su_T)u_R=u_S(u_Tu_R),
\]
hence
\[
  a_{S,T}^{R}a_{ST,R}=a_{T,R}a_{S,TR}.
\]
This is the factor-system relation.

The factor system depends on the choice of representatives. If
\[
  v_S=u_Sc_S,
  \qquad c_S\in K^{\times},
\]
then
\[
  v_Sv_T=v_{ST}a'_{S,T},
\]
where
\[
  a'_{S,T}=a_{S,T}c_S^Tc_Tc_{ST}^{-1}.
\]
The two systems are associated. Association is the natural equivalence relation, because it leaves the isomorphism class of the crossed product unchanged.

A crossed representation of \(G\) with factor system \(a\) is a family of transformations \(C_S\) satisfying the same multiplication law as the representatives:
\[
  C_S^T C_T=C_{ST}a_{S,T}.
\]
If \(a=1\), this is an ordinary representation. If \(a\) is a transformation system, then the crossed representation can be untwisted by a change of basis and again becomes ordinary. Thus the representation theory depends only on the associated factor-system class.

The irreducible crossed representations belonging to \(a\) are precisely the simple modules of the crossed product \((a,K,G)\). If the corresponding algebra is a matrix algebra of degree \(r\) over a division algebra \(D\), then the irreducible crossed representation has degree equal to the degree of \(D\). This is the Schur index. Hence the Schur index is not an extra numerical invariant attached to characters; it is the degree of the division algebra determined by the factor system.

\subsection{Multiplication of factor systems and the Brauer group}

Let \(a=(a_{S,T})\) and \(b=(b_{S,T})\) be factor systems for the same Galois extension \(K/P\). Their product
\[
  (ab)_{S,T}=a_{S,T}b_{S,T}
\]
is again a factor system. If \(A=(a,K,G)\) and \(B=(b,K,G)\), then the tensor product \(A\otimes_PB\) contains a component whose factor system is \(ab\). Passing to similarity classes gives
\[
  [A]+[B]=[(ab,K,G)].
\]
This is the group law.

The identity class is represented by the unit factor system. It is also represented by any transformation system
\[
  a_{S,T}=c_S^Tc_Tc_{ST}^{-1}.
\]
The inverse of \(a\) is represented by the reciprocal factor system, or equivalently by the opposite algebra. Thus the set of factor-system classes split by \(K\) becomes an abelian group, which embeds in the Brauer group of \(P\).

Noether's formulation emphasizes that all this is already contained in ideal theory. A factor system defines a group extension
\[
  1\longrightarrow K^{\times}\longrightarrow G^*\longrightarrow G\longrightarrow1,
\]
where \(G^*\) consists of the cosets \(u_SK^{\times}\). Multiplication in the extension is determined by the factor system. Changing representatives changes the cocycle but not the extension. Multiplication of factor systems corresponds to Baer product of such extensions, and to tensor product of the algebras.

In arithmetic applications the elements \(a_{S,T}\) are replaced by ideals or ideal classes. The same formulas remain meaningful, because ideal multiplication is abelian and the Galois group acts on ideals. This is the bridge to principal genera and to maximal orders.

\subsection{Normal representation and minimal splitting fields}

Assume that \(A\) is a central simple algebra with Galois maximal subfield \(K\). The crossed-product normal form gives a representation of \(A\) on the right \(K\)-module with basis \(u_S\). The action of \(K\) is diagonal through the conjugates, while the action of \(u_S\) permutes the components and multiplies them by factor-system elements. This is the normal representation.

If the factor system is associated to one, the normal representation decomposes into ordinary matrix units. Explicitly choose \(c_S\) with
\[
  a_{S,T}=c_S^Tc_Tc_{ST}^{-1}.
\]
Then \(v_S=u_Sc_S^{-1}\) satisfies
\[
  v_Sv_T=v_{ST}.
\]
Putting
\[
  E=\sum_{S\in G}v_S
\]
one obtains idempotents and matrix units from complementary bases exactly as in the split case. Hence the algebra is a full matrix algebra.

A field \(L\) is a splitting field for \(A\) if
\[
  A\otimes_PL
\]
is a full matrix algebra over \(L\). A detachment field is a field over which the irreducible representation decomposes into absolutely irreducible constituents. Minimal splitting and detachment fields are described by embedding maximal commutative subfields into the simple system. Their degrees are multiples of the product of the degree of the center and the index of the algebra.

If the center is separable, the absolutely irreducible constituents are permuted by the Galois group of the center. If the center is inseparable, one works with the semisimple quotient after extending scalars and accounts for the radical by composition factors. In either case the normal representation encodes the field-theoretic data through the algebra's center and index.

\subsection{Crossed products as algebras and ideals}

The crossed product \(A=(a,K,G)\) can be recovered from the group extension \(G^*\) by adjoining addition. Conversely, from \(A\) one recovers \(G^*\) as the set of regular elements normalizing \(K\):
\[
  G^*=\bigcup_{S\in G}u_SK^{\times}.
\]
The quotient of \(G^*\) by \(K^{\times}\) is \(G\). Thus the multiplicative structure of the algebra already contains the Galois group and the factor system.

Ideals in crossed products are treated by the same method. A left ideal \(L\) of the crossed product gives an endomorphism ring
\[
  \operatorname{End}_A(L),
\]
which is a similar algebra. If \(L\) has the minimal possible absolute length, its endomorphism ring is the division algebra opposite to the one represented by \(A\). This supplies the inverse Brauer class. Products of ideals correspond to multiplication of factor-system classes. The ideal-theoretic view is therefore not separate from the factor-system view; it is its arithmetic realization.

When a maximal order is present, one chooses integral representatives of the factor system. Ideals of the maximal order transform the order into equivalent maximal orders. Locally this is Brandt's theory; globally it becomes the theory used in the principal genus theorem. The reason the same formulas appear in both places is that a factor system is simply the multiplication table of a normalizer, and an ideal class is the arithmetic substitute for a scalar multiplier.

\subsection{Cyclic algebras and norm residues}

For a cyclic extension \(K/P\) of degree \(n\) with generator \(S\), the crossed product has the form
\[
  A=(a,K,S)=K+uK+\cdots+u^{n-1}K,
\]
\[
  zu=uz^S,
\]
\[
  u^n=a.
\]
Changing \(u\) to \(uc\) changes \(a\) by a norm:
\[
  a\mapsto aN(c).
\]
Thus the class depends on the residue class of \(a\) modulo the norm group.

If \(a\) is a norm, say \(a=N(c)^{-1}\), then \(v=uc\) satisfies \(v^n=1\) and the algebra splits. Conversely, if the algebra splits, the unit factor system is associated to the cyclic factor system, so \(a\) is a norm. This is the algebraic form of the norm theorem for cyclic algebras.

Over a finite field every element of the ground field is a norm from a finite extension. Therefore every cyclic algebra over a finite field splits. Since every finite division algebra has a maximal commutative subfield which is finite and cyclic over the center, it follows that the division algebra is equal to its center. This is Wedderburn's theorem that every finite skew field is commutative.

Over a local field the norm group has finite index and is described by valuation and residue data. The cyclic algebra \((a,K,S)\) is split exactly when \(a\) lies in the local norm group. If \(K/P\) is unramified, every unit is a norm and the invariant depends only on the valuation of \(a\). This is the local invariant of the cyclic algebra. It is the same local computation used in the proof of the principal genus theorem: once the local factor system has unit parameter in an unramified extension, the local algebra splits.

\subsection{Closing synthesis of the lecture}

The lecture has passed through the following chain of ideas. A representation is best understood as a module. A module is best understood through its operator ring. Simple operator rings are matrix rings over division rings. Extending coefficients decomposes these rings into simple components. When a Galois group acts on the components, the obstruction to choosing the component isomorphisms multiplicatively is a factor system. A factor system is the multiplication law of a crossed product. Crossed products are exactly the central simple algebras with a chosen Galois maximal subfield. Their arithmetic is the arithmetic of their orders and ideals.

This chain turns the older character-theoretic representation theory into ring theory. It also explains why the same formulas occur in Galois theory, class-field theory, the theory of maximal orders, and the arithmetic of hypercomplex systems. Each occurrence is a different expression of the same mechanism: inner automorphisms replace external automorphisms, factor systems measure their failure to multiply strictly, and ideals measure the arithmetic ambiguity in choosing the representatives.

\end{document}
